% Generated mechanically from frozen input p09/ko/src/spaces-perfect.tex.
% Do not edit this generated file; edit the bound locale source instead.

% OK, start here.
%

\setcounter{chapter}{74}
\chapter{대수공간의 유도 범주}



\phantomsection
\label{spaces-perfect-section-phantom}


\section{서론}
\label{spaces-perfect-section-introduction}

\noindent
이 장에서는 대수공간 위 가군의 유도 범주를 논한다. 이 주제를 다루는 좋은
입문 문헌은 없는 듯하다. 문헌에서는 대수 스택 위 가군의 유도 범주를 다룬
논문을 인용하는 방식으로 이 주제를 다루며, 예를 들어
\cite{olsson_sheaves}를 보라.



\section{규약}
\label{spaces-perfect-section-conventions}

\noindent
$\mathcal{A}$가 아벨 범주이고 $M$이 $\mathcal{A}$의 대상이면, $M$을
$K(\mathcal{A})$ 및/또는 $D(\mathcal{A})$의 대상으로도 쓴다. 이는 $M$이
차수 $0$에 있고 그 밖의 모든 차수에서 영인 복합체에 대응하는 대상이다.

\medskip\noindent
환 $A$가 주어지면 $K(A)$는 $A$-가군 복합체의 호모토피 범주를 나타내고
$D(A)$는 대응하는 유도 범주를 나타낸다. 마찬가지로 환 달린 공간
$(X,\mathcal{O}_X)$가 주어지면 $K(\mathcal{O}_X)$는
$\mathcal{O}_X$-가군 복합체의 호모토피 범주를 나타내고
$D(\mathcal{O}_X)$는 대응하는 유도 범주를 나타낸다.




\section{일반론}
\label{spaces-perfect-section-generalities}

\noindent
이 절에서는 대수공간 위 가군의 비유계 복합체의 코호몰로지에 관한 몇 가지
일반 결과를 제시한다.

\begin{lemma}
\label{spaces-perfect-lemma-restrict-direct-image-open}
$S$를 스킴이라 하자. $f:X\to Y$를 $S$ 위 대수공간의 사상이라 하자.
에탈 사상 $V\to Y$가 주어졌을 때 $U=V\times_YX$라 두고, 사영 사상을
$g:U\to V$라 쓰자. 그러면 $(Rf_*E)|_V=Rg_*(E|_U)$이다. 여기서
$E$는 $D(\mathcal{O}_X)$의 대상이다.
\end{lemma}

\begin{proof}
$E$를 K-단사 복합체 $\mathcal{I}^\bullet$로 나타내자. 이는
$\mathcal{O}_X$-가군의 복합체이다. 그러면
$Rf_*(E)=f_*\mathcal{I}^\bullet$이고, Cohomology on Sites, Lemma
\ref{sites-cohomology-lemma-restrict-K-injective-to-open}에 의해
$Rg_*(E|_U)=g_*(\mathcal{I}^\bullet|_U)$이다. 따라서 결론은
Properties of Spaces, Lemma
\ref{spaces-properties-lemma-pushforward-etale-base-change-modules}에서
따른다.
\end{proof}










\begin{definition}
\label{spaces-perfect-definition-supported-on}
$S$를 스킴이라 하자. $X$를 $S$ 위 대수공간이라 하자.
$E$를 $D(\mathcal{O}_X)$의 대상이라 하자.
$T\subset|X|$를 닫힌 부분집합이라 하자. $E$가 {\it $T$ 위에
지지된다}고 하는 것은 코호몰로지 층 $H^i(E)$가 $T$ 위에 지지될 때이다.
\end{definition}








\section{작은 에탈 사이트 위 준연접 가군의 유도 범주}
\label{spaces-perfect-section-derived-quasi-coherent-etale}

\noindent
$X$를 스킴이라 하자. 이 절에서는 $D_\QCoh(\mathcal{O}_X)$를 작은 에탈
사이트 $X_\etale$로 정의할 수 있음을 보인다. 이는 $X$의 작은 에탈
사이트이다.
$\mathcal{O}_\etale$를 $X_\etale$ 위 구조층이라 쓰자.
환 달린 사이트의 사상
\begin{equation}
\label{spaces-perfect-equation-epsilon}
\epsilon :
(X_\etale, \mathcal{O}_\etale)
\longrightarrow
(X_{Zar}, \mathcal{O}_X).
\end{equation}
을 생각하자. 이는 Descent, Lemma \ref{descent-lemma-compare-sites}에서
$\text{id}_{small, \etale, Zar}$로 쓰는 사상이다.

\begin{lemma}
\label{spaces-perfect-lemma-epsilon-flat}
(\ref{spaces-perfect-equation-epsilon})의 사상 $\epsilon$은 환 달린 사이트의 평탄
사상이다. 특히 함자
$\epsilon^* : \textit{Mod}(\mathcal{O}_X) \to
\textit{Mod}(\mathcal{O}_\etale)$
는 완전하다. 더 나아가 $\epsilon^*\mathcal{F}=0$이면
$\mathcal{F}=0$이다.
\end{lemma}

\begin{proof}
사상 $\epsilon$의 평탄성은 Descent, Lemma
\ref{descent-lemma-compare-etale-zariski-flat}이다.
다른 증명을 주자. $\mathcal{O}_\etale$가 평탄한
$\epsilon^{-1}\mathcal{O}_X$-가군임을 보여야 한다. 이를 위해서는 사상
$\mathcal{O}_{X,x}\to\mathcal{O}_{\etale,\overline{x}}$가 임의의
기하점 $\overline{x}$가 $X$의 점일 때 평탄임을
확인하면 충분하다. Modules on Sites, Lemma
\ref{sites-modules-lemma-check-flat-stalks}, Sites, Lemma
\ref{sites-lemma-point-morphism-sites}, 그리고 \'Etale Cohomology,
Remarks \ref{etale-cohomology-remarks-enough-points}를 보라.
\'Etale Cohomology, Lemma
\ref{etale-cohomology-lemma-describe-etale-local-ring}에 의해
$\mathcal{O}_{\etale,\overline{x}}$는 $\mathcal{O}_{X,x}$의 강
Hensel화이다. 따라서
$\mathcal{O}_{X,x}\to\mathcal{O}_{\etale,\overline{x}}$는
More on Algebra, Lemma
\ref{more-algebra-lemma-dumb-properties-henselization}에 의해 충실
평탄이다.

\medskip\noindent
$\epsilon^*$의 완전성은 Modules on Sites, Lemma
\ref{sites-modules-lemma-flat-pullback-exact}에 의해 $\epsilon$의
평탄성에서 따른다.

\medskip\noindent
$\mathcal{F}$를 $\mathcal{O}_X$-가군이라 하자.
$\epsilon^*\mathcal{F}=0$이라면 위의 기호 아래
$$
0 = \epsilon^*\mathcal{F}_{\overline{x}} =
\mathcal{F}_x \otimes_{\mathcal{O}_{X, x}} \mathcal{O}_{\etale, \overline{x}}
$$
이다(Modules on Sites, Lemma
\ref{sites-modules-lemma-pullback-stalk}). 이는 모든 기하점
$\overline{x}$에 대해 성립한다. 사상
$\mathcal{O}_{X,x}\to\mathcal{O}_{\etale,\overline{x}}$의 충실
평탄성에 의해 $\mathcal{F}_x=0$임을 모든 $x\in X$에 대해 얻는다.
\end{proof}

\noindent
$X$를 스킴이라 하고 (\ref{spaces-perfect-equation-epsilon})의 기호를 쓰자.
Descent, Proposition
\ref{descent-proposition-equivalence-quasi-coherent}과 Remark
\ref{descent-remark-change-topologies-ringed-sites}에 의해
$\epsilon^*:\QCoh(\mathcal{O}_X)\to\QCoh(\mathcal{O}_\etale)$가
동치임을 상기하자. 더 나아가 Descent, Lemma
\ref{descent-lemma-equivalence-quasi-coherent-limits}에 의해
$\QCoh(\mathcal{O}_\etale)$는
$\textit{Mod}(\mathcal{O}_\etale)$의 Serre 부분범주를 이룬다.
따라서 $D_\QCoh(\mathcal{O}_\etale)$를
$D(\mathcal{O}_\etale)$의 삼각 부분범주로 둘 수 있으며, 그 대상은
준연접 코호몰로지 층을 갖는 복합체들이다. Derived Categories,
Section \ref{derived-section-triangulated-sub}을 보라. 함자
$\epsilon^*$는 완전하므로(Lemma \ref{spaces-perfect-lemma-epsilon-flat})
$$\epsilon^* :  D(\mathcal{O}_X) \to D(\mathcal{O}_\etale)$$
를 유도한다. 또한 준연접 가군의 당김은 준연접이므로
$$\epsilon^* : D_\QCoh(\mathcal{O}_X) \to
D_\QCoh(\mathcal{O}_\etale)$$
도 유도한다.

\begin{lemma}
\label{spaces-perfect-lemma-derived-quasi-coherent-small-etale-site}
$X$를 스킴이라 하자. 위에서 정의한 함자
$\epsilon^*:D_\QCoh(\mathcal{O}_X)\to
D_\QCoh(\mathcal{O}_\etale)$는 동치이다.
\end{lemma}

\begin{proof}
함자
$R\epsilon_*:D(\mathcal{O}_\etale)\to D(\mathcal{O}_X)$가
준역함자를 유도함을 보여 증명하겠다. 여기서는 $\epsilon_*$가
$X_{Zar}\subset X_\etale$로의 제한으로 주어진다는 사실과 Descent,
Lemma \ref{descent-lemma-compare-sites}에서의
$\epsilon^*=\text{id}_{small,\etale,Zar}^*$의 기술을 자유롭게 쓴다.

\medskip\noindent
준연접 $\mathcal{O}_X$-가군 $\mathcal{F}$에 대한 수반 사상
$\mathcal{F}\to\epsilon_*\epsilon^*\mathcal{F}$는 동형이다. 실제로
$\mathcal{F}^a$가 층이라는 사실은 Descent, Definition
\ref{descent-definition-structure-sheaf}과 Descent, Lemma
\ref{descent-lemma-sheaf-condition-holds}에서 증명되었다.
반대로 모든 준연접 $\mathcal{O}_\etale$-가군 $\mathcal{H}$는
$\epsilon^*\mathcal{F}$의 꼴이다. 여기서 $\mathcal{O}_X$-가군
$\mathcal{F}$는 준연접이다. Descent, Proposition
\ref{descent-proposition-equivalence-quasi-coherent}을 보라.
위에서 말한 바에 의해 $\mathcal{F}=\epsilon_*\mathcal{H}$이고, 따라서
수반 사상 $\epsilon^*\epsilon_*\mathcal{H}\to\mathcal{H}$는 모든
준연접 $\mathcal{O}_\etale$-가군 $\mathcal{H}$에 대해 동형이다.

\medskip\noindent
$E$를 $D_\QCoh(\mathcal{O}_\etale)$의 대상이라 하고,
$\mathcal{H}^q=H^q(E)$를 그 $q$차 코호몰로지 층이라 쓰자. $\mathcal{B}$를
$X_\etale$의 아핀 대상들의 집합이라 하자.
$H^p(U,\mathcal{H}^q)=0$이 모든 $p>0$, 모든 $q\in\mathbf{Z}$, 모든
$U\in\mathcal{B}$에 대해 성립한다. Descent, Proposition
\ref{descent-proposition-same-cohomology-quasi-coherent}과
Cohomology of Schemes, Lemma
\ref{coherent-lemma-quasi-coherent-affine-cohomology-zero}를 보라.
Cohomology on Sites, Lemma
\ref{sites-cohomology-lemma-cohomology-over-U-trivial}에 의해 이는
$$
H^q(U, E) = H^0(U, \mathcal{H}^q)
$$
를 뜻하며, 모든 $U\in\mathcal{B}$에 대해 성립한다. 특히 아핀 열린집합
$U\subset X$에 대해 성립한다. 따라서 $q$차 코호몰로지는
$R\epsilon_*E$에 대해 $U$ 위에서 층 $\epsilon_*\mathcal{H}^q$의 $U$에서의
값이다. 층화하면
$$
H^q(R\epsilon_*E)=\epsilon_*\mathcal{H}^q
$$
를 얻는다. 특히 $R\epsilon_*$는 함자
$D_\QCoh(\mathcal{O}_\etale)\to D_\QCoh(\mathcal{O}_X)$를 유도한다.
$\epsilon^*$가 완전하므로
$H^q(\epsilon^*R\epsilon_*E)=\epsilon^*\epsilon_*\mathcal{H}^q=
\mathcal{H}^q$를 얻는다(위의 논의를 보라). 따라서 수반 사상
$\epsilon^*R\epsilon_*E\to E$는 동형이다.

\medskip\noindent
반대로 $F\in D_\QCoh(\mathcal{O}_X)$에 대한 수반 사상
$F\to R\epsilon_*\epsilon^*F$도 같은 이유로 동형이다. 즉
$R\epsilon_*\epsilon^*F$의 코호몰로지 층들은
$\epsilon_*H^m(\epsilon^*F)=\epsilon_*\epsilon^*H^m(F)=H^m(F)$와
동형이다.
\end{proof}











\section{준연접 가군의 유도 범주}
\label{spaces-perfect-section-derived-quasi-coherent}

\noindent
$S$를 스킴이라 하자. Lemma
\ref{spaces-perfect-lemma-derived-quasi-coherent-small-etale-site}에 의해 범주
$D_\QCoh(\mathcal{O}_S)$는 $\mathcal{O}_S$-가군의 복합체로 스킴 $S$
위에서 정의할 수도 있고, $\mathcal{O}$-가군의 복합체로 $S$의 작은 에탈
사이트 위에서 정의할 수도 있다. 따라서 다음 정의는 스킴의 경우의 정의와
양립한다.

\begin{definition}
\label{spaces-perfect-definition-derived-quasi-coherent}
$S$를 스킴이라 하자. $X$를 $S$ 위 대수공간이라 하자.
{\it 준연접 코호몰로지 층을 갖는 $\mathcal{O}_X$-가군의 유도 범주}를
$D_\QCoh(\mathcal{O}_X)$로 쓴다.
\end{definition}

\noindent
이는 Properties of Spaces, Lemma
\ref{spaces-properties-lemma-properties-quasi-coherent}와 Derived
Categories, Section \ref{derived-section-triangulated-sub}에 의해
의미가 있다. 따라서 표준 함자
\begin{equation}
\label{spaces-perfect-equation-compare}
D(\QCoh(\mathcal{O}_X))
\longrightarrow
D_\QCoh(\mathcal{O}_X)
\end{equation}
를 얻는다. Derived Categories, Equation
(\ref{derived-equation-compare})을 보라.

\medskip\noindent
대수공간의 평탄 사상 $f:Y\to X$는 완전 함자
$f^*:\textit{Mod}(\mathcal{O}_X)\to\textit{Mod}(\mathcal{O}_Y)$를
유도한다. Morphisms of Spaces, Lemma
\ref{spaces-morphisms-lemma-flat-morphism-sites}와 Modules on Sites,
Lemma \ref{sites-modules-lemma-flat-pullback-exact}를 보라.
특히 $Lf^*:D(\mathcal{O}_X)\to D(\mathcal{O}_Y)$는 임의의 대표 복합체
위에서 계산된다(Derived Categories, Lemma
\ref{derived-lemma-right-derived-exact-functor}).
$Lf^*=f^*$로 쓰는 것은 $f$가 평탄할 때이며, 이 경우
$H^i(f^*E)=f^*H^i(E)$가 $E$가 $D(\mathcal{O}_X)$에 속할 때 성립한다.
$f$가 에탈일 때 이를 자주 쓸 것이다.
물론 에탈인 경우 당김 함자는 단지 $Y_\etale$로의 제한이다.
Properties of Spaces, Equation
(\ref{spaces-properties-equation-restrict-modules})을 보라.

\begin{lemma}
\label{spaces-perfect-lemma-check-quasi-coherence-on-covering}
$S$를 스킴이라 하자. $X$를 $S$ 위 대수공간이라 하자.
$E$를 $D(\mathcal{O}_X)$의 대상이라 하자. 다음은 서로 동치이다.
\begin{enumerate}
\item $E$는 $D_\QCoh(\mathcal{O}_X)$에 속한다.
\item 임의의 에탈 사상 $\varphi:U\to X$에서 $U$가 아핀 스킴이면
$\varphi^*E$는 $D_\QCoh(\mathcal{O}_U)$의 대상이다.
\item 임의의 에탈 사상 $\varphi:U\to X$에서 $U$가 스킴이면
$\varphi^*E$는 $D_\QCoh(\mathcal{O}_U)$의 대상이다.
\item 전사 에탈 사상 $\varphi:U\to X$가 존재하여 $U$가 스킴이고
$\varphi^*E$가 $D_\QCoh(\mathcal{O}_U)$의 대상이다.
\item 대수공간의 전사 에탈 사상 $f:Y\to X$가 존재하여
$Lf^*E$가 $D_\QCoh(\mathcal{O}_Y)$의 대상이다.
\end{enumerate}
\end{lemma}

\begin{proof}
보조정리 앞의 논의와 Properties of Spaces, Lemma
\ref{spaces-properties-lemma-characterize-quasi-coherent}에서 바로
따른다.
\end{proof}

\begin{lemma}
\label{spaces-perfect-lemma-quasi-coherence-direct-sums}
$S$를 스킴이라 하자. $X$를 $S$ 위 대수공간이라 하자. 그러면
$D_\QCoh(\mathcal{O}_X)$에는 직합이 존재한다.
\end{lemma}

\begin{proof}
Injectives, Lemma \ref{injectives-lemma-derived-products}에 의해 유도 범주
$D(\mathcal{O}_X)$에는 직합이 존재하고, 이는 임의의 대표에서 항별
직합을 취하여 계산된다. 따라서 직합을 취하는 것이 완전 함자이므로
(임의의 Grothendieck 아벨 범주에서) 직합의 코호몰로지 층은 코호몰로지
층들의 직합임이 분명하다. 준연접 층들의 직합은 준연접이므로 보조정리가
따른다. Properties of Spaces, Lemma
\ref{spaces-properties-lemma-properties-quasi-coherent}를 보라.
\end{proof}

\noindent
유도 극한에 관한 몇 가지 정보가 필요하다. 아래 보조정리에서 유도 극한은
일반적으로 $D_\QCoh$의 대상이 아님을 주의한다.

\begin{lemma}
\label{spaces-perfect-lemma-Rlim-quasi-coherent}
$S$를 스킴이라 하자. $X$를 $S$ 위 대수공간이라 하자.
$(K_n)$을 $D_\QCoh(\mathcal{O}_X)$의 역계라 하고, 그 유도 극한을
$K=R\lim K_n$이라 하자. 이는 $D(\mathcal{O}_X)$ 안에서 취한다.
$H^q(K_{n+1})\to H^q(K_n)$이 모든 $q\in\mathbf{Z}$와 $n\geq1$에
대해 전사라고 가정하자. 그러면
\begin{enumerate}
\item $H^q(K)=\lim H^q(K_n)$이다.
\item $R\lim H^q(K_n)=\lim H^q(K_n)$이다.
\item 모든 아핀 열린집합 $U\subset X$에 대해
$H^p(U,\lim H^q(K_n))=0$이고 이는 $p>0$일 때 성립한다.
\end{enumerate}
\end{lemma}

\begin{proof}
$\mathcal{B}\subset\Ob(X_\etale)$를 아핀 대상들의 집합이라 하자.
$H^q(K_n)$이 준연접이므로 $H^p(U,H^q(K_n))=0$이
$U\in\mathcal{B}$에 대해 성립한다. Cohomology of Spaces, Section
\ref{spaces-cohomology-section-higher-direct-image}의 논의와 Cohomology
of Schemes, Lemma
\ref{coherent-lemma-quasi-coherent-affine-cohomology-zero}를 보라.
더 나아가 같은 이유로 사상
$H^0(U,H^q(K_{n+1}))\to H^0(U,H^q(K_n))$은
$U\in\mathcal{B}$에 대해 전사이다. 방금 확인한 조건을 Cohomology on
Sites, Lemma
\ref{sites-cohomology-lemma-derived-limit-suitable-system}에 적용하면
(1)이 따른다. (2)와 (3)은 Cohomology on Sites, Lemma
\ref{sites-cohomology-lemma-inverse-limit-is-derived-limit}에서 따른다.
\end{proof}

\begin{lemma}
\label{spaces-perfect-lemma-quasi-coherence-pullback}
$S$를 스킴이라 하자.
$f:Y\to X$를 $S$ 위 대수공간의 사상이라 하자.
함자 $Lf^*$는 $D_\QCoh(\mathcal{O}_X)$를
$D_\QCoh(\mathcal{O}_Y)$ 안으로 보낸다.
\end{lemma}

\begin{proof}
다음 도식을 택하자.
$$
\xymatrix{
U \ar[d]_a \ar[r]_h & V \ar[d]^b \\
X \ar[r]^f & Y
}
$$
여기서 $U$와 $V$는 스킴이고 수직 화살표들은 에탈이며,
$a$는 전사이다.% P09-ERR-0393
$a^*\circ Lf^*=Lh^*\circ b^*$이므로, 결론은 Lemma
\ref{spaces-perfect-lemma-check-quasi-coherence-on-covering}과 스킴의 경우, 즉
Derived Categories of Schemes, Lemma
\ref{perfect-lemma-quasi-coherence-pullback}에서 따른다.
\end{proof}

\begin{lemma}
\label{spaces-perfect-lemma-quasi-coherence-tensor-product}
$S$를 스킴이라 하자. $X$를 $S$ 위 대수공간이라 하자.
대상 $K,L$을 $D_\QCoh(\mathcal{O}_X)$에서 택하면 유도 텐서곱
$K\otimes^\mathbf{L}L$은 $D_\QCoh(\mathcal{O}_X)$에 속한다.
\end{lemma}

\begin{proof}
$\varphi:U\to X$를 스킴 $U$로부터의 전사 에탈 사상이라 하자.
$$
\varphi^*(K \otimes_{\mathcal{O}_X}^\mathbf{L} L) =
\varphi^*K \otimes_{\mathcal{O}_U}^\mathbf{L} \varphi^*L
$$
이므로 Lemma \ref{spaces-perfect-lemma-check-quasi-coherence-on-covering}에 의해 스킴의
경우에서 결론이 따른다. 스킴의 경우는 Derived Categories of Schemes,
Lemma \ref{perfect-lemma-quasi-coherence-tensor-product}이다.
\end{proof}

\noindent
다음 보조정리는 $D_\QCoh(\mathcal{O}_X)$의 대상 위에서 우유도 함자를
계산하는 데 도움이 된다.

\begin{lemma}
\label{spaces-perfect-lemma-nice-K-injective}
$S$를 스킴이라 하자. $X$를 $S$ 위 대수공간이라 하자. $E$를
$D_\QCoh(\mathcal{O}_X)$의 대상이라 하자. 그러면 Derived Categories,
Remark \ref{derived-remark-map-into-derived-limit-truncations}의 사상
$E\to R\lim\tau_{\geq-n}E$는 동형이다.\footnote{특히 $E$에는 K-단사
대표가 존재한다. Derived Categories, Lemma
\ref{derived-lemma-difficulty-K-injectives}를 보라.}
\end{lemma}

\begin{proof}
$\mathcal{H}^i=H^i(E)$를 $i$차 코호몰로지 층, 즉 $E$의 코호몰로지
층이라 하자.
$\mathcal{B}$를 $X_\etale$의 아핀 대상들의 집합이라 하자.
$H^p(U,\mathcal{H}^i)=0$이 모든 $p>0$, 모든 $i\in\mathbf{Z}$, 모든
$U\in\mathcal{B}$에 대해 성립한다. 이는 $U$가 아핀 스킴이기 때문이다.
Cohomology of Spaces, Section
\ref{spaces-cohomology-section-higher-direct-image}의 논의와 Cohomology
of Schemes, Lemma
\ref{coherent-lemma-quasi-coherent-affine-cohomology-zero}를 보라.
따라서 $d=0$인 Cohomology on Sites, Lemma
\ref{sites-cohomology-lemma-is-limit-dimension}에서 보조정리가 따른다.
\end{proof}

\begin{lemma}
\label{spaces-perfect-lemma-application-nice-K-injective}
$S$를 스킴이라 하자. $X$를 $S$ 위 대수공간이라 하자.
$F:\textit{Mod}(\mathcal{O}_X)\to\textit{Ab}$를 함자라 하고
$N\geq0$을 정수라 하자. 다음을 가정하자.
\begin{enumerate}
\item $F$는 좌완전이다.
\item $F$는 가산 직적과 가환한다.
\item $R^pF(\mathcal{F})=0$이 모든 $p\geq N$과 준연접
$\mathcal{F}$에 대해 성립한다.
\end{enumerate}
그러면 $E\in D_\QCoh(\mathcal{O}_X)$에 대해 다음이 성립한다.
\begin{enumerate}
\item $H^i(RF(\tau_{\leq a}E)\to H^i(RF(E))$는 $i\leq a$일 때
동형이다.% P09-ERR-0394
\item $H^i(RF(E))\to H^i(RF(\tau_{\geq b-N+1}E))$는 $i\geq b$일 때
동형이다.
\item $H^i(E)=0$이 $i\not\in[a,b]$에 대해 성립한다고 하자. 여기서 어떤
$-\infty\leq a\leq b\leq\infty$를 택한다. 그러면
$H^i(RF(E))=0$이 $i\not\in[a,b+N-1]$에 대해 성립한다.
\end{enumerate}
\end{lemma}

\begin{proof}
(1)은 Derived Categories, Lemma
\ref{derived-lemma-negative-vanishing}이다.

\medskip\noindent
(2)를 증명하자. $E_n=\tau_{\geq-n}E$라 쓰자. Lemma
\ref{spaces-perfect-lemma-nice-K-injective}에 의해 $E=R\lim E_n$이다. 따라서
Injectives, Lemma \ref{injectives-lemma-RF-commutes-with-Rlim}에 의해
$RF(E)=R\lim RF(E_n)$가 $D(\textit{Ab})$에서 성립한다. 그러므로 모든
$i\in\mathbf{Z}$에 대해 짧은 완전열
$$
0 \to R^1\lim H^{i - 1}(RF(E_n)) \to H^i(RF(E)) \to \lim H^i(RF(E_n)) \to 0
$$
이 있다. More on Algebra, Remark
\ref{more-algebra-remark-compare-derived-limit}을 보라. (2)를 증명하기
위해 왼쪽 항이 영이고 오른쪽 항이
$H^i(RF(E_{-b+N-1})$와 같음을 보이겠다.% P09-ERR-0395
여기서 $b$는 $i\geq b$를 만족하는 임의의 수이다.

\medskip\noindent
모든 $n$에 대해 구별 삼각형
$$
H^{-n}(E)[n] \to E_n \to E_{n - 1} \to H^{-n}(E)[n + 1]
$$
이 $D(\mathcal{O}_X)$ 안에 있다(Derived Categories, Remark
\ref{derived-remark-truncation-distinguished-triangle}).
$H^{-n}(E)$가 준연접이므로
$$
H^i(RF(H^{-n}(E)[n])) = R^{i + n}F(H^{-n}(E)) = 0
$$
이 $i+n\geq N$일 때 성립하고,
$$
H^i(RF(H^{-n}(E)[n + 1])) = R^{i + n + 1}F(H^{-n}(E)) = 0
$$
이 $i+n+1\geq N$일 때 성립한다. 따라서
$$
H^i(RF(E_n)) \to H^i(RF(E_{n - 1}))
$$
은 $n\geq N-i$일 때 동형이다. 그러므로 계 $H^i(RF(E_n))$들은 모두
ML 조건을 만족하고, 위 짧은 완전열의 $R^1\lim$ 항은 영이다(More on
Algebra, Section \ref{more-algebra-section-Rlim}의 논의를 보라).
더 나아가 계 $H^i(RF(E_n))$은 원하는 대로 $n=N-i-1$부터 상수이다.

\medskip\noindent
(3)을 증명하자. $E$에 대한 가정 아래 $\tau_{\leq a-1}E=0$이고, (1)에
의해 $H^i(RF(E))$의 소멸을 $i\leq a-1$에 대해 얻는다. 마찬가지로
$\tau_{\geq b+1}E=0$이므로 (2)에 의해 $H^i(RF(E))$의 소멸을
$i\geq b+n$에 대해 얻는다.
% P09-ERR-0396
\end{proof}










\section{전체 직접상}
\label{spaces-perfect-section-total-direct-image}

\noindent
다음 보조정리는 Cohomology of Spaces, Lemma
\ref{spaces-cohomology-lemma-vanishing-higher-direct-images}에 대응한다.

\begin{lemma}
\label{spaces-perfect-lemma-quasi-coherence-direct-image}
$S$를 스킴이라 하자. $f:X\to Y$를 $S$ 위 대수공간의 준분리이고
준콤팩트인 사상이라 하자.
\begin{enumerate}
\item 함자 $Rf_*$는 $D_\QCoh(\mathcal{O}_X)$를
$D_\QCoh(\mathcal{O}_Y)$ 안으로 보낸다.
\item $Y$가 준콤팩트이면 정수 $N=N(X,Y,f)$가 존재하여, 대상 $E$를
$D_\QCoh(\mathcal{O}_X)$에서 택하고 $H^m(E)=0$이 $m>0$에 대해
성립할 때 $H^m(Rf_*E)=0$이 $m\geq N$에 대해 성립한다.
\item 실제로 $Y$가 준콤팩트이면 $N=N(X,Y,f)$를 택하여, 대수공간의
모든 사상 $Y'\to Y$에 대해 같은 결론이 함자 $R(f')_*$에 대해
성립하게 할 수 있다. 여기서 $f':X'\to Y'$는 $f$의 밑변환이다.
\end{enumerate}
\end{lemma}

\begin{proof}
$E$를 $D_\QCoh(\mathcal{O}_X)$의 대상이라 하자. (1)을 증명하려면
$Rf_*E$가 준연접 코호몰로지 층을 가짐을 보여야 한다. 이 문제는 $Y$
위에서 국소적이므로 $Y$가 준콤팩트라고 가정해도 된다. Cohomology of
Spaces, Lemma
\ref{spaces-cohomology-lemma-vanishing-higher-direct-images}에서처럼
$N=N(X,Y,f)$를 택하자. 그러면 $R^pf_*\mathcal{F}=0$이 모든 준연접
$\mathcal{O}_X$-가군 $\mathcal{F}$와 모든 $p\geq N$에 대해 성립한다.
또한 $R^pf_*\mathcal{F}$는 모든 $p$에 대해 준연접이다(Cohomology of
Spaces, Lemma \ref{spaces-cohomology-lemma-higher-direct-image}).
이 명제들은 밑변환 뒤에도 참이다.

\medskip\noindent
먼저 $E$가 아래로 유계라고 가정하자. 이렇게 주어진 $E$에 대해서는
우리가 택한 $N$으로 (1), (2), (3)이 성립함을 보이겠다. 이 경우 예를
들어 스펙트럼 열
$$
R^pf_*H^q(E) \Rightarrow R^{p + q}f_*E
$$
(Derived Categories, Lemma
\ref{derived-lemma-two-ss-complex-functor})과
$R^pf_*H^q(E)$의 준연접성 및 $R^pf_*H^q(E)$가 $p\geq N$일 때
소멸한다는 사실을 사용하면 (1), (2), (3)이 성립함을 알 수 있다.

\medskip\noindent
다음으로 (2)와 (3)을 증명하자. $H^m(E)=0$이 $m>0$에 대해 성립한다고
하자. $V$를 $Y_\etale$의 아핀 대상이라 하자.
$H^p(V\times_YX,\mathcal{F})=0$이 $p\geq N$에 대해 성립한다.
Cohomology of Spaces, Lemma
\ref{spaces-cohomology-lemma-quasi-coherence-higher-direct-images-application}
을 보라. 따라서 함자 $\Gamma(V\times_YX,-)$에 Lemma
\ref{spaces-perfect-lemma-application-nice-K-injective}를 적용하여
$$
R\Gamma(V, Rf_*E) = R\Gamma(V \times_Y X, E)
$$
의 코호몰로지가 차수 $\geq N$에서 소멸함을 얻는다. 이는 모든 아핀
$V$가 $Y_\etale$에 속할 때 성립하므로, $H^m(Rf_*E)=0$이
$m\geq N$에 대해 성립함을 얻는다.

\medskip\noindent
이제 일반적인 경우에 (1)을 증명하자. 구별 삼각형
$$
\tau_{\leq -n - 1}E \to E \to \tau_{\geq -n}E \to
(\tau_{\leq -n - 1}E)[1]
$$
이 $D(\mathcal{O}_X)$ 안에 있음을 상기하자. Derived Categories, Remark
\ref{derived-remark-truncation-distinguished-triangle}을 보라.
(2)에 의해 $Rf_*\tau_{\leq-n-1}E$의 코호몰로지 층은 차수
$\geq-n+N$에서 소멸한다. 따라서 정수 $q$가 주어지면
$R^qf_*E$가 $R^qf_*\tau_{\geq-n}E$와 같도록 하는 어떤 $n$이 존재하고,
위의 결과를 적용할 수 있다.
\end{proof}

\begin{lemma}
\label{spaces-perfect-lemma-quasi-coherence-pushforward-direct-sums}
$S$를 스킴이라 하자. $f:X\to Y$를 $S$ 위 대수공간의 준분리이고
준콤팩트인 사상이라 하자. 그러면
$Rf_*:D_\QCoh(\mathcal{O}_X)\to D_\QCoh(\mathcal{O}_Y)$는 직합과
가환한다.
\end{lemma}

\begin{proof}
$E_i$를 $D_\QCoh(\mathcal{O}_X)$의 대상족이라 하고
$E=\bigoplus E_i$라 두자. 사상
$$
\bigoplus Rf_*E_i \longrightarrow Rf_*E
$$
이 동형임을 보이고자 한다. 이것이 차수 $0$의 코호몰로지 층 위에서
동형을 유도함을 보이면 보조정리가 따른다. Lemma
\ref{spaces-perfect-lemma-quasi-coherence-direct-image}에서처럼 정수 $N$을 택하자.
인용한 보조정리에 의해 $R^0f_*E=R^0f_*\tau_{\geq-N}E$이고
$R^0f_*E_i=R^0f_*\tau_{\geq-N}E_i$이다. 또한
$\tau_{\geq-N}E=\bigoplus\tau_{\geq-N}E_i$이다. 따라서 모든 $E_i$의
코호몰로지 층이 차수 $<-N$에서 소멸한다고 가정해도 된다. 이제
스펙트럼 열
$$
R^pf_*H^q(E) \Rightarrow R^{p + q}f_*E
\quad\text{and}\quad
R^pf_*H^q(E_i) \Rightarrow R^{p + q}f_*E_i
$$
(Derived Categories, Lemma
\ref{derived-lemma-two-ss-complex-functor})을 사용하면 문제는 준연접
층들의 직합인 경우로 환원된다. 이 경우는 Cohomology of Spaces, Lemma
\ref{spaces-cohomology-lemma-colimit-cohomology}에서 다룬다.
\end{proof}

\begin{remark}
\label{spaces-perfect-remark-match-total-direct-images}
$S$를 스킴이라 하자. $f:X\to Y$를 표현가능한 대수공간 $X$와 $Y$의
사상이라 하자. 이 두 공간은 $S$ 위에 있다. $f_0:X_0\to Y_0$를
$f$를 나타내는 스킴의 사상이라 하자(어색하지만 임시적인 기호이다).
그러면 도식
$$
\xymatrix{
D_\QCoh(\mathcal{O}_{X_0})
\ar@{=}[rrrrrr]_{\text{Lemma
\ref{spaces-perfect-lemma-derived-quasi-coherent-small-etale-site}}}
& & & & & &
D_\QCoh(\mathcal{O}_X) \\
D_\QCoh(\mathcal{O}_{Y_0})
\ar[u]^{Lf^*_0}
\ar@{=}[rrrrrr]^{\text{Lemma
\ref{spaces-perfect-lemma-derived-quasi-coherent-small-etale-site}}}
& & & & & &
D_\QCoh(\mathcal{O}_Y) \ar[u]_{Lf^*}
}
$$
은 가환한다(Lemma \ref{spaces-perfect-lemma-quasi-coherence-pullback}과 Derived
Categories of Schemes, Lemma
\ref{perfect-lemma-quasi-coherence-pullback}). 실제로 Lemma
\ref{spaces-perfect-lemma-derived-quasi-coherent-small-etale-site}의 동치
$D_\QCoh(\mathcal{O}_{X_0})\to D_\QCoh(\mathcal{O}_X)$와
$D_\QCoh(\mathcal{O}_{Y_0})\to D_\QCoh(\mathcal{O}_Y)$는 환 달린
사이트의 평탄 사상
$\epsilon:X_\etale\to X_{0,Zar}$와
$\epsilon:Y_\etale\to Y_{0,Zar}$를 따라 당겨오면서 생기며,
환 달린 사이트의 도식
$$
\xymatrix{
X_{0, Zar} \ar[d]_{f_0} & X_\etale \ar[l]^\epsilon \ar[d]^f \\
Y_{0, Zar} & Y_\etale \ar[l]_\epsilon
}
$$
은 가환하기 때문이다(세부사항은 생략한다). $f$가 준콤팩트이고
준분리라는 것, 동치로 $f_0$가 준콤팩트이고 준분리라고 하자. 그러면
$$
\xymatrix{
D_\QCoh(\mathcal{O}_{X_0})
\ar[d]_{Rf_{0, *}} \ar@{=}[rrrrrr]_{\text{Lemma
\ref{spaces-perfect-lemma-derived-quasi-coherent-small-etale-site}}}
& & & & & &
D_\QCoh(\mathcal{O}_X) \ar[d]^{Rf_*} \\
D_\QCoh(\mathcal{O}_{Y_0})
\ar@{=}[rrrrrr]^{\text{Lemma
\ref{spaces-perfect-lemma-derived-quasi-coherent-small-etale-site}}}
& & & & & &
D_\QCoh(\mathcal{O}_Y)
}
$$
도 가환한다고 주장한다(Lemma \ref{spaces-perfect-lemma-quasi-coherence-direct-image}와
Derived Categories of Schemes, Lemma
\ref{perfect-lemma-quasi-coherence-direct-image}).
이 또한 위에 표시한 사이트의 가환도식에서 따른다. 실제로 Lemma
\ref{spaces-perfect-lemma-derived-quasi-coherent-small-etale-site}의 증명은 함자
$R\epsilon_*$가 동치
$D_\QCoh(\mathcal{O}_X)\to D_\QCoh(\mathcal{O}_{X_0})$와
$D_\QCoh(\mathcal{O}_Y)\to D_\QCoh(\mathcal{O}_{Y_0})$를 줌을
보여준다.
\end{remark}

\begin{lemma}
\label{spaces-perfect-lemma-affine-morphism}
$S$를 스킴이라 하자. $f:X\to Y$를 $S$ 위 대수공간의 아핀 사상이라
하자. 그러면
$Rf_*:D_\QCoh(\mathcal{O}_X)\to D_\QCoh(\mathcal{O}_Y)$는 동형을
반영한다.
\end{lemma}

\begin{proof}
이 명제는 사상 $\alpha:E\to F$가 $D_\QCoh(\mathcal{O}_X)$에 속하고
$Rf_*\alpha$가 동형이면 이 사상 자체가 동형이라는 뜻이다. 이는 코호몰로지
층 위에서 확인할 수 있다. 특히 문제는 $Y$ 위에서 에탈 국소적이다.
따라서 $Y$와 그에 따라 $X$가 아핀이라고 가정해도 된다. 이 경우 문제는
스킴의 경우(Derived Categories of Schemes, Lemma
\ref{perfect-lemma-affine-morphism})로 환원된다. 이 환원에는 Lemma
\ref{spaces-perfect-lemma-derived-quasi-coherent-small-etale-site}와 Remark
\ref{spaces-perfect-remark-match-total-direct-images}를 사용한다.
\end{proof}










\begin{lemma}
\label{spaces-perfect-lemma-affine-morphism-pull-push}
$S$를 스킴이라 하자. $f:X\to Y$를 $S$ 위 대수공간의 아핀 사상이라
하자. $E$를 $D_\QCoh(\mathcal{O}_Y)$에서 택하면
$Rf_*Lf^*E=E\otimes^\mathbf{L}_{\mathcal{O}_Y}f_*\mathcal{O}_X$이다.
\end{lemma}

\begin{proof}
$f$가 아핀이므로 사상
$f_*\mathcal{O}_X\to Rf_*\mathcal{O}_X$는 동형이다(Cohomology of
Spaces, Lemma
\ref{spaces-cohomology-lemma-affine-vanishing-higher-direct-images}).
표준 사상
$E\otimes^\mathbf{L}f_*\mathcal{O}_X=
E\otimes^\mathbf{L}Rf_*\mathcal{O}_X\to Rf_*Lf^*E$
가 존재한다. 이는 사상
$$
Lf^*(E \otimes^\mathbf{L} Rf_*\mathcal{O}_X) =
Lf^*E \otimes^\mathbf{L} Lf^*Rf_*\mathcal{O}_X \longrightarrow
Lf^* E \otimes^\mathbf{L} \mathcal{O}_X = Lf^* E
$$
에 수반되고, 이 사상은 $1:Lf^*E\to Lf^*E$와 표준 사상
$Lf^*Rf_*\mathcal{O}_X\to\mathcal{O}_X$에서 생긴다. 이렇게 구성한
사상이 동형임을 확인하려면 $Y$ 위에서 국소적으로 작업해도 된다.
따라서 $Y$와 그에 따라 $X$가 아핀이라고 가정해도 된다. 이 경우 문제는
스킴의 경우(Derived Categories of Schemes, Lemma
\ref{perfect-lemma-affine-morphism-pull-push})로 환원된다. 이 환원에는
Lemma \ref{spaces-perfect-lemma-derived-quasi-coherent-small-etale-site}와 Remark
\ref{spaces-perfect-remark-match-total-direct-images}를 사용한다.
\end{proof}

\section{기저 위에서 고유함}
\label{spaces-perfect-section-proper-over-base}

\noindent
이 절은 Cohomology of Schemes, Section
\ref{coherent-section-proper-over-base}에 대응한다. 점들의 집합 위 위상과
관련된 내용에서 늘 그렇듯, 이를 대수공간으로 옮길 때에는 주의해야 한다.

\begin{lemma}
\label{spaces-perfect-lemma-closed-proper-over-base}
$S$를 스킴이라 하자. $f:X\to Y$를 $S$ 위 대수공간의 국소 유한형
사상이라 하자. $T\subset|X|$를 닫힌 부분집합이라 하자. 다음은 서로
동치이다.
\begin{enumerate}
\item 사상 $Z\to Y$는, $Z$가 $T$ 위의 유도된 축약 대수공간 구조이면
고유이다(Properties of Spaces, Definition
\ref{spaces-properties-definition-reduced-induced-space}).
\item 어떤 닫힌 부분공간 $Z\subset X$가 $|Z|=T$를 만족하고 사상
$Z\to Y$가 고유이다.
\item 임의의 닫힌 부분공간 $Z\subset X$가 $|Z|=T$를 만족하면 사상
$Z\to Y$는 고유이다.
\end{enumerate}
\end{lemma}

\begin{proof}
(3) $\Rightarrow$ (1)과 (1) $\Rightarrow$ (2)는 바로 따른다. 따라서
(2)가 (3)을 함의함을 보이면 충분하다. 독자에게 이 사실의 증명을 직접
찾아보기를 권한다. $Z'$와 $Z''$를 $T=|Z'|=|Z''|$를 만족하는 닫힌
부분공간이라 하고, $Z'\to Y$가 대수공간의 고유 사상이라고 하자.
$Z''\to Y$도 고유임을 보여야 한다. $Z'''=Z'\cup Z''$를 스킴 이론적
합집합이라 하자. Morphisms of Spaces, Definition
\ref{spaces-morphisms-definition-scheme-theoretic-intersection-union}을
보라. 그러면 $Z'''$은 $|Z'''|=T$를 만족하는 또 다른 닫힌 부분공간이다.
이는 예를 들어 Morphisms of Spaces, Lemma
\ref{spaces-morphisms-lemma-scheme-theoretic-union}의 스킴 이론적
합집합의 기술에서 따른다.
$Z''\to Z'''$이 닫힌 몰입이므로 $Z'''\to Y$가 고유임을 보이면
충분하다(Morphisms of Spaces, Lemmas
\ref{spaces-morphisms-lemma-closed-immersion-proper}와
\ref{spaces-morphisms-lemma-composition-proper}를 보라).
사상 $Z'\to Z'''$은 전단사 닫힌 몰입이고, 특히 전사이며 보편적으로
닫혀 있다. 이제 $Z'\to Y$가 분리라는 사실에 의해 $Z'''\to Y$도
분리이다(Morphisms of Spaces, Lemma
\ref{spaces-morphisms-lemma-image-universally-closed-separated}).
또한 $Z'''\to Y$도 국소 유한형이다. 이는 $X\to Y$가 국소 유한형이기
때문이다
(Morphisms of Spaces, Lemmas
\ref{spaces-morphisms-lemma-immersion-locally-finite-type}와
\ref{spaces-morphisms-lemma-composition-finite-type}).
$Z'\to Y$가 준콤팩트이고 $Z'\to Z'''$이 보편 위상동형이므로
$Z'''\to Y$가 준콤팩트임을 알 수 있다. 마지막으로 $Z'\to Y$가
보편적으로 닫혀 있으므로 Morphisms of Spaces, Lemma
\ref{spaces-morphisms-lemma-image-proper-is-proper}에 의해
$Z'''\to Y$도 그러하다. 이로써 증명이 끝난다.
\end{proof}

\begin{definition}
\label{spaces-perfect-definition-proper-over-base}
$S$를 스킴이라 하자. $f:X\to Y$를 $S$ 위 대수공간의 국소 유한형
사상이라 하자. $T\subset|X|$를 닫힌 부분집합이라 하자. Lemma
\ref{spaces-perfect-lemma-closed-proper-over-base}의 서로 동치인 조건들이 성립하면
{\it $T$가 $Y$ 위에서 고유하다}고 한다.
\end{definition}

\noindent
위 정의에 사용한 보조정리는 사상 $f:X\to Y$가 국소 유한형이 아니면
거짓이다. 따라서 $f$가 국소 유한형이 아닐 때에는 이 용어를 사용하지
말기를 권한다.

\begin{lemma}
\label{spaces-perfect-lemma-closed-closed-proper-over-base}
$S$를 스킴이라 하자. $f:X\to Y$를 $S$ 위 대수공간의 국소 유한형
사상이라 하자. $T'\subset T\subset|X|$를 닫힌 부분집합이라 하자.
$T$가 $Y$ 위에서 고유이면 $T'$도 그러하다.
\end{lemma}

\begin{proof}
생략한다.
\end{proof}

\begin{lemma}
\label{spaces-perfect-lemma-base-change-closed-proper-over-base}
$S$를 스킴이라 하자. $S$ 위 대수공간의 Cartesian 도식
$$
\xymatrix{
X' \ar[d]_{f'} \ar[r]_{g'} & X \ar[d]^f \\
Y' \ar[r]^g & Y
}
$$
을 생각하자. 여기서 $f$는 국소 유한형이다. $T$가 $|X|$의 닫힌
부분집합이고 $Y$ 위에서 고유이면, $|g'|^{-1}(T)$는 $|X'|$의 닫힌
부분집합이고 $Y'$ 위에서 고유이다.
\end{lemma}

\begin{proof}
Morphisms of Spaces, Lemma
\ref{spaces-morphisms-lemma-base-change-finite-type}에 의해 $f'$가 국소
유한형이므로 명제는 의미가 있다. $Z\subset X$를 $T$ 위의 유도된 축약
닫힌 부분공간 구조라 하자. 스킴 이론적 역상을
$Z'=(g')^{-1}(Z)$라 쓰자. 그러면
$Z'=X'\times_XZ=(Y'\times_YX)\times_XZ=Y'\times_YZ$는 $Y'$ 위에서
고유이다. 이는 $Z$를 $Y$ 위에서 밑변환한 것이기 때문이다(Morphisms of Spaces, Lemma
\ref{spaces-morphisms-lemma-base-change-proper}). 한편
$T'=|Z'|$이다.% P09-ERR-0397
따라서 보조정리가 성립한다.
\end{proof}

\begin{lemma}
\label{spaces-perfect-lemma-functoriality-closed-proper-over-base}
$S$를 스킴이라 하자. $B$를 $S$ 위 대수공간이라 하자. $f:X\to Y$를
$B$ 위 국소 유한형인 대수공간들의 사상이라 하자.
\begin{enumerate}
\item $Y$가 $B$ 위에서 분리이고 $T\subset|X|$가 $B$ 위에서 고유인
닫힌 부분집합이면, $|f|(T)$는 $|Y|$의 닫힌 부분집합이고 $B$ 위에서
고유이다.
\item $f$가 보편적으로 닫혀 있고 $T\subset|X|$가 $B$ 위에서 고유인
닫힌 부분집합이면, $|f|(T)$는 $Y$의 닫힌 부분집합이고 $B$ 위에서
고유이다.% P09-ERR-0398
\item $f$가 고유이고 $T\subset|Y|$가 $B$ 위에서 고유인 닫힌
부분집합이면, $|f|^{-1}(T)$는 $|X|$의 닫힌 부분집합이고 $B$ 위에서
고유이다.
\end{enumerate}
\end{lemma}

\begin{proof}
(1)을 증명하자. $Y$가 $B$ 위에서 분리이고 $T\subset|X|$가 $B$
위에서 고유인 닫힌 부분집합이라고 가정하자. $Z$를 $T$ 위의 유도된
축약 닫힌 부분공간 구조라 하고, Morphisms of Spaces, Lemma
\ref{spaces-morphisms-lemma-scheme-theoretic-image-is-proper}를
$Z\to Y$에 $B$ 위에서 적용하면 된다.

\medskip\noindent
(2)를 증명하자. $f$가 보편적으로 닫혀 있고 $T\subset|X|$가 $B$
위에서 고유인 닫힌 부분집합이라고 가정하자. $Z$를 $T$ 위의 유도된
축약 닫힌 부분공간 구조라 하고, $Z'$를 $|f|(T)$ 위의 유도된 축약 닫힌
부분공간 구조라 하자. 유도된 사상 $Z\to Z'$를 얻는다.
$Z''=f^{-1}(Z')$를 스킴 이론적 역상이라 쓰자. 그러면
$Z''\to Z'$는 $f$의 밑변환이므로 보편적으로 닫혀 있다(Morphisms of
Spaces, Lemma \ref{spaces-morphisms-lemma-base-change-proper}).
따라서 $Z\to Z'$는 닫힌 몰입 $Z\to Z''$과 $Z''\to Z'$의 합성이므로
보편적으로 닫혀 있다(Morphisms of Spaces, Lemmas
\ref{spaces-morphisms-lemma-closed-immersion-proper}와
\ref{spaces-morphisms-lemma-composition-proper}).
Morphisms of Spaces, Lemma
\ref{spaces-morphisms-lemma-image-universally-closed-separated}에 의해
$Z'\to B$가 분리임을 얻는다. $Z\to B$가 준콤팩트이고 $Z\to Z'$가
전사이므로 $Z'\to B$도 준콤팩트이다. $Z'\to B$는 $Z'\to Y$와
$Y\to B$의 합성이므로 $Z'\to B$가 국소 유한형이다(Morphisms of Spaces, Lemmas
\ref{spaces-morphisms-lemma-immersion-locally-finite-type}와
\ref{spaces-morphisms-lemma-composition-finite-type}).
마지막으로 $Z\to B$가 보편적으로 닫혀 있으므로 Morphisms of Spaces,
Lemma \ref{spaces-morphisms-lemma-image-proper-is-proper}에 의해
$Z'\to B$도 그러하다. 이로써 증명이 끝난다.

\medskip\noindent
(3)을 증명하자. $f$가 고유이고 $T\subset|Y|$가 $B$ 위에서 고유인
닫힌 부분집합이라고 가정하자. $Z$를 $T$ 위의 유도된 축약 닫힌 부분공간
구조라 하자. $Z'=f^{-1}(Z)$를 스킴 이론적 역상이라 쓰자.
그러면 $Z'\to Z$는 $f$의 밑변환이므로 고유이다(Morphisms of Spaces,
Lemma \ref{spaces-morphisms-lemma-base-change-proper}).
따라서 $Z'\to B$는 $Z'\to Z$와 $Z\to B$의 합성이므로 고유이다
(Morphisms of Spaces, Lemma
\ref{spaces-morphisms-lemma-composition-proper}). 이로써 증명이
끝난다.
\end{proof}
\begin{lemma}
\label{spaces-perfect-lemma-union-closed-proper-over-base}
$S$를 스킴이라 하자. $f : X \to Y$를 $S$ 위 대수공간의 국소 유한형
사상이라 하자. $T_i \subset |X|$, $i = 1, \ldots, n$을 닫힌
부분집합이라 하자. $T_i$, $i = 1, \ldots, n$이 $Y$ 위에서
고유이면 $T_1 \cup \ldots \cup T_n$도 그러하다.
\end{lemma}

\begin{proof}
$Z_i$를 $T_i$ 위의 유도된 축약 닫힌 부분스킴 구조라 하자.
% P09-ERR-0399
사상
$$
Z_1 \amalg \ldots \amalg Z_n \longrightarrow X
$$
은 Morphisms of Spaces, Lemmas
\ref{spaces-morphisms-lemma-closed-immersion-finite}와
\ref{spaces-morphisms-lemma-finite-union-finite}에 의해 유한이다.
유한 사상은 보편적으로 닫혀 있고
(Morphisms of Spaces, Lemma \ref{spaces-morphisms-lemma-finite-proper}),
$Z_1 \amalg \ldots \amalg Z_n$은 $S$ 위에서 고유이므로
% P09-ERR-0400
Lemma \ref{spaces-perfect-lemma-functoriality-closed-proper-over-base}의 (2)에 의해
상 $Z_1 \cup \ldots \cup Z_n$은 $S$ 위에서 고유이다.
\end{proof}

\noindent
$S$를 스킴이라 하자. $f : X \to Y$를 $S$ 위 대수공간의 국소 유한형
사상이라 하자. $\mathcal{F}$를 유한형 준연접 $\mathcal{O}_X$-가군이라
하자. 그러면 지지 $\text{Supp}(\mathcal{F})$, 즉 $\mathcal{F}$의 지지는
$|X|$의 닫힌 부분집합이다. Morphisms of Spaces, Lemma
\ref{spaces-morphisms-lemma-support-finite-type}를 보라. 따라서
``$\mathcal{F}$의 지지는 $Y$ 위에서 고유이다''라고 말하는 것이
의미가 있다.

\begin{lemma}
\label{spaces-perfect-lemma-module-support-proper-over-base}
$S$를 스킴이라 하자. $f : X \to Y$를 $S$ 위 대수공간의 국소 유한형
사상이라 하자. $\mathcal{F}$를 유한형 준연접
$\mathcal{O}_X$-가군이라 하자. 다음은 서로 동치이다.
\begin{enumerate}
\item $\mathcal{F}$의 지지가 $Y$ 위에서 고유이다.
\item $\mathcal{F}$의 스킴 이론적 지지
(Morphisms of Spaces, Definition
\ref{spaces-morphisms-definition-scheme-theoretic-support})가
$Y$ 위에서 고유이다.
\item 닫힌 부분공간 $Z \subset X$와 유한형 준연접
$\mathcal{O}_Z$-가군 $\mathcal{G}$가 존재하여 (a) $Z \to Y$가
고유이고, (b) $(Z \to X)_*\mathcal{G} = \mathcal{F}$이다.
\end{enumerate}
\end{lemma}

\begin{proof}
지지 $\text{Supp}(\mathcal{F})$, 즉 $\mathcal{F}$의 지지는 $|X|$의 닫힌
부분집합이다. Morphisms of Spaces, Lemma
\ref{spaces-morphisms-lemma-support-finite-type}를 보라. 따라서
Definition \ref{spaces-perfect-definition-proper-over-base}를 적용할 수 있다.
$\mathcal{F}$의 스킴 이론적 지지는 그 기저 닫힌 부분집합이
$\text{Supp}(\mathcal{F})$인 닫힌 부분공간이므로 Definition
\ref{spaces-perfect-definition-proper-over-base}에 의해 (1)과 (2)는 서로 동치이다.
(2)가 (3)을 함의하는 것은 명백하다. 반대로 (3)이 참이면
$\text{Supp}(\mathcal{F}) \subset |Z|$이고, 따라서 예를 들어 Lemma
\ref{spaces-perfect-lemma-closed-closed-proper-over-base}에 의해
$\text{Supp}(\mathcal{F})$는 $Y$ 위에서 고유이다.
\end{proof}

\begin{lemma}
\label{spaces-perfect-lemma-base-change-module-support-proper-over-base}
$S$를 스킴이라 하자. $S$ 위 대수공간의 Cartesian 도식
$$
\xymatrix{
X' \ar[d]_{f'} \ar[r]_{g'} & X \ar[d]^f \\
Y' \ar[r]^g & Y
}
$$
을 생각하자. 여기서 $f$는 국소 유한형이다. $\mathcal{F}$를 유한형
준연접 $\mathcal{O}_X$-가군이라 하자. $\mathcal{F}$의 지지가
$Y$ 위에서 고유이면 $(g')^*\mathcal{F}$의 지지는 $Y'$ 위에서
고유이다.
\end{lemma}

\begin{proof}
Modules on Sites, Lemma \ref{sites-modules-lemma-local-pullback}에 의해
$(g')*\mathcal{F}$가 유한형이므로 이 명제는 의미가 있다.
% P09-ERR-0401
Morphisms of Spaces, Lemma
\ref{spaces-morphisms-lemma-support-finite-type}에 의해
$\text{Supp}((g')^*\mathcal{F}) = |g'|^{-1}(\text{Supp}(\mathcal{F}))$이다.
따라서 이 보조정리는 Lemma
\ref{spaces-perfect-lemma-base-change-closed-proper-over-base}에서 따른다.
\end{proof}

\begin{lemma}
\label{spaces-perfect-lemma-cat-module-support-proper-over-base}
$S$를 스킴이라 하자. $f : X \to Y$를 $S$ 위 대수공간의 국소 유한형
사상이라 하자. $\mathcal{F}$, $\mathcal{G}$를 유한형 준연접
$\mathcal{O}_X$-가군이라 하자.
\begin{enumerate}
\item $\mathcal{F}$와 $\mathcal{G}$의 지지가 $Y$ 위에서 고유이면,
$\mathcal{F} \oplus \mathcal{G}$, $\mathcal{G}$를 $\mathcal{F}$로
확대한 임의의 확대, $\Im(u)$와 $\Coker(u)$ (임의의
$\mathcal{O}_X$-가군 사상 $u : \mathcal{F} \to \mathcal{G}$에서
나오는 것),
그리고 $\mathcal{F}$ 또는 $\mathcal{G}$의 임의의 준연접 몫의
지지도 그러하다.
\item $Y$가 국소 Noether이면 연접 $\mathcal{O}_X$-가군들 가운데
지지가 $Y$ 위에서 고유인 것들의 범주는 연접 $\mathcal{O}_X$-가군들의
아벨 범주의 Serre 부분범주이다(Homology, Definition
\ref{homology-definition-serre-subcategory}).
\end{enumerate}
\end{lemma}

\begin{proof}
(1)의 증명. $T$, $T'$을 각각 $\mathcal{F}$와 $\mathcal{G}$의
지지라 하자. 그러면 (1)에서 언급한 모든 층의 지지는
$T \cup T'$에 포함된다. 따라서 이 층들이 유한형이고 준연접임을
확인하면, 주장 자체는 Lemmas
\ref{spaces-perfect-lemma-closed-closed-proper-over-base}와
\ref{spaces-perfect-lemma-union-closed-proper-over-base}에서 명백하다.
준연접성에 대해서는 Properties of Spaces, Section
\ref{spaces-properties-section-quasi-coherent}를 보라. ``유한형''에
대해서는 Properties of Spaces, Section
\ref{spaces-properties-section-properties-modules}를 보라.

\medskip\noindent
(2)의 증명. 증명은 (1)의 증명과 같다. Morphisms of Spaces, Lemma
\ref{spaces-morphisms-lemma-locally-finite-type-locally-noetherian}과
Cohomology of Spaces, Section
\ref{spaces-cohomology-section-coherent}의 연접 가군 범주에 대한
설명에 의해 $X$가 국소 Noether이므로 각 주장이 의미가 있음에
유의하라.
\end{proof}

\begin{lemma}
\label{spaces-perfect-lemma-support-proper-over-base-pushforward}
$S$를 스킴이라 하자. $f : X \to Y$를 $S$ 위 대수공간의 사상이라
하자. $f$는 국소 유한형이고 $Y$는 국소 Noether이라고 가정하자.
$\mathcal{F}$를 연접 $\mathcal{O}_X$-가군이라 하고 그 지지가
$Y$ 위에서 고유이라고 하자. 그러면 $R^pf_*\mathcal{F}$는 연접
$\mathcal{O}_Y$-가군이다. 이는 모든 $p \geq 0$에 대해 성립한다.
\end{lemma}

\begin{proof}
Lemma \ref{spaces-perfect-lemma-module-support-proper-over-base}에 의해 닫힌 몰입
$i : Z \to X$가 존재하여 $g = f \circ i : Z \to Y$가 고유이고,
$\mathcal{F} = i_*\mathcal{G}$가 되는 어떤 연접 가군
$\mathcal{G}$가 $Z$ 위에 존재한다. Cohomology of Spaces, Lemma
\ref{spaces-cohomology-lemma-proper-pushforward-coherent}에 의해
$R^pg_*\mathcal{G}$는 $S$ 위에서 연접이다.% P09-ERR-0402
한편 $R^qi_*\mathcal{G} = 0$이다($q > 0$에 대해)
(Cohomology of Spaces, Lemma
\ref{spaces-cohomology-lemma-finite-pushforward-coherent}).
Cohomology on Sites, Lemma
\ref{sites-cohomology-lemma-relative-Leray}에 의해
$R^pf_*\mathcal{F} = R^pg_*\mathcal{G}$를 얻고 보조정리가 따른다.
\end{proof}






\section{연접 가군의 유도 범주}
\label{spaces-perfect-section-derived-coherent}

\noindent
$S$를 스킴이라 하고 $X$를 $S$ 위의 국소 Noether 대수공간이라 하자.
이 경우 범주
$\textit{Coh}(\mathcal{O}_X) \subset \textit{Mod}(\mathcal{O}_X)$,
즉 연접 $\mathcal{O}_X$-가군들의 범주는
약한 Serre 부분범주이다. Homology, Section
\ref{homology-section-serre-subcategories}와 Cohomology of Spaces,
Lemma \ref{spaces-cohomology-lemma-coherent-abelian-Noetherian}을 보라.
코호몰로지 층들이 연접인 복합체들의 부분범주를
$$
D_{\textit{Coh}}(\mathcal{O}_X) \subset D(\mathcal{O}_X)
$$
로 나타낸다. Derived Categories, Section
\ref{derived-section-triangulated-sub}를 보라. 따라서 표준 함자
\begin{equation}
\label{spaces-perfect-equation-compare-coherent}
D(\textit{Coh}(\mathcal{O}_X))
\longrightarrow
D_{\textit{Coh}}(\mathcal{O}_X)
\end{equation}
를 얻는다. Derived Categories, Equation
(\ref{derived-equation-compare})을 보라.

\begin{lemma}
\label{spaces-perfect-lemma-direct-image-coherent}
$S$를 스킴이라 하자. $f : X \to Y$를 $S$ 위 대수공간의 사상이라
하자. $f$가 국소 유한형이고 $Y$가 Noether이라고 가정하자.
$E$를 $D^b_{\textit{Coh}}(\mathcal{O}_X)$의 대상이라 하고,
$H^i(E)$의 지지가 $Y$ 위에서 고유이라고 하자(모든 $i$에 대해).
그러면 $Rf_*E$는 $D^b_{\textit{Coh}}(\mathcal{O}_Y)$의 대상이다.
\end{lemma}

\begin{proof}
스펙트럼 열
$$
R^pf_*H^q(E) \Rightarrow R^{p + q}f_*E
$$
을 생각하자. Derived Categories, Lemma
\ref{derived-lemma-two-ss-complex-functor}를 보라. 가정과 Lemma
\ref{spaces-perfect-lemma-support-proper-over-base-pushforward}에 의해 층
$R^pf_*H^q(E)$는 연접이다. 따라서 $R^{p + q}f_*E$는 연접이다.
즉, $E \in D_{\textit{Coh}}(\mathcal{O}_Y)$이다.% P09-ERR-0403
아래로 유계인 것은 명백하다. 위로 유계인 것은 Cohomology of Spaces,
Lemma \ref{spaces-cohomology-lemma-vanishing-higher-direct-images} 또는
Lemma \ref{spaces-perfect-lemma-quasi-coherence-direct-image}에서 따른다.
\end{proof}

\begin{lemma}
\label{spaces-perfect-lemma-direct-image-coherent-bdd-below}
$S$를 스킴이라 하자. $f : X \to Y$를 $S$ 위 대수공간의 사상이라
하자. $f$가 국소 유한형이고 $Y$가 Noether이라고 가정하자.
$E$를 $D^+_{\textit{Coh}}(\mathcal{O}_X)$의 대상이라 하고,
$H^i(E)$의 지지가 $S$ 위에서 고유이라고 하자(모든 $i$에 대해).
% P09-ERR-0404
그러면 $Rf_*E$는 $D^+_{\textit{Coh}}(\mathcal{O}_Y)$의 대상이다.
\end{lemma}

\begin{proof}
증명은 Lemma \ref{spaces-perfect-lemma-direct-image-coherent}의 증명과 같다.
또는 완전 함자 $Rf_*$가 구별 삼각형
$\tau_{\leq a}E \to E \to \tau_{\geq a + 1}E \to \tau_{\leq a}E[1]$
에 무엇을 하는지 생각하여 Lemma
\ref{spaces-perfect-lemma-direct-image-coherent}에서 이를 도출할 수도 있다.
\end{proof}

\begin{lemma}
\label{spaces-perfect-lemma-coherent-internal-hom}
$S$를 스킴이라 하자. $X$를 $S$ 위의 국소 Noether 대수공간이라 하자.
$L$이 $D^+_{\textit{Coh}}(\mathcal{O}_X)$에 속하고 $K$가
$D^-_{\textit{Coh}}(\mathcal{O}_X)$에 속하면,
$R\SheafHom(K, L)$은 $D^+_{\textit{Coh}}(\mathcal{O}_X)$에 속한다.
\end{lemma}

\begin{proof}
$D(\mathcal{O}_X)$의 대상이 $D_{\textit{Coh}}(\mathcal{O}_X)$에
속하는지는 $X$ 위에서 \`etale 국소적으로 확인할 수 있다. Cohomology
of Spaces, Lemma \ref{spaces-cohomology-lemma-coherent-Noetherian}을 보라.
따라서 이 보조정리는 스킴의 경우에서 따른다. Derived Categories of
Schemes, Lemma \ref{perfect-lemma-coherent-internal-hom}을 보라.
\end{proof}

\begin{lemma}
\label{spaces-perfect-lemma-ext-finite}
$A$를 Noether 환이라 하자. $X$를 $A$ 위의 고유 대수공간이라 하자.
$L$이 $D^+_{\textit{Coh}}(\mathcal{O}_X)$에 속하고 $K$가
$D^-_{\textit{Coh}}(\mathcal{O}_X)$에 속하면 $A$-가군
$\Ext_{\mathcal{O}_X}^n(K, L)$은 유한이다.
\end{lemma}

\begin{proof}
다음을 상기하자.
$$
\Ext_{\mathcal{O}_X}^n(K, L) =
H^n(X, R\SheafHom_{\mathcal{O}_X}(K, L)) =
H^n(\Spec(A), Rf_*R\SheafHom_{\mathcal{O}_X}(K, L))
$$
Cohomology on Sites, Lemma
\ref{sites-cohomology-lemma-section-RHom-over-U}와 Cohomology on Sites,
Section \ref{sites-cohomology-section-leray}를 보라. 따라서 결과는
Lemmas \ref{spaces-perfect-lemma-coherent-internal-hom}과
\ref{spaces-perfect-lemma-direct-image-coherent-bdd-below}에서 따른다.
\end{proof}




\section{귀납 원리}
\label{spaces-perfect-section-induction}

\noindent
이 절에서는 스킴에 대한 Cohomology of Schemes, Lemma
\ref{coherent-lemma-induction-principle}와 유사한 대수공간의 귀납
원리를 논한다. 이를 정식화하기 위해 {\it 기본 구별 정사각형}의
개념을 도입한다. 이 용어는 \cite{MV}에서 빌려왔다. 여기서 정식화한
원리는 Raynaud와 Gruson의 논문 \cite{GruRay}에 암묵적으로 들어 있다.
대수 스택에 대한 관련 원리는 David Rydh의
\cite[Theorem D]{rydh_etale_devissage}이다.

\begin{definition}
\label{spaces-perfect-definition-elementary-distinguished-square}
$S$를 스킴이라 하자. 가환 도식
$$
\xymatrix{
U \times_W V \ar[r] \ar[d] & V \ar[d]^f \\
U \ar[r]^j & W
}
$$
을 생각하자. 이 도식이 $S$ 위 대수공간의 도식이고 다음을 만족할 때
이를 {\it 기본 구별 정사각형}이라 한다.
\begin{enumerate}
\item $U$는 $W$의 열린 부분공간이고 $j$는 포함 사상이다.
\item $f$는 \`etale이다.
\item $T = W \setminus U$라 두면(유도된 축약 부분공간 구조를
부여하여) 사상 $f^{-1}(T) \to T$는 동형이다.
\end{enumerate}
이를 ``$(U \subset W, f : V \to W)$를 기본 구별 정사각형이라
하자''라고 나타내겠다.
\end{definition}

\noindent
$(U \subset W, f : V \to W)$가 기본 구별 정사각형이면
$W = U \cup f(V)$임에 유의하라. 따라서
$\{U \to W, V \to W\}$는 $W$의 \`etale 피복이다. 이 \`etale
피복들은 좋은 성질을 가지며, 어떤 의미에서는 그러한 피복이
``충분히 많이'' 있음이 밝혀진다.

\begin{lemma}
\label{spaces-perfect-lemma-make-more-elementary-distinguished-squares}
$S$를 스킴이라 하자. $(U \subset W, f : V \to W)$를 $S$ 위
대수공간의 기본 구별 정사각형이라 하자.
\begin{enumerate}
\item $V' \subset V$와 $U \subset U' \subset W$가 열린
부분공간이고 $W' = U' \cup f(V')$이면
$(U' \subset W', f|_{V'} : V' \to W')$는 기본 구별 정사각형이다.
\item $p : W' \to W$가 대수공간의 사상이면
$(p^{-1}(U) \subset W', V \times_W W' \to W')$는 기본 구별
정사각형이다.
\item $S' \to S$가 스킴의 사상이면
$(S' \times_S U \subset S' \times_S W, S' \times_S V \to S' \times_S W)$는
기본 구별 정사각형이다.
\end{enumerate}
\end{lemma}

\begin{proof}
생략한다.
\end{proof}

\begin{lemma}
\label{spaces-perfect-lemma-induction-principle}
$S$를 스킴이라 하자. $X$를 $S$ 위의 준콤팩트 준분리 대수공간이라
하자. $P$를 $X_{spaces, \etale}$의 준콤팩트 준분리 대상들의
성질이라 하자. 다음을 가정하자.
\begin{enumerate}
\item $P$는 $X_{spaces, \etale}$의 모든 아핀 대상에 대해 성립한다.
\item 다음 조건을 만족하는 모든 기본 구별 정사각형
$(U \subset W, f : V \to W)$에 대해,
\begin{enumerate}
\item $W$는 $X_{spaces, \etale}$의 준콤팩트 준분리 대상이다.
\item $U$는 준콤팩트이다.
\item $V$는 아핀이다.
\item $P$는 $U$, $V$, $U \times_W V$에 대해 성립한다.
\end{enumerate}
그러면 $P$는 $W$에 대해 성립한다.
\end{enumerate}
그러면 $P$는 $X_{spaces, \etale}$의 모든 준콤팩트 준분리 대상에
대해 성립하며, 특히 $X$에 대해 성립한다.
\end{lemma}

\begin{proof}
먼저 $P$가 $X_{spaces, \etale}$의 표현 가능한 모든 준콤팩트 준분리
대상에 대해 성립한다고 주장한다. 실제로 $U \to X$가 \`etale이고
$U$가 준콤팩트 준분리 스킴이라고 하자. 가정 (1)에 의해 성질 $P$는
$U$의 모든 아핀 열린집합에 대해 성립한다. 더욱이
$W, V \subset U$가 준콤팩트 열린집합이고 $V$가 아핀이며 $P$가
$W$, $V$, $W \cap V$에 대해 성립하면, (2)에 의해 $P$는
$W \cup V$에 대해 성립한다(쌍
$(W \subset W \cup V, V \to W \cup V)$가 기본 구별 정사각형이기
때문이다). 따라서 스킴에 대한 귀납 원리, 즉 Cohomology of Schemes,
Lemma \ref{coherent-lemma-induction-principle}에 의해 $P$는 $U$에
대해 성립한다.

\medskip\noindent
증명을 끝내려면 $P$가 $X$에 대해 성립함을 증명하면 충분하다
($X$를 결과를 증명하고 싶은 $X_{spaces, \etale}$의 임의의 준콤팩트
준분리 대상으로 바꾸면 되기 때문이다). Decent Spaces, Lemma
\ref{decent-spaces-lemma-filter-quasi-compact-quasi-separated}의 여과
$$
\emptyset = U_{n + 1} \subset
U_n \subset U_{n - 1} \subset \ldots \subset U_1 = X
$$
와 사상 $f_p : V_p \to U_p$를 사용하겠다. $P$가 $U_p$에 대해
성립함을 $p$에 대한 하향 귀납법으로 증명한다. 공 대수공간은 아핀이므로
(1)에 의해 $P$는 $U_{n + 1}$에 대해 성립함에 유의하라.
$P$가 $U_{p + 1}$에 대해 성립한다고 가정하자.
$(U_{p + 1} \subset U_p, f_p : V_p \to U_p)$는 기본 구별 정사각형이지만,
$V_p$가 아핀이 아닐 수 있으므로 (2)를 적용하지 못할 수 있다.
그러나 $V_p$는 준콤팩트 스킴이므로 유한 아핀 열린 피복
$V_p = V_{p, 1} \cup \ldots \cup V_{p, m}$을 택할 수 있다.
$W_{p, 0} = U_{p + 1}$이라 두고
$$
W_{p, i} = U_{p + 1} \cup f_p(V_{p, 1} \cup \ldots \cup V_{p, i})
$$
라 두자($i = 1, \ldots, m$에 대해). 이들은 $X$의 준콤팩트 열린
부분공간이다. 그러면
$$
U_{p + 1} = W_{p, 0} \subset
W_{p, 1} \subset \ldots \subset
W_{p, m} = U_p
$$
이고, 쌍들
$$
(W_{p, 0} \subset W_{p, 1}, f_p|_{V_{p, 1}}),
(W_{p, 1} \subset W_{p, 2}, f_p|_{V_{p, 2}}),\ldots,
(W_{p, m - 1} \subset W_{p, m}, f_p|_{V_{p, m}})
$$
은 Lemma \ref{spaces-perfect-lemma-make-more-elementary-distinguished-squares}에 의해
기본 구별 정사각형이다. $P$는 각 $V_{p, 1}$에 대해 성립하고(아핀
스킴이므로), $W_{p, i} \times_{W_{p, i + 1}} V_{p, i + 1}$에 대해서도
성립한다. 후자는 $V_{p, i + 1}$의 준콤팩트 열린집합이므로 이 증명의
첫 문단에 의해 $P$가 성립하기 때문이다. 따라서 이들 각각에 (2)를
적용하고 귀납적으로 $P$가
$W_{p, 1}, \ldots, W_{p, m} = U_p$에 대해 성립함을 얻는다.
\end{proof}

\begin{lemma}
\label{spaces-perfect-lemma-induction-principle-separated}
$S$를 스킴이라 하자. $X$를 $S$ 위의 준콤팩트 준분리 대수공간이라
하자. $\mathcal{B} \subset \Ob(X_{spaces, \etale})$라 하고 $P$를
$\mathcal{B}$의 원소들의 성질이라 하자. 다음을 가정하자.
\begin{enumerate}
\item 모든 $W \in \mathcal{B}$는 준콤팩트이고 준분리이다.
\item $W \in \mathcal{B}$이고 $U \subset W$가 준콤팩트
열린집합이면 $U \in \mathcal{B}$이다.
\item $V \in \Ob(X_{spaces, \etale})$가 아핀이면 (a)
$V \in \mathcal{B}$이고 (b) $P$는 $V$에 대해 성립한다.
\item 다음 조건을 만족하는 모든 기본 구별 정사각형
$(U \subset W, f : V \to W)$에 대해,
\begin{enumerate}
\item $W \in \mathcal{B}$이다.
\item $U$는 준콤팩트이다.
\item $V$는 아핀이다.
\item $P$는 $U$, $V$, $U \times_W V$에 대해 성립한다.
\end{enumerate}
그러면 $P$는 $W$에 대해 성립한다.
\end{enumerate}
그러면 $P$는 모든 $W \in \mathcal{B}$에 대해 성립한다.
\end{lemma}

\begin{proof}
이는 Lemma \ref{spaces-perfect-lemma-induction-principle}의 증명과 정확히 같은
방식으로 증명된다. ((4)(d)가 의미를 가지는 것은
$U \times_W V$가 $V$의 준콤팩트 열린집합이고, 따라서 조건 (2)와
(3)에 의해 $\mathcal{B}$의 원소이기 때문이다.)
\end{proof}

\begin{remark}
\label{spaces-perfect-remark-how-to}
$\mathcal{B}$라는 모음을 Lemma
\ref{spaces-perfect-lemma-induction-principle-separated}에서 어떻게 택할 것인가?
몇 가지 예는 다음과 같다.
\begin{enumerate}
\item $X$가 준콤팩트이고 분리이면, $\mathcal{B}$를
$X_{spaces, \etale}$의 준콤팩트 분리 대상들의 집합으로 택할 수 있다.
그러면 $X \in \mathcal{B}$이고 $\mathcal{B}$는 (1), (2), (3)(a)를
만족한다. 이 $\mathcal{B}$의 선택으로 Lemma
\ref{spaces-perfect-lemma-induction-principle-separated}는 Lemma
\ref{spaces-perfect-lemma-induction-principle}를 재현한다.
\item $X$가 $\mathbf{Z}$ 위에서 아핀 대각사상을 갖는 준콤팩트
공간이면(Properties of Spaces, Definition
\ref{spaces-properties-definition-separated}에서와 같이),
$\mathcal{B}$를 $X_{spaces, \etale}$의 대상들 가운데 준콤팩트이고
$\mathbf{Z}$ 위에서 아핀 대각사상을 갖는 것들의 집합으로 택할 수
있다. 다시 $X \in \mathcal{B}$이고 $\mathcal{B}$는 (1), (2),
(3)(a)를 만족한다.
\item $X$가 준콤팩트이고 준분리이면, 가장 작은 부분집합
$\mathcal{B}$, 즉 $X$를 포함하고 (1), (2), (3)(a)를 만족하는 것은
다음 규칙으로
주어진다. $W \in \mathcal{B}$일 필요충분조건은 $W$가 $X$의
준콤팩트 열린 부분공간이거나, $W$가 $X_{spaces, \etale}$의 어떤
아핀 대상의 준콤팩트 열린집합인 것이다.
\end{enumerate}
\end{remark}

\noindent
다음은 열린집합에서 더 큰 열린집합으로 참을 확장하는 변형이다.

\begin{lemma}
\label{spaces-perfect-lemma-induction-principle-enlarge}
$S$를 스킴이라 하자. $X$를 $S$ 위의 준콤팩트 준분리 대수공간이라
하자. $W \subset X$를 준콤팩트 열린 부분공간이라 하자. $P$를
$X$의 준콤팩트 열린 부분공간들의 성질이라 하자. 다음을 가정하자.
\begin{enumerate}
\item $P$는 $W$에 대해 성립한다.
\item 모든 기본 구별 정사각형
$(W_1 \subset W_2, f : V \to W_2)$에 대해
% P09-ERR-0406
다음이 성립한다고 하자.
\begin{enumerate}
\item $W_1$, $W_2$는 $X$의 준콤팩트 열린 부분공간이다.
\item $W \subset W_1$이다.
\item $V$는 아핀이다.
\item $P$는 $W_1$에 대해 성립한다.
\end{enumerate}
그러면 $P$는 $W_2$에 대해 성립한다.
\end{enumerate}
그러면 $P$는 $X$에 대해 성립한다.
\end{lemma}

\begin{proof}
이를 Lemma \ref{spaces-perfect-lemma-induction-principle-separated}에서 도출할 수
있지만, 대신 Lemma \ref{spaces-perfect-lemma-induction-principle}의 증명을 명시적으로
다시 수행하여 직접 논증하겠다. Decent Spaces, Lemma
\ref{decent-spaces-lemma-filter-quasi-compact-quasi-separated}의 여과
$$
\emptyset = U_{n + 1} \subset
U_n \subset U_{n - 1} \subset \ldots \subset U_1 = X
$$
와 사상 $f_p : V_p \to U_p$를 사용하겠다. $P$가
$W_p = W \cup U_p$에 대해 성립함을 $p$에 대한 하향 귀납법으로
증명하겠다. $W_1 = X$이므로 이것으로 증명이 끝난다. (1)에 의해
$P$는 $W_{n + 1} = W \cap U_{n + 1} = W$에 대해 성립한다.
% P09-ERR-0407
$P$가 $W_{p + 1}$에 대해 성립한다고 가정하자.
$W_p \setminus W_{p + 1}$은(유도된 축약 부분공간 구조를 부여하여)
$U_p \setminus U_{p + 1}$의 닫힌 부분공간이다.
$(U_{p + 1} \subset U_p, f_p : V_p \to U_p)$가 기본 구별
정사각형이므로 $(W_{p + 1} \subset W_p, f_p : V_p \to W_p)$도
그러하다. 그러나 $V_p$가 아핀이 아닐 수 있으므로 (2)를 적용하지
못할 수 있다. 하지만 $V_p$는 준콤팩트 스킴이므로 유한 아핀 열린 피복
$V_p = V_{p, 1} \cup \ldots \cup V_{p, m}$을 택할 수 있다.
$W_{p, 0} = W_{p + 1}$이라 두고
$$
W_{p, i} = W_{p + 1} \cup f_p(V_{p, 1} \cup \ldots \cup V_{p, i})
$$
라 두자($i = 1, \ldots, m$에 대해). 이들은 $X$의 준콤팩트 열린
부분공간이며 $W$를 포함한다. 그러면
$$
W_{p + 1} = W_{p, 0} \subset
W_{p, 1} \subset \ldots \subset
W_{p, m} = W_p
$$
이고, 쌍들
$$
(W_{p, 0} \subset W_{p, 1}, f_p|_{V_{p, 1}}),
(W_{p, 1} \subset W_{p, 2}, f_p|_{V_{p, 2}}),\ldots,
(W_{p, m - 1} \subset W_{p, m}, f_p|_{V_{p, m}})
$$
은 Lemma \ref{spaces-perfect-lemma-make-more-elementary-distinguished-squares}에 의해
기본 구별 정사각형이다. 이제 이들 각각에 (2)를 적용하고 귀납적으로
$P$가 $W_{p, 1}, \ldots, W_{p, m} = W_p$에 대해 성립함을 얻는다.
\end{proof}



\section{Mayer--Vietoris}
\label{spaces-perfect-section-mayer-vietoris}

\noindent
이 절에서는 기본 구별 삼각형이 여러 Mayer--Vietoris 열을 준다는
것을 증명한다.% P09-ERR-0408

\medskip\noindent
$S$를 스킴이라 하자. $U \to X$를 $S$ 위 대수공간의 \`etale
사상이라 하자. Properties of Spaces, Section
\ref{spaces-properties-section-localize}에서 구조층과 양립하는 등식
$U_{spaces, \etale} = X_{spaces, \etale}/U$을 증명하였다. 따라서 이
상황에서는 흔히 사상 $j_U : U \to X$를 국소화 사상으로 생각한다
(Modules on Sites, Definition
\ref{sites-modules-definition-localize-ringed-site}를 보라). 특히 당김
$j_U^*$을 $U$로의 제한으로 생각하고 흔히 ${}|_U$로 나타낸다.
이는 Properties of Spaces, Equation
(\ref{spaces-properties-equation-restrict-modules})과 양립한다. 특히
다음을 얻는다.
\begin{equation}
\label{spaces-perfect-equation-stalk-restriction}
(\mathcal{F}|_U)_{\overline{u}} = \mathcal{F}_{\overline{x}}
\end{equation}
여기서 $\overline{u}$는 $U$의 기하점이고 $\overline{x}$는
$\overline{u}$의 $X$에서의 상이다. 더욱이 제한은 완전한 왼쪽
수반함자 $j_{U!}$를 갖는다. Modules on Sites, Lemmas
\ref{sites-modules-lemma-extension-by-zero}와
\ref{sites-modules-lemma-extension-by-zero-exact}를 보라. 마지막으로
$\mathcal{G}$가 $\mathcal{O}_X$-가군이면
\begin{equation}
\label{spaces-perfect-equation-stalk-j-shriek}
(j_{U!}\mathcal{G})_{\overline{x}} =
\bigoplus\nolimits_{\overline{u}} \mathcal{G}_{\overline{u}}
\end{equation}
이다. 이는 임의의 기하점 $\overline{x} : \Spec(k) \to X$에 대해
성립하며, 직합은 사상 $\overline{u} : \Spec(k) \to U$들 가운데
$j_U \circ \overline{u} = \overline{x}$를 만족하는 것들에 걸쳐 취한다. Modules on
Sites, Lemma \ref{sites-modules-lemma-stalk-j-shriek}와 Properties of
Spaces, Lemma \ref{spaces-properties-lemma-points-small-etale-site}를
보라.

\begin{lemma}
\label{spaces-perfect-lemma-exact-sequence-lower-shriek}
$S$를 스킴이라 하자. $(U \subset X, V \to X)$를 $S$ 위
대수공간의 기본 구별 정사각형이라 하자.
\begin{enumerate}
\item $\mathcal{O}_X$-가군의 층 $\mathcal{F}$에 대해 짧은 완전열
$$
0 \to j_{U \times_X V!}\mathcal{F}|_{U \times_X V} \to
j_{U!}\mathcal{F}|_U \oplus j_{V!}\mathcal{F}|_V \to \mathcal{F} \to 0
$$
이 있다.
\item 대상 $E$가 $D(\mathcal{O}_X)$에 속할 때 구별 삼각형
$$
j_{U \times_X V!}E|_{U \times_X V} \to
j_{U!}E|_U \oplus j_{V!}E|_V \to E \to 
j_{U \times_X V!}E|_{U \times_X V}[1]
$$
이 $D(\mathcal{O}_X)$ 안에 존재한다.
\end{enumerate}
\end{lemma}

\begin{proof}
(1)의 열이 완전임을 보이기 위해 Properties of Spaces, Theorem
\ref{spaces-properties-theorem-exactness-stalks}에 의해 기하점에서의
줄기를 확인하면 된다. $\overline{x}$를 $X$의 기하점이라 하자.
Equations (\ref{spaces-perfect-equation-stalk-restriction})과
(\ref{spaces-perfect-equation-stalk-j-shriek})에 의해 $\overline{x}$에서 줄기를
취하면 다음 열을 얻는다.
$$
0 \to
\bigoplus\nolimits_{(\overline{u}, \overline{v})} \mathcal{F}_{\overline{x}}
\to
\bigoplus\nolimits_{\overline{u}} \mathcal{F}_{\overline{x}}
\oplus
\bigoplus\nolimits_{\overline{v}} \mathcal{F}_{\overline{x}}
\to
\mathcal{F}_{\overline{x}} \to 0
$$
이 열은 완전이다. 실제로 모든 $\overline{x}$에 대해 정확히 하나의
$\overline{u}$가 $\overline{x}$로 가거나, $\overline{u}$는 없고
정확히 하나의 $\overline{v}$가 $\overline{x}$로 간다.

\medskip\noindent
(2)의 증명. Cohomology on Sites, Section
\ref{sites-cohomology-section-properties-K-injective}에서 유도 범주의
제한 함자와 영 연장 함자는 임의의 복합체에 그 함자를 단순히 적용하여
계산됨을 보았다. $\mathcal{E}^\bullet$을 $\mathcal{O}_X$-가군의
복합체라 하되 $E$를 나타내는 것으로 택하자. 보조정리의 구별 삼각형은
Derived Categories, Section
\ref{derived-section-canonical-delta-functor}, 특히 Lemma
\ref{derived-lemma-derived-canonical-delta-functor}에 의해
$\mathcal{O}_X$-가군 복합체의 짧은 완전열
$$
0 \to j_{U \times_X V!}\mathcal{E}^\bullet|_{U \times_X V} \to
j_{U!}\mathcal{E}^\bullet|_U \oplus j_{V!}\mathcal{E}^\bullet|_V
\to \mathcal{E}^\bullet \to 0
$$
에 결부된 구별 삼각형이다. 이 열은 (1)에 의해 짧은 완전열이다.
\end{proof}

\begin{lemma}
\label{spaces-perfect-lemma-exact-sequence-j-star}
$S$를 스킴이라 하자. $(U \subset X, V \to X)$를 $S$ 위
대수공간의 기본 구별 정사각형이라 하자.
\begin{enumerate}
\item 모든 $\mathcal{O}_X$-가군의 층 $\mathcal{F}$에 대해 짧은 완전열
$$
0 \to \mathcal{F} \to
j_{U, *}\mathcal{F}|_U \oplus j_{V, *}\mathcal{F}|_V \to
j_{U \times_X V, *}\mathcal{F}|_{U \times_X V} \to 0
$$
이 있다.
\item 임의의 대상 $E$가 $D(\mathcal{O}_X)$에 속할 때 구별 삼각형
$$
E \to 
Rj_{U, *}E|_U \oplus Rj_{V, *}E|_V \to
Rj_{U \times_X V, *}E|_{U \times_X V} \to
E[1]
$$
이 $D(\mathcal{O}_X)$ 안에 존재한다.
\end{enumerate}
\end{lemma}

\begin{proof}
$W$를 $X_\etale$의 대상이라 하자. 열
$$
0 \to
\mathcal{F}(W) \to
\mathcal{F}(W \times_X U) \oplus \mathcal{F}(W \times_X V) \to
\mathcal{F}(W \times_X U \times_X V)
$$
이 완전이고 마지막 군의 원소가 $W$ 위에서 국소적으로 가운데 군의
원소로 올려진다고 주장한다. Lemma
\ref{spaces-perfect-lemma-make-more-elementary-distinguished-squares}에 의해 쌍
$(W \times_X U \subset W, V \times_X W \to W)$는 기본 구별
정사각형이다. 따라서 $W = X$라고 가정할 수 있고 다음 열에 대해
같은 것을 증명하면 충분하다.
$$
0 \to
\mathcal{F}(X) \to
\mathcal{F}(U) \oplus \mathcal{F}(V) \to
\mathcal{F}(U \times_X V)
$$
Lemma \ref{spaces-perfect-lemma-exact-sequence-lower-shriek}에서
$$
0 \to j_{U \times_X V!}\mathcal{O}_{U \times_X V}
\to j_{U!}\mathcal{O}_U \oplus
j_{V!}\mathcal{O}_V \to
\mathcal{O}_X \to 0
$$
이 $\mathcal{O}_X$-가군의 완전열임을 보았다.% P09-ERR-0409
오른쪽 완전 함자 $\Hom_{\mathcal{O}_X}(- , \mathcal{F})$를 적용하면
위의 열을 얻는다. 이는 또한 $s \in \mathcal{F}(U \times_X V)$를
$\mathcal{F}(U) \oplus \mathcal{F}(V)$의 원소로 올리는 데 대한
장애물이
$\Ext^1_{\mathcal{O}_X}(\mathcal{O}_X, \mathcal{F}) =
H^1(X, \mathcal{F})$에 놓인다는 뜻이다. 코호몰로지의 국소성
(Cohomology on Sites, Lemma
\ref{sites-cohomology-lemma-kill-cohomology-class-on-covering})에 의해
이 장애물은 $X$ 위에서 \`etale 국소적으로 사라지며, (1)의 증명이
완료된다.

\medskip\noindent
(2)의 증명. K-단사 복합체 $\mathcal{I}^\bullet$을 택하여 $E$를
나타내게 하되, 그 항 $\mathcal{I}^n$은
$\textit{Mod}(\mathcal{O}_X)$의 단사 대상이 되게 하자.
Injectives, Theorem
\ref{injectives-theorem-K-injective-embedding-grothendieck}를 보라.
그러면 $\mathcal{I}^\bullet|U$는 K-단사 복합체이다(Cohomology on
Sites, Lemma
\ref{sites-cohomology-lemma-restrict-K-injective-to-open}). 따라서
$Rj_{U, *}E|_U$는 $j_{U, *}\mathcal{I}^\bullet|_U$로 나타내어진다.
$V$와 $U \times_X V$에 대해서도 마찬가지이다. 따라서 보조정리의
구별 삼각형은 Derived Categories, Section
\ref{derived-section-canonical-delta-functor}, 특히 Lemma
\ref{derived-lemma-derived-canonical-delta-functor}에 의해 복합체의
짧은 완전열
$$
0 \to
\mathcal{I}^\bullet \to
j_{U, *}\mathcal{I}^\bullet|_U \oplus j_{V, *}\mathcal{I}^\bullet|_V \to
j_{U \times_X V, *}\mathcal{I}^\bullet|_{U \times_X V} \to
0.
$$
에 결부된 구별 삼각형이다. 이 열은 (1)에 의해 완전이다.
\end{proof}

\begin{lemma}
\label{spaces-perfect-lemma-unbounded-relative-mayer-vietoris}
$S$를 스킴이라 하자. $f : X \to Y$를 $S$ 위 대수공간의 사상이라
하자. $(U \subset X, V \to X)$를 기본 구별 정사각형이라 하자.
제한 사상들을 $a = f|_U : U \to Y$, $b = f|_V : V \to Y$, 그리고
$c = f|_{U \times_X V} : U \times_X V \to Y$로 나타내자.
모든 대상 $E$가 $D(\mathcal{O}_X)$에 속할 때 구별 삼각형
$$
Rf_*E \to
Ra_*(E|_U) \oplus Rb_*(E|_V) \to
Rc_*(E|_{U \times_X V}) \to
Rf_*E[1]
$$
이 $D(\mathcal{O}_Y)$ 안에 존재한다. 이 삼각형은 $E$에 대해
함자적이다.
\end{lemma}

\begin{proof}
K-단사 복합체 $\mathcal{I}^\bullet$을 택하여 $E$를 나타내게 하자.
$\mathcal{I}^n$이 $\textit{Mod}(\mathcal{O}_X)$의 단사 대상이라고
가정할 수 있으며, 이는 모든 $n$에 대해 그러하다.
Injectives, Theorem
\ref{injectives-theorem-K-injective-embedding-grothendieck}를 보라.
그러면 $Rf_*E$는 $f_*\mathcal{I}^\bullet$으로 계산된다.
$U$, $V$, $U \cap V$에 대해서도 Cohomology on Sites, Lemma
\ref{sites-cohomology-lemma-restrict-K-injective-to-open}에 의해
마찬가지이다.% P09-ERR-0410
따라서 보조정리의 구별 삼각형은 Derived Categories, Section
\ref{derived-section-canonical-delta-functor}, 특히 Lemma
\ref{derived-lemma-derived-canonical-delta-functor}에 의해 복합체의
짧은 완전열
$$
0 \to
f_*\mathcal{I}^\bullet \to
a_*\mathcal{I}^\bullet|_U \oplus b_*\mathcal{I}^\bullet|_V \to
c_*\mathcal{I}^\bullet|_{U \times_X V} \to
0.
$$
에 결부된 구별 삼각형이다. 이것이 복합체의 짧은 완전열임을 보이기
위해 다음과 같이 논증하자. $\mathcal{I}$를
$\textit{Mod}(\mathcal{O}_X)$의 단사 대상으로 택한다. $f_*$를
짧은 완전열
$$
0 \to \mathcal{I} \to
j_{U, *}\mathcal{I}|_U \oplus j_{V, *}\mathcal{I}|_V \to
j_{U \times_X V, *}\mathcal{I}|_{U \times_X V} \to 0
$$
에 적용하자. 이 열은 Lemma \ref{spaces-perfect-lemma-exact-sequence-j-star}의
열이다. 이제 $R^1f_*\mathcal{I} = 0$임을 사용하면 짧은
완전열
$$
0 \to f_*\mathcal{I} \to
f_*j_{U, *}\mathcal{I}|_U \oplus f_*j_{V, *}\mathcal{I}|_V \to
f_*j_{U \times_X V, *}\mathcal{I}|_{U \times_X V} \to 0
$$
을 얻는다. 이제 $a_* = f_*j_{U, *}$이고 $b_*$와 $c_*$에 대해서도
마찬가지임을 관찰하면 증명이 끝난다.
\end{proof}

\begin{lemma}
\label{spaces-perfect-lemma-mayer-vietoris-hom}
$S$를 스킴이라 하자. $(U \subset X, V \to X)$를 $S$ 위
대수공간의 기본 구별 정사각형이라 하자. $E$, $F$를
$D(\mathcal{O}_X)$의 대상이라 하면 Mayer--Vietoris 열
$$
\xymatrix{
& \ldots \ar[r] &
\Ext^{-1}(E_{U \times_X V}, F_{U \times_X V}) \ar[lld] \\
\Hom(E, F) \ar[r] &
\Hom(E_U, F_U) \oplus
\Hom(E_V, F_V) \ar[r] &
\Hom(E_{U \times_X V}, F_{U \times_X V})
}
$$
을 얻는다. 아래첨자는 해당 열린집합으로의 제한을 뜻하고,
$\Hom$들은 해당 유도 범주에서 취한다.
\end{lemma}

\begin{proof}
Lemma \ref{spaces-perfect-lemma-exact-sequence-lower-shriek}의 구별 삼각형을 사용하여
$\Hom$들의 긴 완전열을 얻고(Derived Categories, Lemma
\ref{derived-lemma-representable-homological}에서), Cohomology on Sites,
Lemma \ref{sites-cohomology-lemma-adjoint-lower-shriek-restrict}에 의해
$\Hom(j_{U!}E|_U, F) = \Hom(E|_U, F|_U)$임을 사용하라.
\end{proof}

\begin{lemma}
\label{spaces-perfect-lemma-unbounded-mayer-vietoris}
$S$를 스킴이라 하자. $(U \subset X, V \to X)$를 $S$ 위
대수공간의 기본 구별 정사각형이라 하자. 대상 $E$가
$D(\mathcal{O}_X)$에 속할 때 구별 삼각형
$$
R\Gamma(X, E) \to R\Gamma(U, E) \oplus R\Gamma(V, E) \to
R\Gamma(U \times_X V, E) \to R\Gamma(X, E)[1]
$$
과, 특히 코호몰로지의 긴 완전열
$$
\ldots \to
H^n(X, E) \to
H^n(U, E) \oplus H^n(V, E) \to
H^n(U \times_X V, E) \to
H^{n + 1}(X, E) \to \ldots
$$
을 얻는다. 구별 삼각형과 긴 완전열의 구성은 $E$에 대해 함자적이다.
\end{lemma}

\begin{proof}
K-단사 복합체 $\mathcal{I}^\bullet$을 택하여 $E$를 나타내게 하되,
그 항 $\mathcal{I}^n$은 $\textit{Mod}(\mathcal{O}_X)$의 단사
대상이 되게 하자. Injectives, Theorem
\ref{injectives-theorem-K-injective-embedding-grothendieck}를 보라.
Lemma \ref{spaces-perfect-lemma-exact-sequence-j-star}의 증명에서 복합체의 짧은 완전열
$$
0 \to \mathcal{I}^\bullet \to
j_{U, *}\mathcal{I}^\bullet|_U \oplus j_{V, *}\mathcal{I}^\bullet|_V \to
j_{U \times_X V, *}\mathcal{I}^\bullet|_{U \times_X V} \to 0
$$
을 얻었다. $H^1(X, \mathcal{I}^n) = 0$이므로 대역 단면을 취하면
복합체의 완전열
$$
0 \to \Gamma(X, \mathcal{I}^\bullet) \to
\Gamma(U, \mathcal{I}^\bullet) \oplus
\Gamma(V, \mathcal{I}^\bullet) \to
\Gamma(U \times_X V, \mathcal{I}^\bullet) \to 0
$$
을 얻는다. 이 복합체들은
$R\Gamma(X, E)$, $R\Gamma(U, E)$, $R\Gamma(V, E)$, 그리고
$R\Gamma(U \times_X V, E)$를 나타내므로 Derived Categories, Section
\ref{derived-section-canonical-delta-functor}, 특히 Lemma
\ref{derived-lemma-derived-canonical-delta-functor}에 의해 구별 삼각형을
얻는다.
\end{proof}

\begin{lemma}
\label{spaces-perfect-lemma-restrict-lower-shriek}
$S$를 스킴이라 하자. $j : U \to X$를 $S$ 위 대수공간의 \`etale
사상이라 하자. \`Etale 사상 $V \to Y$가 주어졌을 때
% P09-ERR-0411
$W = V \times_X U$라 두고 사영 사상을 $j_W : W \to V$로 나타내자.
그러면 $(j_!E)|_V = j_{W!}(E|_W)$이다. 여기서 $E$는
$D(\mathcal{O}_U)$에 속한다.
\end{lemma}

\begin{proof}
이는 $(j_!\mathcal{F})|_V = j_{W!}(\mathcal{F}|_W)$이기 때문이다.
이 등식은 $\mathcal{O}_X$-가군 $\mathcal{F}$에 대해 성립하고,
함자 $j_!$와 $j_{W!}$의 구성에서 곧바로 따른다. Modules on Sites,
Lemma \ref{sites-modules-lemma-extension-by-zero}를 보라.
\end{proof}

\begin{lemma}
\label{spaces-perfect-lemma-pushforward-with-support-in-open}
$S$를 스킴이라 하자. $(U \subset X, j : V \to X)$를 $S$ 위
대수공간의 기본 구별 정사각형이라 하자.
$T = |X| \setminus |U|$라 두자.
\begin{enumerate}
\item $E$가 $D(\mathcal{O}_X)$의 대상이고 $T$ 위에 지지되면,
(a) $E \to Rj_*(E|_V)$와 (b) $j_!(E|_V) \to E$는 동형이다.
\item $F$가 $D(\mathcal{O}_V)$의 대상이고 $j^{-1}T$ 위에
지지되면, (a) $F \to (j_!F)|_V$, (b) $(Rj_*F)|_V \to F$, 그리고
(c) $j_!F \to Rj_*F$는 동형이다.
\end{enumerate}
\end{lemma}

\begin{proof}
$E$를 $D(\mathcal{O}_X)$의 대상이라 하고 그 코호몰로지 층들이
$T$ 위에 지지된다고 하자. 그러면 $E|_U = 0$이고
$E|_{U \times_X V} = 0$임을 알 수 있다. $T$는 $U$와 만나지 않고
$j^{-1}T$는 $U \times_X V$와 만나지 않기 때문이다. 따라서 (1)(a)는
Lemma \ref{spaces-perfect-lemma-exact-sequence-j-star}에서 따른다. 정확히 같은
방식으로 (1)(b)는 Lemma
\ref{spaces-perfect-lemma-exact-sequence-lower-shriek}에서 따른다.

\medskip\noindent
$F$를 $D(\mathcal{O}_V)$의 대상이라 하고 그 코호몰로지 층들이
$j^{-1}T$ 위에 지지된다고 하자. Lemma
\ref{spaces-perfect-lemma-restrict-direct-image-open}에 의해
$(Rj_*F)|_U = Rj_{W, *}(F|_W) = 0$이다. 이는 가정에 의해
$F|_W = 0$이기 때문이다. 마찬가지로 Lemma
\ref{spaces-perfect-lemma-restrict-lower-shriek}에 의해
$(j_!F)|_U = j_{W!}(F|_W) = 0$이다. 따라서 $j_!F$와 $Rj_*F$는
$T$ 위에 지지되고, $(j_!F)|_V$와 $(Rj_*F)|_V$는
$j^{-1}(T)$ 위에 지지된다. 유도 범주에서 사상 (2)(a), (b), (c)가
동형인지 확인하려면, 이 사상들이 $T$와 $j^{-1}(T)$의 기하점에서
코호몰로지 층의 줄기에 동형을 유도하는지 확인하면 충분하다.% P09-ERR-0412
Properties of Spaces, Theorem
\ref{spaces-properties-theorem-exactness-stalks}를 보라. 이를 위해
$X$를 $V$로, $U$를 $U \times_X V$로, $V$를 $V \times_X V$로,
$F$를 $F|_{V \times_X V}$로 바꿀 수 있다(첫째 사영에 의한 제한).
Lemmas \ref{spaces-perfect-lemma-restrict-direct-image-open},
\ref{spaces-perfect-lemma-restrict-lower-shriek}, 그리고
\ref{spaces-perfect-lemma-make-more-elementary-distinguished-squares}를 보라.
$V \times_X V \to V$가 단면을 가지므로 (2)는
$j : V \to X$가 단면을 가지는 경우로 환원된다.

\medskip\noindent
$j$가 단면 $\sigma : X \to V$를 가진다고 가정하자.
$V' = \sigma(X)$라 두자. 이는 $V$의 열린 부분공간이다.
$U' = j^{-1}(U)$라 두자. 이는 $V$의 또 다른 열린 부분공간이다.
그러면 $(U' \subset V, V' \to V)$는 기본 구별 정사각형이다.
$F|_{U'} = 0$이고 $F|_{V' \cap U'} = 0$임을 관찰하라. 이는 $F$가
$j^{-1}(T)$ 위에 지지되기 때문이다. 열린 몰입을
$j' : V' \to V$라 하고, $j_{V'} : V' \to X$를 합성
$V' \to V \to X$로 나타내자. 이 합성은 $\sigma$의 역이다.
$F' = \sigma^*F$라 두자.
Lemmas \ref{spaces-perfect-lemma-exact-sequence-lower-shriek}와
\ref{spaces-perfect-lemma-exact-sequence-j-star}의 구별 삼각형에 의해
$F = j'_!(F|_{V'})$이고 $F = Rj'_*(F|_{V'})$이다. 따라서
$j_!F = j_!j'_!(F|_{V'}) = j_{V'!}F = F'$이다. 실제로
$j_{V'} : V' \to X$는 동형이고 $\sigma$의 역이다. 마찬가지로
$Rj_*F = Rj_*Rj'_*F = Rj_{V', *}F = F'$이다. 이것으로 (2)(c)가
증명된다. (2)(a)와 (2)(b)를 증명하려면 $F = F'|_V$임을 보이면
충분하다. 이는 명백하다. 실제로 $F$와 $F'|_V$는 모두 $U'$와
$U' \cap V'$ 위에서 영으로 제한되고, $V'$ 위에서는 같은 대상으로
제한된다.
\end{proof}

\noindent
복합체를 접합할 수 있다!

\begin{lemma}
\label{spaces-perfect-lemma-glue}
$S$를 스킴이라 하자. $(U \subset X, V \to X)$를 $S$ 위
대수공간의 기본 구별 정사각형이라 하자. 다음이 주어졌다고 하자.
\begin{enumerate}
\item 대상 $A$가 $D(\mathcal{O}_U)$에 속하고,
\item 대상 $B$가 $D(\mathcal{O}_V)$에 속하며, 그리고
\item 동형 $c : A|_{U \times_X V} \to B|_{U \times_X V}$.
\end{enumerate}
그러면 대상 $F$가 $D(\mathcal{O}_X)$에 존재하고 동형
$f : F|_U \to A$, $g : F|_V \to B$가 존재하여
$c = g|_{U \times_X V} \circ f^{-1}|_{U \times_X V}$이다. 더욱이 다음이
주어졌다고 하자.
\begin{enumerate}
\item 대상 $E$가 $D(\mathcal{O}_X)$에 속하고,
\item 사상 $a : A \to E|_U$가 $D(\mathcal{O}_U)$에 속하며,
\item 사상 $b : B \to E|_V$가 $D(\mathcal{O}_V)$에 속하고,
\end{enumerate}
그리고 다음이 성립한다고 하자.
$$
a|_{U \times_X V}  = b|_{U \times_X V} \circ c.
$$
그러면 사상 $F \to E$가 $D(\mathcal{O}_X)$ 안에 존재하며, 그
$U$로의 제한은 $a \circ f$이고 $V$로의 제한은 $b \circ g$이다.
\end{lemma}

\begin{proof}
사상들을 $j_U$, $j_V$, $j_{U \times_X V}$로 나타내자. 이들은
대응하는 $X$로 향하는 사상들이다. 구별 삼각형
$$
F \to Rj_{U, *}A \oplus Rj_{V, *}B \to
Rj_{U \times_X V, *}(B|_{U \times_X V}) \to F[1]
$$
을 택하자. 여기서 사상
$Rj_{V, *}B \to Rj_{U \times_X V, *}(B|_{U \times_X V})$는 명백한
것이다. 사상
$Rj_{U, *}A \to Rj_{U \times_X V, *}(B|_{U \times_X V})$는
$Rj_{U, *}A \to Rj_{U \times_X V, *}(A|_{U \times_X V})$와
$Rj_{U \times_X V, *}c$의 합성이다. $U$로 제한하면
$$
F|_U \to A \oplus (Rj_{V, *}B)|_U \to
(Rj_{U \times_X V, *}(B|_{U \times_X V}))|_U \to F|_U[1]
$$
을 얻는다. $j : U \times_X V \to U$라 하자. 제한과 전체 직접상의
양립성(Lemma \ref{spaces-perfect-lemma-restrict-direct-image-open})에 의해
$(Rj_{V, *}B)|_U$와
$(Rj_{U \times_X V, *}(B|_{U \times_X V}))|_U$는 둘 다 표준적으로
$Rj_*(B|_{U \times_X V})$와 동형이다. 따라서 마지막 표시식의 둘째
화살표는 단면을 가지며, 사상 $F|_U \to A$가 동형임을 얻는다.

\medskip\noindent
사상 $F|_V \to B$가 동형임을 보기 위해 한 가지 기법을 사용하겠다.
즉, 구별 삼각형
$$
F|_V \to B \to B' \to F[1]|_V
$$
을 $D(\mathcal{O}_V)$ 안에서 택하자. $F|_U \to A$가 동형이고 동형
$c : A|_{U \times_X V} \to B|_{U \times_X V}$가 있으므로,
$F|_V \to B$의 제한은 $U \times_X V$ 위에서 동형이다. 따라서
$B'$는 $j_V^{-1}(T)$ 위에 지지되며, 여기서
$T = |X| \setminus |U|$이다. 한편 구별 삼각형의 사상
$$
\xymatrix{
F \ar[r] \ar[d] &
Rj_{U, *}F|_U \oplus Rj_{V, *}F|_V \ar[r] \ar[d] &
Rj_{U \times_X V, *}F|_{U \times_X V} \ar[r] \ar[d] &
F[1] \ar[d] \\
F \ar[r] &
Rj_{U, *}A \oplus Rj_{V, *}B \ar[r] &
Rj_{U \times_X V, *}(B|_{U \times_X V}) \ar[r] &
F[1]
}
$$
이 있다. 이 도식의 모든 수직 사상은
$Rj_{V, *}F|_V \to Rj_{V, *}B$를 제외하고 동형이므로, 그 사상도
동형이다(Derived Categories, Lemma
\ref{derived-lemma-third-isomorphism-triangle}).% P09-ERR-0413
이는 $Rj_{V, *}B' = 0$임을 함의한다. 따라서 Lemma
\ref{spaces-perfect-lemma-pushforward-with-support-in-open}에 의해 $B' = 0$이다.

\medskip\noindent
사상 $F \to E$의 존재는 $\Hom$에 대한 Mayer--Vietoris 열, 즉
Lemma \ref{spaces-perfect-lemma-mayer-vietoris-hom}에서 따른다.
\end{proof}







\section{코히레이터}
\label{spaces-perfect-section-coherator}

\noindent
$S$를 스킴이라 하자. $X$를 $S$ 위 대수공간이라 하자.
{\it 코히레이터(coherator)}는 함자
$$
Q_X :
\textit{Mod}(\mathcal{O}_X)
\longrightarrow
\QCoh(\mathcal{O}_X)
$$
로서 포함 함자
$\QCoh(\mathcal{O}_X) \to \textit{Mod}(\mathcal{O}_X)$의 오른쪽
수반이다. 이는 모든 대수공간 $X$에 대해 존재하며, 더욱이 수반 사상
$Q_X(\mathcal{F}) \to \mathcal{F}$는 모든 준연접 가군
$\mathcal{F}$에 대해 동형이다. Properties of Spaces,
Proposition \ref{spaces-properties-proposition-coherator}를 보라.
$Q_X$는 (오른쪽 수반이므로) 왼쪽 완전이어서 그 오른쪽 유도 확장
$$
RQ_X :
D(\mathcal{O}_X)
\longrightarrow
D(\QCoh(\mathcal{O}_X)).
$$
을 생각할 수 있다. $Q_X$가 포함 함자
$\QCoh(\mathcal{O}_X) \to \textit{Mod}(\mathcal{O}_X)$의 오른쪽
수반이므로, Derived Categories, Lemma
\ref{derived-lemma-derived-adjoint-functors}에 의해 $RQ_X$는 표준 함자
$D(\QCoh(\mathcal{O}_X)) \to D(\mathcal{O}_X)$의 오른쪽 수반이다.

\medskip\noindent
이 절에서는 함자 $RQ_X$를 연구한다. Section
\ref{spaces-perfect-section-better-coherator}에서는 포함 함자
$D_\QCoh(\mathcal{O}_X) \to D(\mathcal{O}_X)$의 (존재할 경우)
오른쪽 수반을 연구한다. 이는 전자와 밀접하게 관련되어 있다.

\begin{lemma}
\label{spaces-perfect-lemma-affine-pushforward}
$S$를 스킴이라 하자. $f : X \to Y$를 $S$ 위 대수공간의 아핀
사상이라 하자. 그러면 $f_*$는 유도 함자
$f_* : D(\QCoh(\mathcal{O}_X)) \to D(\QCoh(\mathcal{O}_Y))$를 정의한다.
이 함자에 대해 도식
$$
\xymatrix{
D(\QCoh(\mathcal{O}_X)) \ar[d]_{f_*} \ar[r] &
D_\QCoh(\mathcal{O}_X) \ar[d]^{Rf_*} \\
D(\QCoh(\mathcal{O}_Y)) \ar[r] &
D_\QCoh(\mathcal{O}_Y)
}
$$
은 가환이다.
\end{lemma}

\begin{proof}
함자 $f_* : \QCoh(\mathcal{O}_X) \to \QCoh(\mathcal{O}_Y)$는
완전이다. Cohomology of Spaces, Lemma
\ref{spaces-cohomology-lemma-affine-vanishing-higher-direct-images}를 보라.
따라서 Derived Categories, Lemma
\ref{derived-lemma-right-derived-exact-functor}에 의해 $f_*$는 유도 함자
$f_* : D(\QCoh(\mathcal{O}_X)) \to D(\QCoh(\mathcal{O}_Y))$를 정의하며,
이는 임의의 대표 복합체에 $f_*$를 단순히 적용하여 얻는다.
임의의 $\mathcal{O}_X$-가군 복합체 $\mathcal{F}^\bullet$에 대해 표준
사상 $f_*\mathcal{F}^\bullet \to Rf_*\mathcal{F}^\bullet$이 있다.
증명을 끝내기 위해 $\mathcal{F}^\bullet$이 각
$\mathcal{F}^n$이 준연접인 복합체일 때 이 사상이 준동형임을 보이자.
명제는 $Y$ 위에서 \`etale 국소적이므로 $Y$가 아핀이라고 가정할
수 있다. 아핀 사상은 표현 가능하므로 Remark
\ref{spaces-perfect-remark-match-total-direct-images}의 양립성에 의해 스킴의 경우로
환원된다. 스킴의 경우는 Derived Categories of Schemes, Lemma
\ref{perfect-lemma-affine-pushforward}이다.
\end{proof}

\begin{lemma}
\label{spaces-perfect-lemma-flat-pushforward-coherator}
$S$를 스킴이라 하자. $f : X \to Y$를 $S$ 위 대수공간의 사상이라
하자. $f$가 준콤팩트, 준분리이고 평탄이라고 가정하자. 다음으로
나타낸다고 하자.
$$
\Phi : D(\QCoh(\mathcal{O}_X)) \to D(\QCoh(\mathcal{O}_Y))
$$
이는 함자
$f_* : \QCoh(\mathcal{O}_X) \to \QCoh(\mathcal{O}_Y)$의 오른쪽
유도 함자이다. 그러면 $RQ_Y \circ Rf_* = \Phi \circ RQ_X$이다.
\end{lemma}

\begin{proof}
$RQ_Y \circ Rf_*$와 $\Phi \circ RQ_X$가 같은 함자
$D(\QCoh(\mathcal{O}_Y)) \to D(\mathcal{O}_X)$의 오른쪽 수반임을
보여 이를 증명하겠다.

\medskip\noindent
$f$가 준콤팩트이고 준분리이므로 $f_*$는 준연접성을 보존한다.
Morphisms of Spaces, Lemma
\ref{spaces-morphisms-lemma-pushforward}를 보라.
$\QCoh(\mathcal{O}_X)$가 Grothendieck 아벨 범주임을 상기하자
(Properties of Spaces, Proposition
\ref{spaces-properties-proposition-coherator}). 따라서 임의의 $K$가
$D(\QCoh(\mathcal{O}_X))$에 속하면 K-단사 복합체
$\mathcal{I}^\bullet$로 나타낼 수 있으며, 이 복합체는
$\QCoh(\mathcal{O}_X)$ 안에 있다. Injectives, Theorem
\ref{injectives-theorem-K-injective-embedding-grothendieck}를 보라.
이제 $\Phi(K) = f_*\mathcal{I}^\bullet$로 정의할 수 있다.

\medskip\noindent
$f$가 평탄하므로 함자 $f^*$는 완전이다. 따라서 $f^*$는
$f^* : D(\mathcal{O}_Y) \to D(\mathcal{O}_X)$와
$f^* : D(\QCoh(\mathcal{O}_Y)) \to D(\QCoh(\mathcal{O}_X))$를
정의한다. 함자 $f^* = Lf^* : D(\mathcal{O}_Y) \to D(\mathcal{O}_X)$는
$Rf_* : D(\mathcal{O}_X) \to D(\mathcal{O}_Y)$의 왼쪽 수반이다.
Cohomology on Sites, Lemma \ref{sites-cohomology-lemma-adjoint}를 보라.
마찬가지로 함자
$f^* : D(\QCoh(\mathcal{O}_Y)) \to D(\QCoh(\mathcal{O}_X))$는
Derived Categories, Lemma \ref{derived-lemma-derived-adjoint-functors}에
의해
$\Phi : D(\QCoh(\mathcal{O}_X)) \to D(\QCoh(\mathcal{O}_Y))$의 왼쪽
수반이다.

\medskip\noindent
$A$를 $D(\QCoh(\mathcal{O}_Y))$의 대상이라 하고, $E$를
$D(\mathcal{O}_X)$의 대상이라 하자. 그러면
\begin{align*}
\Hom_{D(\QCoh(\mathcal{O}_Y))}(A, RQ_Y(Rf_*E))
& =
\Hom_{D(\mathcal{O}_Y)}(A, Rf_*E) \\
& =
\Hom_{D(\mathcal{O}_X)}(f^*A, E) \\
& =
\Hom_{D(\QCoh(\mathcal{O}_X))}(f^*A, RQ_X(E)) \\
& =
\Hom_{D(\QCoh(\mathcal{O}_Y))}(A, \Phi(RQ_X(E)))
\end{align*}
이다. 이는 원하는 바를 함의한다.
\end{proof}

\begin{lemma}
\label{spaces-perfect-lemma-affine-coherator}
$S$를 스킴이라 하자. $X$를 $S$ 위의 아핀 대수공간이라 하자.
$A = \Gamma(X, \mathcal{O}_X)$라 두자. 그러면
\begin{enumerate}
\item $Q_X : \textit{Mod}(\mathcal{O}_X) \to \QCoh(\mathcal{O}_X)$는
$\mathcal{F}$를 준연접 $\mathcal{O}_X$-가군으로 보내는 함자이다.
이 가군은 $A$-가군 $\Gamma(X, \mathcal{F})$에 결부된 것이다.
\item $RQ_X : D(\mathcal{O}_X) \to D(\QCoh(\mathcal{O}_X))$는
$E$를 준연접 $\mathcal{O}_X$-가군의 복합체로 보내는 함자이다.
이 복합체는 대상 $R\Gamma(X, E)$에 결부되며, 이 대상은 $D(A)$에
속한다.
\item $D_\QCoh(\mathcal{O}_X)$로 제한하면 함자 $RQ_X$는
(\ref{spaces-perfect-equation-compare})의 준역함자를 정의한다.
\end{enumerate}
\end{lemma}

\begin{proof}
$X_0 = \Spec(A)$를 $X$를 나타내는 아핀 스킴이라 하자. 환 달린
사이트의 사상 $\epsilon : X_\etale \to X_{0, Zar}$가 존재하여 동치
$$
\xymatrix{
\QCoh(\mathcal{O}_X) \ar@<1ex>[r]^{{\epsilon_*}} &
\QCoh(\mathcal{O}_{X_0}) \ar@<1ex>[l]^{{\epsilon^*}}
}
$$
를 유도함을 상기하자. Lemma
\ref{spaces-perfect-lemma-derived-quasi-coherent-small-etale-site}를 보라. 따라서
수반함자의 유일성에 의해
$Q_X = \epsilon^* \circ Q_{X_0} \circ \epsilon_*$이다. 그러므로 (1)은
Derived Categories of Schemes, Lemma
\ref{perfect-lemma-affine-coherator}의 $Q_{X_0}$에 대한 설명과 등식
$\Gamma(X_0, \epsilon_*\mathcal{F}) = \Gamma(X, \mathcal{F})$에서 따른다.
(2)는 (1)과 $A$-가군에서 준연접 $\mathcal{O}_X$-가군으로 가는
함자가 완전이라는 사실에서 따른다. 이제 셋째 주장은 스킴에 대한 결과
(Derived Categories of Schemes, Lemma
\ref{perfect-lemma-affine-coherator})와 Lemma
\ref{spaces-perfect-lemma-derived-quasi-coherent-small-etale-site}에서 따른다.
\end{proof}

\noindent
다음으로 함자
$D(\QCoh(\mathcal{O}_X)) \to D_\QCoh(\mathcal{O}_X)$가 언제 동치인지
판정하는 기준을 증명한다.

\begin{lemma}
\label{spaces-perfect-lemma-argument-proves}
$S$를 스킴이라 하자. $X$를 $S$ 위의 준콤팩트 준분리 대수공간이라
하자. 모든 \`etale 사상 $j : V \to W$에 대해 다음을 가정하자.
여기서 $W \subset X$는 준콤팩트 열린집합이고 $V$는 아핀이다.
오른쪽 유도 함자
$$
\Phi : D(\QCoh(\mathcal{O}_U)) \to D(\QCoh(\mathcal{O}_W))
$$
가 다음 왼쪽 완전 함자의 유도 함자라고 가정하자.
% P09-ERR-0414
$j_* : \QCoh(\mathcal{O}_V) \to \QCoh(\mathcal{O}_W)$
또한 이 함자가 가환 도식
$$
\xymatrix{
D(\QCoh(\mathcal{O}_V)) \ar[d]_\Phi \ar[r]_{i_V} &
D_\QCoh(\mathcal{O}_V) \ar[d]^{Rj_*} \\
D(\QCoh(\mathcal{O}_W)) \ar[r]^{i_W} &
D_\QCoh(\mathcal{O}_W)
}
$$
에 맞는다고 하자. 그러면 함자 (\ref{spaces-perfect-equation-compare})
$$
D(\QCoh(\mathcal{O}_X))
\longrightarrow
D_\QCoh(\mathcal{O}_X)
$$
는 동치이고, 그 준역함자는 $RQ_X$로 주어진다.
\end{lemma}

\begin{proof}
먼저 귀납 원리를 사용하여 $i_X$가 충실충만임을 증명한다. 더
정확히는 Lemma \ref{spaces-perfect-lemma-induction-principle-enlarge}를 사용하겠다.
$(U \subset W, V \to W)$를 기본 구별 정사각형이라 하고, 여기서
$V$는 아핀이며 $U, W$는 $X$의 준콤팩트 열린집합이라고 하자.
$i_U$가 충실충만이라고 가정한다. $i_W$가 충실충만임을 보여야 한다.
$X$를 $W$로 바꿀 수 있다. 즉, $W = X$라고 가정할 수 있다
(이는 표기를 단순화하기 위한 것이며, 보조정리의 명제에 있는 조건은
이 연산 아래 보존됨에 유의하라).

\medskip\noindent
$A, B$가 $D(\QCoh(\mathcal{O}_X))$의 대상이라고 하자. 사상
$$
\Hom_{D(\QCoh(\mathcal{O}_X))}(A, B)
\longrightarrow
\Hom_{D(\mathcal{O}_X)}(i_X(A), i_X(B))
$$
이 전단사임을 보이고자 한다. $T = |X| \setminus |U|$라 두자.

\medskip\noindent
먼저 $i_X(B)$가 $T$ 위에 지지된다고 가정하자. 이 경우 사상
$$
i_X(B) \to Rj_{V, *}(i_X(B)|_V) = Rj_{V, *}(i_V(B|_V))
$$
은 준동형이다(Lemma
\ref{spaces-perfect-lemma-pushforward-with-support-in-open}). 가정에 의해 동형
$i_X(\Phi(B|_V)) \to Rj_{V, *}(i_V(B|_V))$가
$D(\mathcal{O}_X)$ 안에 존재한다. 더욱이 $\Phi$와 ${-}|_V$는
준연접 가군의 유도 범주에서 수반함자이다(Derived Categories, Lemma
\ref{derived-lemma-derived-adjoint-functors}에 의해). 수반 사상
$B \to \Phi(B|_V)$는 $i_X$를 적용한 뒤 동형이 되므로
$D(\QCoh(\mathcal{O}_X))$ 안에서도 동형이다. 따라서
\begin{align*}
\Mor_{D(\QCoh(\mathcal{O}_X))}(A, B)
& =
\Mor_{D(\QCoh(\mathcal{O}_X))}(A, \Phi(B|_V)) \\
& =
\Mor_{D(\QCoh(\mathcal{O}_V))}(A|_V, B|_V) \\
& =
\Mor_{D(\mathcal{O}_V)}(i_V(A|_V), i_V(B|_V)) \\
& =
\Mor_{D(\mathcal{O}_X)}(i_X(A), Rj_{V, *}(i_V(B|_V))) \\
& =
\Mor_{D(\mathcal{O}_X)}(i_X(A), i_X(B))
\end{align*}
이다. 여기서는 $i_V$가 충실충만임을 사용하였다(Lemma
\ref{spaces-perfect-lemma-affine-coherator}).

\medskip\noindent
일반적으로 임의의 복합체 $\mathcal{B}^\bullet$을 준연접
$\mathcal{O}_X$-가군들로 이루어진 것으로 택하여 $B$를 나타내게 하자.
다음으로 임의의 준동형
$s : \mathcal{B}^\bullet|_U \to \mathcal{C}^\bullet$을 택하자. 이는
$U$ 위 준연접 가군의 복합체에 대한 것이다.
$j_U : U \to X$가 준콤팩트이고 준분리이므로 함자 $j_{U, *}$는
준연접 가군을 준연접 가군으로 보낸다(Morphisms of Spaces, Lemma
\ref{spaces-morphisms-lemma-pushforward}). 따라서 표준 사상
$\mathcal{B}^\bullet \to j_{U, *}(\mathcal{B}^\bullet|_U) \to
j_{U, *}\mathcal{C}^\bullet$이 $X$ 위 준연접 가군의 복합체 사이에
존재한다. $B'' = j_{U, *}\mathcal{C}^\bullet$라
$D(\QCoh(\mathcal{O}_X))$에서 두고 구별 삼각형
$$
B \to B'' \to B' \to B[1]
$$
을 $D(\QCoh(\mathcal{O}_X))$ 안에서 택하자. 이 삼각형의 첫 화살표는
$U$ 위로 제한하면 동형이므로 $B'$는 $T$ 위에 지지된다. 따라서 도식
$$
\xymatrix{
\Hom_{D(\QCoh(\mathcal{O}_X))}(A, B'[-1]) \ar[r] \ar[d] &
\Hom_{D(\mathcal{O}_X)}(i_X(A), i_X(B')[-1]) \ar[d] \\
\Hom_{D(\QCoh(\mathcal{O}_X))}(A, B) \ar[r] \ar[d] &
\Hom_{D(\mathcal{O}_X)}(i_X(A), i_X(B)) \ar[d] \\
\Hom_{D(\QCoh(\mathcal{O}_X))}(A, B'') \ar[r] \ar[d] &
\Hom_{D(\mathcal{O}_X)}(i_X(A), i_X(B'')) \ar[d] \\
\Hom_{D(\QCoh(\mathcal{O}_X))}(A, B') \ar[r] &
\Hom_{D(\mathcal{O}_X)}(i_X(A), i_X(B'))
}
$$
에서 열들은 완전이고 맨 위와 맨 아래의 수평 화살표는 전단사이다.
마지막으로 준연접 가군의 복합체 $\mathcal{A}^\bullet$을 택하여
$A$를 나타내게 하자.

\medskip\noindent
$\alpha : i_X(A) \to i_X(B)$를 $D(\mathcal{O}_X)$ 안의 사상이라
하자.% P09-ERR-0415
제한 $\alpha|_U$는 $D(\QCoh(\mathcal{O}_U))$ 안의 사상에서 온다.
이는 $i_U$가 충실충만이기 때문이다. 따라서 위와 같은
$s : \mathcal{B}^\bullet|_U \to \mathcal{C}^\bullet$을 택하되
$\alpha|_U$가 복합체의 실제 사상
$\mathcal{A}^\bullet|_U \to \mathcal{C}^\bullet$으로 나타내어지게
할 수 있다. 이는 복합체의 사상
$\mathcal{A} \to j_{U, *}\mathcal{C}^\bullet$에 대응한다. 다시 말해
$\alpha$의
$\Hom_{D(\mathcal{O}_X)}(i_X(A), i_X(B''))$에서의 상은
$\Hom_{D(\QCoh(\mathcal{O}_X))}(A, B'')$의 한 원소에서 온다.
이제 도식 추적을 하면 $\alpha$가 사상 $A \to B$에서 옴을 알 수
있다. 이 사상은 $D(\QCoh(\mathcal{O}_X))$ 안에 있다. 마지막으로
$a : A \to B$가 $D(\QCoh(\mathcal{O}_X))$의 사상이고
$D(\mathcal{O}_X)$ 안에서는 영이 된다고 하자.
$\mathcal{B}^\bullet$을 알맞게 택한 뒤 $a$가 복합체의 사상
$a^\bullet : \mathcal{A}^\bullet \to \mathcal{B}^\bullet$으로
나타내어진다고 가정할 수 있다. $i_U$가 충실충만이므로
% P09-ERR-0416
제한 $a^\bullet|_U$는 $D(\QCoh(\mathcal{O}_U))$ 안에서 영이다.
따라서 $s$를 택하여
$s \circ a^\bullet|_U : \mathcal{A}^\bullet|_U \to \mathcal{C}^\bullet$이
영과 호모토픽이 되게 할 수 있다. 함자 $j_{U, *}$를 적용하면
$\mathcal{A}^\bullet \to j_{U, *}\mathcal{C}^\bullet$이 영과
호모토픽임을 얻는다. 따라서 $a$는
$\Hom_{D(\QCoh(\mathcal{O}_X))}(A, B'')$에서 영으로 간다.
그러므로 $a$가 어떤 원소
$b \in \Hom_{D(\QCoh(\mathcal{O}_X))}(A, B'[-1])$의 상이라고
가정할 수 있다. $b$의
$\Hom_{D(\mathcal{O}_X)}(i_X(A), i_X(B')[-1])$에서의 상은 어떤
$\gamma \in \Hom_{D(\mathcal{O}_X)}(A, B''[-1])$에서 온다.
% P09-ERR-0417
(이는 $a$가 오른쪽 군에서 영으로 가기 때문이다.) 위에서 수평
화살표들이 전사임을 보았으므로 $\gamma$는 어떤 $c$에서 오며, 이는
$\Hom_{D(\QCoh(\mathcal{O}_X))}(A, B''[-1])$에 속한다. 이것은
원하는 대로 $a = 0$을 함의한다.

\medskip\noindent
이제 $i_X$가 원래의 $X$에 대해 충실충만임을 안다. $RQ_X$가 그
오른쪽 수반이므로
$RQ_X \circ i_X = \text{id}$이다(Categories, Lemma
\ref{categories-lemma-adjoint-fully-faithful}). 증명을 끝내기 위해
임의의 $E$가 $D_\QCoh(\mathcal{O}_X)$에 속할 때 사상
$i_X(RQ_X(E)) \to E$가 동형임을 보이자. 구별 삼각형
$$
i_X(RQ_X(E)) \to E \to E' \to i_X(RQ_X(E))[1]
$$
을 $D_\QCoh(\mathcal{O}_X)$ 안에서 택하자. 위의 결과를 사용하는
형식적 논증으로 $i_X(RQ_X(E')) = 0$을 얻는다. 따라서
$E \in D_\QCoh(\mathcal{O}_X)$에 대해 조건
$i_X(RQ_X(E)) = 0$이 $E = 0$을 함의함을 증명하면 충분하다.
\`Etale 사상 $j : V \to X$를 생각하되 $V$가 아핀이라고 하자. Lemmas
\ref{spaces-perfect-lemma-affine-coherator}와
\ref{spaces-perfect-lemma-flat-pushforward-coherator}, 그리고 우리의 가정에 의해
$$
Rj_*(E|_V) = Rj_*(i_V(RQ_V(E|_V))) = i_X(\Phi(RQ_V(E|_V))) =
i_X(RQ_X(Rj_*(E|_V)))
$$
이다. 구별 삼각형
$$
E \to Rj_*(E|_V) \to E' \to E[1]
$$
을 택하자. $RQ_X$를 적용하여 구별 삼각형
$$
0 \to RQ_X(Rj_*(E|_V)) \to RQ_X(E') \to 0[1]
$$
을 얻는다. 다시 말해 가운데 사상은 동형이다. 이를 위의 등식들과
결합하면 첫 구별 삼각형이 구별 삼각형
$$
E \to i_X(RQ_X(E')) \to E' \to E[1]
$$
이 됨을 알 수 있다. 여기서 가운데 사상은 수반 사상이다. 그러나 합성
$E \to E'$는 영이므로 수반성에 의해
$E \to i_X(RQ_X(E'))$도 영이다! 이 사상은 사상
$E \to Rj_*(E|_V)$와 동형이고, 후자는
$\text{id} : E|_V \to E|_V$에 수반한다. 따라서 $E|_V$가 영임을
얻는다. 이는 모든 아핀 $V$가 $X$ 위에서 \`etale일 때 성립한다. 따라서 원하는 대로
$E$가 영임을 얻는다.
\end{proof}

\begin{proposition}
\label{spaces-perfect-proposition-quasi-compact-affine-diagonal}
$S$를 스킴이라 하자. $X$를 $S$ 위의 준콤팩트 대수공간이라 하고,
$\mathbf{Z}$ 위에서 아핀 대각사상을 갖는다고 하자(Properties of
Spaces, Definition \ref{spaces-properties-definition-separated}에서와
같이). 그러면 함자 (\ref{spaces-perfect-equation-compare})
$$
D(\QCoh(\mathcal{O}_X))
\longrightarrow
D_\QCoh(\mathcal{O}_X)
$$
는 동치이고, 그 준역함자는 $RQ_X$로 주어진다.
\end{proposition}

\begin{proof}
$V \to W$를 \`etale 사상이라 하자. 여기서 $V$는 아핀이고 $W$는
$X$의 준콤팩트 열린 부분공간이다. 그러면 사상 $V \to W$는
아핀이다. 실제로 $W$는 $\mathbf{Z}$ 위에서 아핀 대각사상을 가지고
$V$는 아핀이다(Morphisms of
Spaces, Lemma \ref{spaces-morphisms-lemma-affine-permanence}). 그러면
Lemma \ref{spaces-perfect-lemma-affine-pushforward}가 Lemma
\ref{spaces-perfect-lemma-argument-proves}의 가정이 성립함을 보장한다. 따라서 결론을
얻는다.
\end{proof}

\begin{lemma}
\label{spaces-perfect-lemma-direct-image-coherator}
$S$를 스킴이라 하고 $f : X \to Y$를 $S$ 위 대수공간의 사상이라
하자. $X$와 $Y$가 준콤팩트이고 $\mathbf{Z}$ 위에서 아핀
대각사상을 갖는다고 가정하자(Properties of Spaces, Definition
\ref{spaces-properties-definition-separated}에서와 같이). 다음으로
나타낸다고 하자.
$$
\Phi : D(\QCoh(\mathcal{O}_X)) \to D(\QCoh(\mathcal{O}_Y))
$$
이는 다음 함자의 오른쪽 유도 함자이다.
$f_* : \QCoh(\mathcal{O}_X) \to \QCoh(\mathcal{O}_Y)$
% P09-ERR-0418
도식
$$
\xymatrix{
D(\QCoh(\mathcal{O}_X)) \ar[d]_\Phi \ar[r] &
D_\QCoh(\mathcal{O}_X) \ar[d]^{Rf_*} \\
D(\QCoh(\mathcal{O}_Y)) \ar[r] &
D_\QCoh(\mathcal{O}_Y)
}
$$
은 가환이다.
\end{lemma}

\begin{proof}
도식의 수평 화살표들은 Proposition
\ref{spaces-perfect-proposition-quasi-compact-affine-diagonal}에 의해 범주의
동치임을 관찰하라. 따라서 이 범주들을 동일시할 수 있으며(아핀
대각사상을 갖는 다른 준콤팩트 대수공간에 대해서도 마찬가지이다),
그러면 보조정리의 명제는 표준 사상
$\Phi(K) \to Rf_*(K)$가 모든 $K$에 대해 동형이라는 것이다.
여기서 그 대상은 $D(\QCoh(\mathcal{O}_X))$에 속한다. 만일
$K_1 \to K_2 \to K_3 \to K_1[1]$이
$D(\QCoh(\mathcal{O}_X))$ 안의 구별 삼각형이고 세 대상 가운데 둘에
대해 명제가 참이면 셋째 대상에 대해서도 참임에 유의하라.

\medskip\noindent
$\mathcal{B} \subset \Ob(X_{spaces, \etale})$를 준콤팩트이고 아핀
대각사상을 갖는 대상들의 집합이라 하자. $U \in \mathcal{B}$이고
임의의 사상 $g : U \to Z$가 주어졌다고 하자. 여기서 $Z$는 $S$
위의, 아핀 대각사상을 갖는 준콤팩트 대수공간이다. 다음으로 나타내자.
$$
\Phi_g : D(\QCoh(\mathcal{O}_U)) \to D(\QCoh(\mathcal{O}_Z))
$$
이는 $g_*$의 유도 확장이다. $P(U) =$ ``임의의 $K$가
$D(\QCoh(\mathcal{O}_U))$에 속하고 위와 같은 임의의
$g : U \to Z$가 주어졌을 때 사상 $\Phi_g(K) \to Rg_*K$가 동형이다''라고
두자. Remark \ref{spaces-perfect-remark-how-to}에 의해 Lemma
\ref{spaces-perfect-lemma-induction-principle-separated}의 조건 (1), (2), (3)(a)는
성립하며, (3)(b)와 (4)를 증명하는 일만 남는다.

\medskip\noindent
조건 (3)(b)를 확인하자. $U$를 아핀 스킴이라 하고 $X$ 위에서
\`etale이라고 하자. $g : U \to Z$를 위와 같은 사상이라 하자.
$Z$의 대각사상이 아핀이므로 $g : U \to Z$는 아핀이다(Morphisms of Spaces, Lemma
\ref{spaces-morphisms-lemma-affine-permanence}). 따라서 Lemma
\ref{spaces-perfect-lemma-affine-pushforward}에 의해 $P(U)$가 성립한다.

\medskip\noindent
조건 (4)를 확인하자. $(U \subset W, V \to W)$를 기본 구별
정사각형이라 하자. 이는 $X_{spaces, \etale}$ 안에 있고
$U, W, V$는 $\mathcal{B}$에 속하며 $V$는 아핀이다. $P$가 $U$,
$V$, $U \times_W V$에 대해 성립한다고 가정한다. $P$가 $W$에 대해
성립함을 보여야 한다. $g : W \to Z$를 아핀 대각사상을 갖는
준콤팩트 대수공간으로 가는 사상이라 하자. $K$를
$D(\QCoh(\mathcal{O}_W))$의 대상이라 하자. 구별 삼각형
$$
K \to Rj_{U, *}K|_U \oplus Rj_{V, *}K|_V \to
Rj_{U \times_W V, *}K|_{U \times_W V} \to K[1]
$$
을 $D(\mathcal{O}_W)$ 안에서 생각하자. 위에서 언급한 세 항 중 둘
성질에 의해
$\Phi_g(Rj_{U, *}K|_U) \to Rg_*(Rj_{U, *}K|_U)$가 동형임을 보이면
충분하며, $V$와 $U \times_W V$에 대해서도 마찬가지이다. 이는 다음
문단에서 논한다.

\medskip\noindent
$j : U \to W$를 사상이라 하자. 이 사상은
$X_{spaces, \etale}$에 속한다.% P09-ERR-0419
$U, W$는 $\mathcal{B}$에 속하고 $P$는 $U$에 대해 성립한다고 하자.
$g : W \to Z$를 아핀 대각사상을 갖는 준콤팩트 대수공간으로 가는
사상이라 하자. 증명을 끝내려면
$\Phi_g(Rj_*K) \to Rg_*(Rj_*K)$가 임의의 $K$에 대해 동형임을
보여야 한다. 여기서 그 대상은 $D(\QCoh(\mathcal{O}_U))$에 속한다.
$\mathcal{I}^\bullet$을 $\QCoh(\mathcal{O}_U)$ 안의 K-단사
복합체라 하여 $K$를 나타내게 하자. $P(U)$를 $j$에 적용하면
$j_*\mathcal{I}^\bullet$이 $Rj_*K$를 나타냄을 알 수 있다.
함자 $j_* : \QCoh(\mathcal{O}_U) \to \QCoh(\mathcal{O}_X)$는 완전한
왼쪽 수반 $j^* : \QCoh(\mathcal{O}_X) \to \QCoh(\mathcal{O}_U)$를
가지므로, $j_*\mathcal{I}^\bullet$은 $\QCoh(\mathcal{O}_W)$ 안의
K-단사 복합체이다.% P09-ERR-0420
Derived Categories, Lemma
\ref{derived-lemma-adjoint-preserve-K-injectives}를 보라. 따라서
$\Phi_g(Rj_*K)$는
$g_*j_*\mathcal{I}^\bullet = (g \circ j)_*\mathcal{I}^\bullet$로
나타내어진다. $P(U)$를 $g \circ j$에 적용하면 이것이
$R_{g \circ j, *}(K) = Rg_*(Rj_*K)$를 나타냄을 알 수 있다. 이것으로
증명이 끝난다.
\end{proof}










\section{Noether 공간의 코히레이터}
\label{spaces-perfect-section-coherator-Noetherian}

\noindent
Noether 대수공간의 경우를 다루려면 단사 가군에 관해 조금 더 알아야 한다.

\begin{lemma}
\label{spaces-perfect-lemma-affine-injective-colimit-direct-sum-pushforwards-artin}
$S$를 Noether 아핀 스킴이라 하자. $\QCoh(\mathcal{O}_S)$의 모든 단사
대상은 다음 꼴의 준연접 층들로 이루어진 여과 쌍대극한
$\colim_i \mathcal{F}_i$이다.
$$
\mathcal{F}_i = (Z_i \to S)_*\mathcal{G}_i
$$
여기서 $Z_i$는 Artin 환의 스펙트럼이고 $\mathcal{G}_i$는 $Z_i$ 위의
연접 가군이다.
\end{lemma}

\begin{proof}
$S = \Spec(A)$라 하자. $\mathcal{J}$를 $\QCoh(\mathcal{O}_S)$의 단사
대상이라 하자. $\QCoh(\mathcal{O}_S)$는 $A$-가군들의 범주와
동치이므로, $\mathcal{J}$가 $\widetilde{J}$와 같음을 알 수 있다.
여기서 어떤 단사 $A$-가군을 $J$라 쓴다. Dualizing Complexes, Proposition
\ref{dualizing-proposition-structure-injectives-noetherian}에 의해
$J = \bigoplus E_\alpha$로 쓸 수 있으며, 각 $E_\alpha$는 직분해
불가능하고 따라서 한 점에서의 reside 체의 단사 포락과 동형이다.% P09-ERR-0421
그러므로 Artin 스킴들의 유한 분리합은 Artin이므로, $J$가
$\kappa(\mathfrak p)$의 단사 포락이라고 가정해도 된다. 여기서
$\mathfrak p$는 $A$의 소 아이디얼이다. 그러면
$J = \bigcup J[\mathfrak p^n]$이고, $J[\mathfrak p^n]$은
$\kappa(\mathfrak p)$의
$A_\mathfrak/\mathfrak p^nA_\mathfrak p$ 위 단사 포락이다.% P09-ERR-0422
Dualizing Complexes, Lemma \ref{dualizing-lemma-union-artinian}를 보라.
따라서 $\widetilde{J}$는 층
$(Z_n \to X)_*\mathcal{G}_n$들의 쌍대극한이다.% P09-ERR-0423
여기서 $Z_n = \Spec(A_\mathfrak p/\mathfrak p^nA_\mathfrak p)$이고
$\mathfrak G_n$은 유한
$A_\mathfrak/\mathfrak p^nA_\mathfrak p$-가군 $J[\mathfrak p^n]$에
연관된 연접 층이다.% P09-ERR-0422 P09-ERR-0424
유한성은 Dualizing Complexes, Lemma
\ref{dualizing-lemma-finite}에서 따른다.
\end{proof}

\begin{lemma}
\label{spaces-perfect-lemma-injective-colimit-direct-sum-pushforwards-artin}
$S$를 아핀 스킴이라 하자. $X$를 $S$ 위의 Noether 대수공간이라 하자.
$\QCoh(\mathcal{O}_X)$의 모든 단사 대상은 다음 꼴의 준연접 층들로
이루어진 여과 쌍대극한 $\colim_i \mathcal{F}_i$의 직합 성분이다.
$$
\mathcal{F}_i = (Z_i \to X)_*\mathcal{G}_i
$$
여기서 $Z_i$는 Artin 환의 스펙트럼이고 $\mathcal{G}_i$는 $Z_i$ 위의
연접 가군이다.
\end{lemma}

\begin{proof}
아핀 스킴 $U$와 전사 \`etale 사상 $j : U \to X$를 택하자(Properties
of Spaces, Lemma
\ref{spaces-properties-lemma-quasi-compact-affine-cover}). 그러면 $U$는
Noether 아핀 스킴이다. 단사 대상 $\mathcal{J}'$을
$\QCoh(\mathcal{O}_U)$ 안에서 택하여 단사 사상
$\mathcal{J}|_U \to \mathcal{J}'$이 존재하게 하자. 그러면
$$
\mathcal{J} \to j_*\mathcal{J}'
$$
은 $\QCoh(\mathcal{O}_X)$ 안의 단사 사상이므로, $\mathcal{J}$를
$j_*\mathcal{J}'$의 직합 성분과 동일시한다. 따라서 결과는 Lemma
\ref{spaces-perfect-lemma-affine-injective-colimit-direct-sum-pushforwards-artin}에서
증명한 $\mathcal{J}'$에 대한 해당 결과로부터 따른다.
\end{proof}

\begin{lemma}
\label{spaces-perfect-lemma-flat-pullback-injective-quasi-coherent}
$S$를 스킴이라 하자. $f : X \to Y$를 $S$ 위 대수공간들의 평탄하고
준콤팩트이며 준분리인 사상이라 하자. $\mathcal{J}$가
$\QCoh(\mathcal{O}_X)$의 단사 대상이면, $f_*\mathcal{J}$는
$\QCoh(\mathcal{O}_Y)$의 단사 대상이다.
\end{lemma}

\begin{proof}
$f$가 준콤팩트이고 준분리이므로, 함자 $f_*$는 준연접 층을 준연접
층으로 보낸다(Morphisms of Spaces, Lemma
\ref{spaces-morphisms-lemma-pushforward}). 함자 $f^*$는 $f_*$의 왼쪽
수반이며 단사 사상을 단사 사상으로 보낸다. 따라서 결과는 Homology,
Lemma \ref{homology-lemma-adjoint-preserve-injectives}에서 따른다% P09-ERR-0425
\end{proof}

\begin{lemma}
\label{spaces-perfect-lemma-injective-pushforward}
$S$를 스킴이라 하자. $X$를 $S$ 위의 Noether 대수공간이라 하자.
$\mathcal{J}$가 $\QCoh(\mathcal{O}_X)$의 단사 대상이면 다음이 성립한다.
\begin{enumerate}
\item $H^p(U, \mathcal{J}|_U) = 0$이다. 여기서 $p > 0$이고,
$U$는 $X$ 위에서 \`etale인 임의의 준콤팩트 준분리 대수공간이다.
\item $f : X \to Y$가 $S$ 위 대수공간들의 임의의 사상이고 $Y$가
준분리이면 $R^pf_*\mathcal{J} = 0$이다. 여기서 $p > 0$이다.
\end{enumerate}
\end{lemma}

\begin{proof}
(1)의 증명. Lemma
\ref{spaces-perfect-lemma-injective-colimit-direct-sum-pushforwards-artin}에서와 같이
$\mathcal{J}$를 $\colim \mathcal{F}_i$의 직합 성분으로 쓰자. 여기서
$\mathcal{F}_i = (Z_i \to X)_*\mathcal{G}_i$이다. 소멸을
$\colim \mathcal{F}_i$에 대해 증명하면 충분함은 분명하다. 당김은
쌍대극한과 가환하고 $U$는 준콤팩트이고 준분리이므로,
$H^p(U, \mathcal{F}_i|_U) = 0$임을 증명하면 충분하다. 여기서
$p > 0$이다. Cohomology
of Spaces, Lemma \ref{spaces-cohomology-lemma-colimits}를 보라.
$Z_i \to X$가 아핀 사상임을 관찰하라(Morphisms of Spaces, Lemma
\ref{spaces-morphisms-lemma-Artinian-affine}). 따라서
$$
\mathcal{F}_i|_U = (Z_i \times_X U \to U)_*\mathcal{G}'_i =
R(Z_i \times_X U \to U)_*\mathcal{G}'_i
$$
이다. 여기서 $\mathcal{G}'_i$는 $\mathcal{G}_i$를
$Z_i \times_X U$로 당긴 것이다. Cohomology of Spaces, Lemma
\ref{spaces-cohomology-lemma-affine-base-change}를 보라.
$Z_i \times_X U$가 아핀이므로, $\mathcal{G}'_i$는
$Z_i \times_X U$ 위에서 고차 코호몰로지가 없다. Leray 스펙트럼
열에 의해 $\mathcal{F}_i|_U$에도 같은 사실이 성립한다(Cohomology on
Sites, Lemma \ref{sites-cohomology-lemma-apply-Leray}).

\medskip\noindent
(2)의 증명. $f : X \to Y$를 $S$ 위 대수공간들의 사상이라 하자.
$V \to Y$를 \`etale 사상이라 하고 $V$가 아핀이라고 하자. 그러면
$V \times_Y X \to X$는 \`etale 사상이고, $V \times_Y X$는 $X$ 위에서
\`etale인 준콤팩트 준분리 대수공간이다(세부사항은 생략한다). 따라서
(1)에 의해 $H^p(V \times_Y X, \mathcal{J})$는 영이다.
$R^pf_*\mathcal{J}$는 준층
$V \mapsto H^p(V \times_Y X, \mathcal{J})$에 연관된 층이므로 결과가
증명된다.
\end{proof}

\begin{lemma}
\label{spaces-perfect-lemma-Noetherian-pushforward}
$S$를 스킴이라 하자. $f : X \to Y$를 $S$ 위 Noether 대수공간들의
사상이라 하자. 그러면 준연접 층 위의 $f_*$는 오른쪽 유도 확장
$\Phi : D(\QCoh(\mathcal{O}_X)) \to D(\QCoh(\mathcal{O}_Y))$를
가지며, 도식
$$
\xymatrix{
D(\QCoh(\mathcal{O}_X)) \ar[d]_{\Phi} \ar[r] &
D_\QCoh(\mathcal{O}_X) \ar[d]^{Rf_*} \\
D(\QCoh(\mathcal{O}_Y)) \ar[r] &
D_\QCoh(\mathcal{O}_Y)
}
$$
은 가환이다.
\end{lemma}

\begin{proof}
$X$와 $Y$가 Noether이므로 이 사상은 준콤팩트이고 준분리이다
(Morphisms of Spaces, Lemma
\ref{spaces-morphisms-lemma-quasi-compact-quasi-separated-permanence}를
보라). 따라서 $f_*$는 준연접성을 보존한다.% P09-ERR-0426
Morphisms of Spaces, Lemma \ref{spaces-morphisms-lemma-pushforward}를 보라.
다음으로 $K$를 $D(\QCoh(\mathcal{O}_X))$의 대상이라 하자.
$\QCoh(\mathcal{O}_X)$는 Grothendieck 아벨 범주이므로(Properties of
Spaces, Proposition \ref{spaces-properties-proposition-coherator}),
$K$를 K-단사 복합체 $\mathcal{I}^\bullet$로 나타낼 수 있고, 각
$\mathcal{I}^n$은 $\QCoh(\mathcal{O}_X)$의 단사 대상이 되게 할 수 있다.
Injectives, Theorem
\ref{injectives-theorem-K-injective-embedding-grothendieck}을 보라.
따라서 함자 $\Phi$는 다음과 같이 정의된다.
$$
\Phi(K) = f_*\mathcal{I}^\bullet
$$
여기서 오른쪽 항은 $D(\QCoh(\mathcal{O}_Y))$의 대상으로 본다.
보조정리의 증명을 끝내려면 표준 사상
$$
f_*\mathcal{I}^\bullet \longrightarrow Rf_*\mathcal{I}^\bullet
$$
이 $D(\mathcal{O}_Y)$ 안에서 동형임을 보이면 충분하다. 이를 위해서는
그 사상이 코호몰로지 층들 위에 동형을 유도함을 보이면 충분하다.
임의의 $m \in \mathbf{Z}$를 택하자. $N = N(X, Y, f)$를 Lemma
\ref{spaces-perfect-lemma-quasi-coherence-direct-image}에서와 같이 두자. 다음 짧은 완전열
$$
0 \to \sigma_{\geq m - N - 1}\mathcal{I}^\bullet \to
\mathcal{I}^\bullet \to \sigma_{\leq m - N - 2}\mathcal{I}^\bullet \to 0
$$
을 생각하자. 이것은 $X$ 위 준연접 층들의 복합체로 이루어진 열이다.
Lemma \ref{spaces-perfect-lemma-quasi-coherence-direct-image}에 의해
$Rf_*\sigma_{\leq m - N - 2}\mathcal{I}^\bullet$의 코호몰로지 층들은
$\geq m - 1$ 차수에서 영이다. 따라서
$R^mf_*\mathcal{I}^\bullet$은
$R^mf_*\sigma_{\geq m - N - 1}\mathcal{I}^\bullet$과 동형이다. 다시
말해 $\mathcal{I}^\bullet$이 $\QCoh(\mathcal{O}_X)$의 단사 대상들로
이루어진 아래로 유계인 복합체라고 가정해도 된다. 이 경우는 Leray의
비순환성 보조정리(Derived Categories, Lemma
\ref{derived-lemma-leray-acyclicity})에서 따르며, 필요한 소멸은 Lemma
\ref{spaces-perfect-lemma-injective-pushforward}가 준다.% P09-ERR-0427
\end{proof}

\begin{proposition}
\label{spaces-perfect-proposition-Noetherian}
$S$를 스킴이라 하자. $X$를 $S$ 위의 Noether 대수공간이라 하자.
그러면 함자 (
\ref{spaces-perfect-equation-compare})
$$
D(\QCoh(\mathcal{O}_X))
\longrightarrow
D_\QCoh(\mathcal{O}_X)
$$
는 동치이고, 그 준역함자는 $RQ_X$로 주어진다.
\end{proposition}

\begin{proof}
Lemmas \ref{spaces-perfect-lemma-Noetherian-pushforward}와
\ref{spaces-perfect-lemma-argument-proves}에서 곧바로 따른다.
\end{proof}














\section{유사연접 복합체와 완전 복합체}
\label{spaces-perfect-section-spell-out}

\noindent
이 절에서는 대수공간의 \`etale 사이트에 대해 Cohomology on Sites,
Sections \ref{sites-cohomology-section-strictly-perfect},
\ref{sites-cohomology-section-pseudo-coherent},
\ref{sites-cohomology-section-tor}, 그리고
\ref{sites-cohomology-section-perfect}에서 정의한 일반 개념들을
연구한다. 특히 이를 스킴에서 일어나는 현상과 맞춘다.

\medskip\noindent
먼저 스킴 위의 유사연접 복합체 개념과 그에 연관된 작은 \`etale
사이트 위의 개념을 비교한다.

\begin{lemma}
\label{spaces-perfect-lemma-descend-finite-type}
$X$를 스킴이라 하자. $\mathcal{F}$를 $\mathcal{O}_X$-가군이라 하자.
다음은 동치이다.
\begin{enumerate}
\item $\mathcal{F}$는 $\mathcal{O}_X$-가군으로서 유한형이다.
\item $\epsilon^*\mathcal{F}$는 $\mathcal{O}_\etale$-가군으로서
유한형이다. 여기서 가군은 $X$의 작은 \`etale 사이트 위에 있다.
\end{enumerate}
여기서 $\epsilon$은 (\ref{spaces-perfect-equation-epsilon})에서와 같다.
\end{lemma}

\begin{proof}
(1) $\Rightarrow$ (2)는 일반적인 사실이다. Modules on Sites, Lemma
\ref{sites-modules-lemma-local-pullback}을 보라. (2)를 가정하자.
가정에 의해 \`etale 덮개
$\{f_i : X_i \to X\}$가 존재하여
$\epsilon^*\mathcal{F}|_{(X_i)_\etale}$가 유한 개의 단면으로 생성된다.
$x \in X$라 하자. $\mathcal{F}$가 $x$의 근방에서 유한 개의 단면으로
생성됨을 보이겠다. $x$가 $X_i \to X$의 상에 속한다고 하고
$X' = X_i$라 쓰자. 생성 단면들을
$s_1, \ldots, s_n \in
\Gamma(X', \epsilon^*\mathcal{F}|_{X'_\etale})$라 하자. 등식
$$
\epsilon^*\mathcal{F} =
\epsilon^{-1}\mathcal{F} \otimes_{\epsilon^{-1}\mathcal{O}_X}
\mathcal{O}_\etale
$$
에 의해 \`etale 사상 $X'' \to X'$을 찾을 수 있어서, $x$가
$X' \to X$의 상에 속하며% P09-ERR-0428
$s_i|_{X''} = \sum s_{ij} \otimes a_{ij}$가 성립한다. 여기서 단면들은
$s_{ij} \in \epsilon^{-1}\mathcal{F}(X'')$와
$a_{ij} \in \mathcal{O}_\etale(X'')$이다.
$U \subset X$를 $X'' \to X$의 상이라 하자. $f'' : X'' \to X$가
\`etale이므로 이는 열린 부분스킴이다(Morphisms, Lemma
\ref{morphisms-lemma-etale-open}). 필요하다면 $X''$을 더 줄여서,
$s_{ij}$들이 원소 $t_{ij} \in \mathcal{F}(U)$에서 온다고 가정할 수
있다. 이는 역상 함자 $\epsilon^{-1}$의 구성에서 따른다. 이제
$t_{ij}$들이 $\mathcal{F}|_U$를 생성한다고 주장하면 보조정리의
증명이 끝난다. 즉, 해당 사상
$\mathcal{O}_U^{\oplus N} \to \mathcal{F}|_U$는 $f''$에 의해
$X''$으로 당겼을 때 전사이다. $f'' : X'' \to U$가 스킴들의 전사
평탄 사상이므로, 이는
$\mathcal{O}_U^{\oplus N} \to \mathcal{F}|_U$가 전사임을 함의한다.
실제로 줄기들을 살펴보고
$\mathcal{O}_{U, f''(z)} \to \mathcal{O}_{X'', z}$가 모든
$z \in X''$에 대해 충실히 평탄함을 사용하면 된다.
\end{proof}

\noindent
위 상황에서 사이트의 사상 $\epsilon$은 평탄하므로 가군의 복합체에
대한 당김을 정의한다.

\begin{lemma}
\label{spaces-perfect-lemma-descend-pseudo-coherent}
$X$를 스킴이라 하자. $E$를 $D(\mathcal{O}_X)$의 대상이라 하자.
다음은 동치이다.
\begin{enumerate}
\item $E$는 $m$-유사연접이다.
\item $\epsilon^*E$는 $m$-유사연접이고, 이는 $X$의 작은 \`etale
사이트 위에서의 명제이다.
\end{enumerate}
여기서 $\epsilon$은 (\ref{spaces-perfect-equation-epsilon})에서와 같다.
\end{lemma}

\begin{proof}
(1) $\Rightarrow$ (2)는 일반적인 사실이다. Cohomology on Sites, Lemma
\ref{sites-cohomology-lemma-pseudo-coherent-pullback}을 보라.
$\epsilon^*E$가 $m$-유사연접이라고 가정하자. 이하에서는
$\epsilon^*$가 완전 함자이고 따라서
$$
\epsilon^*H^i(E) = H^i(\epsilon^*E).
$$
임을 따로 언급하지 않고 사용한다. $E$가 $m$-유사연접임을 보이기
위해 $X$ 위에서 국소적으로 작업할 수 있으므로, $X$가 준콤팩트라고
가정할 수 있다(예를 들어 아핀이라고 할 수 있다). $X$가 준콤팩트이므로
모든 \`etale 덮개 $\{U_i \to X\}$에는 유한 세분이 있다. 따라서
$\epsilon^*E$는 $D^{-}(\mathcal{O}_\etale)$의 대상이다. Cohomology on
Sites, Definition
\ref{sites-cohomology-definition-pseudo-coherent} 뒤의 설명을 보라.
Lemma \ref{spaces-perfect-lemma-epsilon-flat}에 의해 $E$는
$D^-(\mathcal{O}_X)$의 대상이다.

\medskip\noindent
$n \in \mathbf{Z}$를 $H^n(E)$가 영이 아닌 가장 큰 정수라 하자. 그러면
$n$은 $H^n(\epsilon^*E)$가 영이 아닌 가장 큰 정수이기도 하다.
$n - m$에 대한 귀납법으로 보조정리를 증명하겠다. $n < m$이면
보조정리는 분명히 참이다. $n \geq m$이면
$H^n(\epsilon^*E)$는 유한 $\mathcal{O}_\etale$-가군이다. Cohomology on
Sites, Lemma \ref{sites-cohomology-lemma-finite-cohomology}를 보라.
따라서 Lemma \ref{spaces-perfect-lemma-descend-finite-type}에 의해 $H^n(E)$는 유한
$\mathcal{O}_X$-가군이다. $X$를 열린 덮개의 원소들로 바꾼 뒤에는 전사
$\mathcal{O}_X^{\oplus t} \to H^n(E)$가 존재한다고 가정할 수 있다.
$X$ 위에서 국소적으로 이를 복합체의 사상
$\alpha : \mathcal{O}_X^{\oplus t}[-n] \to E$로 올릴 수 있다
(세부사항은 생략한다). 구별 삼각형
$$
\mathcal{O}_X^{\oplus t}[-n] \to E \to C \to \mathcal{O}_X^{\oplus t}[-n + 1]
$$
을 택하자. 그러면 $C$의 코호몰로지는 $\geq n$ 차수에서 영이다.
한편 복합체 $\epsilon^*C$는 $m$-유사연접이다. Cohomology on Sites,
Lemma \ref{sites-cohomology-lemma-cone-pseudo-coherent}를 보라. 따라서
귀납 가정에 의해 $C$는 $m$-유사연접이다. Cohomology on Sites, Lemma
\ref{sites-cohomology-lemma-cone-pseudo-coherent}를 한 번 더 적용하면
결론을 얻는다.
\end{proof}

\begin{lemma}
\label{spaces-perfect-lemma-descend-tor-amplitude}
$X$를 스킴이라 하자. $E$를 $D(\mathcal{O}_X)$의 대상이라 하자. 그러면
\begin{enumerate}
\item $E$가 $[a, b]$ 안의 tor 진폭을 가질 필요충분조건은
$\epsilon^*E$가 $[a, b]$ 안의 tor 진폭을 갖는 것이다.
\item $E$가 유한 tor 차원을 가질 필요충분조건은 $\epsilon^*E$가
유한 tor 차원을 갖는 것이다.
\end{enumerate}
여기서 $\epsilon$은 (\ref{spaces-perfect-equation-epsilon})에서와 같다.
\end{lemma}

\begin{proof}
쉬운 함의는 Cohomology on Sites, Lemma
\ref{sites-cohomology-lemma-tor-amplitude-pullback}에서 따른다. 역으로
$\epsilon^*E$가 $[a, b]$ 안의 tor 진폭을 갖는다고 가정하자.
$\mathcal{F}$를 $\mathcal{O}_X$-가군이라 하자. $\epsilon$은 환 달린
사이트의 평탄 사상이므로(Lemma \ref{spaces-perfect-lemma-epsilon-flat}),
$$
\epsilon^*(E \otimes^\mathbf{L}_{\mathcal{O}_X} \mathcal{F})
=
\epsilon^*E
\otimes^\mathbf{L}_{\mathcal{O}_\etale}
\epsilon^*\mathcal{F}
$$
이다. 따라서 오른쪽 항의 코호몰로지 층이 소멸한다는 가정은 Lemma
\ref{spaces-perfect-lemma-epsilon-flat}에 의해
$E \otimes^\mathbf{L}_{\mathcal{O}_X} \mathcal{F}$의 코호몰로지 층에
대한 원하는 소멸을 함의한다.
\end{proof}

\begin{lemma}
\label{spaces-perfect-lemma-tor-dimension-rel}
$f : X \to Y$를 스킴의 사상이라 하자. $E$를
$D(\mathcal{O}_X)$의 대상이라 하자. 그러면
\begin{enumerate}
\item $E$가 $D(f^{-1}\mathcal{O}_Y)$의 대상으로서 $[a, b]$ 안의 tor
진폭을 가질 필요충분조건은 $\epsilon^*E$가 $[a, b]$ 안의 tor 진폭을
$D(f_{small}^{-1}\mathcal{O}_{Y_\etale})$의 대상으로서 갖는 것이다.
\item $E$가 국소적으로 $D(f^{-1}\mathcal{O}_Y)$의 대상으로서 유한 tor
차원을 가질 필요충분조건은 $\epsilon^*E$가 국소적으로
$D(f_{small}^{-1}\mathcal{O}_{Y_\etale})$의 대상으로서 유한 tor 차원을
갖는 것이다.
\end{enumerate}
여기서 $\epsilon$은 (\ref{spaces-perfect-equation-epsilon})에서와 같다.
\end{lemma}

\begin{proof}
(1)의 쉬운 방향은 Cohomology on Sites, Lemma
\ref{sites-cohomology-lemma-tor-amplitude-pullback}에서 따른다.
$x \in X$를 한 점이라 하고, $\overline{x}$를 $x$ 위에 놓인 기하적
점이라 하자. $y = f(x)$라 두고, $\overline{y}$를 $Y$의 기하적
점이라 하되 $\overline{x}$에서 오는 점으로 택하자. 그러면
$(f^{-1}\mathcal{O}_Y)_x = \mathcal{O}_{Y, y}$이고(Sheaves, Lemma
\ref{sheaves-lemma-stalk-pullback}),
$$
(f_{small}^{-1}\mathcal{O}_{Y_\etale})_{\overline{x}} =
\mathcal{O}_{Y_\etale, \overline{y}} =
\mathcal{O}_{Y, y}^{sh}
$$
는 엄밀 Hensel화이다(\`Etale Cohomology, Lemmas
\ref{etale-cohomology-lemma-stalk-pullback}와
\ref{etale-cohomology-lemma-describe-etale-local-ring}에 의해).
$\mathcal{O}_{X_\etale}$의 $X$에서의 줄기는% P09-ERR-0429
$\mathcal{O}_{X, x}^{sh}$이므로, 당김의 추이성에 의해
$$
(\epsilon^*E)_{\overline{x}} =
E_x \otimes_{\mathcal{O}_{X, x}}^\mathbf{L} \mathcal{O}_{X, x}^{sh}
$$
를 얻는다. $\epsilon^*E$가 $[a, b]$ 안의 tor 진폭을 가지며 이를
$f_{small}^{-1}\mathcal{O}_{Y_\etale}$-가군의 복합체로 보면,
$(\epsilon^*E)_{\overline{x}}$는 $[a, b]$ 안의 tor 진폭을 가지며, 이를
$\mathcal{O}_{Y, y}^{sh}$-가군의 복합체로 보아도 그러하다(줄기를 취하는
것은 당김이고 위에서 인용한 보조정리를 사용한다). More on Flatness, Lemma
\ref{flat-lemma-tor-amplitude-up-down-henselization}에 의해
$(\epsilon^*E)_{\overline{x}}$를 $\mathcal{O}_{Y, y}$-가군의 복합체로
보아도 tor 진폭은 $[a, b]$ 안에 있다.
$\mathcal{O}_{X, x} \to \mathcal{O}_{X, x}^{sh}$는 충실히 평탄하고(More
on Algebra, Lemma
\ref{more-algebra-lemma-dumb-properties-henselization}), 또한
$(\epsilon^*E)_{\overline{x}} =
E_x \otimes_{\mathcal{O}_{X, x}}^\mathbf{L} \mathcal{O}_{X, x}^{sh}$이므로,
More on Algebra, Lemma
\ref{more-algebra-lemma-no-change-tor-amplitude}를 적용하면 $E_x$를
$\mathcal{O}_{Y, y}$-가군의 복합체로 본 tor 진폭이 $[a, b]$ 안에
있음을 얻는다. Cohomology, Lemma
\ref{cohomology-lemma-tor-amplitude-stalk}에 의해 $E$를
$D(f^{-1}\mathcal{O}_Y)$의 대상으로 본 tor 진폭은 $[a, b]$ 안에 있다.
이것이 (1)의 역함의를 준다. (2)는 (1)에서 형식적으로 따른다.
\end{proof}

\begin{lemma}
\label{spaces-perfect-lemma-descend-perfect}
$X$를 스킴이라 하자. $E$를 $D(\mathcal{O}_X)$의 대상이라 하자.
$E$가 $D(\mathcal{O}_X)$의 완전 대상일 필요충분조건은
$\epsilon^*E$가 $D(\mathcal{O}_\etale)$의 완전 대상인 것이다.
여기서 $\epsilon$은 (\ref{spaces-perfect-equation-epsilon})에서와 같다.
\end{lemma}

\begin{proof}
쉬운 함의는 Cohomology on Sites, Lemma
\ref{sites-cohomology-lemma-perfect-pullback}에 포함된 일반 결과에서
따른다. 역으로 Cohomology on Sites, Lemma
\ref{sites-cohomology-lemma-perfect}의 동치와, 유사연접 복합체 및 유한
tor 차원 복합체에 대한 해당 결과인 Lemmas
\ref{spaces-perfect-lemma-descend-pseudo-coherent}와
\ref{spaces-perfect-lemma-descend-tor-amplitude}를 사용할 수 있다. 일부 세부사항은
생략한다.
\end{proof}

\begin{lemma}
\label{spaces-perfect-lemma-pseudo-coherent}
$S$를 스킴이라 하자. $X$를 $S$ 위의 대수공간이라 하자. $E$가
$m$-유사연접인 $D(\mathcal{O}_X)$의 대상이면,
$H^i(E)$는 준연접 $\mathcal{O}_X$-가군이다. 여기서 $i > m$이다.
$E$가 유사연접이면
$E$는 $D_\QCoh(\mathcal{O}_X)$의 대상이다.
\end{lemma}

\begin{proof}
국소적으로 $H^i(E)$는 $H^i(\mathcal{E}^\bullet)$과 동형이고,
$\mathcal{E}^\bullet$은 엄밀 완전 복합체이다. 층
$\mathcal{E}^i$들은 유한 자유 가군의 직합 성분이므로 준연접이다.
보조정리가 따른다.
\end{proof}

\begin{lemma}
\label{spaces-perfect-lemma-identify-pseudo-coherent-noetherian}
$S$를 스킴이라 하자. $X$를 $S$ 위의 Noether 대수공간이라 하자.
$E$를 $D_\QCoh(\mathcal{O}_X)$의 대상이라 하자. $m \in \mathbf{Z}$에
대해 다음은 동치이다.
\begin{enumerate}
\item $H^i(E)$는 $i \geq m$일 때 연접이고 $i \gg 0$일 때 영이다.
\item $E$는 $m$-유사연접이다.
\end{enumerate}
특히 $E$가 유사연접일 필요충분조건은 $E$가
$D^-_{\textit{Coh}}(\mathcal{O}_X)$의 대상인 것이다.
\end{lemma}

\begin{proof}
$X$가 준콤팩트이므로 아핀 스킴 $U$와 전사 \`etale 사상
$U \to X$를 찾을 수 있다(Properties of Spaces, Lemma
\ref{spaces-properties-lemma-quasi-compact-affine-cover}). $U$가
Noether임을 관찰하라. $E$가 $m$-유사연접일 필요충분조건은
$E|_U$가 $m$-유사연접인 것이다(정의 또는 Cohomology on Sites,
Lemma
\ref{sites-cohomology-lemma-pseudo-coherent-independent-representative}에서
따른다). 마찬가지로 $H^i(E)$가 연접일 필요충분조건은
$H^i(E)|_U = H^i(E|_U)$가 연접인 것이다(Cohomology of Spaces, Lemma
\ref{spaces-cohomology-lemma-coherent-Noetherian}를 보라). 따라서 $X$가
표현 가능하다고 가정할 수 있다.

\medskip\noindent
$X$가 스킴 $X_0$로 표현되면(Lemma
\ref{spaces-perfect-lemma-derived-quasi-coherent-small-etale-site}),
$E = \epsilon^*E_0$로 쓸 수 있다. 여기서 $E_0$는
$D_\QCoh(\mathcal{O}_{X_0})$의 대상이고,
$\epsilon : X_\etale \to (X_0)_{Zar}$은 (\ref{spaces-perfect-equation-epsilon})에서와
같다. 이 경우 Lemma \ref{spaces-perfect-lemma-descend-pseudo-coherent}에 의해
$E$가 $m$-유사연접일 필요충분조건은 $E_0$가 그러한 것이다.
마찬가지로 Lemma \ref{spaces-perfect-lemma-descend-finite-type}에 의해
$H^i(E_0)$가 유한형(즉 연접)일 필요충분조건은 $H^i(E)$가 그러한
것이다. 마지막으로 Lemma \ref{spaces-perfect-lemma-epsilon-flat}에 의해
$H^i(E_0) = 0$일 필요충분조건은 $H^i(E) = 0$인 것이다. 따라서
스킴의 경우로 환원되며, 이는 Derived Categories of Schemes, Lemma
\ref{perfect-lemma-identify-pseudo-coherent-noetherian}이다.
\end{proof}

\begin{lemma}
\label{spaces-perfect-lemma-tor-qc-qs}
$S$를 스킴이라 하자. $X$를 $S$ 위의 준분리 대수공간이라 하자.
$E$를 $D_\QCoh(\mathcal{O}_X)$의 대상이라 하자. $a \leq b$라 하자.
다음은 동치이다.
\begin{enumerate}
\item $E$는 $[a, b]$ 안의 tor 진폭을 갖는다.
\item 모든 $\mathcal{F}$가 $\QCoh(\mathcal{O}_X)$에 속할 때
$H^i(E \otimes_{\mathcal{O}_X}^\mathbf{L} \mathcal{F}) = 0$이고, 여기서
$i \not \in [a, b]$이다.
\end{enumerate}
\end{lemma}

\begin{proof}
(1)이 (2)를 함의함은 분명하다. (2)를 가정하자. $j : U \to X$를
\`etale 사상이라 하고 $U$가 아핀이라고 하자. $X$가 준분리이므로
$j$는 준콤팩트이고 분리이다. 따라서 $j_*$는 준연접 가군을 준연접
가군으로 보낸다(Morphisms of Spaces, Lemma
\ref{spaces-morphisms-lemma-pushforward}). 임의의 준연접 가군
$\mathcal{G}$를 $U$ 위에서 택하자. 가정에 의해
$H^i(E \otimes_{\mathcal{O}_X}^\mathbf{L} j_*\mathcal{G})$는
$i \not \in [a, b]$일 때 소멸한다. 평탄 사상 $j$로 당기면
$H^i(E|_U \otimes_{\mathcal{O}_U}^\mathbf{L} j^*j_*\mathcal{G})$가
$i \not \in [a, b]$일 때 소멸함을 얻는다. Cohomology of Spaces, Lemma
\ref{spaces-cohomology-lemma-etale-pull-push-split}에 의해
$\mathcal{G}$는 $j^*j_*\mathcal{G}$의 직합 성분이다. 따라서 조건
(2)는 $H^i(E|_U \otimes_{\mathcal{O}_U}^\mathbf{L} \mathcal{G})$가
$i \not \in [a, b]$일 때 소멸함을 함의하며, 이는 모든 준연접
$\mathcal{O}_U$-가군 $\mathcal{G}$에 대해 성립한다. $E|_U$가 $[a, b]$ 안의
tor 진폭을 가짐을 증명하면 충분하므로, $X$가 표현 가능한 경우로
환원된다.

\medskip\noindent
$X$가 스킴 $X_0$로 표현되면(Lemma
\ref{spaces-perfect-lemma-derived-quasi-coherent-small-etale-site}),
$E = \epsilon^*E_0$로 쓸 수 있다. 여기서 $E_0$는
$D_\QCoh(\mathcal{O}_{X_0})$의 대상이고,
$\epsilon : X_\etale \to (X_0)_{Zar}$은 (\ref{spaces-perfect-equation-epsilon})에서와
같다. 모든 준연접 가군 $\mathcal{F}_0$에 대해, 이를 $X_0$ 위에서 보면
$\epsilon^*\mathcal{F}_0$는 $X$ 위에서 준연접이고
$$
H^i(E \otimes_{\mathcal{O}_X}^\mathbf{L} \epsilon^*\mathcal{F}_0)
=
\epsilon^*H^i(E_0 \otimes_{\mathcal{O}_{X_0}}^\mathbf{L} \mathcal{F}_0)
$$
이다. 이는 $\epsilon$이 평탄하기 때문이다(Lemma
\ref{spaces-perfect-lemma-epsilon-flat}). 또한 이 층들이 $i \not \in [a, b]$일 때
소멸하면, 같은 보조정리에 의해
$H^i(E_0 \otimes_{\mathcal{O}_{X_0}}^\mathbf{L} \mathcal{F}_0)$도
그러하다. 따라서 문제를 스킴의 경우로 환원했으며, 이는 Derived
Categories of Schemes, Lemma \ref{perfect-lemma-tor-qc-qs}에서 다룬다.
\end{proof}

\begin{lemma}
\label{spaces-perfect-lemma-descend-RHom}
$X$를 스킴이라 하자. $E, F$를 $D(\mathcal{O}_X)$의 대상이라 하자.
다음 가운데 하나를 가정하자.
\begin{enumerate}
\item $E$는 유사연접이고 $F$는 $D^+(\mathcal{O}_X)$에 속한다.
\item $E$는 완전이고 $F$는 임의적이다.
\end{enumerate}
그러면 표준 동형
$$
\epsilon^*R\SheafHom(E, F) \longrightarrow R\SheafHom(\epsilon^*E, \epsilon^*F)
$$
이 존재한다. 여기서 $\epsilon$은 (\ref{spaces-perfect-equation-epsilon})에서와 같다.
\end{lemma}

\begin{proof}
$\epsilon$이 평탄하고(Lemma \ref{spaces-perfect-lemma-epsilon-flat}), 따라서
$\epsilon^* = L\epsilon^*$임을 상기하라. Cohomology on Sites, Remark
\ref{sites-cohomology-remark-prepare-fancy-base-change}에 의해 왼쪽에서
오른쪽으로 가는 표준 사상이 있다. 이것이 동형임을 보이기 위해
국소적으로 작업할 수 있으므로, $X$가 아핀 스킴이라고 가정할 수 있다.

\medskip\noindent
(1)의 경우 $E$를 위로 유계인 복합체 $\mathcal{E}^\bullet$로 나타낼 수
있고, 이 복합체는 유한 자유 $\mathcal{O}_X$-가군들로 이루어진다. Derived Categories
of Schemes, Lemma \ref{perfect-lemma-lift-pseudo-coherent}를 보라. 또한
$F$를 아래로 유계인 복합체 $\mathcal{F}^\bullet$로 나타낼 수 있고,
이 복합체는 $\mathcal{O}_X$-가군들로 이루어진다. Cohomology, Lemma
\ref{cohomology-lemma-Rhom-complex-of-direct-summands-finite-free}를
적용하면 $R\SheafHom(E, F)$는 다음을 항으로 갖는 복합체로 나타내어진다.
$$
\bigoplus\nolimits_{n = - p + q}
\SheafHom_{\mathcal{O}_X}(\mathcal{E}^p, \mathcal{F}^q)
$$
Cohomology on Sites, Lemma
\ref{sites-cohomology-lemma-Rhom-complex-of-direct-summands-finite-free}를
적용하면 $R\SheafHom(\epsilon^*E, \epsilon^*F)$는 다음을 항으로 갖는
복합체로 나타내어진다.
$$
\bigoplus\nolimits_{n = - p + q}
\SheafHom_{\mathcal{O}_\etale}
(\epsilon^*\mathcal{E}^p, \epsilon^*\mathcal{F}^q)
$$
따라서 보조정리의 명제는 표준 사상
$$
\epsilon^*\SheafHom_{\mathcal{O}_X}(\mathcal{E}, \mathcal{F})
\longrightarrow
\SheafHom_{\mathcal{O}_\etale}
(\epsilon^*\mathcal{E}, \epsilon^*\mathcal{F})
$$
이 임의의 $\mathcal{O}_X$-가군 $\mathcal{F}$와 유한 자유
$\mathcal{O}_X$-가군 $\mathcal{E}$에 대해 동형이라는 참인 사실로
귀착된다.

\medskip\noindent
(2)의 경우 $E$를 엄밀 완전 복합체 $\mathcal{E}^\bullet$로 나타낼 수
있고, 이 복합체는 $\mathcal{O}_X$-가군들로 이루어진다. Derived Categories of Schemes,
Lemmas \ref{perfect-lemma-affine-compare-bounded}와
\ref{perfect-lemma-perfect-affine}, 그리고 가군의 완전 복합체가 유한
생성 사영 가군들로 이루어진 유한 복합체로 나타내어진다는 사실을
사용하라. 그러면 위와 정확히 같은 증명을 할 수 있는데, Cohomology,
Lemma \ref{cohomology-lemma-Rhom-complex-of-direct-summands-finite-free}의
인용을 Cohomology, Lemma
\ref{cohomology-lemma-Rhom-strictly-perfect}의 인용으로 바꾸면 된다.
\end{proof}

\begin{lemma}
\label{spaces-perfect-lemma-quasi-coherence-internal-hom}
$S$를 스킴이라 하자. $X$를 $S$ 위의 대수공간이라 하자. $L, K$를
$D(\mathcal{O}_X)$의 대상이라 하자. 다음 가운데 하나가 성립한다고
하자.
\begin{enumerate}
\item $L$은 $D^+_\QCoh(\mathcal{O}_X)$에 속하고 $K$는 유사연접이다.% P09-ERR-0430
\item $L$은 $D_\QCoh(\mathcal{O}_X)$에 속하고 $K$는 완전이다.% P09-ERR-0430
\end{enumerate}
그러면 $R\SheafHom(K, L)$은 $D_\QCoh(\mathcal{O}_X)$에 속한다.
\end{lemma}

\begin{proof}
이는 스킴에 대한 유사한 명제(Derived Categories of Schemes, Lemma
\ref{perfect-lemma-quasi-coherence-internal-hom})에서 따른다. 여기에는
Lemma \ref{spaces-perfect-lemma-check-quasi-coherence-on-covering}의 판정법, Lemmas
\ref{spaces-perfect-lemma-descend-pseudo-coherent}와 \ref{spaces-perfect-lemma-descend-perfect}의 판정법,
그리고 Lemma \ref{spaces-perfect-lemma-descend-RHom}의 결과를 사용한다.
\end{proof}

\begin{lemma}
\label{spaces-perfect-lemma-internal-hom-evaluate-tensor-isomorphism}
$S$를 스킴이라 하자. $X$를 $S$ 위의 대수공간이라 하자. $K, L, M$을
$D_\QCoh(\mathcal{O}_X)$의 대상이라 하자. 사상
$$
K \otimes_{\mathcal{O}_X}^\mathbf{L} R\SheafHom(M, L)
\longrightarrow
R\SheafHom(M, K \otimes_{\mathcal{O}_X}^\mathbf{L} L)
$$
은 Cohomology on Sites, Lemma
\ref{sites-cohomology-lemma-internal-hom-diagonal-better}에서 주어지며,
다음 경우에 동형이다.
\begin{enumerate}
\item $M$이 완전인 경우,
\item $K$가 완전인 경우,
\item $M$이 유사연접이고 $L \in D^+(\mathcal{O}_X)$이며 $K$가 유한
tor 차원을 갖는 경우.
\end{enumerate}
\end{lemma}

\begin{proof}
사상이 동형인지 여부는 \`etale 국소적으로 확인할 수 있으므로 $X$가
아핀 스킴이라고 가정할 수 있다. 그러면
$K, L, M$을 $\epsilon^*K_0, \epsilon^*L_0, \epsilon^*M_0$로 쓸 수
있다. 여기서 어떤 $K_0, L_0, M_0$는
$D_\QCoh(\mathcal{O}_X)$에 속한다. 이는 Lemma
\ref{spaces-perfect-lemma-derived-quasi-coherent-small-etale-site}에서 따른다. 그러면
Lemma \ref{spaces-perfect-lemma-descend-RHom}이 (1)과 (3)을 스킴의
경우로 환원하며, 이는 Derived Categories of Schemes, Lemma
\ref{perfect-lemma-internal-hom-evaluate-tensor-isomorphism}이다.
$K$가 완전이지만 다른 가정이 없으면 화살표의 어느 쪽도
$D_\QCoh(\mathcal{O}_X)$에 속한다는 것을 알 수 없다. 그러나 더
국소화한 뒤 $K$가 유한 자유 가군들로 이루어진 유한 복합체로
나타내어질 것이므로 결과는 여전히 참이며, 이 경우에는 분명하다.
\end{proof}










\section{완전 복합체에 의한 근사}
\label{spaces-perfect-section-approximation}

\noindent
이 절에서는 Derived Categories of Schemes, Section
\ref{perfect-section-approximation}에서 시작한 논의를 계속한다.

\begin{definition}
\label{spaces-perfect-definition-approximation-holds}
$S$를 스킴이라 하자. $X$를 $S$ 위의 대수공간이라 하자. 다음 조건의
삼중항 $(T, E, m)$을 생각하자.
\begin{enumerate}
\item $T \subset |X|$는 닫힌 부분집합이다.
\item $E$는 $D_\QCoh(\mathcal{O}_X)$의 대상이다.
\item $m \in \mathbf{Z}$이다.
\end{enumerate}
삼중항 $(T, E, m)$에 대해 {\it 근사가 성립한다}고 하는 것은 완전
대상 $P$가 $D(\mathcal{O}_X)$에 속하고 $T$ 위에 지지되며, 사상
$\alpha : P \to E$가 존재하여 동형 $H^i(P) \to H^i(E)$를
$i > m$일 때 유도하고 전사 $H^m(P) \to H^m(E)$를 유도할 때이다.
\end{definition}

\noindent
모든 삼중항에 대해 근사가 성립할 수는 없다. 그 이유는 Derived
Categories of Schemes, Definition
\ref{perfect-definition-approximation-holds} 뒤의 설명을 읽어 보라.

\begin{definition}
\label{spaces-perfect-definition-approximation}
$S$를 스킴이라 하자. $X$를 $S$ 위의 대수공간이라 하자. $X$ 위에서
{\it 완전 복합체에 의한 근사가 성립한다}고 하는 것은, 임의의 닫힌
부분집합 $T \subset |X|$에 대해 사상 $X \setminus T \to X$가
준콤팩트이면 다음 성질의 정수 $r$이 존재한다고 하자. 즉, Definition
\ref{spaces-perfect-definition-approximation-holds}에서와 같은 모든 삼중항
$(T, E, m)$ 가운데 다음을 만족하는 것에 대해
\begin{enumerate}
\item $E$는 $(m - r)$-유사연접이다.
\item $H^i(E)$는 $T$ 위에 지지되며, 여기서 $i \geq m - r$이다.
\end{enumerate}
근사가 성립하는 것이다.
\end{definition}

\begin{lemma}
\label{spaces-perfect-lemma-pushforward-perfect}
$S$를 스킴이라 하자. $(U \subset X, j : V \to X)$를 $S$ 위 대수공간의
기본 구별 정사각형이라 하자. $E$를 $D(\mathcal{O}_V)$의 완전 대상으로
하되, $j^{-1}(T)$ 위에 지지된다고 하자. 여기서
$T = |X| \setminus |U|$이다. 그러면 $Rj_*E$는
$D(\mathcal{O}_X)$의 완전 대상이다.
\end{lemma}

\begin{proof}
완전성은 $X_\etale$ 위에서 국소적인 성질이다. 따라서
$Rj_*E$를 $U$와 $V$에 제한했을 때 완전임을 확인하면 충분하다.
Lemma \ref{spaces-perfect-lemma-pushforward-with-support-in-open}에 의해
$Rj_*E|_V = E$이고 이는 완전이다. 또한 $Rj_*E|_U = 0$인데, 이는
$E|_{V \setminus j^{-1}(T)} = 0$이기 때문이다(Lemma
\ref{spaces-perfect-lemma-restrict-direct-image-open}을 사용하라).
\end{proof}

\begin{lemma}
\label{spaces-perfect-lemma-open}
$S$를 스킴이라 하자. $(U \subset X, j : V \to X)$를 $S$ 위
대수공간들의 기본 구별 정사각형이라 하자. $T$를
$|X| \setminus |U|$의 닫힌 부분집합이라 하고, $(T, E, m)$을 Definition
\ref{spaces-perfect-definition-approximation-holds}에서와 같은 삼중항이라 하자.
다음을 가정하자.
\begin{enumerate}
\item $(j^{-1}T, E|_V, m)$에 대해 근사가 성립한다.
\item $H^i(E)$는 $i \geq m$일 때 $T$ 위에 지지된다.
\end{enumerate}
그러면 $(T, E, m)$에 대해 근사가 성립한다.
\end{lemma}

\begin{proof}
$P \to E|_V$를 삼중항 $(j^{-1}T, E|_V, m)$의 근사라 하자. 이 근사는
$V$ 위에 있다.
그러면 Lemma \ref{spaces-perfect-lemma-pushforward-perfect}에 의해 $Rj_*P$는
$D(\mathcal{O}_X)$의 완전 대상이다. 한편 Lemma
\ref{spaces-perfect-lemma-pushforward-with-support-in-open}에 의해 $Rj_*P = j_!P$이다.
예를 들어 (\ref{spaces-perfect-equation-stalk-j-shriek})에 의해 $j_!P$는 $T$ 위에
지지된다. 따라서 근사
$Rj_*P = j_!P \to j_!(E|_V) \to E$를 얻는다.
\end{proof}

\begin{lemma}
\label{spaces-perfect-lemma-approximation-affine}
$S$를 스킴이라 하자. $X$를 아핀 스킴으로 표현되는 $S$ 위의
대수공간이라 하자. Definition \ref{spaces-perfect-definition-approximation-holds}에서와
같은 모든 삼중항 $(T, E, m)$ 가운데 다음 조건을 만족하는 정수
$r \geq 0$가 존재하는 것에 대해 근사가 성립한다.
\begin{enumerate}
\item $E$는 $m$-유사연접이다.
\item $H^i(E)$는 $T$ 위에 지지되며, 여기서 $i \geq m - r + 1$이다.
\item $X \setminus T$는 $r$개의 아핀 열린집합의 합집합이다.
\end{enumerate}
특히 아핀 스킴에 대해서는 완전 복합체에 의한 근사가 성립한다.
\end{lemma}

\begin{proof}
$X_0$를 $X$를 표현하는 아핀 스킴이라 하자. $T_0 \subset X_0$를
$T$에 대응하는 닫힌 부분집합이라 하자.% P09-ERR-0431
$\epsilon : X_\etale \to X_{0, Zar}$을 (\ref{spaces-perfect-equation-epsilon})의
사상이라 하자. $E = \epsilon^*E_0$로 쓸 수 있다. 여기서 어떤 대상
$E_0$는 $D_\QCoh(\mathcal{O}_{X_0})$에 속한다. Lemma
\ref{spaces-perfect-lemma-derived-quasi-coherent-small-etale-site}를 보라. 그러면
$E_0$는 $m$-유사연접이다(Lemma
\ref{spaces-perfect-lemma-descend-pseudo-coherent}를 보라). 코호몰로지 층의 줄기를
비교하면(Lemma \ref{spaces-perfect-lemma-epsilon-flat}의 증명을 보라),
$H^i(E_0)$는 $T_0$ 위에 지지되며 여기서 $i \geq m - r + 1$이다.
Derived Categories of Schemes, Lemma
\ref{perfect-lemma-approximation-affine}에 의해 근사 $P_0 \to E_0$가
존재하며, 이는 $(T_0, E_0, m)$의 근사이다. Lemma
\ref{spaces-perfect-lemma-descend-perfect}에 의해
$P = \epsilon^*P_0$는 $D(\mathcal{O}_X)$의 완전 대상이다. 당기면
원하는 근사 $P = \epsilon^*P_0 \to \epsilon^*E_0 = E$를 얻는다.
\end{proof}

\begin{lemma}
\label{spaces-perfect-lemma-induction-step}
$S$를 스킴이라 하자. $(U \subset X, j : V \to X)$를 $S$ 위
대수공간들의 기본 구별 정사각형이라 하자. $U$는 준콤팩트이고 $V$는
아핀이며 $U \times_X V$는 준콤팩트이라고 가정하자. 완전 복합체에
의한 근사가 $U$ 위에서 성립하면, 완전 복합체에 의한 근사는 $X$
위에서도 성립한다.
\end{lemma}

\begin{proof}
$T \subset |X|$를 닫힌 부분집합이라 하고 $X \setminus T \to X$가
준콤팩트이라고 하자. $r_U$를 Definition \ref{spaces-perfect-definition-approximation}의
정수로서 쌍 $(U, T \cap |U|)$에 맞는 것으로 택하자.
$T' = T \setminus |U|$라 두고 $T'$에 유도된 기약 부분공간 구조를
주자. $|T'|$가 $|X| \setminus |U|$에 포함되므로
$j^{-1}(T') \to T'$는 동형이다. 또한
$V \setminus j^{-1}(T')$는 준콤팩트이다. 실제로 이것은
$U \times_X V$와 $X \setminus T$의 $X$ 위 올곱이고, 우리는
$U \times_X V$가 준콤팩트이며 $X \setminus T \to X$가
준콤팩트이라고 가정했다. $r'$을 $V \setminus j^{-1}(T')$를 덮는 데
필요한 아핀들의 개수라 하자. $r = \max(r_U, r')$가 쌍 $(X, T)$에
대해 적합하다고 주장한다.

\medskip\noindent
이를 보기 위해 삼중항 $(T, E, m)$을 택하되, $E$가
$(m - r)$-유사연접이고 $H^i(E)$가 $T$ 위에 지지되며
$i \geq m - r$이라고 하자. $t$를 $H^t(E)|_U$가 영이 아닌 가장 큰
정수라 하자. 이러한 정수는 $U$가 준콤팩트이고 $E|_U$가
$(m - r)$-유사연접이므로 존재한다. $E$를 근사할 수 있음을 $t$에
대한 귀납법으로 증명하겠다.

\medskip\noindent
기초 단계: $t \leq m - r'$라 하자. 이는 $H^i(E)$가 $T'$ 위에
지지되며 $i \geq m - r'$임을 뜻한다. 따라서 Lemma
\ref{spaces-perfect-lemma-approximation-affine}는 근사 $P \to E|_V$가 존재함을
보장한다. 이는 삼중항 $(T', E|_V, m)$의 근사이며 $V$ 위에 있다. Lemma
\ref{spaces-perfect-lemma-open}을 적용하면 $(T', E, m)$을 근사할 수 있다. 그러한
근사는 $(T, E, m)$의 근사이기도 하다.

\medskip\noindent
귀납 단계. 근사 $P \to E|_U$를
$(T \cap |U|, E|_U, m)$의 근사로 택하자. 이는 특히 전사
$H^t(P) \to H^t(E|_U)$를 준다. 이하의 증명에서는 표현 가능한
대수공간 $V$와 $U \times_X V$에 대해 Lemma
\ref{spaces-perfect-lemma-derived-quasi-coherent-small-etale-site}의 동치와 Remark
\ref{spaces-perfect-remark-match-total-direct-images}의 호환성을 사용하겠다. 또한
Zariski 사이트와 \'etale 사이트 위에서 $(m - r)$-유사연접성 및 각각
완전성이 일치한다는 사실을 사용하겠다. Lemmas
\ref{spaces-perfect-lemma-descend-pseudo-coherent}와 \ref{spaces-perfect-lemma-descend-perfect}를 보라.
따라서 열린 매장 $U \times_X V \subset V$에 대해 Derived Categories
of Schemes, Section \ref{perfect-section-lift}의 결과를 사용할 수 있다.
이 방법으로 Derived Categories of Schemes, Lemma
\ref{perfect-lemma-lift-perfect-complex-plus-shift-support}는 완전 대상
$Q$가 $D(\mathcal{O}_V)$에 속하고 $j^{-1}(T)$ 위에 지지되며, 동형
$Q|_{U \times_X V} \to (P \oplus P[1])|_{U \times_X V}$가 존재함을
함의한다. Derived Categories of Schemes, Lemma
\ref{perfect-lemma-lift-map}에 의해 $Q$를
$Q \otimes^\mathbf{L} I$로 바꾸어 사상
$$
Q|_{U \times_X V} \longrightarrow
(P \oplus P[1])|_{U \times_X V} \longrightarrow
P|_{U \times_X V} \longrightarrow E|_{U \times_X V}
$$
이 $Q \to E|_V$로 올라간다고 가정할 수 있다. Lemma
\ref{spaces-perfect-lemma-glue}에 의해 사상 $a : R \to E$를
$D(\mathcal{O}_X)$ 안에서 찾을 수 있어서,
$a|_U$는 $P \oplus P[1] \to E|_U$와 동형이고
$a|_V$는 $Q \to E|_V$와 동형이다.% P09-ERR-0432
따라서 $R$은 완전이고 $T$ 위에 지지되며, 사상
$H^t(R) \to H^t(E)$는 $U$로 제한하면 전사이다. 구별 삼각형
$$
R \to E \to E' \to R[1]
$$
을 택하자. 그러면 $E'$은 $(m - r)$-유사연접이고(Cohomology on Sites,
Lemma \ref{sites-cohomology-lemma-cone-pseudo-coherent}),
$H^i(E')|_U = 0$은 $i \geq t$일 때 성립하며, $H^i(E')$는 $T$ 위에
지지되고 여기서 $i \geq m - r$이다. 귀납 가정에 의해 근사
$R' \to E'$을 찾을 수 있고, 이는 $(T, E', m)$의 근사이다. 합성
$R' \to E' \to R[1]$을 구별 삼각형
$R \to R'' \to R' \to R[1]$에 넣고, 사상 $R' \to E'$과
$R[1] \to R[1]$을 TR3을 사용하여 다음 구별 삼각형의 사상으로
확장하자.
$$
\xymatrix{
R \ar[r] \ar[d] & R'' \ar[d] \ar[r] & R' \ar[d] \ar[r] & R[1] \ar[d] \\
R \ar[r] &  E \ar[r] & E' \ar[r] & R[1]
}
$$
그러면 $R''$은 완전 복합체이고(Cohomology on Sites, Lemma
\ref{sites-cohomology-lemma-two-out-of-three-perfect}) $T$ 위에
지지된다. 쉬운 도식 추적을 하면 $R'' \to E$가 원하는 근사임을 알 수
있다.
\end{proof}

\begin{theorem}
\label{spaces-perfect-theorem-approximation}
$S$를 스킴이라 하자. $X$를 $S$ 위의 준콤팩트 준분리 대수공간이라
하자. 그러면 $X$ 위에서 완전 복합체에 의한 근사가 성립한다.
\end{theorem}

\begin{proof}
이는 Lemma \ref{spaces-perfect-lemma-induction-principle}의 귀납 원리와 Lemmas
\ref{spaces-perfect-lemma-induction-step} 및 \ref{spaces-perfect-lemma-approximation-affine}에서 따른다.
\end{proof}






\section{유도 범주의 생성}
\label{spaces-perfect-section-generating}

\noindent
이 절은 Derived Categories of Schemes, Section
\ref{perfect-section-generating}에 해당한다. 그러나 먼저 Derived
Categories of Schemes, Lemma
\ref{perfect-lemma-lift-map-from-perfect-complex-with-support}에 해당하는
다음 보조정리를 증명한다.

\begin{lemma}
\label{spaces-perfect-lemma-lift-map-from-perfect-complex-with-support}
$S$를 스킴이라 하자. $X$를 $S$ 위의 준콤팩트 준분리 대수공간이라
하자. $W \subset X$를 준콤팩트 열린집합이라 하자. $T \subset |X|$를
닫힌 부분집합이라 하고 $X \setminus T \to X$가 준콤팩트 사상이라고
하자. $E$를 $D_\QCoh(\mathcal{O}_X)$의 대상이라 하자.
$\alpha : P \to E|_W$를 사상이라 하되, $P$가
$D(\mathcal{O}_W)$의 완전 대상이고 $T \cap W$ 위에 지지된다고 하자.
그러면 사상 $\beta : R \to E$가 존재하여, $R$은
$D(\mathcal{O}_X)$의 완전 대상이고 $T$ 위에 지지되며, $P$는
$R|_W$의 직합 성분이고 이는 $D(\mathcal{O}_W)$ 안에서의 명제이며 이 포함은
$\alpha$ 및 $\beta|_W$와 호환된다.% P09-ERR-0433
\end{lemma}

\begin{proof}
Lemma \ref{spaces-perfect-lemma-induction-principle-enlarge}의 귀납 원리를 사용해 이를
증명하겠다. 따라서 기본 구별 정사각형
$(W \subset X, f : V \to X)$이 있고, $V$가 아핀이며
$P \to E|_W$가 보조정리의 명제에서와 같은 경우로 즉시 환원된다.
이하의 증명에서는 표현 가능한 대수공간 $V$와 $W \times_X V$에 대해
Lemma \ref{spaces-perfect-lemma-derived-quasi-coherent-small-etale-site}와 Remark
\ref{spaces-perfect-remark-match-total-direct-images}의 호환성을 사용하겠다. 또한
Zariski 사이트와 \'etale 사이트 위의 완전성이 일치한다는 사실을
사용한다. Lemma \ref{spaces-perfect-lemma-descend-perfect}를 보라.

\medskip\noindent
Derived Categories of Schemes, Lemma
\ref{perfect-lemma-lift-perfect-complex-plus-shift-support}에 의해 완전 대상
$Q$를 $D(\mathcal{O}_V)$ 안에서 택할 수 있어서, 이는 $f^{-1}T$ 위에
지지되고 동형
$Q|_{W \times_X V} \to (P \oplus P[1])|_{W \times_X V}$가 존재한다.
Derived Categories of Schemes, Lemma \ref{perfect-lemma-lift-map}에 의해
$Q$를 $Q \otimes^\mathbf{L} I$로 바꾸어도 되며, 여전히 $f^{-1}T$ 위에
지지된다. 그리고 사상
$$
Q|_{W \times_X V} \to (P \oplus P[1])|_{W \times V}
\longrightarrow P|_{W \times_X V}
\longrightarrow
E|_{W \times_X V}
$$
% P09-ERR-0434
이 $Q \to E|_V$로 올라간다고 가정할 수 있다. Lemma
\ref{spaces-perfect-lemma-glue}에 의해 사상 $a : R \to E$를
$D(\mathcal{O}_X)$ 안에서 찾을 수 있어서, $a|_W$는
$P \oplus P[1] \to E|_W$와 동형이고 $a|_V$는
$Q \to E|_V$와 동형이다.% P09-ERR-0435
따라서 $R$은 원하는 대로 완전이고 $T$ 위에 지지된다.
\end{proof}

\begin{remark}
\label{spaces-perfect-remark-addendum}
Lemma \ref{spaces-perfect-lemma-lift-map-from-perfect-complex-with-support}의 증명은
$$
R|_W = P \oplus P^{\oplus n_1}[1] \oplus \ldots \oplus P^{\oplus n_m}[m]
$$
임을 보인다. 여기서 $m \geq 0$이고 $n_j \geq 0$이다. 따라서
$R|_W$의 최고 차수 코호몰로지 층은 $P$의 것과 같다. 사상
$P^{\oplus n_1}[1] \oplus \ldots \oplus P^{\oplus n_m}[m] \to R|_W$에
대해 이 구성을 반복하고, 원뿔을 취하고, 귀납법을 사용하면 임의로
주어진 차수보다 위에서 $R|_W$와 $P$의 코호몰로지 층들이 같게 할 수
있다.
\end{remark}

\begin{lemma}
\label{spaces-perfect-lemma-direct-summand-of-a-restriction}
$S$를 스킴이라 하자. $X$를 $S$ 위의 준콤팩트 준분리 대수공간이라
하자. $W$를 $X$의 준콤팩트 열린 부분공간이라 하자. $P$를
$D(\mathcal{O}_W)$의 완전 대상이라 하자. 그러면 $P$는
$D(\mathcal{O}_X)$의 어떤 완전 대상을 제한한 것의 직합 성분이다.
\end{lemma}

\begin{proof}
Lemma \ref{spaces-perfect-lemma-lift-map-from-perfect-complex-with-support}의 특수한
경우이다.
\end{proof}

\begin{theorem}
\label{spaces-perfect-theorem-bondal-van-den-Bergh}
$S$를 스킴이라 하자. $X$를 $S$ 위의 준콤팩트 준분리 대수공간이라
하자. 범주 $D_\QCoh(\mathcal{O}_X)$는 하나의 완전 대상으로 생성될 수
있다. 더 정확히는 완전 대상 $P$가 $D(\mathcal{O}_X)$에 존재하여,
$E \in D_\QCoh(\mathcal{O}_X)$에 대해 다음은 동치이다.
\begin{enumerate}
\item $E = 0$이다.
\item $\Hom_{D(\mathcal{O}_X)}(P[n], E) = 0$이고, 이는 모든
$n \in \mathbf{Z}$에 대해 성립한다.
\end{enumerate}
\end{theorem}

\begin{proof}
Lemma \ref{spaces-perfect-lemma-induction-principle}의 귀납 원리를 사용하여 이를
증명하겠다.

\medskip\noindent
$X$가 아핀이면 $\mathcal{O}_X$는 완전 생성자이다. 이는 Lemma
\ref{spaces-perfect-lemma-derived-quasi-coherent-small-etale-site}와 Derived Categories
of Schemes, Lemma \ref{perfect-lemma-affine-compare-bounded}에서 따른다.

\medskip\noindent
$(U \subset X, f : V \to X)$가 기본 구별 정사각형이라고 가정하자.
$U$는 준콤팩트이고 정리는 $U$에 대해 성립하며 $V$는 아핀 스킴이라고
하자. $P$를 $D(\mathcal{O}_U)$의 완전 대상이라 하고
$D_\QCoh(\mathcal{O}_U)$의 생성자라고 하자. Lemma
\ref{spaces-perfect-lemma-direct-summand-of-a-restriction}을 사용하여 완전 대상 $Q$를
$D(\mathcal{O}_X)$ 안에서 택할 수 있으며, 이를 $U$에 제한하면
직합이고 그 직합 성분 가운데 하나가 $P$이다. $V = \Spec(A)$라 하자.
$Z \subset V$를 $X \setminus U$의 역상인 기약 닫힌 부분스킴이라 하자.
이는 그 역상에 동형으로 사상한다(Definition
\ref{spaces-perfect-definition-elementary-distinguished-square}을 보라). 이는 $V$의
역콤팩트 닫힌 부분집합이다. 다음을 만족하는
$f_1, \ldots, f_r \in A$를 택하자.
$Z = V(f_1, \ldots, f_r)$. $K \in D(\mathcal{O}_V)$를
$f_1, \ldots, f_r$에 대한 $A$ 위의 Koszul 복합체에 대응하는 완전 대상이라
하자. $K$는 $Z$ 위에 지지되므로 직접상
$K' = Rf_*K$는 $D(\mathcal{O}_X)$의 완전 대상이고, 이를 $V$에
제한하면 $K$이다(Lemmas \ref{spaces-perfect-lemma-pushforward-perfect}와
\ref{spaces-perfect-lemma-pushforward-with-support-in-open}을 보라). 우리는
$Q \oplus K'$가 $D_\QCoh(\mathcal{O}_X)$의 생성자라고 주장한다.

\medskip\noindent
$E$를 $D_\QCoh(\mathcal{O}_X)$의 대상이라 하고, $Q \oplus K'$의 어떤
이동에서도 $E$로 가는 영이 아닌 사상이 없다고 하자. Lemma
\ref{spaces-perfect-lemma-pushforward-with-support-in-open}에 의해 $K' =  f_! K$이고,
따라서
$$
\Hom_{D(\mathcal{O}_X)}(K'[n], E) = \Hom_{D(\mathcal{O}_V)}(K[n], E|_V)
$$
이다. 그러므로 Derived Categories of Schemes, Lemma
\ref{perfect-lemma-orthogonal-koszul-complex}와 Lemma
\ref{spaces-perfect-lemma-derived-quasi-coherent-small-etale-site}에 의해, 이 군들의
소멸은 $E|_V$가
$R(U \times_X V \to V)_*E|_{U \times_X V}$와 동형임을 함의한다.
이는 $E = R(U \to X)_*E|_U$임을 함의한다(작은 세부사항은 생략한다).
이 경우
$$
\Hom_{D(\mathcal{O}_X)}(Q[n], E) = \Hom_{D(\mathcal{O}_U)}(Q|_U[n], E|_U)
$$
이고, 이는 $\Hom_{D(\mathcal{O}_U)}(P[n], E|_U)$를 직합 성분으로
포함한다. 따라서 $P$의 선택에 의해 이 군들의 소멸은 $E|_U$가
영임을 함의한다. 그러므로 $E$는 영이다.
\end{proof}
\begin{remark}
\label{spaces-perfect-remark-pullback-generator}
$S$를 스킴이라 하자. $f : X \to Y$를 $S$ 위의 준콤팩트 준분리
대수공간들의 사상이라 하자. $E \in D_\QCoh(\mathcal{O}_Y)$를
생성자라 하자(Theorem \ref{spaces-perfect-theorem-bondal-van-den-Bergh}를 보라).
그러면 다음은 동치이다.
\begin{enumerate}
\item $K \in D_\QCoh(\mathcal{O}_X)$에 대해 $Rf_*K = 0$일
필요충분조건은 $K = 0$인 것이다.
\item $Rf_* : D_\QCoh(\mathcal{O}_X) \to D_\QCoh(\mathcal{O}_Y)$는
동형을 반영한다.
\item $Lf^*E$는 $D_\QCoh(\mathcal{O}_X)$의 생성자이다.
\end{enumerate}
(1)과 (2)의 동치는
$Rf_* : D_\QCoh(\mathcal{O}_X) \to D_\QCoh(\mathcal{O}_Y)$가 삼각
범주의 완전 함자라는 사실의 형식적 귀결이다. 마찬가지로 (1)과 (3)의
동치는 $Lf^*$가 $Rf_*$의 왼쪽 수반이라는 사실에서 형식적으로 따른다.
이 조건들은 $f$가 아핀일 때(Lemma \ref{spaces-perfect-lemma-affine-morphism}), 또는
$f$가 열린 매장일 때, 또는 $f$가 그러한 사상들의 합성일 때 성립한다.
\end{remark}

\noindent
다음 결과는 Theorem \ref{spaces-perfect-theorem-bondal-van-den-Bergh}의 강화이고% P09-ERR-0436
정확히 같은 방법으로 증명된다. $T \subset |X|$를 닫힌 부분집합이라
하자. 여기서 $X$는 대수공간이다. $D_T(\mathcal{O}_X)$로 코호몰로지
층들이 $T$ 위에 지지되는 복합체들로 이루어진 엄밀히 충만하고 포화된
삼각 부분범주를 나타내자.

\begin{lemma}
\label{spaces-perfect-lemma-generator-with-support}
$S$를 스킴이라 하자. $X$를 $S$ 위의 준콤팩트 준분리 대수공간이라
하자. $T \subset |X|$를 닫힌 부분집합이라 하고
$|X| \setminus T$가 준콤팩트이라고 하자. 위 표기 아래 범주
$D_{\QCoh, T}(\mathcal{O}_X)$는 하나의 완전 대상으로 생성된다.
\end{lemma}

\begin{proof}
Lemma \ref{spaces-perfect-lemma-induction-principle}의 귀납 원리를 사용하여 이를
증명하겠다. 이 성질은 사이트 $X_{spaces, \etale}$의 표현 가능한
준콤팩트 준분리 대상들에 대해 참이다. 이는 Derived Categories of
Schemes, Lemma \ref{perfect-lemma-generator-with-support}에서 따른다.

\medskip\noindent
$(U \subset X, f : V \to X)$가 기본 구별 정사각형이라고 가정하자.
보조정리는 $U$에 대해 성립하고 $V$는 아핀이라고 하자. 증명을
끝내려면 결과가 $X$에 대해 성립함을 보여야 한다. $P$를
$D(\mathcal{O}_U)$의 완전 대상이라 하되 $T \cap U$ 위에 지지되고
$D_{\QCoh, T \cap U}(\mathcal{O}_U)$의 생성자라고 하자. Lemma
\ref{spaces-perfect-lemma-lift-map-from-perfect-complex-with-support}를 사용하여 완전
대상 $Q$를 $D(\mathcal{O}_X)$ 안에서 택할 수 있으며, 이는 $T$ 위에
지지되고 $U$로 제한하면 직합이고 그 직합 성분 가운데 하나가
$P$이다. $V = \Spec(B)$라 쓰자. $Z = X \setminus U$라 하자. 그러면
$f^{-1}Z$는 $V$의 닫힌 부분집합이고 $V \setminus f^{-1}Z$는
준콤팩트이다. $X$가 준분리이므로
$f^{-1}Z \cap f^{-1}T = f^{-1}(Z \cap T)$는 $V$의 닫힌 부분집합이고
$W = V \setminus f^{-1}(Z \cap T)$는 준콤팩트이다. 따라서
$g_1, \ldots, g_s \in B$를 택하여
$f^{-1}(Z \cap T) = V(g_1, \ldots, g_r)$가 되게 할 수 있다.% P09-ERR-0437
$K \in D(\mathcal{O}_V)$를 $g_1, \ldots, g_s$에 대한 $B$ 위의 Koszul
복합체에 대응하는 완전 대상이라 하자. $K$는 닫힌 부분집합
$f^{-1}(Z \cap T) \subset V$ 위에 지지되므로,% P09-ERR-0438
직접상 $K' = R(V \to X)_*K$는 $D(\mathcal{O}_X)$의 완전 대상이고,
이를 $V$에 제한하면 $K$이다(Lemmas \ref{spaces-perfect-lemma-pushforward-perfect}와
\ref{spaces-perfect-lemma-pushforward-with-support-in-open}을 보라). 우리는
$Q \oplus K'$가 $D_{\QCoh, T}(\mathcal{O}_X)$의 생성자라고 주장한다.

\medskip\noindent
$E$를 $D_{\QCoh, T}(\mathcal{O}_X)$의 대상이라 하고, $Q \oplus K'$의
어떤 이동에서도 $E$로 가는 영이 아닌 사상이 없다고 하자. Lemma
\ref{spaces-perfect-lemma-pushforward-with-support-in-open}에 의해
$K' =  R(V \to X)_! K$이고, 따라서
$$
\Hom_{D(\mathcal{O}_X)}(K'[n], E) = \Hom_{D(\mathcal{O}_V)}(K[n], E|_V)
$$
이다. 그러므로 Derived Categories of Schemes, Lemma
\ref{perfect-lemma-orthogonal-koszul-complex}에 의해
$E|_V = Rj_*E|_W$이다. 여기서 $j : W \to V$는 포함 사상이다. 그림은
다음과 같다.
$$
\xymatrix{
W \ar[r]_j & V & Z \cap T \ar[l] \ar[d] \\
V \setminus f^{-1}Z \ar[u]^{j'} \ar[ru]_{j''} & & Z \ar[lu]
}
$$
$E$는 $T$ 위에 지지되므로 $E|_W$는
$f^{-1}T \cap W = f^{-1}T \cap (V \setminus f^{-1}Z)$ 위에 지지되고,
이는 $W$ 안에서 닫혀 있다. 따라서
$$
E|_V = Rj_*(E|_W) = Rj_*(Rj'_*(E|_{U \cap V})) = Rj''_*(E|_{U \cap V})
$$
이다. 둘째 등식은 Cohomology, Lemma
\ref{cohomology-lemma-pushforward-restriction}의 (1)이며, $V$가
스킴이고 $E$가 준연접 코호몰로지 층을 가져 준콤팩트 열린 매장
$j'$을 따른 직접상이 바탕 스킴 위의 직접상과 일치하므로 적용할 수
있다. Remark \ref{spaces-perfect-remark-match-total-direct-images}를 보라. 이는
$E = R(U \to X)_*E|_U$임을 함의한다(작은 세부사항은 생략한다). 이
경우
$$
\Hom_{D(\mathcal{O}_X)}(Q[n], E) = \Hom_{D(\mathcal{O}_U)}(Q|_U[n], E|_U)
$$
이고, 이는 $\Hom_{D(\mathcal{O}_U)}(P[n], E|_U)$를 직합 성분으로
포함한다. 따라서 $P$의 선택에 의해 이 군들의 소멸은 $E|_U$가
영임을 함의한다. 그러므로 $E$는 영이다.
\end{proof}











\section{콤팩트 대상과 완전 대상}
\label{spaces-perfect-section-compact}

\noindent
이 절은 Derived Categories of Schemes, Section
\ref{perfect-section-compact}에 해당한다.

\begin{proposition}
\label{spaces-perfect-proposition-compact-is-perfect}
$S$를 스킴이라 하자. $X$를 $S$ 위의 준콤팩트 준분리 대수공간이라
하자. $D_\QCoh(\mathcal{O}_X)$의 대상이 콤팩트일 필요충분조건은
그것이 완전인 것이다.
\end{proposition}

\begin{proof}
$K$가 $D(\mathcal{O}_X)$의 완전 대상이고 그 쌍대가 $K^\vee$이면
(Cohomology on Sites, Lemma
\ref{sites-cohomology-lemma-dual-perfect-complex}), 다음이 성립한다.
$$
\Hom_{D(\mathcal{O}_X)}(K, M) =
H^0(X, K^\vee \otimes_{\mathcal{O}_X}^\mathbf{L} M)
$$
이는 $M$에 함자적이다.
$K^\vee \otimes_{\mathcal{O}_X}^\mathbf{L} -$는 직합과 가환하고,
Lemma \ref{spaces-perfect-lemma-quasi-coherence-pushforward-direct-sums}에 의해
$H^0(X, -)$도 $D_\QCoh(\mathcal{O}_X)$ 위에서 직합과 가환한다.
따라서 $K$는 $D_\QCoh(\mathcal{O}_X)$ 안에서 콤팩트이다.

\medskip\noindent
역으로 $K$를 $D_\QCoh(\mathcal{O}_X)$의 콤팩트 대상이라 하자. $K$가
완전임을 보이려면 $K|_U$가 완전임을 보이면 충분하다. 여기서 $U$는
$X$ 위에서 \'etale인 임의의 아핀 스킴이다. Cohomology on Sites, Lemma
\ref{sites-cohomology-lemma-perfect-independent-representative}를 보라.
$j : U \to X$가 준콤팩트 분리 사상임을 관찰하라. 따라서
$$
Rj_* : D_\QCoh(\mathcal{O}_U) \to D_\QCoh(\mathcal{O}_X)
$$
는 직합과 가환한다. Lemma
\ref{spaces-perfect-lemma-quasi-coherence-pushforward-direct-sums}를 보라. 그러므로
$U$로의 제한과 $Rj_*$의 수반성은 $K|_U$가
$D_\QCoh(\mathcal{O}_U)$의 완전 대상임을 함의한다. 따라서 $X$가
아핀인 경우, 특히 준콤팩트 준분리 스킴인 경우로 환원된다. Lemmas
\ref{spaces-perfect-lemma-derived-quasi-coherent-small-etale-site}와
\ref{spaces-perfect-lemma-descend-perfect}를 통해 스킴의 경우, 즉 Derived Categories
of Schemes, Proposition \ref{perfect-proposition-compact-is-perfect}로
환원된다.
\end{proof}

\begin{remark}
\label{spaces-perfect-remark-classical-generator}
$S$를 스킴이라 하자. $X$를 $S$ 위의 준콤팩트 준분리 대수공간이라
하자. $G$를 $D(\mathcal{O}_X)$의 완전 대상이자
$D_\QCoh(\mathcal{O}_X)$의 생성자라고 하자. Theorem
\ref{spaces-perfect-theorem-bondal-van-den-Bergh}에 의해 이러한 대상이 적어도 하나
있다. Lemma \ref{spaces-perfect-lemma-quasi-coherence-direct-sums}, Proposition
\ref{spaces-perfect-proposition-compact-is-perfect}, 그리고 Derived Categories,
Proposition \ref{derived-proposition-generator-versus-classical-generator}를
결합하면 $G$가 $D_{perf}(\mathcal{O}_X)$의 고전적 생성자임을 알 수
있다.
\end{remark}

\noindent
다음 결과는 Proposition \ref{spaces-perfect-proposition-compact-is-perfect}의 강화이다.
$T \subset |X|$를 닫힌 부분집합이라 하자. 여기서 $X$는 대수공간이다.
앞에서와 같이 $D_T(\mathcal{O}_X)$는 코호몰로지 층이 $T$ 위에
지지되는 복합체들로 이루어진 엄밀히 충만하고 포화된 삼각 부분범주를
나타낸다. 직합을 취하는 것은 코호몰로지 층을 취하는 것과 가환하므로
$D_T(\mathcal{O}_X)$에는 직합이 있고, 이 직합은
$D(\mathcal{O}_X)$ 안의 직합과 같다.

\begin{lemma}
\label{spaces-perfect-lemma-compact-is-perfect-with-support}
$S$를 스킴이라 하자. $X$를 $S$ 위의 준콤팩트 준분리 대수공간이라
하자. $T \subset |X|$를 닫힌 부분집합이라 하고
$|X| \setminus T$가 준콤팩트이라고 하자.
$D_{\QCoh, T}(\mathcal{O}_X)$의 대상이 콤팩트일 필요충분조건은
$D(\mathcal{O}_X)$의 대상으로서 완전인 것이다.
\end{lemma}

\begin{proof}
보조정리 앞의 remark에 의해 $D_{\QCoh, T}(\mathcal{O}_X)$는 직합을
갖는 삼각 범주이다. Proposition \ref{spaces-perfect-proposition-compact-is-perfect}에
의해 완전 대상들은 $D(\mathcal{O}_X)$의 콤팩트 대상을 정의하므로,
더욱이 직합을 취하는 것 아래 보존되는 임의의 부분범주의 콤팩트
대상을 정의한다. 역으로 생성자
$E \in D_{\QCoh, T}(\mathcal{O}_X)$가 존재하고 이는 완전
$\mathcal{O}_X$-가군 복합체라는 사실을 사용하겠다.% P09-ERR-0439
Lemma \ref{spaces-perfect-lemma-generator-with-support}를 보라. 따라서 위에 의해 $E$는
콤팩트이다. 그러면 Derived Categories, Proposition
\ref{derived-proposition-generator-versus-classical-generator}에 의해
$E$는 $D_{\QCoh, T}(\mathcal{O}_X)$의 콤팩트 대상들로 이루어진 충만
부분범주의 고전적 생성자이다. 따라서 임의의 콤팩트 대상은 $E$에서
다음 연산들로 이루어진 유한한 열을 통해 구성할 수 있다. (a) 이동,
(b) 유한 직합, (c) 원뿔, (d) 직합 성분. 이 연산들은 각각 완전성을
보존하므로 결과가 따른다.
\end{proof}

\begin{remark}
\label{spaces-perfect-remark-classical-generator-with-support}
$S$를 스킴이라 하자. $X$를 $S$ 위의 준콤팩트 준분리 대수공간이라
하자. $T \subset |X|$를 닫힌 부분집합이라 하고
$|X| \setminus T$가 준콤팩트이라고 하자. $G$를
$D_{\QCoh, T}(\mathcal{O}_X)$의 완전 대상이자
$D_{\QCoh, T}(\mathcal{O}_X)$의 생성자라고 하자. Lemma
\ref{spaces-perfect-lemma-generator-with-support}에 의해 이러한 대상이 적어도 하나
있다. $D_{\QCoh, T}(\mathcal{O}_X)$가 직합을 갖는다는 사실과 Lemma
\ref{spaces-perfect-lemma-compact-is-perfect-with-support} 및 Derived Categories,
Proposition \ref{derived-proposition-generator-versus-classical-generator}를
결합하면 $G$가 $D_{perf, T}(\mathcal{O}_X)$의 고전적 생성자임을 알
수 있다.
\end{remark}

\noindent
다음 보조정리는 앞 보조정리의 증명에 들어가는 발상들의 응용이다.

\begin{lemma}
\label{spaces-perfect-lemma-map-from-pseudo-coherent-to-complex-with-support}
$S$를 스킴이라 하자. $X$를 $S$ 위의 준콤팩트 준분리 대수공간이라
하자. $T \subset |X|$를 닫힌 부분집합이라 하고, 그 여집합
$U \subset X$가 준콤팩트이라고 하자.
$\alpha : P \to E$를 $D_\QCoh(\mathcal{O}_X)$의 사상이라 하고 다음
가운데 하나를 가정하자.
\begin{enumerate}
\item $P$는 완전이고 $E$는 $T$ 위에 지지된다.
\item $P$는 유사연접이고 $E$는 $T$ 위에 지지되며 $E$는 아래로
유계이다.
\end{enumerate}
그러면 완전 $\mathcal{O}_X$-가군 복합체 $I$와 사상
$I \to \mathcal{O}_X[0]$가 존재하여, $I \otimes^\mathbf{L} P \to E$는
영이고 $I|_U \to \mathcal{O}_U[0]$는 동형이다.
\end{lemma}

\begin{proof}
$\mathcal{D} = D_{\QCoh, T}(\mathcal{O}_X)$라 두자. 두 경우 모두
복합체 $K = R\SheafHom(P, E)$는 $\mathcal{D}$의 대상이다. 준연접성은
Lemma \ref{spaces-perfect-lemma-quasi-coherence-internal-hom}을 보라. $K$가 $T$ 위에
지지됨은 분명하다. 이는 $R\SheafHom$의 구성이 열린집합으로의 제한과
가환하기 때문이다. 사상 $\alpha$는
$H^0(K) = \Hom_{D(\mathcal{O}_X)}(\mathcal{O}_X[0], K)$의 한 원소를
정의한다. 그러면 사상 $\alpha : \mathcal{O}_X[0] \to K$에 대해
결과를 증명하면 충분하다.

\medskip\noindent
$E \in \mathcal{D}$를 완전 생성자라 하자. Lemma
\ref{spaces-perfect-lemma-generator-with-support}를 보라. 다음처럼 쓰자.
$$
K = \text{hocolim} K_n
$$
이는 Derived Categories, Lemma \ref{derived-lemma-write-as-colimit}에서
생성자 $E$를 사용한 표현이다. 함자 $\mathcal{D} \to D(\mathcal{O}_X)$가 직합과
가환하므로 $K = \text{hocolim} K_n$은 $D(\mathcal{O}_X)$ 안에서도
성립한다. $\mathcal{O}_X$는 $D(\mathcal{O}_X)$의 콤팩트 대상이므로,
어떤 $n$과 사상 $\alpha_n : \mathcal{O}_X \to K_n$을 찾을 수 있고,
이 사상은 $\alpha$를 준다. Derived Categories, Lemma
\ref{derived-lemma-commutes-with-countable-sums}를 보라. Derived
Categories, Lemma \ref{derived-lemma-factor-through}을 사상
$\mathcal{O}_X[0] \to K_n$에 적용하자. 이는 주위 범주
$D(\mathcal{O}_X)$ 안의 사상이다. 그러면
$\alpha_n$이 $\mathcal{O}_X[0] \to Q \to K_n$으로 분해됨을 알 수
있다. 여기서 $Q$는 $\langle E \rangle$의 대상이다. 따라서 $Q$는
$T$ 위에 지지되는 완전 복합체이다.

\medskip\noindent
구별 삼각형
$$
I \to \mathcal{O}_X[0] \to Q \to I[1]
$$
을 택하자. 구성에 의해 $I$는 완전이고, 사상
$I \to \mathcal{O}_X[0]$는 $U$ 위로 제한하면 동형이며, 합성
$I \to K$는 $\alpha$가 $Q$를 통해 분해되므로 영이다. 이것으로
보조정리가 증명된다.
\end{proof}










\section{가군 범주로서의 유도 범주}
\label{spaces-perfect-section-derived-is-dga}

\noindent
이 절은 Derived Categories of Schemes, Section
\ref{perfect-section-derived-is-dga}에 해당한다.

\begin{lemma}
\label{spaces-perfect-lemma-tensor-with-QCoh-complex}
$S$를 스킴이라 하자. $X$를 $S$ 위의 대수공간이라 하자.
$K^\bullet$을 코호몰로지 층들이 준연접인 $\mathcal{O}_X$-가군의
복합체라 하자. 미분 등급 자기준동형 대수를
$(E, d) = \Hom_{\text{Comp}^{dg}(\mathcal{O}_X)}(K^\bullet, K^\bullet)$라
하자. 그러면 Differential Graded Algebra, Lemma
\ref{dga-lemma-tensor-with-complex-derived}의 함자
$$
- \otimes_E^\mathbf{L} K^\bullet :
D(E, \text{d}) \longrightarrow D(\mathcal{O}_X)
$$
의 상은 $D_\QCoh(\mathcal{O}_X)$에 포함된다.
\end{lemma}

\begin{proof}
$P$를 미분 등급 $E$-가군이라 하되 성질 $P$를 갖는다고 하자.
$F_\bullet$을 Differential Graded Algebra, Section
\ref{dga-section-P-resolutions}에서와 같이 $P$ 위의 여과라 하자.
그러면
$$
P \otimes_E K^\bullet = \text{hocolim}\ F_iP \otimes_E K^\bullet
$$
이다. 각 $F_iP$에는 유한 여과가 있고 그 등급 조각들은 $E[k]$들의
직합이다. 결과는 쉽게 따른다.
\end{proof}

\noindent
다음 보조정리는 강화할 수 있다 (영이 아닌 코호몰로지 층들이
고정된 범위에만 놓이는 모든 $L$에 대하여 소멸을 일양적으로
택할 수 있다).

\begin{lemma}
\label{spaces-perfect-lemma-ext-from-perfect-into-bounded-QCoh}
$S$를 스킴이라 하자. $X$를 $S$ 위의 준콤팩트 준분리 대수공간이라
하자. $K$를 $D(\mathcal{O}_X)$의 완전 대상이라 하자. 그러면
\begin{enumerate}
\item 정수 $a \leq b$가 존재하여 다음이 성립한다:
$\Hom_{D(\mathcal{O}_X)}(K, L) = 0$은
$L \in D_\QCoh(\mathcal{O}_X)$가 $H^i(L) = 0$을 모든
$i \in [a, b]$에 대해 만족할 때 성립하고,
\item $L$이 유계이면 $\Ext^n_{D(\mathcal{O}_X)}(K, L)$은 유한 개의
경우를 제외한 모든 $n$에 대해 영이다.
\end{enumerate}
\end{lemma}

\begin{proof}
$\Ext^n_{D(\mathcal{O}_X)}(K, L) =
\Hom_{D(\mathcal{O}_X)}(K, L[n])$이므로 (2)는 (1)에서 따른다.
(1)을 증명하자. $K$가 완전이므로
$$
\Ext^i_{D(\mathcal{O}_X)}(K, L) =
H^i(X, K^\vee \otimes_{\mathcal{O}_X}^\mathbf{L} L)
$$
이다. 여기서 $K^\vee$는 $K$의 “쌍대” 완전 복합체이다. Cohomology
on Sites, Lemma \ref{sites-cohomology-lemma-dual-perfect-complex}을 보라.
Lemmas \ref{spaces-perfect-lemma-quasi-coherence-tensor-product} and
\ref{spaces-perfect-lemma-pseudo-coherent}에 의해
$P = K^\vee \otimes_{\mathcal{O}_X}^\mathbf{L} L$은 $D_\QCoh(X)$에
속한다 (완전 복합체가 준연접 코호몰로지 층들을 가짐을 보려면
뒤의 보조정리를 사용한다). $K^\vee$가 $[a, b]$ 안의 tor 진폭을
갖는다고 하자. 그러면 스펙트럴 열
$$
E_1^{p, q} = H^p(K^\vee \otimes_{\mathcal{O}_X}^\mathbf{L} H^q(L))
\Rightarrow
H^{p + q}(K^\vee \otimes_{\mathcal{O}_X}^\mathbf{L} L)
$$
은 $H^j(K^\vee \otimes_{\mathcal{O}_X}^\mathbf{L} L)$이
$H^q(L) = 0$이 모든 $q \in [j - b, j - a]$에 대해 성립할 때
영임을 보인다.
$N$을 Cohomology of Spaces, Lemma
\ref{spaces-cohomology-lemma-vanishing-quasi-separated}의 정수
$\max(d_p + p)$라 하자. 그러면
$H^0(X, K^\vee \otimes_{\mathcal{O}_X}^\mathbf{L} L)$은 코호몰로지 층들
$$
H^{-N}(K^\vee \otimes_{\mathcal{O}_X}^\mathbf{L} L),
\ H^{-N + 1}(K^\vee \otimes_{\mathcal{O}_X}^\mathbf{L} L),
\ \ldots,
\ H^0(K^\vee \otimes_{\mathcal{O}_X}^\mathbf{L} L)
$$
이 영이면 영이다.
실제로 인용한 보조정리와
Lemma \ref{spaces-perfect-lemma-application-nice-K-injective}에 의해
$$
H^0(X, K^\vee \otimes_{\mathcal{O}_X}^\mathbf{L} L) =
H^0(X, \tau_{\geq -N}(K^\vee \otimes_{\mathcal{O}_X}^\mathbf{L} L))
$$
이고, 코호몰로지 층들의 소멸에 의해 이는
$H^0(X, \tau_{\geq 1}(K^\vee \otimes_{\mathcal{O}_X}^\mathbf{L} L))$와
같다. 후자는 Derived Categories, Lemma
\ref{derived-lemma-negative-vanishing}에 의해 영이다. 따라서
$\Hom_{D(\mathcal{O}_X)}(K, L)$은 $H^i(L) = 0$이 모든
$i \in [-b - N, -a]$에 대해 성립하면 영이다.
\end{proof}

\noindent
다음은 Derived Categories of Schemes, Theorem
\ref{perfect-theorem-DQCoh-is-Ddga}의 유사이다.

\begin{theorem}
\label{spaces-perfect-theorem-DQCoh-is-Ddga}
$S$를 스킴이라 하자. $X$를 $S$ 위의 준콤팩트 준분리
대수공간이라 하자. 그러면 미분 등급 대수 $(E, \text{d})$가
존재하여 그 비영 코호몰로지 군 $H^i(E)$은 유한 개뿐이고,
$D_\QCoh(\mathcal{O}_X)$는 $D(E, \text{d})$와 동치이다.
\end{theorem}

\begin{proof}
$K^\bullet$을 $\mathcal{O}$-가군의 K-단사 복합체로서 완전이고
$D_\QCoh(\mathcal{O}_X)$를 생성하는 것이라 하자. % P09-ERR-0441
그런 것은 Theorem \ref{spaces-perfect-theorem-bondal-van-den-Bergh}와
K-단사 분해의 존재성에 의해 존재한다. 다음을 택하면 정리가
성립함을 보이겠다.
$$
(E, \text{d}) = \Hom_{\text{Comp}^{dg}(\mathcal{O}_X)}(K^\bullet, K^\bullet)
$$
여기서 $\text{Comp}^{dg}(\mathcal{O}_X)$는 $\mathcal{O}$-가군 복합체의
미분 등급 범주이다. Differential Graded Algebra, Section
\ref{dga-section-variant-base-change}을 보라. $K^\bullet$이
K-단사이므로
\begin{equation}
\label{spaces-perfect-equation-E-is-OK}
H^n(E) = \Ext^n_{D(\mathcal{O}_X)}(K^\bullet, K^\bullet)
\end{equation}
은 모든 $n \in \mathbf{Z}$에 대해 성립한다. Lemma
\ref{spaces-perfect-lemma-ext-from-perfect-into-bounded-QCoh}에 의해 이 Ext들 가운데
비영인 것은 유한 개뿐이다. Differential Graded Algebra, Lemma
\ref{dga-lemma-tensor-with-complex-derived}의 함자
$$
- \otimes_E^\mathbf{L} K^\bullet :
D(E, \text{d}) \longrightarrow D(\mathcal{O}_X)
$$
를 생각하자. $K^\bullet$은 완전이므로 $D(\mathcal{O}_X)$의 콤팩트
대상을 정의한다. Proposition \ref{spaces-perfect-proposition-compact-is-perfect}을
보라. % P09-ERR-0442
이를 (\ref{spaces-perfect-equation-E-is-OK})와 합치면, 위 함자는 Differential
Graded Algebra, Lemma
\ref{dga-lemma-fully-faithful-in-compact-case}에 의해 완전충실하다.
이 함자는 Differential Graded Algebra, Lemma
\ref{dga-lemma-tensor-with-complex-hom-adjoint}에 의해 % P09-ERR-0443
우수반
$$
R\Hom(K^\bullet, - ) : D(\mathcal{O}_X) \longrightarrow D(E, \text{d})
$$
을 가지며, 수반 함자에 관한 일반론에 의해 이는 좌준역함자이다.
한편 Lemma \ref{spaces-perfect-lemma-tensor-with-QCoh-complex}에 의해
$$
- \otimes_E^\mathbf{L} K^\bullet :
D(E, \text{d}) \longrightarrow D_\QCoh(\mathcal{O}_X)
$$
를 얻는다. $K^\bullet$을 $D_\QCoh(\mathcal{O}_X)$의 생성자로
택했으므로, $D_\QCoh(\mathcal{O}_X)$에 제한한 수반의 핵은 영이다.
형식적 논증은 원하는 동치를 얻음을 보인다. Derived Categories,
Lemma \ref{derived-lemma-fully-faithful-adjoint-kernel-zero}를 보라.
\end{proof}

\begin{remark}[지지를 갖는 변형]
\label{spaces-perfect-remark-DQCoh-is-Ddga-with-support}
$S$를 스킴이라 하자. % P09-ERR-0444
$X$를 준콤팩트 준분리 대수공간이라 하자.
$T \subset |X|$를 $|X| \setminus T$가 준콤팩트인 닫힌 부분집합이라
하자. Theorem \ref{spaces-perfect-theorem-DQCoh-is-Ddga}의 유사는
$D_{\QCoh, T}(\mathcal{O}_X)$에 대해 성립한다. 이는 정리의 증명과
완전히 같은 논증에서 Lemmas \ref{spaces-perfect-lemma-generator-with-support}와
\ref{spaces-perfect-lemma-compact-is-perfect-with-support}, 그리고 지지를 갖는
Lemma \ref{spaces-perfect-lemma-tensor-with-QCoh-complex}의 변형을 사용하면 따른다.
이 결과가 필요하게 되면 여기서 이를 정확히 진술하고 상세한 증명을
제시하겠다.
\end{remark}

\begin{remark}[미분 등급 대수의 유일성]
\label{spaces-perfect-remark-independence-choice}
$X$를 환 $R$ 위의 준콤팩트 준분리 대수공간이라 하자. Theorem
\ref{spaces-perfect-theorem-DQCoh-is-Ddga}의 증명에 쓰인 구성에 의해, 미분 등급
대수 $(A, \text{d})$가 $R$ 위에 존재하여 삼각 범주로서
$D_\QCoh(X)$는 $R$-선형으로 $D(A, \text{d})$와 동치이다.
다음과 같이 물을 수 있다: $(A, \text{d})$는 얼마나 유일한가?
그 답은 $(A, \text{d})$가 유도 동치까지 잘 정의된다고 말하는
것보다 (단지) 조금 더 낫다. 실제로 $(B, \text{d})$가 두 번째
그러한 쌍이라고 하자. 그러면
$$
(A, \text{d}) = \Hom_{\text{Comp}^{dg}(\mathcal{O}_X)}(K^\bullet, K^\bullet)
$$
이고
$$
(B, \text{d}) = \Hom_{\text{Comp}^{dg}(\mathcal{O}_X)}(L^\bullet, L^\bullet)
$$
이다. 여기서 $K^\bullet$과 $L^\bullet$은
$\mathcal{O}_X$-가군의 K-단사 복합체들로서
$D_\QCoh(\mathcal{O}_X)$의 완전 생성자들에 대응한다. 다음처럼 두자.
$$
\Omega = \Hom_{\text{Comp}^{dg}(\mathcal{O}_X)}(K^\bullet, L^\bullet)
\quad
\Omega' = \Hom_{\text{Comp}^{dg}(\mathcal{O}_X)}(L^\bullet, K^\bullet)
$$
그러면 $\Omega$는 미분 등급 $B^{opp} \otimes_R A$-가군이고,
$\Omega'$는 미분 등급 $A^{opp} \otimes_R B$-가군이다. 또한 동치
$$
D(A, \text{d}) \to D_\QCoh(\mathcal{O}_X) \to
D(B, \text{d})
$$
는 함자 $- \otimes_A^\mathbf{L} \Omega'$로 주어지고, 준역함자도
마찬가지이다. 따라서 우리는 Differential Graded Algebra, Remark
\ref{dga-remark-hochschild-cohomology}의 상황에 있다. 이 비고가
필요하게 되면 여기서 상세한 증명과 함께 정확한 진술을 제시하겠다.
\end{remark}

\section{유사연접 복합체의 특성화, I}
\label{spaces-perfect-section-pseudo-coherent-hocolim}

\noindent
이 내용은 More on Morphisms of Spaces, Section
\ref{spaces-more-morphisms-section-characterize-pseudo-coherent}에서
계속된다. 유사연접 대상을 완전 대상에 의한 근사들의 유도 호모토피
극한으로 특성화할 수 있다. % P09-ERR-0446

\begin{lemma}
\label{spaces-perfect-lemma-pseudo-coherent-hocolim}
$S$를 스킴이라 하자. $X$를 $S$ 위의 준콤팩트 준분리
대수공간이라 하자. $K \in D(\mathcal{O}_X)$라 하자. 다음은
동치이다.
\begin{enumerate}
\item $K$는 유사연접이고,
\item $K = \text{hocolim} K_n$이며, $K_n$은 완전이고
$\tau_{\geq -n}K_n \to \tau_{\geq -n}K$는 모든 $n$에 대해 동형이다.
\end{enumerate}
\end{lemma}

\begin{proof}
함의 (2) $\Rightarrow$ (1)은 임의의 환 달린 사이트에서 참이다.
실제로 (2)가 성립한다고 하자. 유도 범주의 완전 대상은
유사연접임을 상기하자. Cohomology on Sites, Lemma
\ref{sites-cohomology-lemma-perfect}을 보라. 그러면 정의에서
$\tau_{\geq -n}K_n$은 $(-n + 1)$-유사연접이고, 따라서
$\tau_{\geq -n}K$도 $(-n + 1)$-유사연접이며, 따라서 $K$는
$(-n + 1)$-유사연접임이 따른다. 이는 모든 $n$에 대해 참이므로
$K$는 유사연접이다. Cohomology on Sites, Definition
\ref{sites-cohomology-definition-pseudo-coherent}을 보라.

\medskip\noindent
(1)을 가정하자. 먼저 $K_1 \to K$를 택한다. 이는
$(X, K, -2)$의 근사이고 $K_1$은 완전 복합체이다. Definitions
\ref{spaces-perfect-definition-approximation-holds}와 \ref{spaces-perfect-definition-approximation},
그리고 Theorem \ref{spaces-perfect-theorem-approximation}을 보라. 귀납적으로
$$
K_1 \to K_2 \to \ldots \to K_n \to K
$$
를 얻었다고 하자. 여기서 $K_i$는 완전이고
$\tau_{\geq -i}K_i \to \tau_{\geq -i}K$가 모든
$1 \leq i \leq n$에 대해 동형이다.
% P09-ERR-0447
그러면 Lemma \ref{spaces-perfect-lemma-ext-from-perfect-into-bounded-QCoh}에서와
같이 $a \leq b$를 완전 대상 $K_n$에 대해 택한다.
$K_{n + 1} \to K$를 택하자. 이는
$(X, K, \min(a - 1, -n - 1))$의 근사이다. 구별 삼각형
$$
K_{n + 1} \to K \to C \to K_{n + 1}[1]
$$
을 택한다. 그러면 $C \in D_\QCoh(\mathcal{O}_X)$이고
$H^i(C) = 0$은 $i \geq a$에 대해 성립한다. 따라서
$a, b$의 선택에 의해
$\Hom_{D(\mathcal{O}_X)}(K_n, C) = 0$이다. 그러므로 합성
$K_n \to K \to C$는 영이다. 따라서 Derived Categories, Lemma
\ref{derived-lemma-representable-homological}에 의해
$K_n \to K$를 $K_{n + 1}$을 통해 분해할 수 있다. 이로써
귀납 단계가 증명된다.

\medskip\noindent
이제 $K = \text{hocolim} K_n$임을 증명해야 한다. 이는 함자
$H^i( - ) : D(\mathcal{O}_X) \to \textit{Mod}(\mathcal{O}_X)$와
$K_n$의 선택에 Derived Categories, Lemma
\ref{derived-lemma-cohomology-of-hocolim}을 적용하면 따른다.
\end{proof}

\begin{lemma}
\label{spaces-perfect-lemma-pseudo-coherent-hocolim-with-support}
$X$를 준콤팩트 준분리 스킴이라 하자.
$T \subset X$를 $X \setminus T$가 준콤팩트인 닫힌 부분집합이라
하자. $K \in D(\mathcal{O}_X)$가 $T$ 위에 지지된다고 하자.
다음은 동치이다.
\begin{enumerate}
\item $K$는 유사연접이고,
\item $K = \text{hocolim} K_n$이며, $K_n$은 완전이고 $T$ 위에
지지되며, $\tau_{\geq -n}K_n \to \tau_{\geq -n}K$는 모든 $n$에
대해 동형이다.
\end{enumerate}
\end{lemma}

\begin{proof}
이 보조정리의 증명은 Lemma \ref{spaces-perfect-lemma-pseudo-coherent-hocolim}의
증명과 완전히 같다. 다만 근사들을 택할 때 삼중항 $(T, K, m)$을
사용한다.
\end{proof}

\section{코히레이터의 재고찰}
\label{spaces-perfect-section-better-coherator}

\noindent
Section \ref{spaces-perfect-section-coherator}에서 우리는 우수반 $RQ_X$를
구성하고 연구하였다. 이는 표준 함자
$D(\QCoh(\mathcal{O}_X)) \to D(\mathcal{O}_X)$의 우수반이다.
그것은 코히레이터
$Q_X : \textit{Mod}(\mathcal{O}_X) \to \QCoh(\mathcal{O}_X)$의
우유도 연장으로 구성되었다. 이 절에서는 포함 함자
$$
D_\QCoh(\mathcal{O}_X) \longrightarrow D(\mathcal{O}_X)
$$
가 언제 우수반을 갖는지 연구한다. 이 우수반이 존재하면 이를
\footnote{이는 아마 표준적인 표기가 아닐 것이다. 그러나 우리는 이미
코히레이터에 $Q_X$를, 그 유도 연장에 $RQ_X$를 사용하였다.}
$$
DQ_X :
D(\mathcal{O}_X) \longrightarrow D_\QCoh(\mathcal{O}_X)
$$
로 표기하겠다. 준콤팩트 준분리 대수공간은 그러한 우수반을 갖는다.

\begin{lemma}
\label{spaces-perfect-lemma-better-coherator}
$S$를 스킴이라 하자. $X$를 $S$ 위의 준콤팩트 준분리
대수공간이라 하자. 포함 함자
$D_\QCoh(\mathcal{O}_X) \to D(\mathcal{O}_X)$는 우수반을 갖는다.
\end{lemma}

\begin{proof}[첫째 증명]
이를 증명하기 위해 Lemma \ref{spaces-perfect-lemma-induction-principle}의 귀납
원리를 사용하겠다. $D(\QCoh(\mathcal{O}_X)) \to
D_\QCoh(\mathcal{O}_X)$가 동치라면, Section
\ref{spaces-perfect-section-coherator}의 함자 $RQ_X$가 함자
$D(\QCoh(\mathcal{O}_X)) \to D(\mathcal{O}_X)$의 우수반이므로
보조정리는 참이다. 특히 Lemma \ref{spaces-perfect-lemma-affine-coherator}에 의해
보조정리는 아핀 대수공간에 대해 참이다. 따라서 다음을 보이면
충분하다: $(U \subset X, f : V \to X)$가 기본 구별 정사각형이고
$U$는 준콤팩트이며 $V$는 아핀이고, 보조정리가 $U$, $V$,
$U \times_X V$에 대해 성립한다면, 보조정리는 $X$에 대해서도
성립한다.

\medskip\noindent
우수반이 존재할 필요충분조건은 $K$가
$D(\mathcal{O}_X)$의 임의의 대상일 때 다음 구별 삼각형을
찾을 수 있는 것이다.
% P09-ERR-0448
$$
E' \to E \to K \to E'[1]
$$
이는 $D(\mathcal{O}_X)$ 안의 삼각형이고, $E'$는
$D_\QCoh(\mathcal{O}_X)$에 속하며, $\Hom(M, K) = 0$이 모든
$M$ 가운데 $D_\QCoh(\mathcal{O}_X)$에 속하는 것에 대해 성립한다.
Derived Categories, Lemma \ref{derived-lemma-right-adjoint}를 보라.
구별 삼각형
$$
E \to Rj_{U, *}E|_U \oplus Rj_{V, *}E|_V \to
Rj_{U \times_X V, *}E|_{U \times_X V} \to E[1]
$$
을 $D(\mathcal{O}_X)$ 안에서 생각하자. 이는 Lemma
\ref{spaces-perfect-lemma-exact-sequence-j-star}의 삼각형이다. Derived Categories,
Lemma
\ref{derived-lemma-prepare-adjoint}에 의해
$Rj_{U, *}E|_U$, $Rj_{V, *}E|_V$와
$Rj_{U \times_X V, *}E|_{U \times_X V}$에 대해 원하는 구별
삼각형들을 구성하면 충분하다. 이로써 다음 문단에서 논의하는
명제로 환원된다.

\medskip\noindent
$j : U \to X$를 에탈 사상이라 하되, $U$는 준콤팩트 준분리이고
보조정리는 $U$에 대해 참이라고 하자. % P09-ERR-0449
$L$을 $D(\mathcal{O}_U)$의 대상이라 하자. 그러면 구별 삼각형
$$
E' \to Rj_*L \to K \to E'[1]
$$
이 $D(\mathcal{O}_X)$ 안에 존재하여, $E'$는
$D_\QCoh(\mathcal{O}_X)$에 속하고, $\Hom(M, K) = 0$이 모든
$M$ 가운데 $D_\QCoh(\mathcal{O}_X)$에 속하는 것에 대해 성립한다.
이를 보이기 위해 구별 삼각형
$$
L' \to L \to Q \to L'[1]
$$
을 $D(\mathcal{O}_U)$ 안에서 택하되, $L'$는
$D_\QCoh(\mathcal{O}_U)$에 속하고, $\Hom(N, Q) = 0$이 모든
$N$ 가운데 $D_\QCoh(\mathcal{O}_U)$에 속하는 것에 대해 성립한다고
하자.
Derived Categories, Lemma \ref{derived-lemma-right-adjoint}의 명제가
필요충분조건이므로 이는 가능하다. 다음 구별 삼각형을 얻는다.
$$
Rj_*L' \to Rj_*L \to Rj_*Q \to Rj_*L'[1]
$$
이는 $D(\mathcal{O}_X)$ 안의 삼각형이다.
$Rj_*L'$는 Lemma \ref{spaces-perfect-lemma-quasi-coherence-direct-image}에 의해
$D_\QCoh(\mathcal{O}_X)$에 속한다. 한편
$M$이 $D_\QCoh(\mathcal{O}_X)$에 속하면, % P09-ERR-0450
$$
\Hom(M, Rj_*Q) = \Hom(Lj^*M, Q) = 0
$$
이다. 이는 Lemma \ref{spaces-perfect-lemma-quasi-coherence-pullback}에 의해
$Lj^*M$이 $D_\QCoh(\mathcal{O}_U)$에 속하기 때문이다.
이것으로 증명이 끝난다.
\end{proof}

\begin{proof}[둘째 증명]
우수반은 Derived Categories, Proposition
\ref{derived-proposition-brown}에 의해 존재한다. 가정들은
만족된다. 먼저 $D_\QCoh(\mathcal{O}_X)$는 직합을 가지며 직합은
포함 함자와 가환한다 (Lemma
\ref{spaces-perfect-lemma-quasi-coherence-direct-sums}). 한편
$D_\QCoh(\mathcal{O}_X)$는 콤팩트 생성된다. 실제로 Theorem
\ref{spaces-perfect-theorem-bondal-van-den-Bergh}에 의해 완전 생성자를 가지며
% P09-ERR-0451
Proposition \ref{spaces-perfect-proposition-compact-is-perfect}에 의해 완전 대상은
콤팩트이기 때문이다.
\end{proof}

\begin{lemma}
\label{spaces-perfect-lemma-pushforward-better-coherator}
$S$를 스킴이라 하자. $f : X \to Y$를 $S$ 위 대수공간의
준콤팩트 준분리 사상이라 하자. 우수반 $DQ_X$와 $DQ_Y$가 포함
함자 $D_\QCoh \to D$에 대해 각각 $X$와 $Y$에 존재한다면,
$$
Rf_* \circ DQ_X = DQ_Y \circ Rf_*
$$
이다.
\end{lemma}

\begin{proof}
Lemma \ref{spaces-perfect-lemma-quasi-coherence-direct-image}에 의해 $Rf_*$는
$D_\QCoh(\mathcal{O}_X)$를 $D_\QCoh(\mathcal{O}_Y)$ 안으로
보내므로 이 명제는 뜻이 있다. 마찬가지로 Lemma
\ref{spaces-perfect-lemma-quasi-coherence-pullback}에 의해 $Lf^*$는
$D_\QCoh(\mathcal{O}_Y)$를 $D_\QCoh(\mathcal{O}_X)$ 안으로
보낸다. 따라서 $Rf_* \circ DQ_X$와 $DQ_Y \circ Rf_*$는 둘 다
$Lf^* : D_\QCoh(\mathcal{O}_Y) \to D(\mathcal{O}_X)$의
우수반이므로 명제가 참이다.
\end{proof}

\begin{remark}
\label{spaces-perfect-remark-explain-consequence}
$S$를 스킴이라 하자. $(U \subset X, f : V \to X)$를 $S$ 위
대수공간의 기본 구별 정사각형이라 하자. $X$, $U$, $V$가
준콤팩트 준분리라고 가정하자. Lemma
\ref{spaces-perfect-lemma-better-coherator}에 의해 함자 $DQ_X$, $DQ_U$, $DQ_V$,
$DQ_{U \times_X V}$가 존재한다. 또한 표준 구별 삼각형
$$
DQ_X(K) \to Rj_{U, *}DQ_U(K|_U) \oplus Rj_{V, *}DQ_V(K|_V)
\to Rj_{U \times_X V, *}DQ_{U \times_X V}(K|_{U \times_X V}) \to
$$
이 모든 $K \in D(\mathcal{O}_X)$에 대해 존재한다.
% P09-ERR-0452
이는 Lemma \ref{spaces-perfect-lemma-exact-sequence-j-star}의 구별 삼각형에
완전 함자 $DQ_X$를 적용하고 Lemma
\ref{spaces-perfect-lemma-pushforward-better-coherator}를 세 번 사용하면 따른다.
\end{remark}

\begin{lemma}
\label{spaces-perfect-lemma-boundedness-better-coherator}
$S$를 스킴이라 하자. $X$를 $S$ 위의 준콤팩트 준분리
대수공간이라 하자. Lemma \ref{spaces-perfect-lemma-better-coherator}의 함자
$DQ_X$는 다음 유계성을 갖는다. 정수 $N = N(X)$가 존재하여,
$K$가 $D(\mathcal{O}_X)$에 속하고
$H^i(U, K) = 0$이 $U$가 $X$ 위 아핀 에탈이고
$i \not \in [a, b]$에 대해 성립하면, 코호몰로지 층
$H^i(DQ_X(K))$는 $i \not \in [a, b + N]$에 대해 영이다.
\end{lemma}

\begin{proof}
Lemma \ref{spaces-perfect-lemma-induction-principle}의 귀납 원리를 사용하여 이를
증명하겠다.

\medskip\noindent
$X$가 아핀이면 보조정리는 $N = 0$으로 참이다. 이 경우
$RQ_X = DQ_X$가 $R\Gamma(X, K)$에 결부된 준연접 층의 복합체를
취하는 것으로 주어지기 때문이다. Lemma
\ref{spaces-perfect-lemma-affine-coherator}를 보라.

\medskip\noindent
$(U \subset W, f : V \to W)$를 기본 구별 정사각형이라 하자.
$W$는 준콤팩트 준분리이고, $U \subset W$는 준콤팩트 열린집합이며,
$V$는 아핀이고, 보조정리는 $U$, $V$, $U \times_W V$에 대해
성립한다고 하자. 해당 정수들을 $N(U)$, $N(V)$,
$N(U \times_W V)$라 하자. 이제 $K$가 $D(\mathcal{O}_X)$에 속하고
% P09-ERR-0453
$H^i(W, K) = 0$이 모든 $W$가 $X$ 위 아핀 에탈이고 모든
$i \not \in [a, b]$에 대해 성립한다고 하자. 그러면 $K|_U$,
$K|_V$, $K|_{U \times_W V}$도 같은 성질을 갖는다. 따라서
$RQ_U(K|_U)$와 $RQ_V(K|_V)$, % P09-ERR-0454
$RQ_{U \cap V}(K|_{U \times_W V})$의 코호몰로지 층들은
% P09-ERR-0455
$[a, b + \max(N(U), N(V), N(U \times_W V))$ 밖에서 영이다.
% P09-ERR-0456
함자 $Rj_{U, *}$, $Rj_{V, *}$, $Rj_{U \times_W V, *}$는
Lemma \ref{spaces-perfect-lemma-quasi-coherence-direct-image}에 의해
$D_\QCoh$ 위에서 유한 코호몰로지 차원을 가지므로, 정수 $N$이
존재하여 $Rj_{U, *}DQ_U(K|_U)$, $Rj_{V, *}DQ_V(K|_V)$,
$Rj_{U \cap V, *}DQ_{U \times_W V}(K|_{U \times_W V})$의
코호몰로지 층들은 $[a, b + N]$ 밖에서 영이다. 이제 Remark
\ref{spaces-perfect-remark-explain-consequence}의 구별 삼각형으로 결론을 얻는다.
\end{proof}

\begin{example}
\label{spaces-perfect-example-inverse-limit-quasi-coherent}
$S$를 스킴이라 하자. $X$를 $S$ 위의 준콤팩트 준분리
대수공간이라 하자. $(\mathcal{F}_n)$을 $X$ 위의 준연접 층들의
역계라 하자. $DQ_X$는 우수반이므로 곱과 가환하고, 따라서 유도
극한과도 가환한다. 그러므로
$$
DQ_X(R\lim \mathcal{F}_n) =
(R\lim\text{ in }D_\QCoh(\mathcal{O}_X))(\mathcal{F}_n)
$$
이다. 여기서 첫 번째 $R\lim$은 $D(\mathcal{O}_X)$에서 취한다.
이를 위해 실제로 $K = R\lim \mathcal{F}_n$이라 쓰자. 임의의
$U$가 $X$ 위 아핀 에탈이면
$$
H^i(U, K) =
H^i(R\Gamma(U, R\lim \mathcal{F}_n)) =
H^i(R\lim R\Gamma(U, \mathcal{F}_n)) =
H^i(R\lim \Gamma(U, \mathcal{F}_n))
$$
이다. 이는 코호몰로지가 유도 극한과 가환하고 준연접 층
$\mathcal{F}_n$이 아핀에서 고차 코호몰로지를 갖지 않기 때문이다.
아벨 군의 범주에서 $R\lim$을 계산하면, $H^i(U, K) = 0$은
$i \in [0, 1]$이 아닐 때 성립함을 알 수 있다. 따라서 위의
$R\lim$, 즉
$D_\QCoh(\mathcal{O}_X)$ 안에서 취하여 위에 의해 $DQ_X(K)$인 것은
$D^b_\QCoh(\mathcal{O}_X)$에 속하며,
Lemma \ref{spaces-perfect-lemma-boundedness-better-coherator}에 의해 음의 차수에서
코호몰로지 층들이 영이다.
\end{example}

\section{코호몰로지와 밑변환, IV}
\label{spaces-perfect-section-cohomology-and-base-change-perfect}

\noindent
이 절은 Derived Categories of Schemes, Section
\ref{perfect-section-cohomology-and-base-change-perfect}의 유사이다.

\begin{lemma}
\label{spaces-perfect-lemma-cohomology-base-change}
$S$를 스킴이라 하자. $f : X \to Y$를 $S$ 위 대수공간의
준콤팩트 준분리 사상이라 하자. $E$가
$D_\QCoh(\mathcal{O}_X)$에 속하고 $K$가
$D_\QCoh(\mathcal{O}_Y)$에 속하면
$$
Rf_*(E) \otimes_{\mathcal{O}_Y}^\mathbf{L} K =
Rf_*(E \otimes_{\mathcal{O}_X}^\mathbf{L} Lf^*K)
$$
이다.
\end{lemma}

\begin{proof}
아무 가정 없이도 사상
$Rf_*(E) \otimes_{\mathcal{O}_Y}^\mathbf{L} K \to
Rf_*(E \otimes_{\mathcal{O}_X}^\mathbf{L} Lf^*K)$가 존재한다.
실제로 이는 표준 사상
$$
Lf^*(Rf_*(E) \otimes_{\mathcal{O}_Y}^\mathbf{L} K) =
Lf^*(Rf_*(E)) \otimes_{\mathcal{O}_X}^\mathbf{L} Lf^*K
\longrightarrow
E \otimes_{\mathcal{O}_X}^\mathbf{L} Lf^*K
$$
의 수반인데, 이 표준 사상은 $Lf^*Rf_*E \to E$에서 유도된다.
Cohomology on Sites, Lemmas
\ref{sites-cohomology-lemma-pullback-tensor-product}와
\ref{sites-cohomology-lemma-adjoint}를 보라. 이것이 동형인지 확인할
때 $Y$ 위에서 에탈 국소적으로 작업해도 된다. 따라서 $Y$가 아핀
스킴인 경우로 환원된다.

\medskip\noindent
$K = \bigoplus K_i$가 어떤 복합체들
$K_i \in D_\QCoh(\mathcal{O}_Y)$의 직합이라고 하자. 명제가 각
$K_i$에 대해 성립하면 $K$에 대해서도 성립한다. 실제로 함자
$Lf^*$와 $\otimes^\mathbf{L}$는 구성에 의해 직합을 보존하고,
$Rf_*$는 Lemma \ref{spaces-perfect-lemma-quasi-coherence-pushforward-direct-sums}에
의해 (준연접 코호몰로지 층들을 갖는 복합체들에 대해) 직합과
가환한다. 또한 $K \to L \to M \to K[1]$이 $D_\QCoh(Y)$ 안의
구별 삼각형이라고 하자. 보조정리의 명제가 $K, L, M$ 중 둘에
대해 성립하면 나머지 하나에 대해서도 성립한다 (여기 관련된
함자들은 삼각 범주의 완전 함자들이기 때문이다).

\medskip\noindent
$Y$가 아핀이고 $Y = \Spec(A)$라고 하자. 함자
$\widetilde{\ } : D(A) \to D_\QCoh(\mathcal{O}_Y)$는 Lemma
\ref{spaces-perfect-lemma-derived-quasi-coherent-small-etale-site}와 Derived
Categories of Schemes, Lemma
\ref{perfect-lemma-affine-compare-bounded}에 의해 동치이다.
$T$를 다음과 같은 $K \in D(A)$의 성질이라 하자: 보조정리의 명제가
$\widetilde{K}$에 대해 성립한다. 위 논의와 More on Algebra, Remark
\ref{more-algebra-remark-P-resolution}에 의해 $T$가 $A[k]$에 대해
성립함만 증명하면 충분하다. % P09-ERR-0457
구조층의 이동들에 대해서는 보조정리의 명제가 명백하므로 이것으로
증명이 끝난다.
\end{proof}

\begin{definition}
\label{spaces-perfect-definition-tor-independent}
$S$를 스킴이라 하자. $B$를 $S$ 위의 대수공간이라 하자.
$X$, $Y$를 $B$ 위의 대수공간이라 하자. $X$와 $Y$가
{\it $B$ 위에서 Tor 독립}이라고 하는 것은, 모든 기하점의 가환 도식
$$
\xymatrix{
\Spec(k) \ar[d]_{\overline{y}} \ar[dr]_{\overline{b}} \ar[r]_-{\overline{x}} &
X \ar[d] \\
Y \ar[r] & B
}
$$
에 대하여 환 $\mathcal{O}_{X, \overline{x}}$와
$\mathcal{O}_{Y, \overline{y}}$가
$\mathcal{O}_{B, \overline{b}}$ 위에서 Tor 독립인 것과 같다
(More on Algebra, Definition
\ref{more-algebra-definition-tor-independent}을 보라).
\end{definition}

\noindent
다음 보조정리는 특히 표현가능한 대수공간의 경우에 이 정의가
우리의 정의와 일치함을 보인다.

\begin{lemma}
\label{spaces-perfect-lemma-tor-independent}
$S$를 스킴이라 하자. $B$를 $S$ 위의 대수공간이라 하자.
$X$, $Y$를 $B$ 위의 대수공간이라 하자. 다음은 동치이다.
\begin{enumerate}
\item $X$와 $Y$는 $B$ 위에서 Tor 독립이다.
\item 모든 가환 도식
$$
\xymatrix{
U \ar[d] \ar[r] & W \ar[d] & V \ar[d] \ar[l] \\
X \ar[r] & B & Y \ar[l]
}
$$
에서 수직 화살표들이 에탈이면, $U$와 $V$는 $W$ 위에서 Tor
독립이다. % P09-ERR-0458
\item (2)에서와 같은 어떤 가환 도식에 대하여 (a) $W \to B$는
에탈 전사이고, (b) $U \to X \times_B W$는 에탈 전사이고,
(c) $V \to Y \times_B W$는 에탈 전사이며, 공간 $U$와 $V$는
$W$ 위에서 Tor 독립이다.
\item (3)에서와 같은 어떤 가환 도식에 대하여 $U$, $V$, $W$가
스킴이고, 스킴 $U$와 $V$는 Derived Categories of Schemes,
Definition \ref{perfect-definition-tor-independent}의 의미에서
$W$ 위에서 Tor 독립이다.
\end{enumerate}
\end{lemma}

\begin{proof}
대수공간의 에탈 사상 $\varphi : U \to X$와 기하점
$\overline{u}$에 대하여 국소환의 사상
$\mathcal{O}_{X, \varphi(\overline{u})} \to
\mathcal{O}_{U, \overline{u}}$는 동형이다. 따라서 (1)과 (2)의
동치가 따른다. 함의 (1) $\Rightarrow$ (3)도 마찬가지로 따른다.
(3)을 가정하고 Definition \ref{spaces-perfect-definition-tor-independent}에서와
같은 기하점의 도식을 택하자. 가정에 의해 먼저
$\overline{b}$를 기하점 $\overline{w}$로 올리되 이 점은 $W$에
놓이게 하고, 이어서 기하점 $(\overline{x}, \overline{b})$를
기하점 $\overline{u}$로 올리되 이 점은 $U$에 놓이게 하고,
마지막으로 기하점 $(\overline{y}, \overline{b})$를 기하점
$\overline{v}$로 올리되 이 점은 $V$에 놓이게 할 수 있다.
이러한 올림을 찾으려면 Properties of Spaces,
Lemma \ref{spaces-properties-lemma-geometric-lift-to-usual}를 사용한다.
위의 국소환에 관한 비고를 사용하면 주어진 도식에 대해 정의의
조건이 만족됨을 알 수 있다.

\medskip\noindent
이러한 사전 관찰을 했으므로, (4)는 Definition
\ref{spaces-perfect-definition-tor-independent}이 Derived Categories of Schemes,
Definition \ref{perfect-definition-tor-independent}과 $X$, $Y$,
$B$가 스킴인 경우에 일치한다는 명제로 귀착된다.

\medskip\noindent
$\overline{x}, \overline{b}, \overline{y}$가 Definition
\ref{spaces-perfect-definition-tor-independent}에서와 같고 점 $x, y, b$ 위에
놓인다고 하자. 다음을 상기하자.
$\mathcal{O}_{X, \overline{x}} = \mathcal{O}_{X, x}^{sh}$
(Properties of Spaces, Lemma
\ref{spaces-properties-lemma-describe-etale-local-ring}).
다른 두 환도 마찬가지이다. Algebra, Lemma
\ref{algebra-lemma-strictly-henselian-functorial-improve}에 의해
$\mathcal{O}_{X, \overline{x}}$는
$\mathcal{O}_{X, x} \otimes_{\mathcal{O}_{B, b}}
\mathcal{O}_{B, \overline{b}}$의 엄밀 Hensel화임을 알 수 있다.
특히 환 사상
$$
\mathcal{O}_{X, x} \otimes_{\mathcal{O}_{B, b}} \mathcal{O}_{B, \overline{b}}
\longrightarrow
\mathcal{O}_{X, \overline{x}}
$$
은 평탄하다 (More on Algebra, Lemma
\ref{more-algebra-lemma-dumb-properties-henselization}). More on
Algebra, Lemma \ref{more-algebra-lemma-tor-independent-flat}에 의해
다음을 얻는다.
$$
\text{Tor}_i^{\mathcal{O}_{B, b}}(\mathcal{O}_{X, x}, \mathcal{O}_{Y, y})
\otimes_{\mathcal{O}_{X, x} \otimes_{\mathcal{O}_{B, b}} \mathcal{O}_{Y, y}}
(\mathcal{O}_{X, \overline{x}} \otimes_{\mathcal{O}_{B, \overline{b}}}
\mathcal{O}_{Y, \overline y})
=
\text{Tor}_i^{\mathcal{O}_{B, \overline{b}}}(
\mathcal{O}_{X, \overline{x}}, \mathcal{O}_{Y, \overline{y}})
$$
따라서 스킴으로서 $X$와 $Y$가 $B$ 위에서 Tor 독립이면,
대수공간으로서 $X$와 $Y$도 $B$ 위에서 Tor 독립이다.

\medskip\noindent
역을 위해서는 $X$, $Y$, $B$가 아핀이라고 가정해도 된다.
위의 관찰에 의해 환 사상
$$
\mathcal{O}_{X, x} \otimes_{\mathcal{O}_{B, b}} \mathcal{O}_{Y, y}
\longrightarrow
\mathcal{O}_{X, \overline{x}} \otimes_{\mathcal{O}_{B, \overline{b}}}
\mathcal{O}_{Y, \overline y}
$$
은 평탄하다. 또한 스펙트럼 상의 사상의 상은
$\mathfrak s \subset
\mathcal{O}_{X, x} \otimes_{\mathcal{O}_{B, b}} \mathcal{O}_{Y, y}$이면서
$\mathfrak m_x$와 $\mathfrak m_y$ 위에 놓이는 모든 소아이디얼을
포함한다. 따라서 이것과 위에 표시한 Tor 공식에서, 대수공간으로서
$X$와 $Y$가 $B$ 위에서 Tor 독립이면
$$
\text{Tor}_i^{\mathcal{O}_{B, b}}
(\mathcal{O}_{X, x}, \mathcal{O}_{Y, y})_\mathfrak s = 0
$$
이 모든 $i > 0$과 위와 같은 모든 $\mathfrak s$에 대해 성립함을
알 수 있다. More on Algebra, Lemma
\ref{more-algebra-lemma-tor-independent}을 환 사상
$\Gamma(B, \mathcal{O}_B) \to \Gamma(X, \mathcal{O}_X)$와
$\Gamma(B, \mathcal{O}_B) \to \Gamma(X, \mathcal{O}_X)$에 적용하면,
% P09-ERR-0459
이는 $X$와 $Y$가 $B$ 위에서 Tor 독립임을 함의한다.
\end{proof}

\begin{lemma}
\label{spaces-perfect-lemma-compare-base-change}
$S$를 스킴이라 하자. $g : Y' \to Y$를 $S$ 위 대수공간의
사상이라 하자. $f : X \to Y$를 $S$ 위 대수공간의 준콤팩트
준분리 사상이라 하자. 밑변환 도식
$$
\xymatrix{
X' \ar[r]_{g'} \ar[d]_{f'} &
X \ar[d]^f \\
Y' \ar[r]^g &
Y
}
$$
을 생각하자. $X$와 $Y'$가 $Y$ 위에서 Tor 독립이면, 모든
$E \in D_\QCoh(\mathcal{O}_X)$에 대해
$Rf'_*L(g')^*E = Lg^*Rf_*E$이다.
\end{lemma}

\begin{proof}
임의의 대상 $E$가 $D(\mathcal{O}_X)$에 속할 때 Cohomology on Sites,
Remark \ref{sites-cohomology-remark-base-change}을 사용하여 표준
밑변환 사상 $Lg^*Rf_*E \to Rf'_*L(g')^*E$를 얻는다. 이것이
동형인지 확인할 때 $Y'$ 위에서 에탈 국소적으로 작업해도 된다.
따라서 $g : Y' \to Y$가 아핀 스킴들의 사상이라고 가정해도 된다.
특히 $g$는 아핀이고, 다음 사상이 동형임을 보이면 충분하다.
$$
Rg_*Lg^*Rf_*E \to Rg_*Rf'_*L(g')^*E = Rf_*(Rg'_* L(g')^* E)
$$
Lemma \ref{spaces-perfect-lemma-affine-morphism}를 보라
(대상 $Rf'_*L(g')^*E$와 $Lg^*Rf_*E$가 준연접 코호몰로지 층들을
가짐을 보려면 Lemmas \ref{spaces-perfect-lemma-quasi-coherence-pullback},
\ref{spaces-perfect-lemma-quasi-coherence-tensor-product}, 그리고
\ref{spaces-perfect-lemma-quasi-coherence-direct-image}을 사용한다).
$g'$도 아핀임에 유의하자 (Morphisms of Spaces, Lemma
\ref{spaces-morphisms-lemma-base-change-affine}). Lemma
\ref{spaces-perfect-lemma-affine-morphism-pull-push}에 의해 이 사상은
$$
Rf_*E \otimes_{\mathcal{O}_Y}^\mathbf{L} g_*\mathcal{O}_{Y'}
\longrightarrow
Rf_*(E \otimes_{\mathcal{O}_X}^\mathbf{L} g'_*\mathcal{O}_{X'})
$$
가 된다. $g'_*\mathcal{O}_{X'} = f^*g_*\mathcal{O}_{Y'}$임을
관찰하자. 따라서 Lemma \ref{spaces-perfect-lemma-cohomology-base-change}에 의해
$Lf^*g_*\mathcal{O}_{Y'} = f^*g_*\mathcal{O}_{Y'}$임을 증명하면
충분하다. 이는 $X$와 $Y'$가 $Y$ 위에서 Tor 독립이라는 가정에서
따른다. 실제로 이를 확인할 때 $X$ 위에서 에탈 국소적으로 작업해도
되므로 $X$가 아핀이라고도 가정할 수 있다.
$X = \Spec(A)$, $Y = \Spec(R)$, $Y' = \Spec(R')$라 하자.
우리의 가정은 $A$와 $R'$가 $R$ 위에서 Tor 독립임을 함의한다
(Lemma \ref{spaces-perfect-lemma-tor-independent}와 More on Algebra, Lemma
\ref{more-algebra-lemma-tor-independent}를 보라). 즉
$\text{Tor}_i^R(A, R') = 0$이 $i > 0$에 대해 성립한다. 다시 말해
$A \otimes_R^\mathbf{L} R' = A \otimes_R R'$이고, 이는 정확히
$Lf^*g_*\mathcal{O}_{Y'} = f^*g_*\mathcal{O}_{Y'}$를 뜻한다.
\end{proof}

\noindent
다음 보조정리는 쌍대화 복합체에 관한 장에서 사용될 것이다.

\begin{lemma}
\label{spaces-perfect-lemma-affine-morphism-and-hom-out-of-perfect}
$g : S' \to S$를 아핀 스킴들의 사상이라 하자. 준콤팩트 준분리
대수공간의 카르테시안 정사각형
$$
\xymatrix{
X' \ar[r]_{g'} \ar[d]_{f'} & X \ar[d]^f \\
S' \ar[r]^g & S
}
$$
을 생각하자. $g$와 $f$가 Tor 독립이라고 가정하자.
$S = \Spec(R)$ 및 $S' = \Spec(R')$라 쓰자.
$M, K \in D(\mathcal{O}_X)$에 대하여 표준 사상
$$
R\Hom_X(M, K) \otimes^\mathbf{L}_R R'
\longrightarrow
R\Hom_{X'}(L(g')^*M, L(g')^*K)
$$
은 다음 두 경우에 $D(R')$ 안의 동형이다.
\begin{enumerate}
\item $M \in D(\mathcal{O}_X)$는 완전이고
$K \in D_\QCoh(X)$이다.
\item $M \in D(\mathcal{O}_X)$는 유사연접이고,
$K \in D_\QCoh^+(X)$이며, $R'$는 $R$ 위에서 유한 tor 차원을 갖는다.
\end{enumerate}
\end{lemma}

\begin{proof}
전역 Hom 복합체의 표준 사상
$R\Hom_X(M, K) \to R\Hom_{X'}(L(g')^*M, L(g')^*K)$가
$D(\Gamma(X, \mathcal{O}_X))$ 안에 존재한다.
Cohomology on Sites, Section
\ref{sites-cohomology-section-global-RHom}을 보라. 스칼라를 제한하면
이를 $D(R)$ 안의 사상으로 볼 수 있다. 그러면 제한의 수반성과
$- \otimes_R^\mathbf{L} R'$를 사용하여 보조정리에 표시한 사상을
얻는다. 사상을 정의했으므로, 이것이 아벨 군의 유도 범주에서
동형임을 증명하면 충분하다.

\medskip\noindent
우변은 Lemma \ref{spaces-perfect-lemma-affine-morphism-pull-push}에 의해
$$
R\Hom_X(M, R(g')_*L(g')^*K) =
R\Hom_X(M, K \otimes_{\mathcal{O}_X}^\mathbf{L} g'_*\mathcal{O}_{X'})
$$
와 같다. 두 경우 모두 복합체 $R\SheafHom(M, K)$는 Lemma
\ref{spaces-perfect-lemma-quasi-coherence-internal-hom}에 의해
$D_\QCoh(\mathcal{O}_X)$의 대상이다. 자연스러운 사상
$$
R\SheafHom(M, K) \otimes_{\mathcal{O}_X}^\mathbf{L} g'_*\mathcal{O}_{X'}
\longrightarrow
R\SheafHom(M, K \otimes_{\mathcal{O}_X}^\mathbf{L} g'_*\mathcal{O}_{X'})
$$
이 있고, 이는 두 경우 모두 Lemma
\ref{spaces-perfect-lemma-internal-hom-evaluate-tensor-isomorphism}에 의해
동형이다. % P09-ERR-0461
이 보조정리가 경우 (2)에 적용됨을 보이기 위해
$$
g'_*\mathcal{O}_{X'} = Rg'_*\mathcal{O}_{X'} =
Lf^*g_*\mathcal{O}_X
$$
임에 유의하자. 여기서 두 번째 등식은 Lemma
\ref{spaces-perfect-lemma-compare-base-change}에 의한다. % P09-ERR-0462
Derived Categories of Schemes, Lemma
\ref{perfect-lemma-tor-dimension-affine}, Lemma
\ref{spaces-perfect-lemma-descend-tor-amplitude}, 그리고 Cohomology on Sites,
Lemma \ref{sites-cohomology-lemma-tor-amplitude-pullback}을 사용하면
$g'_*\mathcal{O}_{X'}$가 유한 Tor 차원을 가짐을 알 수 있다.
따라서 두 경우 모두 $K$를 $R\SheafHom(M, K)$로 바꾸면 다음이
동형임을 증명하는 문제로 환원된다.
$$
R\Gamma(X, K) \otimes^\mathbf{L}_A A' \longrightarrow
R\Gamma(X, K \otimes^\mathbf{L}_{\mathcal{O}_X} g'_*\mathcal{O}_{X'})
$$
% P09-ERR-0463
좌변은 Lemma \ref{spaces-perfect-lemma-affine-morphism-pull-push}에 의해
$R\Gamma(X', L(g')^*K)$와 같다. 따라서 결과는 Lemma
\ref{spaces-perfect-lemma-compare-base-change}에서 따른다.
\end{proof}

\begin{remark}
\label{spaces-perfect-remark-multiplication-map}
Lemma \ref{spaces-perfect-lemma-affine-morphism-and-hom-out-of-perfect}의 표기를
사용하자. % P09-ERR-0460
도식
$$
\xymatrix{
R\Hom_X(M, Rg'_*L) \otimes_R^\mathbf{L} R' \ar[r] \ar[d]_\mu &
R\Hom_{X'}(L(g')^*M, L(g')^*Rg'_*L) \ar[d]^a \\
R\Hom_X(M, R(g')_*L) \ar@{=}[r] &
R\Hom_{X'}(L(g')^*M, L)
}
$$
은 가환한다. 여기서 위쪽 수평 화살표는 보조정리의 사상이고,
$\mu$는 곱셈 사상이며, $a$는 수반 사상
$L(g')^*Rg'_*L \to L$에서 유도된다. 곱셈 사상은
$K' \otimes_R^\mathbf{L} R' \to K'$라는 수반 사상이며, 이는 임의의
$K' \in D(R')$에 대해 취한 것이다.
\end{remark}

\begin{lemma}
\label{spaces-perfect-lemma-tor-independence-and-tor-amplitude}
$S$를 스킴이라 하자. 대수공간의 카르테시안 정사각형
$$
\xymatrix{
X' \ar[r]_{g'} \ar[d]_{f'} & X \ar[d]^f \\
Y' \ar[r]^g & Y
}
$$
을 $S$ 위에서 생각하자. $g$와 $f$가 Tor 독립이라고 가정하자.
\begin{enumerate}
\item $E \in D(\mathcal{O}_X)$가 $[a, b]$ 안의 tor 진폭을
$f^{-1}\mathcal{O}_Y$-가군의 복합체로서 가지면,
$L(g')^*E$는 $[a, b]$ 안의 tor 진폭을
$f^{-1}\mathcal{O}_{Y'}$-가군의 복합체로서 갖는다.
% P09-ERR-0464
\item $\mathcal{G}$가 $\mathcal{O}_X$-가군이고 $Y$ 위에서 평탄하면
$L(g')^*\mathcal{G} = (g')^*\mathcal{G}$이다.
\end{enumerate}
\end{lemma}

\begin{proof}
tor 차원은 줄기에서 계산할 수 있다. Cohomology on Sites, Lemma
\ref{sites-cohomology-lemma-tor-amplitude-stalk}와 Properties of
Spaces, Theorem \ref{spaces-properties-theorem-exactness-stalks}를
보라. $\overline{x}'$가 $X'$의 기하점이고 그 상이
$\overline{x}$이며 이 점이 $X$에 놓이면,
$$
(L(g')^*E)_{\overline{x}'} =
E_{\overline{x}}
\otimes_{\mathcal{O}_{X, \overline{x}}}^\mathbf{L}
\mathcal{O}_{X', \overline{x}'}
$$
이다. $\overline{y}'$가 $Y'$에 놓이고 $\overline{y}$가 $Y$에
놓이며 이들이 각각 $\overline{x}'$와 $\overline{x}$의 상이라고
하자. $X$와 $Y'$가
$Y$ 위에서 tor 독립이므로 More on Algebra, Lemma
\ref{more-algebra-lemma-base-change-comparison}을 적용하여 표시된
공식의 우변이
$E_{\overline{x}}
\otimes_{\mathcal{O}_{Y, \overline{y}}}^\mathbf{L}
\mathcal{O}_{Y', \overline{y}'}$와
$D(\mathcal{O}_{Y', \overline{y}'})$ 안에서 같음을 알 수 있다.
따라서 (1)은 More on Algebra, Lemma
\ref{more-algebra-lemma-pull-tor-amplitude}에서 따른다.
(2)를 보려면 $\mathcal{G}$의 평탄성이 $\mathcal{G}[0]$이
$[0, 0]$ 안의 tor 진폭을 갖는 조건과 동치임을 관찰하자.
(1)을 적용하면 결론을 얻는다.
\end{proof}

\section{코호몰로지와 밑변환, V}
\label{spaces-perfect-section-cohomology-base-change}

\noindent
이 절은 Derived Categories of Schemes, Section
\ref{perfect-section-cohomology-base-change}의 유사이다. Section
\ref{spaces-perfect-section-cohomology-and-base-change-perfect}에서 우리는 사상들이
tor 독립이면 밑변환 정리가 성립함을 보았다. 그러한 조건 없이는
아핀인 경우에도 밑변환 정리가 성립할 수 없다. More on Algebra,
Section \ref{more-algebra-section-tor-independence}를 보라. 이 절에서는
“한 번에 하나의 복합체”에 대해 밑변환 결과를 얻을 수 있는 때를
분석한다.

\medskip\noindent
이를 가능하게 하기 위해 $S$를 밑 스킴이라 하고 가환 도식
$$
\xymatrix{
X' \ar[r]_{g'} \ar[d]_{f'} &
X \ar[d]^f \\
Y' \ar[r]^g &
Y
}
$$
을 생각하자. 이는 $S$ 위 대수공간의 도식이다 (보통 우리는 이것이
카르테시안이라고 가정할 것이다). $K \in D_\QCoh(\mathcal{O}_X)$라
하고, $L(g')^*K \to K'$를 $D_\QCoh(\mathcal{O}_{X'})$ 안의
사상이라 하자. 기하점 $\overline{x}'$가 $X'$에 놓일 때
기하점들
$\overline{x} = g'(\overline{x}')$,
$\overline{y}' = f'(\overline{x}')$,
$\overline{y} = f(\overline{x}) = g(\overline{y}')$를 생각하자.
이들은 각각 $X$, $Y'$, $Y$에 놓인다. 그러면 사상들
$$
K_{\overline{x}}
\otimes_{\mathcal{O}_{Y, \overline{y}}}^\mathbf{L}
\mathcal{O}_{Y', \overline{y}'}
\to
K_{\overline{x}}
\otimes_{\mathcal{O}_{X, \overline{x}}}^\mathbf{L}
\mathcal{O}_{X', \overline{x}'} \to
K'_{\overline{x}'}
$$
을 생각할 수 있다. 여기서 첫 화살표는 More on Algebra,
Equation (\ref{more-algebra-equation-comparison-map})이고, 둘째
화살표는
$(L(g')^*K)_{\overline{x}'} =
K_{\overline{x}} \otimes_{\mathcal{O}_{X, \overline{x}}}^\mathbf{L}
\mathcal{O}_{X', \overline{x}'}$와 주어진 사상
$L(g')^*K \to K'$에서 유도된다. 각 $i \in \mathbf{Z}$에 대해
$\mathcal{O}_{X, \overline{x}} \otimes_{\mathcal{O}_{Y, \overline{y}}}
\mathcal{O}_{Y', \overline{y}'}$-가군 구조가
$H^i(K_{\overline{x}}
\otimes_{\mathcal{O}_{Y, \overline{y}}}^\mathbf{L}
\mathcal{O}_{Y', \overline{y}'})$ 위에 생긴다. 모두 합치면 표준
사상
\begin{equation}
\label{spaces-perfect-equation-bc}
H^i(K_{\overline{x}}
\otimes_{\mathcal{O}_{Y, \overline{y}}}^\mathbf{L}
\mathcal{O}_{Y', \overline{y}'})
\otimes_{(\mathcal{O}_{X, \overline{x}}
\otimes_{\mathcal{O}_{Y, \overline{y}}}
\mathcal{O}_{Y', \overline{y}'})}
\mathcal{O}_{X', \overline{x}'}
\longrightarrow
H^i(K'_{\overline{x}'})
\end{equation}
을 $\mathcal{O}_{X', \overline{x}'}$-가군의 사상으로 얻는다.

\begin{lemma}
\label{spaces-perfect-lemma-single-complex-base-change-condition}
$S$를 스킴이라 하자. 대수공간의 카르테시안 도식
$$
\xymatrix{
X' \ar[r]_{g'} \ar[d]_{f'} &
X \ar[d]^f \\
Y' \ar[r]^g &
Y
}
$$
을 $S$ 위에서 생각하자. $K \in D_\QCoh(\mathcal{O}_X)$라 하고
$L(g')^*K \to K'$를 $D_\QCoh(\mathcal{O}_{X'})$ 안의 사상이라
하자. 다음은 동치이다.
\begin{enumerate}
\item 임의의 $x' \in X'$와 $i \in \mathbf{Z}$에 대해 사상
(\ref{spaces-perfect-equation-bc})는 동형이다. % P09-ERR-0465
\item 가환 도식
$$
\xymatrix{
& U \ar[d] \ar[rd]^a \\
V' \ar[r] \ar[rd]^c & V \ar[rd]^b & X \ar[d]^f \\
& Y' \ar[r]^g & Y
}
$$
에서 $a, b, c$가 에탈이고, $U, V, V'$가 스킴이며
$U' = V' \times_V U$일 때, Derived Categories of Schemes, Lemma
\ref{perfect-lemma-single-complex-base-change-condition}의 동치인
조건들이 $(U \to X)^*K$와 $(U' \to X')^*K'$에 대해 성립한다.
\item (2)에서와 같은 어떤 도식에서 $U' \to X'$가 전사이다.
\end{enumerate}
\end{lemma}

\begin{proof}
(1)은 $X'$ 위에서 에탈 국소적임을 관찰하자. 알려진 형식적
함의들을 따라가면 다음을 증명하면 충분함을 알 수 있다:
$X, Y, Y', X'$가 각각 스킴 $X_0, Y_0, Y'_0, X'_0$으로
표현될 때, 이 보조정리의 조건 (1)은 Derived Categories of
Schemes, Lemma
\ref{perfect-lemma-single-complex-base-change-condition}의 조건
(1)과 동치이다. $f_0, g_0, g'_0, f'_0$로 이 스킴들 사이의
사상들 가운데 각각 $f, g, g', f'$에 대응하는 것을 표기하자.
% P09-ERR-0466
Lemma \ref{spaces-perfect-lemma-derived-quasi-coherent-small-etale-site}에 의해
$K = \epsilon^*K_0$ 및 $K' = \epsilon^*K'_0$라고 가정해도 된다.
여기서 $K_0 \in D_\QCoh(\mathcal{O}_{X_0})$와
$K'_0 \in D_\QCoh(\mathcal{O}_{X'_0})$는 어떤 대상들이다.
또한 사상 $Lg^*K \to K'$는 Remark
\ref{spaces-perfect-remark-match-total-direct-images}의 표기에서 어떤 사상
$L(g_0)^*K_0 \to K'_0$의 당김이다. % P09-ERR-0467
$\mathcal{O}_{X, \overline{x}}$가
$\mathcal{O}_{X, x}$의 엄밀 Hensel화임을 상기하자 (Properties of
Spaces, Lemma \ref{spaces-properties-lemma-describe-etale-local-ring}).
또한
$$
K_{\overline{x}} =
K_{0, x} \otimes_{\mathcal{O}_{X, x}}^\mathbf{L} \mathcal{O}_{X, \overline{x}}
\quad\text{and}\quad
K'_{\overline{x}'} =
K'_{0, x'} \otimes_{\mathcal{O}_{X', x'}}^\mathbf{L}
\mathcal{O}_{X', \overline{x}'}
$$
이다 (Properties of Spaces, Lemma
\ref{spaces-properties-lemma-stalk-quasi-coherent}와 비슷하다).
가환 도식
$$
\xymatrix{
H^i(K_{\overline{x}}
\otimes_{\mathcal{O}_{Y, \overline{y}}}^\mathbf{L}
\mathcal{O}_{Y', \overline{y}'})
\otimes_{(\mathcal{O}_{X, \overline{x}}
\otimes_{\mathcal{O}_{Y, \overline{y}}}
\mathcal{O}_{Y', \overline{y}'})}
\mathcal{O}_{X', \overline{x}'}
\ar[r] &
H^i(K'_{\overline{x}'}) \\
H^i(K_{0, x} \otimes_{\mathcal{O}_{Y, y}}^\mathbf{L} \mathcal{O}_{Y', y'})
\otimes_{(\mathcal{O}_{X, x} \otimes_{\mathcal{O}_{Y, y}} \mathcal{O}_{Y', y'})}
\mathcal{O}_{X', x'}
\ar[u] \ar[r] &
H^i(K'_{0, x'}) \ar[u]
}
$$
을 생각하자. 아래쪽 수평 화살표가 동형일 필요충분조건이 위쪽
수평 화살표가 동형인 것임을 보여야 한다.
$\mathcal{O}_{X', x'} \to \mathcal{O}_{X', \overline{x}'}$가 충실히
평탄하므로 (More on Algebra, Lemma
\ref{more-algebra-lemma-dumb-properties-henselization}), 위쪽 화살표가
이 사상에 의한 아래쪽 화살표의 밑변환임을 보이면 충분하다.
오른쪽의 대상들에 대해서는 위에서 준 줄기 사이의 관계에서 곧바로
따른다. 왼쪽의 대상들에 대해서는 다음을 보이면 충분하다.
\begin{align*}
&
H^i\left(
(K_{0, x}
\otimes_{\mathcal{O}_{X, x}}^\mathbf{L} \mathcal{O}_{X, \overline{x}})
\otimes_{\mathcal{O}_{Y, \overline{y}}}^\mathbf{L}
\mathcal{O}_{Y', \overline{y}'}\right) \\
& =
H^i(K_{0, x} \otimes_{\mathcal{O}_{Y, y}}^\mathbf{L} \mathcal{O}_{Y', y'})
\otimes_{(\mathcal{O}_{X, x} \otimes_{\mathcal{O}_{Y, y}} \mathcal{O}_{Y', y'})}
(\mathcal{O}_{X, \overline{x}}
\otimes_{\mathcal{O}_{Y, \overline{y}}}
\mathcal{O}_{Y', \overline{y}'})
\end{align*}
이는 More on Algebra, Lemma
\ref{more-algebra-lemma-lemma-tor-independent-flat-compare}에서 따른다.
이 보조정리의 평탄성 가정들은 위에서 말한 사실들과 Algebra, Lemma
\ref{algebra-lemma-strictly-henselian-functorial-improve}에 의해
성립한다. 마지막 보조정리에 의해
$\mathcal{O}_{X, \overline{x}}$는
$\mathcal{O}_{X, x} \otimes_{\mathcal{O}_{Y, y}}
\mathcal{O}_{Y, \overline{y}}$의 엄밀 Hensel화이고,
$\mathcal{O}_{Y', \overline{y}'}$는
$\mathcal{O}_{Y', y'} \otimes_{\mathcal{O}_{Y, y}}
\mathcal{O}_{Y, \overline{y}}$의 엄밀 Hensel화이다.
\end{proof}

\begin{lemma}
\label{spaces-perfect-lemma-single-complex-base-change}
$S$를 스킴이라 하자. 대수공간의 카르테시안 도식
$$
\xymatrix{
X' \ar[r]_{g'} \ar[d]_{f'} &
X \ar[d]^f \\
Y' \ar[r]^g &
Y
}
$$
을 $S$ 위에서 생각하자. $K \in D_\QCoh(\mathcal{O}_X)$라 하고
$L(g')^*K \to K'$를 $D_\QCoh(\mathcal{O}_{X'})$ 안의 사상이라
하자. 다음을 가정하자.
\begin{enumerate}
\item Lemma \ref{spaces-perfect-lemma-single-complex-base-change-condition}의 동치인
조건들이 성립한다.
\item $f$는 준콤팩트 준분리이다.
\end{enumerate}
그러면 합성
$Lg^*Rf_*K \to Rf'_*L(g')^*K \to Rf'_*K'$는 동형이다.
\end{lemma}

\begin{proof}
사상이 동형인지 확인할 때 $Y'$ 위에서 에탈 국소적으로 작업해도
된다. 따라서 $g : Y' \to Y$가 아핀 스킴들의 사상이라고 가정해도
된다. 이 경우 Lemma \ref{spaces-perfect-lemma-induction-principle}의 귀납 원리를
사용하여 다음을 증명하겠다. $U$가 $X$ 위 에탈인 준콤팩트 준분리
대수공간이면, 같은 방식으로 구성한 사상
$Lg^*R(U \to Y)_*K|_U \to R(U' \to Y')_*K'|_{U'}$는 동형이다.
여기서
$U' = X' \times_{g', X} U = Y' \times_{g, Y} U$이다.

\medskip\noindent
$U$가 스킴이면 (예를 들어 아핀이면) 결과가 성립한다. 실제로 이
경우 $Y, Y', U, U'$는 스킴이고, $K$와 $K'$는 Lemma
\ref{spaces-perfect-lemma-derived-quasi-coherent-small-etale-site}에 의해 밑에 놓인
스킴들의 유도 범주의 대상에서 유도된다. 또한 Derived Categories
of Schemes, Lemma
\ref{perfect-lemma-single-complex-base-change-condition}의 조건은
Lemma \ref{spaces-perfect-lemma-single-complex-base-change-condition}에 의해 이
복합체들에 대해 성립한다. 따라서 Remark
\ref{spaces-perfect-remark-match-total-direct-images}에서 설명한 호환성에 의해,
스킴의 경우의 결과인 Derived Categories of Schemes, Lemma
\ref{perfect-lemma-single-complex-base-change}를 적용할 수 있다.

\medskip\noindent
귀납 단계를 보자. $(U \subset W, V \to W)$를 기본 구별 정사각형이라
하자. 여기서 $W$는 $X$ 위 에탈인 준콤팩트 준분리 대수공간이고,
$U$는 준콤팩트이며, $V$는 아핀이고, 결과는 $U$, $V$,
$U \times_W V$에 대해 성립한다고 하자. 표기를 간단히 하기 위해
$W$를 $X$로 바꾸어 쓰겠다 (이 시점에서는 허용된다).
% P09-ERR-0468
$a : U \to Y$, $b : V \to Y$, $c : U \times_X V \to Y$로 자명한
사상들을 표기하자. % P09-ERR-0469
$a' : U' \to Y'$, $b' : V' \to Y'$와
$c' : U' \times_{X'} V' \to Y'$를 각각 $a$, $b$, $c$의
밑변환이라 하자. 상대 Mayer--Vietoris의 구별 삼각형들
(Lemma \ref{spaces-perfect-lemma-unbounded-relative-mayer-vietoris})을 사용하면
가환 도식
$$
\xymatrix{
Lg^*Rf_*K \ar[r] \ar[d] & Rf'_*K' \ar[d] \\
Lg^*Ra_*K|_U \oplus Lg^*Rb_*K|_V \ar[r] \ar[d] &
Ra'_* K'|_{U'} \oplus Rb'_* K'|_{V'} \ar[d] \\
Lg^*Rc_*K|_{U \times_X V} \ar[r] \ar[d] &
Rc'_*K'|_{U' \times_{X'} V'} \ar[d] \\
Lg^*Rf_* K[1] \ar[r] &
Rf'_* K'[1]
}
$$
을 얻는다. 둘째와 셋째 수평 화살표가 동형이므로 첫째도 동형이다
(Derived Categories, Lemma
\ref{derived-lemma-third-isomorphism-triangle}). 이것으로 보조정리의
증명이 끝난다.
\end{proof}

\begin{lemma}
\label{spaces-perfect-lemma-single-complex-base-change-condition-inherited}
$S$를 스킴이라 하자. 대수공간의 카르테시안 도식
$$
\xymatrix{
X' \ar[r]_{g'} \ar[d]_{f'} &
X \ar[d]^f \\
S' \ar[r]^g &
S
}
$$
을 $S$ 위에서 생각하자. $K \in D_\QCoh(\mathcal{O}_X)$라 하고
$L(g')^*K \to K'$를 $D_\QCoh(\mathcal{O}_{X'})$ 안의 사상이라
하자. Lemma \ref{spaces-perfect-lemma-single-complex-base-change-condition}의 동치인
조건들이 성립하면 다음이 참이다.
\begin{enumerate}
\item $E \in D_\QCoh(\mathcal{O}_X)$에 대하여 Lemma
\ref{spaces-perfect-lemma-single-complex-base-change-condition}의 동치인 조건들은
$L(g')^*(E \otimes^\mathbf{L} K) \to
L(g')^*E \otimes^\mathbf{L} K'$에 대해 성립한다.
\item $E$가 $D(\mathcal{O}_X)$ 안의 완전 대상이면, Lemma
\ref{spaces-perfect-lemma-single-complex-base-change-condition}의 동치인 조건들은
$L(g')^*R\SheafHom(E, K) \to R\SheafHom(L(g')^*E, K')$에 대해
성립한다. % P09-ERR-0470
\item $K$가 아래로 유계이고 $E$가 $D(\mathcal{O}_X)$ 안의
유사연접 대상이면, Lemma
\ref{spaces-perfect-lemma-single-complex-base-change-condition}의 동치인 조건들은
$L(g')^*R\SheafHom(E, K) \to R\SheafHom(L(g')^*E, K')$에 대해
성립한다. % P09-ERR-0471
\end{enumerate}
\end{lemma}

\begin{proof}
관련된 복합체들은 Lemmas
\ref{spaces-perfect-lemma-quasi-coherence-pullback},
\ref{spaces-perfect-lemma-quasi-coherence-tensor-product},
\ref{spaces-perfect-lemma-quasi-coherence-internal-hom}과 Cohomology on Sites,
Lemmas \ref{sites-cohomology-lemma-pseudo-coherent-pullback},
\ref{sites-cohomology-lemma-perfect-pullback}에 의해 준연접
코호몰로지 층들을 가지므로 명제는 뜻이 있다. 이제 사상들
(\ref{spaces-perfect-equation-bc})이 (1)의 경우에 동형임은
(\ref{spaces-perfect-equation-bc})의 원천과 표적을 텐서곱의 추이성을 사용하여
계산하면 확인할 수 있다. More on Algebra, Lemma
\ref{more-algebra-lemma-triple-tensor-product}을 보라.
(2)와 (3)의 사상은 합성
$$
L(g')^*R\SheafHom(E, K) \to R\SheafHom(L(g')^*E, L(g')^*K)
\to R\SheafHom(L(g')^*E, K')
$$
이다. 여기서 첫 화살표는 Cohomology on Sites, Remark
\ref{sites-cohomology-remark-prepare-fancy-base-change}이고, 둘째
화살표는 주어진 사상 $L(g')^*K \to K'$에서 유도된다. 사상들
(\ref{spaces-perfect-equation-bc})이 동형임을 증명하기 위해, 경우 (2)에는
$E_x$를 유한 생성 사영 $\mathcal{O}_{X. x}$-가군의 유계 복합체로
표현하고, 경우 (3)에는 유한 생성 자유 가군의 위로 유계 복합체로
표현한 뒤 화살표의 원천과 표적을 계산한다. % P09-ERR-0472
몇몇 세부사항은 생략한다.
\end{proof}

\begin{lemma}
\label{spaces-perfect-lemma-base-change-tensor}
$S$를 스킴이라 하자. $f : X \to Y$를 $S$ 위 대수공간의
준콤팩트 준분리 사상이라 하자.
$E \in D_\QCoh(\mathcal{O}_X)$라 하자.
$\mathcal{G}^\bullet$을 준연접 $\mathcal{O}_X$-가군의 위로 유계
복합체로서 $Y$ 위에서 평탄한 것이라 하자. 그러면
$$
Rf_*(E \otimes^\mathbf{L}_{\mathcal{O}_X} \mathcal{G}^\bullet)
$$
의 형성은 임의의 밑변환과 가환한다 (정확한 진술은 증명을 보라).
\end{lemma}

\begin{proof}
명제의 뜻은 다음과 같다. $g : Y' \to Y$를 대수공간의 사상이라
하고 밑변환 도식
$$
\xymatrix{
X' \ar[r]_{g'} \ar[d]_{f'} &
X \ar[d]^f \\
Y' \ar[r]^g &
Y
}
$$
을 생각하자. 다시 말해 $X' = Y' \times_Y X$이다. 보조정리는
$$
Lg^*Rf_*(E \otimes^\mathbf{L}_{\mathcal{O}_X} \mathcal{G}^\bullet)
\longrightarrow
Rf'_*(L(g')^*E \otimes^\mathbf{L}_{\mathcal{O}_{X'}} (g')^*\mathcal{G}^\bullet)
$$
가 동형이라고 주장한다. 우변에서는
$\mathcal{G}^\bullet$에 유도 당김을 사용하지 \textbf{않음}을
관찰하자. 이를 증명하기 위해 Lemmas
\ref{spaces-perfect-lemma-single-complex-base-change}와
\ref{spaces-perfect-lemma-single-complex-base-change-condition-inherited}을 적용하면,
표준 사상
$$
L(g')^*\mathcal{G}^\bullet \to (g')^*\mathcal{G}^\bullet
$$
이 Lemma \ref{spaces-perfect-lemma-single-complex-base-change-condition}의 동치인
조건들을 만족함을 증명하면 충분하다. 이는 줄기에서 조건을
확인하면 따른다. 그곳에서는
$\mathcal{G}^\bullet_{\overline{x}}
\otimes_{\mathcal{O}_{Y, \overline{y}}}
\mathcal{O}_{Y', \overline{y}'}$가 유도 텐서곱을 계산한다는
사실에서 곧바로 따른다. 이는 복합체 $\mathcal{G}^\bullet$에 대한
가정 때문이다.
\end{proof}

\begin{lemma}
\label{spaces-perfect-lemma-base-change-RHom}
$S$를 스킴이라 하자. $f : X \to Y$를 $S$ 위 대수공간의
준콤팩트 준분리 사상이라 하자. $E$를 $D(\mathcal{O}_X)$의
대상이라 하자. $\mathcal{G}^\bullet$을 준연접
$\mathcal{O}_X$-가군의 복합체라 하자. 다음 중 하나를 가정하자.
\begin{enumerate}
\item $E$는 완전이고, $\mathcal{G}^\bullet$은 위로 유계이며,
$\mathcal{G}^n$은 $Y$ 위에서 평탄하다. % P09-ERR-0473
\item $E$는 유사연접이고, $\mathcal{G}^\bullet$은 유계이며,
$\mathcal{G}^n$은 $Y$ 위에서 평탄하다.
\end{enumerate}
그러면
$$
Rf_*R\SheafHom(E, \mathcal{G}^\bullet)
$$
의 형성은 임의의 밑변환과 가환한다 (정확한 진술은 증명을 보라).
\end{lemma}

\begin{proof}
명제의 뜻은 다음과 같다. $g : Y' \to Y$를 대수공간의 사상이라
하고 밑변환 도식
$$
\xymatrix{
X' \ar[r]_h \ar[d]_{f'} &
X \ar[d]^f \\
Y' \ar[r]^g &
Y
}
$$
을 생각하자. % P09-ERR-0474
다시 말해 $X' = Y' \times_Y X$이다. 보조정리는
$$
Lg^*Rf_*R\SheafHom(E, \mathcal{G}^\bullet)
\longrightarrow
R(f')_*R\SheafHom(L(g')^*E, (g')^*\mathcal{G}^\bullet)
$$
가 동형이라고 주장한다. 우변에서는
$\mathcal{G}^\bullet$에 유도 당김을 사용하지 \textbf{않음}을
관찰하자. 이를 증명하기 위해 Lemmas
\ref{spaces-perfect-lemma-single-complex-base-change}와
\ref{spaces-perfect-lemma-single-complex-base-change-condition-inherited}을 적용하면,
표준 사상
$$
L(g')^*\mathcal{G}^\bullet \to (g')^*\mathcal{G}^\bullet
$$
이 Lemma \ref{spaces-perfect-lemma-single-complex-base-change-condition}의 동치인
조건들을 만족함을 증명하면 충분하다. 이는 Lemma
\ref{spaces-perfect-lemma-base-change-tensor}의 증명에서 보였다.
\end{proof}















\section{완전 복합체의 구성}
\label{spaces-perfect-section-producing-perfect}

\noindent
다음 보조정리는 완전 복합체를 구성하는 데 쓰이는 주된 기술적
도구이다. 이 결과의 후속 판본들은 뇌터 근사를 통해 이 결과로
귀착될 것이다.

\begin{lemma}
\label{spaces-perfect-lemma-perfect-direct-image}
$S$를 스킴이라 하자. $Y$를 $S$ 위의 뇌터 대수공간이라 하자.
$f : X \to Y$를 국소 유한형이고 준분리인 대수공간의 사상이라
하자. $E \in D(\mathcal{O}_X)$가 다음 조건들을 만족한다고 하자.
\begin{enumerate}
\item $E \in D^b_{\textit{Coh}}(\mathcal{O}_X)$이다.
\item $H^i(E)$의 지지가 $Y$ 위에서 고유인 것은 모든 $i$에 대하여
성립한다.
\item $E$는 $D(f^{-1}\mathcal{O}_Y)$의 대상으로서 유한 Tor 차원을
가진다.
\end{enumerate}
그러면 $Rf_*E$는 $D(\mathcal{O}_Y)$의 완전 대상이다.
\end{lemma}

\begin{proof}
Lemma \ref{spaces-perfect-lemma-direct-image-coherent}에 의해 $Rf_*E$는
$D^b_{\textit{Coh}}(\mathcal{O}_Y)$의 대상이다. 따라서
$Rf_*E$는 유사연접이다
(Lemma \ref{spaces-perfect-lemma-identify-pseudo-coherent-noetherian}).
그러므로 $Rf_*E$가 유한 Tor 차원을 가짐을 보이면 충분하다.
Cohomology on Sites, Lemma
\ref{sites-cohomology-lemma-perfect}을 보라.
Lemma \ref{spaces-perfect-lemma-tor-qc-qs}에 의해
$Rf_*(E) \otimes_{\mathcal{O}_Y}^\mathbf{L} \mathcal{F}$가 모든
준연접 층 $\mathcal{F}$에 대하여 $Y$ 위에서 보편적으로 유계인
코호몰로지를 가짐을 확인하면 충분하다.
$Rf_*(E)$가 위로 유계이므로 위로 유계라는 것은 명백하다.
$T \subset |X|$를 $H^i(E)$의 지지들을 모든 $i$에 걸쳐 합친
합집합이라 하자. 그러면 가정 (1), (2)와 Lemma
\ref{spaces-perfect-lemma-union-closed-proper-over-base}에 의해 $T$는 $Y$ 위에서
고유이다. 특히 $X' \subset X$인 준콤팩트 열린 부분공간이
존재하며, 이것은 $T$를 포함한다. $f' = f|_{X'}$로 놓으면
$Rf_*(E) = Rf'_*(E|_{X'})$이다. 이는 $E$가 $X \setminus T$ 위에서
영으로 제한되기 때문이다. 따라서 $X$를 $X'$로 바꾸어
$f$가 준콤팩트라고 가정할 수 있다. $f$가 준분리라고 이미
가정하였다. 그러므로
$$
Rf_*(E) \otimes_{\mathcal{O}_Y}^\mathbf{L} \mathcal{F} =
Rf_*\left(E \otimes_{\mathcal{O}_X}^\mathbf{L} Lf^*\mathcal{F}\right) =
Rf_*\left(E \otimes_{f^{-1}\mathcal{O}_Y}^\mathbf{L} f^{-1}\mathcal{F}\right)
$$
이다. 이는 Lemma \ref{spaces-perfect-lemma-cohomology-base-change}와 Cohomology on
Sites, Lemma
\ref{sites-cohomology-lemma-variant-derived-pullback}에서 따른다.
가정 (3)에 의해 복합체
$E \otimes_{f^{-1}\mathcal{O}_Y}^\mathbf{L} f^{-1}\mathcal{F}$의
코호몰로지 층들은 어떤 주어진 유한 범위, 이를테면 $[a, b]$ 안에
놓인다. 그러면 여기에 $Rf_*$를 적용한 것은 범위 $[a, \infty)$
안에 코호몰로지를 가지며, 이로써 결론을 얻는다.
\end{proof}

\begin{lemma}
\label{spaces-perfect-lemma-tensor-perfect}
$S$를 스킴이라 하자. $B$를 $S$ 위의 뇌터 대수공간이라 하자.
$f : X \to B$를 국소 유한형이고 준분리인 대수공간의 사상이라
하자. $E \in D(\mathcal{O}_X)$를 완전 대상이라 하자.
$\mathcal{G}^\bullet$을 연접 $\mathcal{O}_X$-가군의 유계
복합체로서 $B$ 위에서 평탄하고 지지가 $B$ 위에서 고유인 것이라
하자. 그러면
$K = Rf_*(E \otimes^\mathbf{L}_{\mathcal{O}_X} \mathcal{G}^\bullet)$는
$D(\mathcal{O}_B)$의 완전 대상이다.
\end{lemma}

\begin{proof}
Lemma \ref{spaces-perfect-lemma-perfect-direct-image}에 의해 대상 $K$는 완전이다.
그 보조정리를 적용할 수 있음을 확인하자. 국소적으로 $E$는 유한
생성 자유 $\mathcal{O}_X$-가군의 유한 복합체와 동형이다. 따라서
국소적으로
$E \otimes^\mathbf{L}_{\mathcal{O}_X} \mathcal{G}^\bullet$은 그 항들이
$$
\bigoplus\nolimits_{i = a, \ldots, b} (\mathcal{G}^i)^{\oplus r_i}
$$
꼴인 유한 복합체와 동형이다. 여기서
$a, b, r_a, \ldots, r_b$는 어떤 정수들이다. 이로부터 코호몰로지
층들
$H^i(E \otimes^\mathbf{L}_{\mathcal{O}_X} \mathcal{G})$이 연접임이
곧바로 따른다. % P09-ERR-0475
또한 $\mathcal{G}^i$가 $f^{-1}\mathcal{O}_Y$ 위에서 평탄하므로
Tor 차원에 관한 가정도 따른다. % P09-ERR-0476
\end{proof}

\begin{lemma}
\label{spaces-perfect-lemma-ext-perfect}
$S$를 스킴이라 하자. $B$를 $S$ 위의 뇌터 대수공간이라 하자.
$f : X \to B$를 국소 유한형이고 준분리인 대수공간의 사상이라
하자. $E \in D(\mathcal{O}_X)$를 완전 대상이라 하자.
$\mathcal{G}^\bullet$을 연접 $\mathcal{O}_X$-가군의 유계
복합체로서 $B$ 위에서 평탄하고 지지가 $B$ 위에서 고유인 것이라
하자. 그러면
$K = Rf_*R\SheafHom(E, \mathcal{G})$는 $D(\mathcal{O}_B)$의
완전 대상이다. % P09-ERR-0477
\end{lemma}

\begin{proof}
$E$가 완전 복합체이므로 쌍대 완전 복합체 $E^\vee$가 존재한다.
Cohomology on Sites, Lemma
\ref{sites-cohomology-lemma-dual-perfect-complex}을 보라.
$R\SheafHom(E, \mathcal{G}^\bullet) =
E^\vee \otimes^\mathbf{L}_{\mathcal{O}_X} \mathcal{G}^\bullet$임을
관찰하자. 따라서 $K$의 완전성은 Lemma
\ref{spaces-perfect-lemma-tensor-perfect}에서 따른다.
\end{proof}






\section{Ext에 대한 사영 공식}
\label{spaces-perfect-section-ext}

\noindent
Lemma \ref{spaces-perfect-lemma-compute-ext} (또는 문헌의 유사한 결과들)는 때때로
Quot 함자, 힐베르트 스킴 및 그 밖의 모듈라이 문제에 대한
아르틴의 판정 기준 가운데 하나를 확인하는 데 필요한 장애 이론의
성질들을 확인할 때 유용하다. $f : X \to Y$가 고유이고 평탄하며
유한 표시인 대수공간의 사상이고
$E \in D(\mathcal{O}_X)$가 완전이라고 하자. 여기서 보조정리는
$$
\Ext^i_X(E, f^*\mathcal{F}) =
\Ext^i_Y((Rf_*E^\vee)^\vee, \mathcal{F})
$$
라고 말한다. 단, $\mathcal{F}$는 $Y$ 위에서 준연접이다.
이와 같이 쓰면 Ext에 대한 사영 공식처럼 보이며, 실제로 그 결과는
Lemma \ref{spaces-perfect-lemma-cohomology-base-change}에서 비교적 쉽게 따른다.

\begin{lemma}
\label{spaces-perfect-lemma-compute-tensor-perfect}
가정과 표기는 Lemma \ref{spaces-perfect-lemma-tensor-perfect}에서와 같다. 그러면
함자적인 동형사상들
$$
H^i(B, K \otimes^\mathbf{L}_{\mathcal{O}_B} \mathcal{F})
\longrightarrow
H^i(X, E \otimes^\mathbf{L}_{\mathcal{O}_X}
(\mathcal{G}^\bullet \otimes_{\mathcal{O}_X} f^*\mathcal{F}))
$$
이 존재한다. 여기서 $\mathcal{F}$는 $B$ 위에서 준연접이고,
이 동형사상들은 경계 사상들과 양립한다 (증명을 보라).
\end{lemma}

\begin{proof}
다음 등식들이 성립한다.
$$
\mathcal{G}^\bullet \otimes_{\mathcal{O}_X}^\mathbf{L} Lf^*\mathcal{F} =
\mathcal{G}^\bullet \otimes_{f^{-1}\mathcal{O}_B}^\mathbf{L} f^{-1}\mathcal{F} =
\mathcal{G}^\bullet \otimes_{f^{-1}\mathcal{O}_B} f^{-1}\mathcal{F} =
\mathcal{G}^\bullet \otimes_{\mathcal{O}_X} f^*\mathcal{F}
$$
첫째 등식은 Cohomology on Sites, Lemma
\ref{sites-cohomology-lemma-variant-derived-pullback}에서 따르고,
둘째 등식은 $\mathcal{G}^n$이 평탄한
$f^{-1}\mathcal{O}_B$-가군이기 때문에 성립하며, 셋째 등식은
당김의 정의에서 따른다. 따라서
\begin{align*}
H^i(X, E \otimes^\mathbf{L}_{\mathcal{O}_X}
(\mathcal{G}^\bullet \otimes_{\mathcal{O}_X} f^*\mathcal{F}))
& =
H^i(X, E \otimes^\mathbf{L}_{\mathcal{O}_X} \mathcal{G}^\bullet
\otimes_{\mathcal{O}_X}^\mathbf{L} Lf^*\mathcal{F}) \\
& =
H^i(B,
Rf_*(E \otimes^\mathbf{L}_{\mathcal{O}_X} \mathcal{G}^\bullet
\otimes^\mathbf{L}_{\mathcal{O}_X} Lf^*\mathcal{F})) \\
& =
H^i(B,
Rf_*(E \otimes^\mathbf{L}_{\mathcal{O}_X} \mathcal{G}^\bullet)
\otimes^\mathbf{L}_{\mathcal{O}_B} \mathcal{F}) \\
& =
H^i(B, K \otimes^\mathbf{L}_{\mathcal{O}_B} \mathcal{F})
\end{align*}
를 얻는다. 첫째 등식은 위의 계산에서, 둘째 등식은 Leray
(Cohomology on Sites, Remark
\ref{sites-cohomology-remark-before-Leray})에서, 셋째 등식은 Lemma
\ref{spaces-perfect-lemma-cohomology-base-change}에서 따른다.
경계 사상에 관한 진술의 뜻은 다음과 같다. 짧은 완전열
$0 \to \mathcal{F}_1 \to \mathcal{F}_2 \to \mathcal{F}_3 \to 0$이
주어지면 이 동형사상들은 가환 도식
$$
\xymatrix{
H^i(B, K \otimes^\mathbf{L}_{\mathcal{O}_B} \mathcal{F}_3)
\ar[r] \ar[d]_\delta &
H^i(X, E \otimes^\mathbf{L}_{\mathcal{O}_X}
(\mathcal{G}^\bullet \otimes_{\mathcal{O}_X} f^*\mathcal{F}_3)) \ar[d]^\delta \\
H^{i + 1}(B, K \otimes^\mathbf{L}_{\mathcal{O}_B} \mathcal{F}_1)
\ar[r] &
H^{i + 1}(X, E \otimes^\mathbf{L}_{\mathcal{O}_X}
(\mathcal{G}^\bullet \otimes_{\mathcal{O}_X} f^*\mathcal{F}_1))
}
$$
에 들어맞는다. 여기서 경계 사상들은 구별 삼각형
$$
K \otimes^\mathbf{L}_{\mathcal{O}_B} \mathcal{F}_1 \to
K \otimes^\mathbf{L}_{\mathcal{O}_B} \mathcal{F}_2 \to
K \otimes^\mathbf{L}_{\mathcal{O}_B} \mathcal{F}_3 \to
K \otimes^\mathbf{L}_{\mathcal{O}_B} \mathcal{F}_1[1]
$$
과 $D(\mathcal{O}_X)$에서 짧은 완전열
$$
0 \to
\mathcal{G}^\bullet \otimes_{\mathcal{O}_X} f^*\mathcal{F}_1 \to
\mathcal{G}^\bullet \otimes_{\mathcal{O}_X} f^*\mathcal{F}_2 \to
\mathcal{G}^\bullet \otimes_{\mathcal{O}_X} f^*\mathcal{F}_3 \to 0
$$
에 결부된 구별 삼각형에서 나온다. 이는
복합체들의 짧은 완전열이다. 이 열은 $\mathcal{G}^n$이 $B$
위에서 평탄하기 때문에 완전하다. 위에 표시한 도식의 가환성에
대한 확인은 생략한다.
\end{proof}


\begin{lemma}
\label{spaces-perfect-lemma-compute-ext-perfect}
가정과 표기는 Lemma \ref{spaces-perfect-lemma-ext-perfect}에서와 같다. 그러면
함자적인 동형사상들
$$
H^i(B, K \otimes^\mathbf{L}_{\mathcal{O}_B} \mathcal{F})
\longrightarrow
\Ext^i_{\mathcal{O}_X}(E,
\mathcal{G}^\bullet \otimes_{\mathcal{O}_X} f^*\mathcal{F})
$$
이 존재한다. 여기서 $\mathcal{F}$는 $B$ 위에서 준연접이고,
이 동형사상들은 경계 사상들과 양립한다 (증명을 보라).
\end{lemma}

\begin{proof}
Lemma \ref{spaces-perfect-lemma-ext-perfect}의 증명에서와 같이 $E^\vee$를 쌍대
완전 복합체라 하자. 또한
$K = Rf_*(E^\vee \otimes_{\mathcal{O}_X}^\mathbf{L} \mathcal{G}^\bullet)$
임을 상기하자. 한편
$$
\Ext^i_{\mathcal{O}_X}(E,
\mathcal{G}^\bullet \otimes_{\mathcal{O}_X} f^*\mathcal{F})
=
H^i(X, E^\vee \otimes^\mathbf{L}_{\mathcal{O}_X}
(\mathcal{G}^\bullet \otimes_{\mathcal{O}_X} f^*\mathcal{F}))
$$
이기도 하므로, $E^\vee$의 구성에 의해 동형사상들의 존재는
$E^\vee$와 $\mathcal{G}^\bullet$에 Lemma
\ref{spaces-perfect-lemma-compute-tensor-perfect}을 적용하면 따른다.
경계 사상에 관한 진술의 뜻은 다음과 같다. 짧은 완전열
$0 \to \mathcal{F}_1 \to \mathcal{F}_2 \to \mathcal{F}_3 \to 0$이
주어지면 이 동형사상들은 가환 도식
$$
\xymatrix{
H^i(B, K \otimes^\mathbf{L}_{\mathcal{O}_B} \mathcal{F}_3)
\ar[r] \ar[d]_\delta &
\Ext^i_{\mathcal{O}_X}(E,
\mathcal{G}^\bullet \otimes_{\mathcal{O}_X} f^*\mathcal{F}_3) \ar[d]^\delta \\
H^{i + 1}(B, K \otimes^\mathbf{L}_{\mathcal{O}_B} \mathcal{F}_1)
\ar[r] &
\Ext^{i + 1}_{\mathcal{O}_X}(E,
\mathcal{G}^\bullet \otimes_{\mathcal{O}_X} f^*\mathcal{F}_1)
}
$$
에 들어맞는다. 여기서 경계 사상들은 구별 삼각형
$$
K \otimes^\mathbf{L}_{\mathcal{O}_B} \mathcal{F}_1 \to
K \otimes^\mathbf{L}_{\mathcal{O}_B} \mathcal{F}_2 \to
K \otimes^\mathbf{L}_{\mathcal{O}_B} \mathcal{F}_3 \to
K \otimes^\mathbf{L}_{\mathcal{O}_B} \mathcal{F}_1[1]
$$
과 $D(\mathcal{O}_X)$에서 짧은 완전열
$$
0 \to
\mathcal{G}^\bullet \otimes_{\mathcal{O}_X} f^*\mathcal{F}_1 \to
\mathcal{G}^\bullet \otimes_{\mathcal{O}_X} f^*\mathcal{F}_2 \to
\mathcal{G}^\bullet \otimes_{\mathcal{O}_X} f^*\mathcal{F}_3 \to 0
$$
에 결부된 구별 삼각형에서 나온다. 이는
복합체들의 짧은 완전열이다. 이 열은 $\mathcal{G}^n$이 $B$
위에서 평탄하기 때문에 완전하다. 위에 표시한 도식의 가환성에
대한 확인은 생략한다.
\end{proof}
\begin{lemma}
\label{spaces-perfect-lemma-compute-ext}
$S$를 스킴이라 하자. $f : X \to B$를 $S$ 위 대수공간의
사상이라 하고, $E \in D(\mathcal{O}_X)$라 하며,
$\mathcal{F}^\bullet$을 $\mathcal{O}_X$-가군의 복합체라 하자.
% P09-ERR-0478
다음을 가정하자.
\begin{enumerate}
\item $B$는 뇌터이다.
\item $f$는 국소 유한형이고 준분리이다.
\item $E \in D^-_{\textit{Coh}}(\mathcal{O}_X)$이다.
\item $\mathcal{G}^\bullet$은 연접 $\mathcal{O}_X$-가군의 유계
복합체로서 $B$ 위에서 평탄하고 지지가 $B$ 위에서 고유이다.
% P09-ERR-0479
\end{enumerate}
그러면 다음 두 명제가 참이다.
\begin{enumerate}
\item[(A)] 모든 $m \in \mathbf{Z}$에 대하여 완전 대상 $K$가
$D(\mathcal{O}_B)$에 존재하고, 함자적인 사상들
$$
\alpha^i_\mathcal{F} :
\Ext^i_{\mathcal{O}_X}(E,
\mathcal{G}^\bullet \otimes_{\mathcal{O}_X} f^*\mathcal{F})
\longrightarrow
H^i(B, K \otimes^\mathbf{L}_{\mathcal{O}_B} \mathcal{F})
$$
이 존재한다. 여기서 $\mathcal{F}$는 $B$ 위에서 준연접이고,
사상들은 경계 사상들과 양립하며 (증명을 보라),
$\alpha^i_\mathcal{F}$는 $i \leq m$에 대하여 동형이다.
\item[(B)] 유사연접 대상 $L \in D(\mathcal{O}_B)$과 함자적인
동형사상들
$$
\Ext^i_{\mathcal{O}_B}(L, \mathcal{F}) \longrightarrow
\Ext^i_{\mathcal{O}_X}(E,
\mathcal{G}^\bullet \otimes_{\mathcal{O}_X} f^*\mathcal{F})
$$
이 존재한다. 여기서 $\mathcal{F}$는 $B$ 위에서 준연접이고,
동형사상들은 경계 사상들과 양립한다.
\end{enumerate}
\end{lemma}

\begin{proof}
(A)의 증명. $\mathcal{G}^i$가 $i \in [a, b]$에 대해서만 영이
아니라고 하자. $X$를 $\mathcal{G}^i$의 지지들의 합집합의
준콤팩트 열린 근방으로 바꿀 수 있다. 따라서 $X$가 뇌터라고
가정할 수 있다. 이 경우 $X$와 $f$는 준콤팩트이고 준분리이다.
근사 $P \to E$를 택하되, 이를 완전 복합체 $P$에 의한
$(X, E, -m - 1 + a)$의 근사라 하자
(Theorem \ref{spaces-perfect-theorem-approximation}에 의해 가능하다). 그러면 유도된
사상
$$
\Ext^i_{\mathcal{O}_X}(E,
\mathcal{G}^\bullet \otimes_{\mathcal{O}_X} f^*\mathcal{F})
\longrightarrow
\Ext^i_{\mathcal{O}_X}(P,
\mathcal{G}^\bullet \otimes_{\mathcal{O}_X} f^*\mathcal{F})
$$
은 $i \leq m$에 대하여 동형이다. 실제로 이 사상의 핵과 여핵은
각각
$$
\Ext^i_{\mathcal{O}_X}(C,
\mathcal{G}^\bullet \otimes_{\mathcal{O}_X} f^*\mathcal{F})
\quad\text{resp.}\quad
\Ext^{i + 1}_{\mathcal{O}_X}(C,
\mathcal{G}^\bullet \otimes_{\mathcal{O}_X} f^*\mathcal{F})
$$
의 몫가군과 부분가군이다. 여기서 $C$는 $P \to E$의 원뿔이다.
$C$의 코호몰로지 층들은 차수 $\geq -m - 1 + a$에서 영이므로,
Derived Categories, Lemma \ref{derived-lemma-negative-exts}에 의해
이 $\Ext$-군들은 $i \leq m + 1$에 대하여 영이다. 이로써 $E$가
완전 복합체인 경우로 귀착되며, 그 경우는 Lemma
\ref{spaces-perfect-lemma-compute-ext-perfect}이다. 경계 사상에 관한 진술은
Lemma \ref{spaces-perfect-lemma-compute-ext-perfect}의 증명에서 설명하였다.

\medskip\noindent
(B)의 증명. (A)의 증명에서와 같이 $X$가 뇌터라고 가정할 수
있다. Lemma \ref{spaces-perfect-lemma-identify-pseudo-coherent-noetherian}에 의해
$E$는 유사연접임을 관찰하자. Lemma
\ref{spaces-perfect-lemma-pseudo-coherent-hocolim}에 의해
$E = \text{hocolim} E_n$으로 쓸 수 있다. 여기서 $E_n$은 완전이고
$E_n \to E$는 절단 $\tau_{\geq -n}$ 위에서 동형을 유도한다.
$E_n^\vee$를 쌍대 완전 복합체라 하자
(Cohomology on Sites, Lemma
\ref{sites-cohomology-lemma-dual-perfect-complex}).
그러면 완전 대상들의 역계
$\ldots \to E_3^\vee \to E_2^\vee \to E_1^\vee$를 얻는다. 이로부터
다시 역계
$$
\ldots \to K_3 \to K_2 \to K_1\quad\text{with}\quad
K_n = Rf_*(E_n^\vee \otimes_{\mathcal{O}_X}^\mathbf{L} \mathcal{G}^\bullet)
$$
를 얻으며, 이는 $Y$ 위에서 완전하다. Lemma
\ref{spaces-perfect-lemma-tensor-perfect}을 보라. % P09-ERR-0480
Lemma \ref{spaces-perfect-lemma-compute-ext-perfect}과 그 증명, 그리고 앞 단락의
논증 ($P = E_n$으로 둔다)에 의해, 임의의 준연접
$\mathcal{F}$가 $Y$ 위에 주어졌을 때 함자적인 표준 사상들
$$
\xymatrix{
& \Ext^i_{\mathcal{O}_X}(E,
\mathcal{G}^\bullet \otimes_{\mathcal{O}_X} f^*\mathcal{F})
\ar[ld] \ar[rd] \\
H^i(Y, K_{n + 1} \otimes_{\mathcal{O}_Y}^\mathbf{L} \mathcal{F})
\ar[rr] & &
H^i(Y, K_n \otimes_{\mathcal{O}_Y}^\mathbf{L} \mathcal{F})
}
$$
을 얻는다. 이들은 $i \leq n + a$에 대하여 동형이다.
% P09-ERR-0481
$L_n = K_n^\vee$를 쌍대 완전 복합체라 하자. 그러면
$L_1 \to L_2 \to L_3 \to \ldots$는 $D(\mathcal{O}_Y)$ 안의 완전
대상들의 계이며, 임의의 준연접 $\mathcal{F}$가 $Y$ 위에
주어졌을 때
사상들
$$
\Ext^i_{\mathcal{O}_Y}(L_{n + 1}, \mathcal{F})
\longrightarrow
\Ext^i_{\mathcal{O}_Y}(L_n, \mathcal{F})
$$
은 $i \leq n + a - 1$에 대하여 동형이다. 이는
$L_n \to L_{n + 1}$이 절단 $\tau_{\geq -n - a + 2}$ 위에서
동형을 유도함을 뜻한다 (힌트: $L_n \to L_{n + 1}$의 원뿔을 잡아
그 마지막 영이 아닌 코호몰로지 층을 살펴보라).
따라서 $L = \text{hocolim} L_n$은 유사연접이다. Lemma
\ref{spaces-perfect-lemma-pseudo-coherent-hocolim}을 보라. 호모토피 여극한의
사상 성질에 의해
$\Ext^i_{\mathcal{O}_Y}(L, \mathcal{F}) =
\Ext^i_{\mathcal{O}_Y}(L_n, \mathcal{F})$가
$i \leq n + a - 3$에 대하여 성립하며, 이것으로 증명이 끝난다.
\end{proof}

\begin{remark}
\label{spaces-perfect-remark-base-change-of-L}
Lemma \ref{spaces-perfect-lemma-compute-ext}의 (B)에서 나온 유사연접 복합체 $L$은
이 상황에 표준적으로 결부된다. 예를 들어 (B)에서와 같은 $L$의
형성은 밑변환과 양립한다. 다시 말해 스킴의 카르테시안 도식
$$
\xymatrix{
X' \ar[r]_{g'} \ar[d]_{f'} &
X \ar[d]^f \\
Y' \ar[r]^g &
Y
}
$$
이 주어지면 표준적이고 함자적인 동형사상들
$$
\Ext^i_{\mathcal{O}_{Y'}}(Lg^*L, \mathcal{F}') \longrightarrow
\Ext^i_{\mathcal{O}_X}(L(g')^*E,
(g')^*\mathcal{G}^\bullet \otimes_{\mathcal{O}_{X'}} (f')^*\mathcal{F}')
$$
이 존재한다. 여기서 $\mathcal{F}'$는 $Y'$ 위에서 준연접이다.
% P09-ERR-0482
우변에서는 $\mathcal{G}^\bullet$에 유도 당김을 사용하지
\textbf{않음}을 관찰하자. % P09-ERR-0483
이 결과가 필요해지면 여기에 정확한 결과를 정식화하고 상세한
증명을 제시할 것이다.
\end{remark}








\section{극한과 유도 범주}
\label{spaces-perfect-section-limits}

\noindent
이 절에서는 대수공간들의 역계의 극한인 대수공간의 유도 범주에
관한 몇 가지 결과를 모은다. 더 정확히 말하면 다음 설정에서
작업할 것이다.

\begin{situation}
\label{spaces-perfect-situation-descent}
$S$를 스킴이라 하자. $X = \lim_{i \in I} X_i$를 $S$ 위
대수공간들의 유향계의 극한이라 하며, 그 전이 사상
$f_{i'i} : X_{i'} \to X_i$는 아핀이라고 하자. 사영을
$f_i : X \to X_i$로 나타낸다. $X_i$가 모든 $i \in I$에 대하여
준콤팩트이고 준분리라고 가정한다. 또한 원소 $0 \in I$를 하나
택한다.
\end{situation}

\begin{lemma}
\label{spaces-perfect-lemma-descend-homomorphisms}
Situation \ref{spaces-perfect-situation-descent}의 설정에서
$E_0$와 $K_0$를 $D(\mathcal{O}_{X_0})$의 대상들이라 하자.
$E_i = Lf_{i0}^*E_0$와 $K_i = Lf_{i0}^*K_0$로 놓되 이는
$i \geq 0$에 대하여 정하고, $E = Lf_0^*E_0$와
$K = Lf_0^*K_0$로 놓자. 그러면 사상
$$
\colim_{i \geq 0} \Hom_{D(\mathcal{O}_{X_i})}(E_i, K_i)
\longrightarrow
\Hom_{D(\mathcal{O}_X)}(E, K)
$$
은 다음 중 하나가 성립하면 동형이다.
\begin{enumerate}
\item $E_0$는 완전이고 $K_0 \in D_\QCoh(\mathcal{O}_{X_0})$이다.
\item $E_0$는 유사연접이고
$K_0 \in D_\QCoh(\mathcal{O}_{X_0})$는 유한 Tor 차원을 가진다.
\end{enumerate}
\end{lemma}

\begin{proof}
모든 준콤팩트 준분리 대상 $U_0$를
$(X_0)_{spaces, \etale}$에서 잡고, 조건 $P$를 생각하자. 이는
표준 사상
$$
\colim_{i \geq 0} \Hom_{D(\mathcal{O}_{U_i})}(E_i|_{U_i}, K_i|_{U_i})
\longrightarrow
\Hom_{D(\mathcal{O}_U)}(E|_U, K|_U)
$$
이 동형이라는 조건이다. 여기서
$U = X \times_{X_0} U_0$이고
$U_i = X_i \times_{X_0} U_0$이다. Lemma
\ref{spaces-perfect-lemma-induction-principle}의 귀납 원리를 사용하여 $P$가 모든
$U_0$에 대해 성립함을 증명할 것이다. 이 보조정리의 조건
(2)는 유도 범주에서의 Hom에 대한 Mayer--Vietoris로부터 곧바로
따른다. Lemma \ref{spaces-perfect-lemma-mayer-vietoris-hom}을 보라. 따라서
$X_0$가 아핀인 경우에 보조정리를 증명하면 충분하다.

\medskip\noindent
$X_0$가 아핀이면 결과는 스킴의 경우에서 따른다. Derived
Categories of Schemes, Lemma
\ref{perfect-lemma-descend-homomorphisms}을 보라. 이를 확인하려면
Lemma \ref{spaces-perfect-lemma-derived-quasi-coherent-small-etale-site}의 동치와
Lemmas \ref{spaces-perfect-lemma-descend-pseudo-coherent},
\ref{spaces-perfect-lemma-descend-tor-amplitude},
\ref{spaces-perfect-lemma-descend-perfect}에서 설명한 성질들의 대응을 사용한다.
\end{proof}

\begin{lemma}
\label{spaces-perfect-lemma-perfect-on-limit}
Situation \ref{spaces-perfect-situation-descent}에서 $D(\mathcal{O}_X)$의 완전
대상들의 범주는 $D(\mathcal{O}_{X_i})$의 완전 대상들의 범주들의
여극한이다.
\end{lemma}

\begin{proof}
모든 준콤팩트 준분리 대상 $U_0$를
$(X_0)_{spaces, \etale}$에서 잡고, 조건 $P$를 생각하자. 이는 함자
$$
\colim_{i \geq 0} D_{perf}(\mathcal{O}_{U_i})
\longrightarrow
D_{perf}(\mathcal{O}_U)
$$
가 동치라는 조건이다. 여기서 ${}_{perf}$는 완전
대상들로 이루어진 충만 부분범주를 뜻하며,
$U = X \times_{X_0} U_0$이고
$U_i = X_i \times_{X_0} U_0$이다. Lemma
\ref{spaces-perfect-lemma-induction-principle}의 귀납 원리를 사용하여 $P$가 모든
$U_0$에 대해 성립함을 증명할 것이다. 우선 Lemma
\ref{spaces-perfect-lemma-descend-homomorphisms}에 의해 이 함자가 이미
충만충실임을 알고 있다. 따라서 본질적 전사성을 증명하면
충분하다.

\medskip\noindent
먼저 귀납 원리의 조건 (2)를 확인하자. 따라서 기본 구별
정사각형 $(U_0 \subset X_0, V_0 \to X_0)$이 주어지고,
$P$가 $U_0$, $V_0$, $U_0 \times_{X_0} V_0$에 대해 성립한다고
하자. $E$를 $D(\mathcal{O}_X)$의 완전 대상이라 하자.
$i \geq 0$과 완전 대상 $E_{U, i}$를 $U_i$ 위에서, 완전 대상
$E_{V, i}$를 $V_i$ 위에서 찾아서, 이들을 $U$와 $V$로 당긴 것들이
$E|_U$와 $E|_V$에 각각 동형이 되게 할 수 있다. 다음 사상들을
각각
$$
a : E_{U, i} \to (R(X \to X_i)_*E)|_{U_i}
\quad\text{and}\quad
b : E_{V, i} \to (R(X \to X_i)_*E)|_{V_i}
$$
로 나타내자. 이들은 동형사상
$L(U \to U_i)^*E_{U, i} \to E|_U$와
$L(V \to V_i)^*E_{V, i} \to E|_V$에 수반한 사상들이다.
충만충실성에 의해 $i$를 늘린 뒤, 동형사상
$c : E_{U, i}|_{U_i \times_{X_i} V_i} \to E_{V, i}|_{U_i \times_{X_i} V_i}$
를 찾을 수 있으며, 이것을 당기면 식별
$$
L(U \to U_i)^*E_{U, i}|_{U \times_X V} \to E|_{U \times_X V} \to
L(V \to V_i)^*E_{V, i}|_{U \times_X V}.
$$
이 된다. Lemma \ref{spaces-perfect-lemma-glue}를 적용하여 대상 $E_i$를 $X_i$
위에 얻고, 사상 $d : E_i \to R(X \to X_i)_*E$를 얻자. 이 사상은
사상 $a$와 $b$로 제한되는데, 각각 $U_i$와 $V_i$ 위에서 그러하다.
그러면
$E_i$가 완전이고 $d$가 동형사상
$L(X \to X_i)^*E_i \to E$에 수반한 사상임은 명백하다.

\medskip\noindent
마지막으로 귀납 원리의 조건 (1), 즉 $X_0$가 아핀일 때
보조정리가 성립함을 확인하자. 이는 스킴의 경우에서 따른다.
Derived Categories of Schemes, Lemma
\ref{perfect-lemma-descend-perfect}을 보라. 이를 확인하려면 Lemma
\ref{spaces-perfect-lemma-derived-quasi-coherent-small-etale-site}의 동치와
Lemma \ref{spaces-perfect-lemma-descend-perfect}의 대응을 사용한다.
\end{proof}







\section{코호몰로지와 밑변환, VI}
\label{spaces-perfect-section-cohomology-and-base-change-final}

\noindent
Sections
\ref{spaces-perfect-section-cohomology-and-base-change-perfect},
\ref{spaces-perfect-section-cohomology-base-change},
\ref{spaces-perfect-section-producing-perfect}의 논의를 이어가는 코호몰로지와
밑변환에 관한 마지막 절이다. 이해하기 쉬운 특수한 경우는
Remark \ref{spaces-perfect-remark-explain-perfect-direct-image}에 제시되어 있다.

\begin{lemma}
\label{spaces-perfect-lemma-base-change-tensor-perfect}
$S$를 스킴이라 하자. $f : X \to Y$를 $S$ 위 대수공간들 사이의
유한 표시 사상이라 하자. $E \in D(\mathcal{O}_X)$를 완전 대상이라
하자. $\mathcal{G}^\bullet$을 유한 표시
$\mathcal{O}_X$-가군의 유계 복합체로서 $Y$ 위에서 평탄하고
지지가 $Y$ 위에서 고유인 것이라 하자. 그러면
$$
K = Rf_*(E \otimes_{\mathcal{O}_X}^\mathbf{L} \mathcal{G}^\bullet)
$$
는 $D(\mathcal{O}_Y)$의 완전 대상이고, 그 형성은 임의의
밑변환과 가환한다.
\end{lemma}

\begin{proof}
밑변환에 관한 진술은 Lemma \ref{spaces-perfect-lemma-base-change-tensor}이다.
따라서 $K$가 완전 대상임을 보이면 충분하다. $Y$가 뇌터이면 이는
Lemma \ref{spaces-perfect-lemma-tensor-perfect}에서 따른다. 뇌터 근사를 통해
이 경우로 귀착시킬 것이다. 독자는 이 증명의 나머지 부분을
건너뛰어도 좋다.

\medskip\noindent
문제는 $Y$ 위에서 국소적이므로 $Y$가 아핀이라고 가정할 수 있다.
$Y = \Spec(R)$이라 하자. $R = \colim R_i$를 뇌터 환들 $R_i$의
여과 여극한으로 쓴다. Limits of Spaces, Lemma
\ref{spaces-limits-lemma-descend-finite-presentation}에 의해 어떤
$i$와 대수공간 $X_i$가 존재하며, 이는 $R_i$ 위에서 유한 표시이고
이를 $R$로 밑변환한 것은 $X$이다. Limits of Spaces, Lemma
\ref{spaces-limits-lemma-descend-modules-finite-presentation}에
의해 $i$를 늘린 뒤, 유한 표시
$\mathcal{O}_{X_i}$-가군의 유계 복합체
$\mathcal{G}_i^\bullet$이 존재하고 이를 $X$로 당긴 것은
$\mathcal{G}^\bullet$이라고 가정할 수 있다. $i$를 늘린 뒤
$\mathcal{G}_i^n$이 $R_i$ 위에서 평탄하다고 가정할 수 있다.
Limits of Spaces, Lemma
\ref{spaces-limits-lemma-descend-flat}을 보라. $i$를 늘린 뒤
$\mathcal{G}_i^n$의 지지가 $R_i$ 위에서 고유라고 가정할 수 있다.
Limits of Spaces, Lemma
\ref{spaces-limits-lemma-eventually-proper-support}를 보라. 마지막으로
Lemma \ref{spaces-perfect-lemma-perfect-on-limit}에 의해 $i$를 늘린 뒤 완전 대상
$E_i$가 $D(\mathcal{O}_{X_i})$에 존재하고, 이를 $X$로 당긴 것이
$E$라고 가정할 수 있다. Lemma \ref{spaces-perfect-lemma-tensor-perfect}에 의해
$$
K_i =
Rf_{i, *}(E_i \otimes_{\mathcal{O}_{X_i}}^\mathbf{L} \mathcal{G}_i^\bullet)
$$
는 $\Spec(R_i)$ 위에서 완전이다. 여기서
$f_i : X_i \to \Spec(R_i)$는 구조 사상이다. 밑변환 결과
(Lemma \ref{spaces-perfect-lemma-base-change-tensor})에 의해 $K_i$를
$Y = \Spec(R)$로 당긴 것은 $K$이며, 이로써 결론을 얻는다.
\end{proof}

\begin{remark}
\label{spaces-perfect-remark-explain-perfect-direct-image}
$R$을 환이라 하자. $X$를 $R$ 위에서 유한 표시인 대수공간이라
하자. $\mathcal{G}$를 유한 표시 $\mathcal{O}_X$-가군으로서 $R$
위에서 평탄하고 지지가 $R$ 위에서 고유인 것이라 하자. Lemma
\ref{spaces-perfect-lemma-base-change-tensor-perfect}에 의해 유한 생성 사영
$R$-가군의 유한 복합체 $M^\bullet$이 존재하여
$$
R\Gamma(X_{R'}, \mathcal{G}_{R'}) = M^\bullet \otimes_R R'
$$
이 $R$-대수 $R'$에 대하여 함자적으로 성립한다.
\end{remark}

\begin{lemma}
\label{spaces-perfect-lemma-base-change-tensor-pseudo-coherent}
$S$를 스킴이라 하자.
$f : X \to Y$를 $S$ 위 대수공간들 사이의 유한 표시 사상이라
하자. $E \in D(\mathcal{O}_X)$를 유사연접 대상이라 하자.
$\mathcal{G}^\bullet$을 유한 표시 $\mathcal{O}_X$-가군의 위로
유계인 복합체로서 $Y$ 위에서 평탄하고 지지가 $Y$ 위에서
고유인 것이라 하자. 그러면
$$
K = Rf_*(E \otimes_{\mathcal{O}_X}^\mathbf{L} \mathcal{G}^\bullet)
$$
는 $D(\mathcal{O}_Y)$의 유사연접 대상이고, 그 형성은 임의의
밑변환과 가환한다.
\end{lemma}

\begin{proof}
밑변환에 관한 진술은 Lemma \ref{spaces-perfect-lemma-base-change-tensor}이다.
따라서 $K$가 유사연접 대상임을 보이면 충분하다. 이는 완전
복합체들에 의한 근사를 사용하여 Lemma
\ref{spaces-perfect-lemma-base-change-tensor-perfect}에서 따를 것이다. 독자는
이 증명의 나머지 부분을 건너뛰어도 좋다.

\medskip\noindent
문제는 $Y$ 위에서 에탈 국소적이므로 $Y$가 아핀이라고 가정할 수
있다. 그러면 $X$는 준콤팩트이고 준분리이다. 더욱이 정수 $N$이
존재하여 총직상
$Rf_* : D_\QCoh(\mathcal{O}_X) \to D_\QCoh(\mathcal{O}_Y)$가
코호몰로지 차원 $N$을 가진다. 이는 Lemma
\ref{spaces-perfect-lemma-quasi-coherence-direct-image}에서 설명되어 있다.
정수 $b$를 택하여 $\mathcal{G}^i = 0$이 $i > b$에 대하여
성립하게 하자. $K$가 모든 $m$에 대하여 $m$-유사연접임을 보이면
충분하다. 근사 $P \to E$를 택하되, 이를 완전 복합체 $P$에 의한
$(X, E, m - N - 1 - b)$의 근사라 하자. 이는 Theorem
\ref{spaces-perfect-theorem-approximation}에 의해 가능하다. 구별 삼각형
$$
P \to E \to C \to P[1]
$$
을 $D_\QCoh(\mathcal{O}_X)$에서 택하자. $C$의 코호몰로지 층들은
차수 $\geq m - N - 1 - b$에서 영이다. 따라서
$C \otimes^\mathbf{L} \mathcal{G}^\bullet$의 코호몰로지 층들은
차수 $\geq m - N - 1$에서 영이다. 그러므로
$Rf_*(C \otimes^\mathbf{L} \mathcal{G})$의 코호몰로지 층들은
차수 $\geq m - 1$에서 영이다. % P09-ERR-0484
따라서
$$
Rf_*(P \otimes^\mathbf{L} \mathcal{G}) \to
Rf_*(E \otimes^\mathbf{L} \mathcal{G})
$$
는 차수 $\geq m$의 코호몰로지 층들 위에서 동형이다.
다음으로 $H^i(P) = 0$이 $i > a$에 대하여 성립한다고 하자. 그러면
$
P \otimes^\mathbf{L} \sigma_{\geq m - N - 1 - a}\mathcal{G}^\bullet
\longrightarrow
P \otimes^\mathbf{L} \mathcal{G}^\bullet
$
는 차수 $\geq m - N - 1$의 코호몰로지 층들 위에서 동형이다.
따라서 다시
$$
Rf_*(P \otimes^\mathbf{L} \sigma_{\geq m - N - 1 - a}\mathcal{G}^\bullet) \to
Rf_*(P \otimes^\mathbf{L} \mathcal{G}^\bullet)
$$
가 차수 $\geq m$의 코호몰로지 층들 위에서 동형임을 얻는다.
Lemma \ref{spaces-perfect-lemma-base-change-tensor-perfect}에 의해 원천은 완전
복합체이다. 결론적으로 $K$는 원하는 대로 $m$-유사연접이다.
\end{proof}

\begin{lemma}
\label{spaces-perfect-lemma-flat-proper-perfect-direct-image-general}
$S$를 스킴이라 하자. $f : X \to Y$를 $S$ 위 대수공간의
고유 유한 표시 사상이라 하자.
\begin{enumerate}
\item $E \in D(\mathcal{O}_X)$가 완전이고 $f$가 평탄하다고 하자.
그러면 $Rf_*E$는 $D(\mathcal{O}_Y)$의 완전 대상이고, 그 형성은
임의의 밑변환과 가환한다.
\item $\mathcal{G}$를 유한 표시 $\mathcal{O}_X$-가군으로서 $S$
위에서 평탄한 것이라 하자. 그러면 $Rf_*\mathcal{G}$는
$D(\mathcal{O}_Y)$의 완전 대상이고, 그 형성은 임의의 밑변환과
가환한다. % P09-ERR-0485
\end{enumerate}
\end{lemma}

\begin{proof}
Lemma \ref{spaces-perfect-lemma-base-change-tensor-perfect}을 적용한 특수한
경우들이다. (1)에서는 $\mathcal{G}^\bullet$을
$\mathcal{O}_X$와 같게 두고 이를 차수 $0$에 놓는다. (2)에서는
$E = \mathcal{O}_X$로 두며 $\mathcal{G}^\bullet$은
$\mathcal{G}$만으로 이루어지게 하고 이를 차수 $0$에 놓는다.
\end{proof}

\begin{lemma}
\label{spaces-perfect-lemma-flat-proper-pseudo-coherent-direct-image-general}
$S$를 스킴이라 하자. $f : X \to Y$를 $S$ 위 대수공간의 평탄하고
고유인 유한 표시 사상이라 하자. $E \in D(\mathcal{O}_X)$가
유사연접이라고 하자. 그러면 $Rf_*E$는 $D(\mathcal{O}_Y)$의
유사연접 대상이고, 그 형성은 임의의 밑변환과 가환한다.
\end{lemma}

\noindent
더 일반적으로 $f : X \to Y$가 고유이고 $E$가 $X$ 위에 있으면서
$Y$에 대해 상대적으로 유사연접이면
(More on Morphisms of Spaces, Definition
\ref{spaces-more-morphisms-definition-relative-pseudo-coherence}),
$Rf_*E$는 유사연접이다 (그러나 이 일반성에서는 그 형성이
밑변환과 가환하지 않는다). 스킴에 대한 이 결과는
\cite{Kiehl}에서 증명되어 있다.

\begin{proof}
$\mathcal{G} = \mathcal{O}_X$로 두고 Lemma
\ref{spaces-perfect-lemma-base-change-tensor-pseudo-coherent}을 적용한 특수한
경우이다.
\end{proof}

\begin{lemma}
\label{spaces-perfect-lemma-pullback-and-limits}
$R$을 환이라 하자. $X$를 대수공간이라 하고
$f : X \to \Spec(R)$을 고유이고 평탄하며 유한 표시인 사상이라
하자. $(M_n)$을 전이 사상들이 전사인 $R$-가군의 역계라 하자.
그러면 표준 사상
$$
\mathcal{O}_X \otimes_R (\lim M_n)
\longrightarrow
\lim \mathcal{O}_X \otimes_R M_n
$$
은 원천에서 표적에 $DQ_X$를 적용한 것으로 가는 동형을 유도한다.
\end{lemma}

\begin{proof}
진술의 뜻은 임의의 대상 $E$가
$D_\QCoh(\mathcal{O}_X)$에 주어졌을 때 유도된 사상
$$
\Hom(E, \mathcal{O}_X \otimes_R (\lim M_n))
\longrightarrow
\Hom(E, \lim \mathcal{O}_X \otimes_R M_n)
$$
이 동형이라는 것이다. $D_\QCoh(\mathcal{O}_X)$는 완전 생성자를
가지므로
(Theorem \ref{spaces-perfect-theorem-bondal-van-den-Bergh}), 완전 $E$에 대해서만
이를 확인하면 충분하다. Lemma \ref{spaces-perfect-lemma-Rlim-quasi-coherent}에
의해
$\lim \mathcal{O}_X \otimes_R M_n =
R\lim \mathcal{O}_X \otimes_R M_n$이다. Cohomology on Sites,
Section \ref{sites-cohomology-section-global-RHom}의 완전함자
$R\Hom_X(E, -) : D_\QCoh(\mathcal{O}_X) \to D(R)$는 곱과 가환하고,
따라서 유도 극한과도 가환한다. 그러므로
$$
R\Hom_X(E, \lim \mathcal{O}_X \otimes_R M_n) =
R\lim R\Hom_X(E, \mathcal{O}_X \otimes_R M_n)
$$
이다. $E^\vee$를 쌍대 완전 복합체라 하자. Cohomology on Sites,
Lemma \ref{sites-cohomology-lemma-dual-perfect-complex}을 보라.
다음 등식들이 성립한다.
$$
R\Hom_X(E, \mathcal{O}_X \otimes_R M_n) =
R\Gamma(X, E^\vee \otimes_{\mathcal{O}_X}^\mathbf{L} Lf^*M_n) =
R\Gamma(X, E^\vee) \otimes_R^\mathbf{L} M_n
$$
이는 Lemma \ref{spaces-perfect-lemma-cohomology-base-change}에서 따른다.
Lemma \ref{spaces-perfect-lemma-flat-proper-perfect-direct-image-general}에 의해
$R\Gamma(X, E^\vee)$는 $R$-가군의 완전 복합체이다. 특히 이는
유사연접 복합체이므로 More on Algebra, Lemma
\ref{more-algebra-lemma-pseudo-coherent-tensor-limit}에 의해
$$
R\lim R\Gamma(X, E^\vee) \otimes_R^\mathbf{L} M_n =
R\Gamma(X, E^\vee) \otimes_R^\mathbf{L} \lim M_n
$$
을 얻으며, 이것이 원하는 결과이다.
\end{proof}

\begin{lemma}
\label{spaces-perfect-lemma-perfect-enough}
$A$를 환이라 하자. $X$를 $A$ 위의 준콤팩트 준분리 대수공간이라
하자. $K \in D^-_\QCoh(\mathcal{O}_X)$라 하자.
$R\Gamma(X, E \otimes^\mathbf{L} K)$가 $D(A)$에서 유사연접이라고
가정하되, 이를 모든 완전 대상 $E$가 $D(\mathcal{O}_X)$에
주어졌을 때 가정한다. 그러면
$R\Gamma(X, E \otimes^\mathbf{L} K)$는 $D(A)$에서 유사연접이며,
이는 모든 유사연접 대상 $E$가 $D(\mathcal{O}_X)$에 주어졌을 때
성립한다.
\end{lemma}

\begin{proof}
정수 $N$이 존재하여
$R\Gamma(X, -) : D_\QCoh(\mathcal{O}_X) \to D(A)$가 코호몰로지
차원 $N$을 가진다. 이는 Lemma
\ref{spaces-perfect-lemma-quasi-coherence-direct-image}에서 설명되어 있다.
$b \in \mathbf{Z}$를 $H^i(K) = 0$이 $i > b$에 대해 성립하도록
택하자. $E$가 $X$ 위에서 유사연접이라고 하자.
$R\Gamma(X, E \otimes^\mathbf{L} K)$가 모든 $m$에 대하여
$m$-유사연접임을 보이면 충분하다. 근사 $P \to E$를 택하되,
이를 완전 복합체 $P$에 의한 $(X, E, m - N - 1 - b)$의 근사라
하자. 이는 Theorem \ref{spaces-perfect-theorem-approximation}에 의해 가능하다.
구별 삼각형
$$
P \to E \to C \to P[1]
$$
을 $D_\QCoh(\mathcal{O}_X)$에서 택하자. $C$의 코호몰로지 층들은
차수 $\geq m - N - 1 - b$에서 영이다. 따라서
$C \otimes^\mathbf{L} K$의 코호몰로지 층들은 차수
$\geq m - N - 1$에서 영이다. 그러므로
$R\Gamma(X, C \otimes^\mathbf{L} K)$의 코호몰로지는 차수
$\geq m - 1$에서 영이다. 따라서
$$
R\Gamma(X, P \otimes^\mathbf{L} K) \to R\Gamma(X, E \otimes^\mathbf{L} K)
$$
는 차수 $\geq m$의 코호몰로지 위에서 동형이다. 가정에 의해
원천은 유사연접이다. 결론적으로
$R\Gamma(X, E \otimes^\mathbf{L} K)$는 원하는 대로
$m$-유사연접이다.
\end{proof}

\begin{lemma}
\label{spaces-perfect-lemma-base-change-RHom-perfect}
$S$를 스킴이라 하자.
$f : X \to Y$를 $S$ 위 대수공간들 사이의 유한 표시 사상이라
하자. $E \in D(\mathcal{O}_X)$를 완전 대상이라 하자.
$\mathcal{G}^\bullet$을 유한 표시 $\mathcal{O}_X$-가군의 유계
복합체로서 $Y$ 위에서 평탄하고 지지가 $Y$ 위에서 고유인 것이라
하자. 그러면
$$
K = Rf_*R\SheafHom(E, \mathcal{G}^\bullet)
$$
는 $D(\mathcal{O}_Y)$의 완전 대상이고, 그 형성은 임의의
밑변환과 가환한다.
\end{lemma}

\begin{proof}
밑변환에 관한 진술은 Lemma \ref{spaces-perfect-lemma-base-change-RHom}이다.
따라서 $K$가 완전 대상임을 보이면 충분하다. $Y$가 뇌터이면 이는
Lemma \ref{spaces-perfect-lemma-ext-perfect}에서 따른다. 뇌터 근사를 통해 이
경우로 귀착시킬 것이다. 독자는 이 증명의 나머지 부분을
건너뛰어도 좋다.

\medskip\noindent
문제는 $Y$ 위에서 국소적이므로 $Y$가 아핀이라고 가정할 수 있다.
$Y = \Spec(R)$이라 하자. $R = \colim R_i$를 뇌터 환들 $R_i$의
여과 여극한으로 쓴다. Limits of Spaces, Lemma
\ref{spaces-limits-lemma-descend-finite-presentation}에 의해 어떤
$i$와 대수공간 $X_i$가 존재하며, 이는 $R_i$ 위에서 유한 표시이고
이를 $R$로 밑변환한 것은 $X$이다. Limits of Spaces, Lemma
\ref{spaces-limits-lemma-descend-modules-finite-presentation}에
의해 $i$를 늘린 뒤, 유한 표시
$\mathcal{O}_{X_i}$-가군의 유계 복합체
$\mathcal{G}_i^\bullet$이 존재하고 이를 $X$로 당긴 것은
$\mathcal{G}$이라고 가정할 수 있다. % P09-ERR-0486
$i$를 늘린 뒤 $\mathcal{G}_i^n$이 $R_i$ 위에서 평탄하다고
가정할 수 있다. Limits of Spaces, Lemma
\ref{spaces-limits-lemma-descend-flat}을 보라. $i$를 늘린 뒤
$\mathcal{G}_i^n$의 지지가 $R_i$ 위에서 고유라고 가정할 수 있다.
Limits of Spaces, Lemma
\ref{spaces-limits-lemma-eventually-proper-support}를 보라. 마지막으로
Lemma \ref{spaces-perfect-lemma-descend-perfect}에 의해 $i$를 늘린 뒤 완전 대상
$E_i$가 $D(\mathcal{O}_{X_i})$에 존재하고, 이를 $X$로 당긴 것이
$E$라고 가정할 수 있다. Lemma
\ref{spaces-perfect-lemma-compute-ext-perfect}을
$X_i \to \Spec(R_i)$, $E_i$, $\mathcal{G}_i^\bullet$에 적용하고
앞서 보인 밑변환 성질을 사용하면 결과를 얻는다.
\end{proof}










\section{완전 복합체}
\label{spaces-perfect-section-perfect-complexes}

\noindent
먼저 완전 복합체의 베티 수가 뛰는 자취를 논한다. 우선 베티 수를
정의해야 한다.

\medskip\noindent
$S$를 스킴이라 하자. $X$를 $S$ 위의 대수공간이라 하자.
$E$를 $D(\mathcal{O}_X)$의 대상이라 하자. $x \in |X|$라 하자.
$\beta_i(x) \in \{0, 1, 2, \ldots \} \cup \{\infty\}$를 정의하고자
한다. 이를 위해 사상 $f : \Spec(k) \to X$를 택하되, 이는
$x$의 동치류에 속하게 한다. 그러면 $Lf^*E$는
$D(\Spec(k)_\etale, \mathcal{O})$의 대상이다. \'Etale Cohomology,
Lemma \ref{etale-cohomology-lemma-all-modules-quasi-coherent}와
Theorem \ref{etale-cohomology-theorem-quasi-coherent}에 의해
$D(\Spec(k)_\etale, \mathcal{O}) = D(k)$는 $k$-벡터공간의 유도
범주이다. 따라서 $Lf^*E$는 $k$-벡터공간의 복합체이고
$\beta_i(x) = \dim_k H^i(Lf^*E)$로 둘 수 있다. 이것이 $x$ 안의
대표원의 선택에 의존하지 않음은 쉽게 알 수 있다. 더욱이 $X$가
스킴이면 이는 Derived Categories of Schemes, Section
\ref{perfect-section-perfect-complexes}에서 사용한 개념과 같다.

\begin{lemma}
\label{spaces-perfect-lemma-jump-loci}
$S$를 스킴이라 하자. $X$를 $S$ 위의 대수공간이라 하자.
$E \in D(\mathcal{O}_X)$가 유사연접이라고 하자 (예를 들어
완전이라고 하자). 임의의 $i \in \mathbf{Z}$에 대하여 위에서
정의한 함수
$$
\beta_i : |X| \longrightarrow \{0, 1, 2, \ldots\}
$$
를 생각하자. 그러면 다음이 성립한다.
\begin{enumerate}
\item $\beta_i$의 형성은 임의의 밑변환과 가환한다.
\item 함수들 $\beta_i$는 위 반연속이다.
\item $\beta_i$의 준위집합들은 에탈 국소적으로 구성 가능하다.
\end{enumerate}
\end{lemma}

\begin{proof}
스킴 $U$와 전사 에탈 사상 $\varphi : U \to X$를 택하자. 그러면
$L\varphi^*E$는 스킴 $U$ 위의 유사연접 복합체이고
(Lemma \ref{spaces-perfect-lemma-descend-pseudo-coherent}을 사용한다), 스킴에
대한 결과를 적용할 수 있다. Derived Categories of Schemes,
Lemma \ref{perfect-lemma-jump-loci}를 보라. (3)의 뜻은 준위집합들의
$U$로의 역상이 국소적으로 구성 가능하다는 것이다. Properties of
Spaces, Definition
\ref{spaces-properties-definition-locally-constructible}을 보라.
\end{proof}

\begin{lemma}
\label{spaces-perfect-lemma-jump-loci-geometric}
$Y$를 스킴이라 하고 $X$를 $Y$ 위의 대수공간이라 하자. 구조 사상
$f : X \to Y$가 평탄하고 고유이며 유한 표시라고 하자.
$\mathcal{F}$를 유한 표시 $\mathcal{O}_X$-가군으로서 $Y$ 위에서
평탄한 것이라 하자. 고정된 $i \in \mathbf{Z}$에 대하여 함수
$$
\beta_i : |Y| \to \{0, 1, 2, \ldots\},\quad
y \longmapsto \dim_{\kappa(y)} H^i(X_y, \mathcal{F}_y)
$$
를 생각하자. 그러면 다음이 성립한다.
\begin{enumerate}
\item $\beta_i$의 형성은 임의의 밑변환과 가환한다.
\item 함수들 $\beta_i$는 위 반연속이다.
\item $\beta_i$의 준위집합들은 $Y$에서 국소적으로 구성 가능하다.
\end{enumerate}
\end{lemma}

\begin{proof}
코호몰로지와 밑변환, 더 정확히는 Lemma
\ref{spaces-perfect-lemma-flat-proper-perfect-direct-image-general}에 의해 대상
$K = Rf_*\mathcal{F}$는 $Y$의 유도 범주의 완전 대상이고, 그
형성은 임의의 밑변환과 가환한다. 특히
$$
H^i(X_y, \mathcal{F}_y) = H^i(K \otimes_{\mathcal{O}_Y}^\mathbf{L} \kappa(y))
$$
이다. 따라서 보조정리는 Lemma \ref{spaces-perfect-lemma-jump-loci}에서 따른다.
\end{proof}

\begin{lemma}
\label{spaces-perfect-lemma-chi-locally-constant}
$S$를 스킴이라 하자. $X$를 $S$ 위의 대수공간이라 하자.
$E \in D(\mathcal{O}_X)$가 완전이라고 하자. 함수
$$
\chi_E : |X| \longrightarrow \mathbf{Z},\quad
x \longmapsto \sum (-1)^i \beta_i(x)
$$
는 $X$ 위에서 국소 상수이다.
\end{lemma}

\begin{proof}
생략한다. 힌트: 에탈 국소화를 사용하여 스킴의 경우에서 따른다.
Derived Categories of Schemes, Lemma
\ref{perfect-lemma-chi-locally-constant}를 보라.
\end{proof}

\begin{lemma}
\label{spaces-perfect-lemma-open-where-cohomology-in-degree-i-rank-r}
$S$를 스킴이라 하자. $X$를 $S$ 위의 대수공간이라 하자.
$E \in D(\mathcal{O}_X)$가 완전이라고 하자.
$i, r \in \mathbf{Z}$가 주어지면 다음 성질로 특징지어지는 열린
부분공간 $U \subset X$가 존재한다.
\begin{enumerate}
\item $E|_U \cong H^i(E|_U)[-i]$이고 $H^i(E|_U)$는 국소 자유
$\mathcal{O}_U$-가군이며 그 계수는 $r$이다.
\item 사상 $f : Y \to X$가 $U$를 경유하는 것은 $Lf^*E$가 계수
$r$인 국소 자유 가군을 차수 $i$에 놓은 것과 동형인 것과
동치이다.
\end{enumerate}
\end{lemma}

\begin{proof}
생략한다. 힌트: 에탈 국소화를 사용하여 스킴의 경우에서 따른다.
Derived Categories of Schemes, Lemma
\ref{perfect-lemma-open-where-cohomology-in-degree-i-rank-r}을 보라.
\end{proof}

\begin{lemma}
\label{spaces-perfect-lemma-open-where-cohomology-in-degree-i-rank-r-geometric}
$S$를 스킴이라 하자.
$f : X \to Y$를 $S$ 위 대수공간의 고유하고 평탄한 유한 표시
사상이라 하자. $\mathcal{F}$를 유한 표시
$\mathcal{O}_X$-가군으로서 $Y$ 위에서 평탄한 것이라 하자.
$i, r \in \mathbf{Z}$를 고정하자. 그러면 다음 성질을 갖는 열린
부분공간 $V \subset Y$가 존재한다. 사상 $T \to Y$가 $V$를
경유하는 것은 $Rf_{T, *}\mathcal{F}_T$가 계수 $r$인 유한 국소
자유 가군을 차수 $i$에 놓은 것과 동형인 것과 동치이다.
\end{lemma}

\begin{proof}
코호몰로지와 밑변환
(Lemma \ref{spaces-perfect-lemma-flat-proper-perfect-direct-image-general})에 의해
대상 $K = Rf_*\mathcal{F}$는 $Y$의 유도 범주의 완전 대상이고,
그 형성은 임의의 밑변환과 가환한다. 따라서 이 보조정리는 Lemma
\ref{spaces-perfect-lemma-open-where-cohomology-in-degree-i-rank-r}에서 곧바로
따른다.
\end{proof}

\begin{lemma}
\label{spaces-perfect-lemma-locally-closed-where-H0-locally-free}
$S$를 스킴이라 하자. $X$를 $S$ 위의 대수공간이라 하자.
$E \in D(\mathcal{O}_X)$가 완전이고, Tor-진폭이 $[a, b]$ 안에
있다고 하자. 여기서 $a, b \in \mathbf{Z}$이다. $r \geq 0$이라
하자. 그러면 다음 성질로 특징지어지는
국소 닫힌 부분공간 $j : Z \to X$가 존재한다.
\begin{enumerate}
\item $H^a(Lj^*E)$는 국소 자유 $\mathcal{O}_Z$-가군이며 그
계수는 $r$이다.
\item 사상 $f : Y \to X$가 $Z$를 경유하는 것은 모든 사상
$g : Y' \to Y$에 대하여 $\mathcal{O}_{Y'}$-가군
$H^a(L(f \circ g)^*E)$가 계수 $r$인 국소 자유 가군인 것과
동치이다.
\end{enumerate}
더욱이 $j : Z \to X$는 유한 표시이고 다음이 성립한다.
\begin{enumerate}
\item[(3)] $f : Y \to X$가 $Y \xrightarrow{g} Z \to X$로
분해되면
$H^a(Lf^*E) = g^*H^a(Lj^*E)$이다.
\item[(4)] $\beta_a(x) \leq r$가 모든 $x \in |X|$에 대하여
성립하면
$j$는 닫힌 몰입이고, $f : Y \to X$가 주어졌을 때 다음은
동치이다.
\begin{enumerate}
\item $f : Y \to X$는 $Z$를 경유한다.
\item $H^0(Lf^*E)$는 국소 자유 $\mathcal{O}_Y$-가군이며 그
계수는 $r$이다. % P09-ERR-0487
\end{enumerate}
또한 $r = 1$이면 이들은 다음과도 동치이다.
\begin{enumerate}
\item[(c)]
$\mathcal{O}_Y \to \SheafHom_{\mathcal{O}_Y}(H^0(Lf^*E), H^0(Lf^*E))$
는 단사이다. % P09-ERR-0487
\end{enumerate}
\end{enumerate}
\end{lemma}

\begin{proof}
생략한다. 힌트: 에탈 국소화를 사용하여 스킴의 경우에서 따른다.
Derived Categories of Schemes, Lemma
\ref{perfect-lemma-locally-closed-where-H0-locally-free}를 보라.
\end{proof}

\begin{lemma}
\label{spaces-perfect-lemma-proper-flat-h0}
$S$를 스킴이라 하자.
$f : X \to Y$를 $S$ 위 대수공간의 사상이라 하자. 다음을
가정하자.
\begin{enumerate}
\item $f$는 고유이고 평탄하며 유한 표시이다.
\item 사상 $\Spec(k) \to Y$가 주어지고 $k$가 체이면
$k = H^0(X_k, \mathcal{O}_{X_k})$이다.
\end{enumerate}
그러면 다음이 성립한다.
\begin{enumerate}
\item[(a)] $f_*\mathcal{O}_X = \mathcal{O}_Y$이고, 이는 임의의
밑변환 뒤에도 성립한다.
\item[(b)] $Y$ 위에서 에탈 국소적으로
$$
Rf_*\mathcal{O}_X = \mathcal{O}_Y \oplus P
$$
가 $D(\mathcal{O}_Y)$ 안에서 성립한다. 여기서 $P$는
Tor-진폭이 $[1, \infty)$ 안에 있는 완전 대상이다.
\end{enumerate}
\end{lemma}

\begin{proof}
(a)와 (b)를 $Y$ 위에서 에탈 국소적으로 증명하면 충분하므로
$Y$가 아핀 스킴이라고 가정할 수 있고 실제로 그렇게 하자.
코호몰로지와 밑변환
(Lemma \ref{spaces-perfect-lemma-flat-proper-perfect-direct-image-general})에 의해
복합체 $E = Rf_*\mathcal{O}_X$는 완전이고 그 형성은 임의의
밑변환과 가환한다. 특히 $y \in Y$에 대하여
$$
H^0(E \otimes^\mathbf{L} \kappa(y)) =
H^0(X_y, \mathcal{O}_{X_y}) = \kappa(y)
$$
임을 안다. 따라서 Lemma \ref{spaces-perfect-lemma-jump-loci}의 표기를 쓰면
$\beta_0(y) \leq 1$이 모든 $y \in Y$에 대하여 성립한다.
Lemma \ref{spaces-perfect-lemma-locally-closed-where-H0-locally-free}를
$a = 0$과 $r = 1$로 적용하자. 그러면 보편 닫힌 부분스킴
$j : Z \to Y$를 얻으며, 여기서 $H^0(Lj^*E)$가 가역이다. 이는
그 보조정리의
(4)(a), (b), (c)의 동치로 특징지어진다. $E$의 형성은 밑변환과
가환하므로
$$
Lf^*E = R\text{pr}_{1, *}\mathcal{O}_{X \times_Y X}
$$
이다. 사상 $\text{pr}_1 : X \times_Y X$는 절단면을 가지며,
그것은 대각 사상 $\Delta$이다. 이는 $X$의 $Y$ 위 대각이다.
사상들
$$
\mathcal{O}_X \longrightarrow R\text{pr}_{1, *}\mathcal{O}_{X \times_Y X}
\longrightarrow \mathcal{O}_X
$$
을 $D(\mathcal{O}_X)$에서 얻고 그 합성은 항등이다. 따라서
$R\text{pr}_{1, *}\mathcal{O}_{X \times_Y X} = \mathcal{O}_X \oplus E'$
가 $D(\mathcal{O}_X)$에서 성립한다. 따라서 $\mathcal{O}_X$는
$H^0(Lf^*E)$의 직합 인자이고, 위에서 인용한 보조정리의 (4)(c)와
(4)(a)의 동치에 의해 $X \to Y$가 $Z$를 경유한다고 결론짓는다.
$\{X \to Y\}$가 fppf 덮개이므로 $Z = Y$이다. 따라서
$f_*\mathcal{O}_X$는 가역 $\mathcal{O}_Y$-가군이다. 따라서
$\mathcal{O}_Y \to f_*\mathcal{O}_X$가 동형이라고 결론짓는다.
실제로 환 사상 $A \to B$에서 $B$가 $A$-가군으로 가역이면 그
사상은 동형이다. 가정들은 밑변환 아래에서 보존되므로 (a)가
참이다.

\medskip\noindent
(b)의 증명. 위에서 모든 $y \in Y$에 대하여 사상
$\mathcal{O}_Y \to H^0(E \otimes^\mathbf{L} \kappa(y))$가
전사임을 보였다. 따라서 More on Algebra, Lemma
\ref{more-algebra-lemma-better-cut-complex-in-two}를 적용하면
$y$의 어떤 열린 근방에서 분해
$Rf_*\mathcal{O}_X = \mathcal{O}_Y \oplus P$를 얻는다.
\end{proof}

\begin{lemma}
\label{spaces-perfect-lemma-proper-flat-geom-red-connected}
$S$를 스킴이라 하자.
$f : X \to Y$를 $S$ 위 대수공간의 사상이라 하자. 다음을
가정하자.
\begin{enumerate}
\item $f$는 고유이고 평탄하며 유한 표시이다.
\item $f$의 기하학적 올들은 축약이고 연결이다.
\end{enumerate}
그러면 $f_*\mathcal{O}_X = \mathcal{O}_Y$이고, 이는 임의의
밑변환 뒤에도 성립한다.
\end{lemma}

\begin{proof}
Lemma \ref{spaces-perfect-lemma-proper-flat-h0}에 의해
$k = H^0(X_k, \mathcal{O}_{X_k})$임을 모든 사상
$\Spec(k) \to Y$에 대하여 보이면 충분하다. 여기서 $k$는 체이다.
이는 Spaces
over Fields, Lemma
\ref{spaces-over-fields-lemma-proper-geometrically-reduced-global-sections}와
$X_k$가 기하학적으로 연결이고 기하학적으로 축약이라는 사실에서
따른다.
\end{proof}








\section{그 밖의 응용}
\label{spaces-perfect-section-other-applications}

\noindent
이 절에서는 위에서 전개한 이론으로부터 도출할 수 있는 몇 가지
결과를 진술하고 증명한다.

\begin{lemma}
\label{spaces-perfect-lemma-countable-cohomology}
$S$를 스킴이라 하자. $X$를 $S$ 위의 준콤팩트 준분리 대수공간이라
하자. $K$를 $D_\QCoh(\mathcal{O}_X)$의 대상으로서, 그 코호몰로지
층 $H^i(K)$의 절단들의 집합이 $X$ 위에서 에탈인 아핀 스킴들
위에서 가산인 것이라 하자. 그러면 임의의 준콤팩트 준분리 에탈
사상 $U \to X$와 임의의 완전 대상 $E$가
$D(\mathcal{O}_X)$에 주어졌을 때 집합들
$$
H^i(U, K \otimes^\mathbf{L} E),\quad \Ext^i(E|_U, K|_U)
$$
은 가산이다.
\end{lemma}

\begin{proof}
Cohomology on Sites, Lemma
\ref{sites-cohomology-lemma-dual-perfect-complex}을 사용하면 군들
$H^i(U, K \otimes^\mathbf{L} E)$에 대해 결과를 증명하는 것으로
충분함을 알 수 있다. 귀납 원리를 사용하여 보조정리를 증명할
것이다. Lemma \ref{spaces-perfect-lemma-induction-principle}을 보라.

\medskip\noindent
$U = \Spec(A)$가 아핀이면 결과는 스킴의 경우에서 따른다. Derived
Categories of Schemes, Lemma
\ref{perfect-lemma-countable-cohomology}를 보라.

\medskip\noindent
증명을 끝내려면 다음을 보이면 충분하다. 기본 구별 삼각형
$(U \subset W, V \to W)$이 주어지고 결과가 $U$, $V$,
$U \times_W V$에 대하여 성립하면, 결과는 $W$에 대하여도
성립한다. 이는 Mayer--Vietoris 열의 즉각적인 귀결이다. Lemma
\ref{spaces-perfect-lemma-unbounded-mayer-vietoris}를 보라.
\end{proof}

\begin{lemma}
\label{spaces-perfect-lemma-countable}
$S$를 스킴이라 하자. $X$를 $S$ 위의 준콤팩트 준분리 대수공간이라
하자. $\mathcal{O}_X$의 절단들의 집합이 $X$ 위에서 에탈인 아핀들
위에서 가산이라고 가정하자. $K$를
$D_\QCoh(\mathcal{O}_X)$의 대상이라 하자. 다음은 동치이다.
\begin{enumerate}
\item $K = \text{hocolim} E_n$이고, $E_n$은
$D(\mathcal{O}_X)$의 완전 대상이다.
\item 코호몰로지 층들 $H^i(K)$의 절단들의 집합은 $X$ 위에서
에탈인 아핀들 위에서 가산이다.
\end{enumerate}
\end{lemma}

\begin{proof}
(1)이 참이면 (2)도 참이다. 실제로 호모토피 여극한은 코호몰로지
층을 취하는 것과 가환하고
(Derived Categories, Lemma
\ref{derived-lemma-cohomology-of-hocolim}에 의해), 완전 복합체는
국소적으로 유한 생성 자유 $\mathcal{O}_X$-가군의 유한 복합체와
동형이어서 $X$에 관한 가정에 의해 (2)를 만족한다.

\medskip\noindent
(2)를 가정하자.
K-단사 복합체 $\mathcal{K}^\bullet$을 택하여 $K$를 표현하게 하자.
완전 생성자 $E$를 $D_\QCoh(\mathcal{O}_X)$에서 택하고, 이를
K-단사 복합체 $\mathcal{I}^\bullet$로 표현하자. Theorem
\ref{spaces-perfect-theorem-DQCoh-is-Ddga}와 그 증명에 따르면 삼각 범주들의
동치
$F : D_\QCoh(\mathcal{O}_X) \to D(A, \text{d})$가 존재한다.
여기서 $(A, \text{d})$는 미분 등급 대수
$$
(A, \text{d}) =
\Hom_{\text{Comp}^{dg}(\mathcal{O}_X)}
(\mathcal{I}^\bullet, \mathcal{I}^\bullet)
$$
이다. 이 동치는 $K$를 미분 등급 가군
$$
M = \Hom_{\text{Comp}^{dg}(\mathcal{O}_X)}
(\mathcal{I}^\bullet, \mathcal{K}^\bullet)
$$
으로 보낸다. $H^i(A) = \Ext^i(E, E)$이고
$H^i(M) = \Ext^i(E, K)$임을 주목하자. 더욱이 $F$는 동치이므로
그것과 그 준역은 호모토피 여극한과 가환한다. 따라서 $M$을
$D(A, \text{d})$의 콤팩트 대상들의 호모토피 여극한으로 쓰면
충분하다. Differential Graded Algebra, Lemma
\ref{dga-lemma-countable}에 의해
$\Ext^i(E, E)$와 $\Ext^i(E, K)$가 각 $i$에 대하여 가산임을
보이면 충분하다. 이는 Lemma
\ref{spaces-perfect-lemma-countable-cohomology}에서 따른다.
\end{proof}

\begin{lemma}
\label{spaces-perfect-lemma-computing-sections-as-colim}
$A$를 환이라 하자. $f : U \to X$를 $A$ 위 대수공간의 평탄한
유한 표시 사상이라 하자. 그러면 다음이 성립한다.
\begin{enumerate}
\item 완전 대상들 $L_n$이 $D(\mathcal{O}_X)$에 놓이는 역계가
존재하여
$$
R\Gamma(U, Lf^*K) = \text{hocolim}\ R\Hom_X(L_n, K)
$$
가 $D(A)$ 안에서, $K$가 $D_\QCoh(\mathcal{O}_X)$에 주어졌을 때
함자적으로 성립한다.
\item 완전 대상들 $E_n$이 $D(\mathcal{O}_X)$에 놓이는 계가
존재하여
$$
R\Gamma(U, Lf^*K) = \text{hocolim}\ R\Gamma(X, E_n \otimes^\mathbf{L} K)
$$
가 $D(A)$ 안에서, $K$가 $D_\QCoh(\mathcal{O}_X)$에 주어졌을 때
함자적으로 성립한다.
\end{enumerate}
\end{lemma}

\begin{proof}
Lemma \ref{spaces-perfect-lemma-cohomology-base-change}에 의해
$$
R\Gamma(U, Lf^*K) = R\Gamma(X, Rf_*\mathcal{O}_U \otimes^\mathbf{L} K)
$$
가 $K$에 대해 함자적으로 성립한다. $R\Gamma(X, -)$는 호모토피
여극한과 가환한다. 실제로 Lemma
\ref{spaces-perfect-lemma-quasi-coherence-pushforward-direct-sums}에 의해 직합과
가환한다. 마찬가지로 $- \otimes^\mathbf{L} K$는 유도 여극한과
가환한다. 이는 $- \otimes^\mathbf{L} K$가 직합과 가환하기
때문이다 ($D(\mathcal{O}_X)$의 직합은 대표 복합체들의 직합으로
주어지기 때문이다). 따라서 (2)를 증명하려면
$Rf_*\mathcal{O}_U = \text{hocolim} E_n$을 완전 대상들 $E_n$이
$D(\mathcal{O}_X)$에 놓이는 계로 쓰면 충분하다.
그렇게 한 뒤 $L_n = E_n^\vee$로 놓으면 (1)을 얻는다. Cohomology
on Sites, Lemma
\ref{sites-cohomology-lemma-dual-perfect-complex}을 보라.

\medskip\noindent
$A = \colim A_i$로 쓰되 $A_i$가 $\mathbf{Z}$ 위에서 유한형이게
하자. Limits of Spaces, Lemma
\ref{spaces-limits-lemma-descend-finite-presentation}에 의해 어떤
$i$와 유한 표시 사상들
$U_i \to X_i \to \Spec(A_i)$를 찾을 수 있고, 이들을
$\Spec(A)$로 밑변환하면 $U \to X \to \Spec(A)$를 되찾는다.
$i$를 늘린 뒤 $f_i : U_i \to X_i$가 평탄하다고 가정할 수 있다.
Limits of Spaces, Lemma
\ref{spaces-limits-lemma-descend-flat}을 보라. Lemma
\ref{spaces-perfect-lemma-compare-base-change}에 의해
$Rf_{i, *}\mathcal{O}_{U_i}$를 $g : X \to X_i$로 유도 당긴 것은
$Rf_*\mathcal{O}_U$와 같다. $Lg^*$는 유도 여극한과 가환하므로
$f_i$에 대해 원하는 결과를 증명하면 충분하다. 따라서 $U$와
$X$가 $\mathbf{Z}$ 위에서 유한형이라고 가정할 수 있다.

\medskip\noindent
$f : U \to X$가 $\mathbf{Z}$ 위에서 유한형인 대수공간의
사상이라고 하자. 증명을 끝내기 위해 $Rf_*\mathcal{O}_U$가 완전
복합체들의 호모토피 여극한임을 보일 것이다. 이를 위해 Lemma
\ref{spaces-perfect-lemma-countable}을 적용한다. 따라서
$R^if_*\mathcal{O}_U$의 절단들의 집합이 $X$ 위에서 에탈인
아핀들 위에서 가산임을 보이면 충분하다. 이는 구조층에 Lemma
\ref{spaces-perfect-lemma-countable-cohomology}를 적용하면 따른다.
\end{proof}








\section{해소 성질}
\label{spaces-perfect-section-resolution-property}

\noindent
이 절은 대수공간에 대한 Derived Categories of Schemes, Section
\ref{perfect-section-resolution-property}의 유사물이다. 먼저 그
절을 읽기를 권한다. 체 위의 매끄럽고 고유한 대수공간이 항상
해소 성질을 가지는지, 아니면 그렇지 않은지는 현재 알려져 있지
않다. 이 질문의 답을 알고 있다면
\href{mailto:stacks.project@gmail.com}{stacks.project@gmail.com}으로
전자우편을 보내기 바란다.

\medskip\noindent
일반 대수공간에 대하여 이를 생각하는 것은 별 의미가 없지만
다음 정의를 내릴 수 있다.

\begin{definition}
\label{spaces-perfect-definition-resolution-property}
$S$를 스킴이라 하자. $X$를 $S$ 위의 대수공간이라 하자.
$X$가 {\it 해소 성질}을 가진다는 것은 모든 유한형 준연접
$\mathcal{O}_X$-가군이 유한 국소 자유 $\mathcal{O}_X$-가군의
몫이라는 뜻이다.
\end{definition}

\noindent
$X$가 준콤팩트 준분리 대수공간이면, 유한 표시인 모든
$\mathcal{O}_X$-가군이 (자동으로 준연접이다) 유한 국소 자유
$\mathcal{O}_X$-가군의 몫임을 확인하면 충분하다. % P09-ERR-0488
Limits of Spaces, Lemma
\ref{spaces-limits-lemma-finite-directed-colimit-surjective-maps}를
보라. $X$가 뇌터 대수공간이면 유한형 준연접 가군들은 정확히
연접 $\mathcal{O}_X$-가군들이다. Cohomology of Spaces, Lemma
\ref{spaces-cohomology-lemma-coherent-Noetherian}을 보라.

\begin{lemma}
\label{spaces-perfect-lemma-resolution-property-ample-relative}
$S$를 스킴이라 하자.
$f : X \to Y$를 $S$ 위 대수공간의 사상이라 하자. 다음을
가정하자.
\begin{enumerate}
\item $Y$는 준콤팩트이고 준분리이며 해소 성질을 가진다.
\item $f$-풍부 가역 가군이 $X$ 위에 존재한다
(Divisors on Spaces, Definition
\ref{spaces-divisors-definition-relatively-ample}).
\end{enumerate}
그러면 $X$는 해소 성질을 가진다.
\end{lemma}

\begin{proof}
$\mathcal{F}$를 유한형 준연접 $\mathcal{O}_X$-가군이라 하자.
$\mathcal{L}$을 $f$-풍부 가역 가군이라 하자. 아핀 스킴 $V$와
전사 에탈 사상 $V \to Y$를 택하자. $U = V \times_Y X$로 놓자.
그러면 $\mathcal{L}|_U$는 $U$ 위에서 풍부하다. Properties,
Proposition \ref{properties-proposition-characterize-ample}에 의해
유한 개의 사상
$s_i : \mathcal{L}^{\otimes n_i}|_U \to \mathcal{F}|_U$가 존재하며
이들은 함께 전사이다. 준연접 $\mathcal{O}_Y$-가군들
$$
\mathcal{H}_n =
f_*(\mathcal{F} \otimes_{\mathcal{O}_X} \mathcal{L}^{\otimes n})
$$
을 생각하자. $s_i$를 $V$ 위에서 층 $\mathcal{H}_{-n_i}$의
절단으로 생각할 수 있다. 유한 국소 자유 $\mathcal{O}_Y$-가군
$\mathcal{E}_i$와 사상
$\mathcal{E}_i \to \mathcal{H}_{-n_i}$를 찾아서 $s_i$가 그 상에
속하게 할 수 있다고 하자. 그러면 대응하는 사상들
$$
f^*\mathcal{E}_i
\otimes_{\mathcal{O}_X}
\mathcal{L}^{\otimes n_i} \longrightarrow \mathcal{F}
$$
은 함께 전사이며 보조정리가 증명된다. Limits of Spaces, Lemma
\ref{spaces-limits-lemma-directed-colimit-finite-type}에 의해
각 $i$에 대하여 유한형 준연접 부분가군
$\mathcal{H}'_i \subset  \mathcal{H}_{-n_i}$를 찾을 수 있고, 이것은
절단 $s_i$를 $V$ 위에서 포함한다.
따라서 $Y$의 해소 성질은 전사
$\mathcal{E}_i \to \mathcal{H}'_i$를 주며, 이로써 결론을 얻는다.
\end{proof}

\begin{lemma}
\label{spaces-perfect-lemma-resolution-property-goes-up-affine}
$S$를 스킴이라 하자.
$f : X \to Y$를 $S$ 위 대수공간의 아핀 또는 준아핀 사상이라
하고, $Y$가 준콤팩트이고 준분리라고 하자. $Y$가 해소 성질을
가지면 $X$도 그러하다.
\end{lemma}

\begin{proof}
Divisors on Spaces, Lemma
\ref{spaces-divisors-lemma-quasi-affine-O-ample}에 의해 이는 Lemma
\ref{spaces-perfect-lemma-resolution-property-ample-relative}의 특수한 경우이다.
\end{proof}

\noindent
해소 성질이 아래로 내려감을 증명할 수 있는 한 경우는 다음과
같다.

\begin{lemma}
\label{spaces-perfect-lemma-resolution-property-goes-down-finite-flat}
$S$를 스킴이라 하자.
$f : X \to Y$를 $S$ 위 대수공간의 전사 유한 국소 자유 사상이라
하자. $X$가 해소 성질을 가지면 $Y$도 그러하다.
\end{lemma}

\begin{proof}
이 조건은 $f$가 아핀이고 $f_*\mathcal{O}_X$가 양의 계수를 가진
유한 국소 자유 $\mathcal{O}_Y$-가군이라는 뜻이다.
$\mathcal{G}$를 유한형 준연접 $\mathcal{O}_Y$-가군이라 하자.
가정에 의해 전사 $\mathcal{E} \to f^*\mathcal{G}$가 존재한다.
여기서 $\mathcal{O}_X$-가군 $\mathcal{E}$는 유한 국소 자유이다.
$f_*$가 완전함자이므로
(Cohomology of Spaces, Section
\ref{spaces-cohomology-section-finite-morphisms}), 전사
$$
f_*\mathcal{E}
\longrightarrow
f_*f^*\mathcal{G} = \mathcal{G} \otimes_{\mathcal{O}_Y} f_*\mathcal{O}_X
$$
를 얻는다. 쌍대를 취하면 전사
$$
f_*\mathcal{E}
\otimes_{\mathcal{O}_Y}
\SheafHom_{\mathcal{O}_Y}(f_*\mathcal{O}_X, \mathcal{O}_Y)
\longrightarrow
\mathcal{G}
$$
를 얻는다. $f_*\mathcal{E}$가 유한 국소 자유이므로 결론을
얻는다.
\end{proof}

\noindent
대수공간의 해소 성질에 관해서는 More on Morphisms of Spaces,
Section \ref{spaces-more-morphisms-section-resolution-property}도
보라.












\section{유계성의 검출}
\label{spaces-perfect-section-detecting-boundedness}

\noindent
이 절에서는 Section \ref{spaces-perfect-section-generating}에서 구성한
준콤팩트 준분리 스킴의 $D_\QCoh$의 콤팩트 생성자들이 특별한
성질을 가짐을 보인다. 이 절과 매우 비슷하므로 먼저 그 절을
읽기를 권한다.

\begin{lemma}
\label{spaces-perfect-lemma-bounded-truncation}
$S$를 스킴이라 하자.
$X$를 $S$ 위의 준콤팩트 준분리 대수공간이라 하자.
$P \in D_{perf}(\mathcal{O}_X)$이고
$E \in D_{\QCoh}(\mathcal{O}_X)$라 하자.
$a \in \mathbf{Z}$라 하자. 다음은 서로 동치이다.
\begin{enumerate}
\item $\Hom_{D(\mathcal{O}_X)}(P[-i], E) = 0$이다. 여기서
$i \gg 0$이다.
\item $\Hom_{D(\mathcal{O}_X)}(P[-i], \tau_{\geq a} E) = 0$이다.
여기서 $i \gg 0$이다.
\end{enumerate}
\end{lemma}

\begin{proof}
삼각형 $\tau_{< a} E \to E \to \tau_{\geq a} E \to$를 쓰면,
$$
\Hom_{D(\mathcal{O}_X)}(P[-i], \tau_{< a} E) =
\Hom_{D(\mathcal{O}_X)}(P, (\tau_{< a} E)[i]) = 0
$$
임을 $i \gg 0$에 대하여 보일 수 있으면 이 동치가 따름을 알 수
있다. $P$가 완전하므로
이는 Lemma \ref{spaces-perfect-lemma-ext-from-perfect-into-bounded-QCoh}에 의해
참이다.
\end{proof}

\begin{lemma}
\label{spaces-perfect-lemma-bounded-below-truncation}
$S$를 스킴이라 하자.
$X$를 $S$ 위의 준콤팩트 준분리 대수공간이라 하자.
$P \in D_{perf}(\mathcal{O}_X)$이고
$E \in D_{\QCoh}(\mathcal{O}_X)$라 하자.
$a \in \mathbf{Z}$라 하자. 다음은 서로 동치이다.
\begin{enumerate}
\item $\Hom_{D(\mathcal{O}_X)}(P[-i], E) = 0$이다. 여기서
$i \ll 0$이다.
\item $\Hom_{D(\mathcal{O}_X)}(P[-i], \tau_{\leq a} E) = 0$이다.
여기서 $i \ll 0$이다.
\end{enumerate}
\end{lemma}

\begin{proof}
삼각형 $\tau_{\leq a} E \to E \to \tau_{> a} E \to$를 쓰면,
$$
\Hom_{D(\mathcal{O}_X)}(P[-i], \tau_{> a} E) =
\Hom_{D(\mathcal{O}_X)}(P, (\tau_{> a} E)[i]) = 0
$$
임을 $i \ll 0$에 대하여 보일 수 있으면 이 동치가 따름을 알 수
있다. $P$가 완전하므로
이는 Lemma \ref{spaces-perfect-lemma-ext-from-perfect-into-bounded-QCoh}에 의해
참이다.
\end{proof}

\begin{proposition}
\label{spaces-perfect-proposition-detecting-bounded-above}
$S$를 스킴이라 하자.
$X$를 $S$ 위의 준콤팩트 준분리 대수공간이라 하자.
$G \in D_{perf}(\mathcal{O}_X)$를
$D_\QCoh (\mathcal{O}_X)$를 생성하는 완전 복합체라 하자.
$E \in D_\QCoh (\mathcal{O}_X)$라 하자. 다음은 서로 동치이다.
\begin{enumerate}
\item $E \in D^-_\QCoh (\mathcal{O}_X)$이다.
\item $\Hom_{D(\mathcal{O}_X)}(G[-i], E) = 0$이다. 여기서
$i \gg 0$이다.
\item $\Ext^i_X(G, E) = 0$이다. 여기서 $i \gg 0$이다.
\item $R\Hom_X(G, E)$는 $D^-(\mathbf{Z})$에 속한다.
\item $H^i(X, G^\vee \otimes_{\mathcal{O}_X}^\mathbf{L} E) = 0$이다.
여기서 $i \gg 0$이다.
\item $R\Gamma(X, G^\vee \otimes_{\mathcal{O}_X}^\mathbf{L} E)$는
$D^-(\mathbf{Z})$에 속한다.
\item $P$를 $D(\mathcal{O}_X)$의 임의의 완전 대상으로 두면
\begin{enumerate}
\item (2), (3), (4)의 명제들에서 $G$를 $P$로 바꾼 것이 성립한다.
\item $H^i(X, P \otimes_{\mathcal{O}_X}^\mathbf{L} E) = 0$이다.
여기서 $i \gg 0$이다.
\item $R\Gamma(X, P \otimes_{\mathcal{O}_X}^\mathbf{L} E)$는
$D^-(\mathbf{Z})$에 속한다.
\end{enumerate}
\end{enumerate}
\end{proposition}

\begin{proof}
(1)을 가정하자.
$\Hom_{D(\mathcal{O}_X)}(G[-i], E) = \Hom_{D(\mathcal{O}_X)}(G, E[i])$
이므로 Lemma \ref{spaces-perfect-lemma-ext-from-perfect-into-bounded-QCoh}에 의해
$i \gg 0$일 때 이것은 영이다. 따라서 (1)은 (2)를 함의한다.

\medskip\noindent
Cohomology on Sites, Section
\ref{sites-cohomology-section-global-RHom}의 논의에 의해 (2), (3),
(4)는 서로 동치이다. 정의에 의해
$H^i(X, -) = H^i(R\Gamma(X, -))$이므로 (5)와 (6)은 동치이다.
동치인 조건 (2), (3), (4)는 Cohomology on Sites, Lemma
\ref{sites-cohomology-lemma-dual-perfect-complex}과 등식
$(G[-i])^\vee = G^\vee[i]$에 의해 동치인 조건 (5), (6)과
동치이다.

\medskip\noindent
(7)이 (2)를 함의함은 분명하다. 역으로 동치인 조건 (2)--(6)이
(7)을 함의함을 증명하자. Remark
\ref{spaces-perfect-remark-classical-generator}에 의해 $G$가
$D_{perf}(\mathcal{O}_X)$의 고전적 생성자임을 상기하자.
$P \in D_{perf}(\mathcal{O}_X)$에 대하여 $T(P)$를
$R\Hom_X(P, E)$가 $D^-(\mathbf{Z})$에 속한다는 명제라 하자.
$T$는 분명히 직합에 의해 유전되고, 구별 삼각형에 대한 셋 중 둘
성질을 만족하며, 직합인자에 의해 유전되고, 시프트에 의해
보존된다. 따라서 Derived Categories, Remark
\ref{derived-remark-check-on-generator}에 의해 (4)는 $T$가
$D_{perf}(\mathcal{O}_X)$ 전체에서 성립함을 함의한다. 같은 논증이
다른 모든 성질에도 적용되는데, 성질 (7)(b)와 (7)(c)에 대해서는
$P \mapsto P^\vee$가 $D_{perf}(\mathcal{O}_X)$의 자기동치라는
사실도 사용한다. 작은 세부사항은 생략한다.

\medskip\noindent
동치인 조건 (2)--(7)이 (1)을 함의함을 Lemma
\ref{spaces-perfect-lemma-induction-principle}의 귀납 원리를 사용하여 증명할
것이다.

\medskip\noindent
먼저 (2)--(7) $\Rightarrow$ (1)임을 증명하자. 여기서 $X$는
아핀이라고 가정한다.
이는 스킴의 경우인 Derived Categories of Schemes, Proposition
\ref{perfect-proposition-detecting-bounded-above}에서 따른다.

\medskip\noindent
이제 $(U \subset X, j : V \to X)$가 $S$ 위 준콤팩트 준분리
대수공간들의 기본 구별 사각형이고, (2)--(7) $\Rightarrow$ (1)이
$U$, $V$, $U \times_X V$에 대하여 알려져 있다고 가정하자.
증명을 끝내려면 (2)--(7) $\Rightarrow$ (1)을 $X$에 대하여
보여야 한다. $E \in D_\QCoh(\mathcal{O}_X)$가 (2)--(7)을
만족한다고 하자. Lemma
\ref{spaces-perfect-lemma-direct-summand-of-a-restriction}와 Theorem
\ref{spaces-perfect-theorem-bondal-van-den-Bergh}에 의해 완전 복합체 $Q$가 $X$ 위에
존재하여 $Q|_U$는 $D_\QCoh (\mathcal{O}_U)$를 생성한다.

\medskip\noindent
$V = \Spec(A)$라 하자. $Z \subset V$를 $X \setminus U$의
역상이며 그곳으로 동형으로 사상되는 축약 폐부분스킴이라 하자
(Definition \ref{spaces-perfect-definition-elementary-distinguished-square}을
보라). 이는 $V$의 역콤팩트 폐부분집합이다.
$f_1, \ldots, f_r \in A$를 택하여
$Z = V(f_1, \ldots, f_r)$가 되게 하자.
$K \in D(\mathcal{O}_V)$를 $f_1, \ldots, f_r$에 대한
$A$ 위 코쥘 복합체에 대응하는 완전 대상이라
하자. $K$는 $Z$ 위에서 지지되므로 그 직상
$K' = Rj_*K$는 $D(\mathcal{O}_X)$의 완전 대상이고 $V$로 제한하면
$K$이다 (Lemmas \ref{spaces-perfect-lemma-pushforward-perfect}와
\ref{spaces-perfect-lemma-pushforward-with-support-in-open}을 보라).
가정에 의해 $R\Hom_{\mathcal{O}_X}(Q, E)$와
$R\Hom_{\mathcal{O}_X}(K', E)$는 위로 유계이다.

\medskip\noindent
Lemma \ref{spaces-perfect-lemma-pushforward-with-support-in-open}에 의해
$K' =  j_!K$이고, 따라서
$$
\Hom_{D(\mathcal{O}_X)}(K'[-i], E) = \Hom_{D(\mathcal{O}_V)}(K[-i], E|_V) = 0
$$
이다. 여기서 $i \gg 0$이다. 그러므로 Derived Categories of
Schemes, Lemma
\ref{perfect-lemma-orthogonal-koszul-first-variant}를 $E|_V$에
적용하면 정수 $a$가 존재하여
$\tau_{\geq a}(E|_V) =
\tau_{\geq a} R (U \times_X V \to V)_* (E|_{U \times_X V})$가 된다.
그러면
$\tau_{\geq a} E = \tau_{\geq a} R (U \to X)_* (E |_U)$이다
(표준 사상을 $U$와 $V$로 제한한 뒤 동형임을 확인하라).
따라서 Lemma \ref{spaces-perfect-lemma-bounded-truncation}을 두 번 사용하면
$$
\Hom_{D(\mathcal{O}_U)}(Q|_U [-i], E|_U) =
\Hom_{D(\mathcal{O}_X)}(Q[-i], R (U \to X)_* (E|_U)) = 0
$$
임을 알 수 있다. 여기서 $i \gg 0$이다. 이 Proposition은 $U$와
생성자 $Q|_U$에 대해 성립하므로
$E|_U \in D^-_\QCoh(\mathcal{O}_U)$이다. 그런데 함자
$R (U \to X)_*$는 $D^-_\QCoh$를 보존하므로
(Lemma \ref{spaces-perfect-lemma-quasi-coherence-direct-image}에 의해)
$\tau_{\geq a}E \in D^-_\QCoh(\mathcal{O}_X)$를 얻는다. 따라서
$E \in D^-_\QCoh (\mathcal{O}_X)$이다.
\end{proof}

\begin{proposition}
\label{spaces-perfect-proposition-detecting-bounded-below}
$S$를 스킴이라 하자.
$X$를 $S$ 위의 준콤팩트 준분리 대수공간이라 하자.
$G \in D_{perf}(\mathcal{O}_X)$를
$D_\QCoh (\mathcal{O}_X)$를 생성하는 완전 복합체라 하자.
$E \in D_\QCoh (\mathcal{O}_X)$라 하자. 다음은 서로 동치이다.
\begin{enumerate}
\item $E \in D^+_\QCoh (\mathcal{O}_X)$이다.
\item $\Hom_{D(\mathcal{O}_X)}(G[-i], E) = 0$이다. 여기서
$i \ll 0$이다.
\item $\Ext^i_X(G, E) = 0$이다. 여기서 $i \ll 0$이다.
\item $R\Hom_X(G, E)$는 $D^+(\mathbf{Z})$에 속한다.
\item $H^i(X, G^\vee \otimes_{\mathcal{O}_X}^\mathbf{L} E) = 0$이다.
여기서 $i \ll 0$이다.
\item $R\Gamma(X, G^\vee \otimes_{\mathcal{O}_X}^\mathbf{L} E)$는
$D^+(\mathbf{Z})$에 속한다.
\item $P$를 $D(\mathcal{O}_X)$의 임의의 완전 대상으로 두면
\begin{enumerate}
\item (2), (3), (4)의 명제들에서 $G$를 $P$로 바꾼 것이 성립한다.
\item $H^i(X, P \otimes_{\mathcal{O}_X}^\mathbf{L} E) = 0$이다.
여기서 $i \ll 0$이다.
\item $R\Gamma(X, P \otimes_{\mathcal{O}_X}^\mathbf{L} E)$는
$D^+(\mathbf{Z})$에 속한다.
\end{enumerate}
\end{enumerate}
\end{proposition}

\begin{proof}
(1)을 가정하자.
$\Hom_{D(\mathcal{O}_X)}(G[-i], E) = \Hom_{D(\mathcal{O}_X)}(G, E[i])$
이므로 Lemma \ref{spaces-perfect-lemma-ext-from-perfect-into-bounded-QCoh}에 의해
$i \ll 0$일 때 이것은 영이다. 따라서 (1)은 (2)를 함의한다.

\medskip\noindent
Cohomology on Sites, Section
\ref{sites-cohomology-section-global-RHom}의 논의에 의해 (2), (3),
(4)는 서로 동치이다. 정의에 의해
$H^i(X, -) = H^i(R\Gamma(X, -))$이므로 (5)와 (6)은 동치이다.
동치인 조건 (2), (3), (4)는 Cohomology on Sites, Lemma
\ref{sites-cohomology-lemma-dual-perfect-complex}과 등식
$(G[-i])^\vee = G^\vee[i]$에 의해 동치인 조건 (5), (6)과
동치이다.

\medskip\noindent
(7)이 (2)를 함의함은 분명하다. 역으로 동치인 조건 (2)--(6)이
(7)을 함의함을 증명하자. Remark
\ref{spaces-perfect-remark-classical-generator}에 의해 $G$가
$D_{perf}(\mathcal{O}_X)$의 고전적 생성자임을 상기하자.
$P \in D_{perf}(\mathcal{O}_X)$에 대하여 $T(P)$를
$R\Hom_X(P, E)$가 $D^+(\mathbf{Z})$에 속한다는 명제라 하자.
$T$는 분명히 직합에 의해 유전되고, 구별 삼각형에 대한 셋 중 둘
성질을 만족하며, 직합인자에 의해 유전되고, 시프트에 의해
보존된다. 따라서 Derived Categories, Remark
\ref{derived-remark-check-on-generator}에 의해 (4)는 $T$가
$D_{perf}(\mathcal{O}_X)$ 전체에서 성립함을 함의한다. 같은 논증이
다른 모든 성질에도 적용되는데, 성질 (7)(b)와 (7)(c)에 대해서는
$P \mapsto P^\vee$가 $D_{perf}(\mathcal{O}_X)$의 자기동치라는
사실도 사용한다. 작은 세부사항은 생략한다.

\medskip\noindent
동치인 조건 (2)--(7)이 (1)을 함의함을 Lemma
\ref{spaces-perfect-lemma-induction-principle}의 귀납 원리를 사용하여 증명할
것이다.

\medskip\noindent
먼저 (2)--(7) $\Rightarrow$ (1)임을 증명하자. 여기서 $X$는
아핀이라고 가정한다. 이는 스킴의 경우인 Derived Categories of
Schemes, Proposition
\ref{perfect-proposition-detecting-bounded-below}에서 따른다.

\medskip\noindent
이제 $(U \subset X, j : V \to X)$가 $S$ 위 준콤팩트 준분리
대수공간들의 기본 구별 사각형이고, (2)--(7) $\Rightarrow$ (1)이
$U$, $V$, $U \times_X V$에 대하여 알려져 있다고 가정하자.
증명을 끝내려면 (2)--(7) $\Rightarrow$ (1)을 $X$에 대하여
보여야 한다. $E \in D_\QCoh(\mathcal{O}_X)$가 (2)--(7)을
만족한다고 하자. Lemma
\ref{spaces-perfect-lemma-direct-summand-of-a-restriction}와 Theorem
\ref{spaces-perfect-theorem-bondal-van-den-Bergh}에 의해 완전 복합체 $Q$가 $X$ 위에
존재하여 $Q|_U$는 $D_\QCoh (\mathcal{O}_U)$를 생성한다.

\medskip\noindent
$V = \Spec(A)$라 하자. $Z \subset V$를 $X \setminus U$의
역상이며 그곳으로 동형으로 사상되는 축약 폐부분스킴이라 하자
(Definition \ref{spaces-perfect-definition-elementary-distinguished-square}을
보라). 이는 $V$의 역콤팩트 폐부분집합이다.
$f_1, \ldots, f_r \in A$를 택하여
$Z = V(f_1, \ldots, f_r)$가 되게 하자.
$K \in D(\mathcal{O}_V)$를 $f_1, \ldots, f_r$에 대한
$A$ 위 코쥘 복합체에 대응하는 완전 대상이라 하자. $K$는 $Z$
위에서 지지되므로 그 직상 $K' = Rj_*K$는
$D(\mathcal{O}_X)$의 완전 대상이고 $V$로 제한하면 $K$이다
(Lemmas \ref{spaces-perfect-lemma-pushforward-perfect}와
\ref{spaces-perfect-lemma-pushforward-with-support-in-open}을 보라).
가정에 의해 $R\Hom_{\mathcal{O}_X}(Q, E)$와
$R\Hom_{\mathcal{O}_X}(K', E)$는 아래로 유계이다.

\medskip\noindent
Lemma \ref{spaces-perfect-lemma-pushforward-with-support-in-open}에 의해
$K' =  j_!K$이고, 따라서
$$
\Hom_{D(\mathcal{O}_X)}(K'[-i], E) = \Hom_{D(\mathcal{O}_V)}(K[-i], E|_V) = 0
$$
이다. 여기서 $i \ll 0$이다. 그러므로 Derived Categories of
Schemes, Lemma
\ref{perfect-lemma-orthogonal-koszul-second-variant}를 $E|_V$에
적용하면 정수 $a$가 존재하여
$\tau_{\leq a}(E|_V) =
\tau_{\leq a} R (U \times_X V \to V)_* (E|_{U \times_X V})$가 된다.
그러면
$\tau_{\leq a} E = \tau_{\leq a} R (U \to X)_* (E |_U)$이다
(표준 사상을 $U$와 $V$로 제한한 뒤 동형임을 확인하라).
따라서 Lemma \ref{spaces-perfect-lemma-bounded-below-truncation}을 두 번 사용하면
$$
\Hom_{D(\mathcal{O}_U)}(Q|_U [-i], E|_U) =
\Hom_{D(\mathcal{O}_X)}(Q[-i], R (U \to X)_* (E|_U)) = 0
$$
임을 알 수 있다. 여기서 $i \ll 0$이다. 이 Proposition은 $U$와
생성자 $Q|_U$에 대해 성립하므로
$E|_U \in D^+_\QCoh(\mathcal{O}_U)$이다. 그런데 함자
$R (U \to X)_*$는 $D^+_\QCoh$를 보존하므로
(Lemma \ref{spaces-perfect-lemma-quasi-coherence-direct-image}에 의해)
$\tau_{\leq a}E \in D^+_\QCoh(\mathcal{O}_X)$를 얻는다. 따라서
$E \in D^+_\QCoh (\mathcal{O}_X)$이다.
\end{proof}












\section{유도범주의 준연접 대상}
\label{spaces-perfect-section-QC}

\noindent
$S$를 스킴이라 하자. $X$를 $S$ 위의 대수공간이라 하자.
$X_{affine, \etale}$는 $X_\etale$의 아핀 대상들로 이루어진
범주에 표준 에탈 덮개들이 주는 위상을 부여한 것임을 상기하자.
Properties of Spaces, Definition
\ref{spaces-properties-definition-affine-etale-site}을 보라.
$X_{affine, \etale}$의 토포스는 $X$의 작은 에탈 토포스이다.
Properties of Spaces, Lemma
\ref{spaces-properties-lemma-alternative}를 보라. 자리
$X_\etale$에는 구조층 $\mathcal{O}_X$가 주어져 있고, 그
$X_{affine, \etale}$로의 제한도 $\mathcal{O}_X$로 나타낸다.
그러면 환 달린 토포스들의 동치
$$
(\Sh(X_{affine, \etale}), \mathcal{O}_X)
\longrightarrow
(\Sh(X_\etale), \mathcal{O}_X)
$$
가 있다. Descent on Spaces, Equation
(\ref{spaces-descent-equation-alternative-small-ringed})과 Descent
on Spaces, Section
\ref{spaces-descent-section-alternative-quasi-coherent}의 논의를
보라.

\medskip\noindent
이 절에서는 $X_{affine}$으로 $X_{affine, \etale}$의 바탕 범주에
혼돈 위상을 부여한 것을 나타낸다. 즉, 층과 전층이
일치하도록 한다. 특히 구조층 $\mathcal{O}_X$는
$X_{affine}$ 위에서도 층이 된다. 이로부터 환 달린 자리들의
사상 % P09-ERR-0489
$$
\epsilon :
(X_{affine, \etale}, \mathcal{O}_X)
\longrightarrow
(X_{affine}, \mathcal{O}_X)
$$
을 얻는다. 이는 Cohomology on Sites, Section
\ref{sites-cohomology-section-compare}의 사상이다. 이 절에서는
$D_\QCoh(\mathcal{O}_X)$를 Cohomology on Sites, Section
\ref{sites-cohomology-section-modules-cohomology}에서 도입한 범주
$\mathit{QC}(X_{affine}, \mathcal{O}_X)$와 동일시할 것이다.

\begin{lemma}
\label{spaces-perfect-lemma-DQCoh-alternative-small}
위 상황에서 다음 삼각범주들 사이에는 표준 완전동치들이 있다.
\begin{enumerate}
\item $D_\QCoh(\mathcal{O}_X)$,
\item $D_\QCoh(X_{affine, \etale}, \mathcal{O}_X)$,
\item $D_\QCoh(X_{affine}, \mathcal{O}_X)$,
\item $\mathit{QC}(X_{affine}, \mathcal{O}_X)$.
\end{enumerate}
\end{lemma}

\begin{proof}
$U \to V \to X$가 에탈 사상들이고 $U$와 $V$가 아핀이면, 환
사상 $\mathcal{O}_X(V) \to \mathcal{O}_X(U)$은 평탄하다. 따라서
(3)과 (4)의 동치는 Cohomology on Sites, Lemma
\ref{sites-cohomology-lemma-cartesian-flat-quasi-coherent}의 특수한
경우이다 (그 증명은 명제의 의미도 명확히 한다).

\medskip\noindent
보조정리 앞의 논의는 환 달린 토포스들의 동치
$(\Sh(X_{affine, \etale}), \mathcal{O}_X) \to
(\Sh(X_\etale), \mathcal{O}_X)$가 있고, 따라서 가군들의
아벨 범주 사이에 동치가 있음을 보여준다. 준연접 가군이라는 개념은
내재적이므로 (Modules on Sites, Lemma
\ref{sites-modules-lemma-special-locally-free}), 이 동치는
준연접 가군들의 부분범주를 보존한다. 따라서 (1)과 (2)의
삼각범주 사이에 표준 완전동치를 얻는다.

\medskip\noindent
(2)와 (3)의 삼각범주 사이의 완전동치를 얻기 위해 위 사상
$\epsilon : (X_{affine, \etale}, \mathcal{O}_X) \to
(X_{affine}, \mathcal{O}_X)$에 Cohomology on Sites, Lemma
\ref{sites-cohomology-lemma-compare-topologies-derived-adequate-modules}를
적용할 것이다. $\mathcal{B} = \Ob(X_{affine})$로 놓고,
$\mathcal{A} \subset \textit{PMod}(X_{affine}, \mathcal{O})$를 전층
$\mathcal{F}$들 가운데
$\mathcal{F}(V) \otimes_{\mathcal{O}_X(V)} \mathcal{O}_X(U)
\to \mathcal{F}(U)$가 동형이 되는 것들로 이루어진 충만한
부분범주로 놓는다. Descent on Spaces, Lemma
\ref{spaces-descent-lemma-quasi-coherent-alternative-small}에 의해
$\mathcal{A}$의 대상들은 정확히 에탈 위상에서 층이며
$(X_{affine, \etale}, \mathcal{O}_X)$ 위의 준연접 가군인
것들이다. 한편 Modules on Sites, Lemma
\ref{sites-modules-lemma-chaotic-quasi-coherent}에 의해
$\mathcal{A}$의 대상들은 정확히 $ (X_{affine}, \mathcal{O}_X)$,
즉 혼돈 위상 위의 준연접 가군들이다. 그러므로 Cohomology on
Sites, Lemma
\ref{sites-cohomology-lemma-compare-topologies-derived-adequate-modules}가
적용됨을 보이면, $\epsilon^*$와 $R\epsilon_*$를 사용하여 실제로
(2)와 (3)의 범주 사이의 표준 동치를 얻는다.

\medskip\noindent
다음 네 조건을 확인해야 한다.
\begin{enumerate}
\item $\mathcal{A}$의 모든 대상은 에탈 위상에 대한 층이다.
이는 위에서 보았다.
\item $\mathcal{A}$는
$\textit{Mod}(X_{affine, \etale}, \mathcal{O}_X)$의 약한 세르
부분범주이다. 위에서
$\mathcal{A} = \QCoh(X_{affine, \etale}, \mathcal{O}_X)$임을
보았고, 또한 이들이 동치
$\textit{Mod}(X_{affine, \etale}, \mathcal{O}) =
\textit{Mod}(X_\etale, \mathcal{O}_X)$를 통해 $X$ 위의 준연접
가군들에 대응함을 보았다. 따라서 이 결과는 Properties of Spaces,
Lemma \ref{spaces-properties-lemma-properties-quasi-coherent}와
Homology, Lemma
\ref{homology-lemma-characterize-weak-serre-subcategory}에서
따른다. % P09-ERR-0490
\item $X_{affine}$의 모든 대상은 혼돈 위상에서 원소들이
$\mathcal{B}$에 속하는 덮개를 가진다. 이는 $\mathcal{B}$가 모든
대상을 포함하기 때문에 성립한다.
\item $U$가 $X_{affine}$의 임의의 대상이고 $\mathcal{F}$가
$\mathcal{A}$에 속하면 $H^p_{Zar}(U, \mathcal{F}) = 0$이다.
여기서 $p > 0$이다. 이는 아핀 위 준연접 가군의
코호몰로지가 소멸하기 때문이다. Cohomology of Spaces, Section
\ref{spaces-cohomology-section-higher-direct-image}의 논의와
Cohomology of Schemes, Lemma
\ref{coherent-lemma-quasi-coherent-affine-cohomology-zero}를 보라.
\end{enumerate}
이로써 증명이 끝난다.
\end{proof}

\begin{remark}
\label{spaces-perfect-remark-QC-compare}
$S$를 스킴이라 하자. $X$를 $S$ 위의 대수공간이라 하자. 나중에
$\mathit{QC}((\textit{Aff}/X), \mathcal{O})$도
$D_\QCoh(\mathcal{O}_X)$와 표준적으로 동치임을 보일 것이다.
Sheaves on Stacks, Proposition
\ref{stacks-sheaves-proposition-QC-compare}를 보라.
\end{remark}
